% Native Chinese working draft; public TeX/font dependencies remain required.
\begin{filecontents*}{fga-bibliography.bib}

  @article{Har1966,
  shorthand = {Har1966},
  sortname = {Har1966},
  title = {Connectedness of the {{Hilbert}} scheme},
  author = {Hartshorne, R.},
  year = {1966},
  journal = {Pub. math. de l'IH\'{E}S},
  volume = {29},
  pages = {5--48},
}

    @article{Ses1962,
  shorthand = {Ses1962},
  sortname = {Ses1962},
  title = {La vari\'{e}t\'{e} de Picard d'une vari\'{e}t\'{e} compl\`{e}te},
  author = {Seshadri, C.S.},
  year = {1962},
  journal = {Annali di mat. pura ed app.},
  volume = {57},
  pages = {117--142},
}

    @article{Oor1962,
  shorthand = {Oor1962},
  sortname = {Oor1962},
  title = {Sur le sch\'{e}ma de Picard},
  author = {Oort, F.},
  year = {1962},
  journal = {Bull. Soc. Math. France},
  volume = {90},
  pages = {1--14},
}

    @article{Mum1961,
  shorthand = {Mum1961},
  sortname = {Mum1961},
  title = {An elementary theorem in geometric invariant theory},
  author = {Mumford, D.},
  year = {1961},
  journal = {Bull. Amer. Math. Soc.},
  volume = {67},
  pages = {483--486},
}

    @article{Ser1961,
  shorthand = {Ser1961},
  sortname = {Ser1961},
  title = {Exemples de vari\'{e}t\'{e}s projectives en caract\'{e}ristique $p$ non relevables en caract\'{e}ristique 0},
  author = {Serre, J.-P.},
  year = {1961},
  journal = {Proc. Nat. Acad. Sc. U.S.A.},
  volume = {47},
  pages = {108--109},
}

    @article{Mum1961a,
  shorthand = {Mum1961a},
  sortname = {Mum1961a},
  title = {The topology of normal singularities of an algebraic surface and a criterion for simplicity},
  author = {Mumford, D.},
  year = {1961},
  journal = {Pub. math. de l'IH\'{E}S},
  volume = {9},
  pages = {5--22},
}

    @article{Ser1958,
  shorthand = {Ser1958},
  sortname = {Ser1958},
  title = {Corps locaux et isog\'{e}nies},
  author = {Serre, J.-P.},
  year = {1960},
  journal = {S\'{e}minaire Bourbaki},
  volume = {11},
  pages = {Talk no. 185},
}

    @article{GD1960,
  shorthand = {GD1960},
  sortname = {GD1960},
  title = {{{El\'{e}ments de G\'{e}om\'{e}trie Alg\'{e}brique}}},
  author = {Grothendieck, A. and Dieudonn\'{e}, J.},
  year = {1960/67},
  journal = {Pub. math. de l'IH\'{E}S},
  volume = {4, 8, 11, 17, 20, 24, 28, 32},
}

    @article{Che1960,
  shorthand = {Che1960},
  sortname = {Che1960},
  title = {Sur la th\'{e}orie de Picard},
  author = {Chevalley, C.},
  year = {1960},
  journal = {Amer. J. of Math.},
  volume = {82},
  pages = {435--490},
}

    @book{Gro1960b,
  shorthand = {Gro1960b},
  sortname = {Gro1960b},
  title = {{{S\'{e}minaire de G\'{e}om\'{e}trie Alg\'{e}brique}}},
  author = {Grothendieck, A.},
  year = {1960/61},
  publisher = {{Paris, Institut des Hautes \'{E}tudes Scientifiques}}
}

    @article{Gro1960a,
  shorthand = {Gro1960a},
  sortname = {Gro1960a},
  title = {Techniques de construction en g\'{e}om\'{e}trie analytique, {{I}}--{{X}}},
  author = {Grothendieck, A.},
  year = {1960/61},
  journal = {S\'{e}minaire Cartan},
  volume = {13},
  pages = {Talks no. 7--17},
}

    @inproceedings{Gro1960,
  shorthand = {Gro1960},
  sortname = {Gro1960},
  title = {The cohomology theory of abstract algebraic varieties},
  booktitle = {Proceedings of the {{International Congress}} of {{Mathematicians}} 1958},
  author = {Grothendieck, A},
  year = {1960},
  publisher = {{Cambridge University Press}},
}

    @book{Lan1959,
  shorthand = {Lan1959},
  sortname = {Lan1959},
  title = {Abelian varieties},
  author = {Lang, S.},
  year = {1959},
  publisher = {Interscience Publishers},
}

    @article{Ser1958a,
  shorthand = {Ser1958a},
  sortname = {Ser1958a},
  title = {Espaces fibr\'{e}s alg\'{e}briques},
  author = {Serre, J.-P.},
  year = {1958},
  journal = {S\'{e}minaire Chevalley},
  volume = {3},
  pages = {Talk no. 1},
}

    @article{GR1958,
  shorthand = {GR1958},
  sortname = {GR1958},
  title = {{{Komplexe R\"{a}ume}}},
  author = {Grauert, H. and Remmert, R.},
  year = {1958},
  journal = {Math. Annalen},
  volume = {136},
  pages = {245--318},
}

    @article{BS1958,
  shorthand = {BS1958},
  sortname = {BS1958},
  title = {Le th\'{e}or\`{e}me de {{Riemann}}\textendash{{Roch}}},
  author = {Borel, A. and Serre, J.-P.},
  journal = {Bull. Soc. Math. France},
  year = {1958},
  volume = {86},
  pages = {97--136},
}

    @article{Mur1958,
  shorthand = {Mur1958},
  sortname = {Mur1958},
  title = {On a connectedness theorem for a birational transformation at a simple point},
  author = {Murre, J.P.},
  year = {1958},
  journal = {Amer. J. of Math.},
  volume = {80},
  pages = {3--15},
}

    @inproceedings{Ser1958b,
  shorthand = {Ser1958b},
  sortname = {Ser1958b},
  title = {Sur la topologie des vari\'{e}t\'{e}s alg\'{e}briques en caract\'{e}ristique $p$},
  booktitle = {Symposium internacional de topologia algebraica, {{Mexico}}, 1956},
  author = {Serre, J.-P.},
  year = {1958},
  publisher = {{Universidad nacional autonoma, Mexico}},
  pages = {24--53},
}

    @article{Car1957,
  shorthand = {Car1957},
  sortname = {Car1957},
  title = {Des groups $\mathrm{Ext}^s(A,B)$},
  author = {Cartier, P.},
  year = {1957},
  journal = {S\'{e}minaire A. Grothendieck: Alg\`{e}bre homologique},
  volume = {1},
  pages = {Talk no. 3},
}

    @article{Gro1957,
  shorthand = {Gro1957},
  sortname = {Gro1957},
  title = {Sur quelques points d'alg\`{e}bre homologique},
  author = {Grothendieck, A.},
  year = {1957},
  journal = {Tohoku math. J.},
  volume = {9},
  pages = {119--221},
}

    @article{Mat1957,
  shorthand = {Mat1957},
  sortname = {Mat1957},
  title = {The criteria for algebraic equivalence and the torsion group},
  author = {Matsusaka, T.},
  year = {1957},
  journal = {Amer. J. Math.},
  volume = {79},
  pages = {53--66},
}

    @article{Kod1956,
  shorthand = {Kod1956},
  sortname = {Kod1956},
  title = {Characteristic linear systems of complete continuous systems},
  author = {Kodaira, K.},
  year = {1956},
  journal = {Amer. J. of Math.},
  volume = {78},
  pages = {716--744},
}

    @article{Ser1956,
  shorthand = {Ser1956},
  sortname = {Ser1956},
  title = {G\'{e}om\'{e}trie alg\'{e}brique et g\'{e}om\'{e}trie analytique},
  author = {Serre, J.-P.},
  year = {1956},
  journal = {Ann. Institut Fourier Grenoble},
  volume = {6},
  pages = {1--42},
}

    @inproceedings{Ser1956a,
  shorthand = {Ser1956a},
  sortname = {Ser1956a},
  title = {Sur la dimension homologique des anneaux et des modules {{Noeth\'{e}riens}}},
  booktitle = {Proc. {{Intern}}. {{Symp}}. on Alg. Number Theory, {{Tokyo}} and {{Nikko}}, 1955},
  author = {Serre, J.-P.},
  year = {1956},
  pages = {175--189},
  publisher = {{Science Council of Japan, Tokyo}},
}

    @article{Ser1955,
  shorthand = {Ser1955},
  sortname = {Ser1955},
  title = {Faisceaux alg\'{e}briques coh\'{e}rents},
  author = {Serre, J.-P.},
  year = {1955},
  journal = {Annals of Math.},
  volume = {61},
  pages = {197--278},
}

    @article{Zap1945,
  shorthand = {Zap1945},
  sortname = {Zap1945},
  title = {Sull'esistenza, sopra la superficie algebriche, di sistemi continui completi infiniti, la cui curva e a serie caratteristica incompleta},
  author = {Zappa, G.},
  journal = {Acta Pontif. Acad. Scient.},
  year = {1945},
  volume = {9},
  pages = {91--93}
}

    

\end{filecontents*}
% Chinese working edition. Original book class, geometry, headings and Latin font retained.
\documentclass[a4paper,oneside,fleqn]{book}
\usepackage{geometry}
% Upstream running heads use two lines; preserve the text block's top position.
\setlength{\headheight}{22pt}
\addtolength{\topmargin}{-10pt}
\PassOptionsToPackage{no-math}{fontspec}
\usepackage[fontset=none,scheme=plain]{ctex}
\setCJKmainfont{Noto Serif CJK SC}[ItalicFont=Noto Serif CJK SC,ItalicFeatures={FakeSlant=0.2}]
\setCJKsansfont{Noto Sans CJK SC}
\setCJKmonofont{Noto Sans Mono CJK SC}
\def\mode{WebMode}
\usepackage{ifthen}

% BEGIN preamble.tex
%!TEX root = main.tex
\usepackage[T1]{fontenc}
\usepackage[final]{microtype}
\usepackage{amsmath,amssymb}
\usepackage{amsthm,mathtools}
\usepackage{amscd}
\usepackage{etoolbox}
\usepackage[cal=esstix,scr=rsfs]{mathalfa}



% DOUBLE QUOTES
\usepackage[autostyle=false,style=english]{csquotes}
\MakeOuterQuote{"}



% HEADERS AND FOOTERS
\usepackage{fancyhdr}
\renewcommand{\chaptermark}[1]{\markboth{\thechapter. #1}{}}
\renewcommand{\sectionmark}[1]{\markboth{\thesection. #1}{}}
\fancyhf{}
\ifthenelse{\equal{\mode}{WebMode}}{%
  \fancyhead[C]{\small FGA~\thepart\\\leftmark}%
  \renewcommand{\headrulewidth}{0.5pt}%
}{%
  \fancyhead[LO]{FGA~\thepart}%
  \fancyhead[RE]{\leftmark}%
  \renewcommand{\headrulewidth}{0pt}%
}
\fancyfoot[C]{\thepage}



% BIBLIOGRAPHY
\usepackage[backend=bibtex]{biblatex}
\renewbibmacro{in:}{%
  \ifboolexpr{%
     test {\ifentrytype{article}}%
  }{}{\printtext{\bibstring{in}\intitlepunct}}%
}



% Load titlesec before hyperref so its heading destinations are emitted with the title.
\usepackage{titlesec}

% LINKS
\ifthenelse{\equal{\mode}{WebMode}}{%
  \usepackage{xcolor}%
  \usepackage[colorlinks=true,linkcolor={blue!60!black},citecolor={blue!60!black},urlcolor={purple!60!black}]{hyperref}%
}{%
  \usepackage[colorlinks=false,hidelinks]{hyperref}
}

\usepackage[nameinlink,capitalize]{cleveref}
\crefname{section}{§}{§§}
\crefname{subsection}{§}{§§}
\labelformat{subsection}{#1}



% HEADINGS

\titleformat{\part}[display]{}{\huge FGA~\thepart}{8pt}{\huge}
\titleclass{\part}{top}
\titlespacing{\part}{0pt}{0pt}{*6}

\titleformat{\chapter}[hang]{}{\Large\thechapter.}{8pt}{\Large}
\ifthenelse{\equal{\mode}{WebMode}}{%
  \titleclass{\chapter}{straight}%
  \titlespacing{\chapter}{0pt}{*8}{*4}%
}{%
  \titlespacing{\chapter}{0pt}{0pt}{*4}%
}

\titleformat{\section}[hang]{}{\large\thesection.}{8pt}{\large}
\titlespacing{\section}{0pt}{*8}{*2}


\titleformat{\subsection}[hang]{}{\large\thesubsection.}{8pt}{\large}
\titlespacing{\subsection}{0pt}{*8}{*2}




% ENVIRONMENTS
\makeatletter
\def\skipsinglepar{%
  \@ifnextchar\par
    {\expandafter\skipsinglepar\@gobble}%
    {}%
}
\makeatother


%% Abstract
\newenvironment{abstract}{\par\vspace{\topsep}\noindent\textbf{译者说明。}~}{\par\vspace{\topsep}}


%% Numbering for \ref
\makeatletter
\newcommand{\setnumber}[1]{\phantomsection\def\@currentlabel{#1}}
\makeatother


%% Numbered equations
\newenvironment{eqn}[1]{\vspace{-1em}}{}


%% Original citations
\newenvironment{original}{\par\vspace{\topsep}\hrule\vspace{\topsep}\noindent}{\vspace{\topsep}\hrule\par\vspace{\topsep}}


%% Custom inner environment
\newenvironment{inenv}[4][]{%
  % first argument (optional) is the title
  % second argument (mandatory) is the number
  % third argument (mandatory) is the taxon
  % fourth argument (mandatory) is empty for non-italic
  \par\vspace{\topsep}\noindent%
  \ifstrempty{#2}{\phantomsection}{}%
  \ifstrempty{#2}{%if no number
    \ifstrempty{#1}{%if no title
      \textbf{#3.}%
    }{%if title
      \textbf{#3} (#1).%
    }%
  }{%if number
    \ifstrempty{#1}{%
      \textbf{#3~#2.}%
    }{%
      \textbf{#3~#2} (#1).%
    }%
  }%
  \ifstrempty{#4}{}{\itshape}~}%non-italic
  {\par\vspace{\topsep}}


%% Italicised environments
\newenvironment{theorem}[2][]%
  {\begin{inenv}[#1]{#2}{定理}{italic}}%
  {\end{inenv}}
\newenvironment{lemma}[2][]%
  {\begin{inenv}[#1]{#2}{引理}{italic}}%
  {\end{inenv}}
\newenvironment{proposition}[2][]%
  {\begin{inenv}[#1]{#2}{命题}{italic}}%
  {\end{inenv}}
\newenvironment{corollary}[2][]%
  {\begin{inenv}[#1]{#2}{推论}{italic}}%
  {\end{inenv}}
\newenvironment{statement}[2][]%
  {\begin{inenv}[#1]{#2}{断言}{italic}}%
  {\end{inenv}}
\newenvironment{conjecture}[2][]%
  {\begin{inenv}[#1]{#2}{猜想}{italic}}%
  {\end{inenv}}


%% Upright environments
\newenvironment{definition}[2][]%
  {\begin{inenv}[#1]{#2}{定义}{}}%
  {\end{inenv}}
\newenvironment{remark}[2][]%
  {\begin{inenv}[#1]{#2}{注}{}}%
  {\end{inenv}}
\newenvironment{remarks}[2][]%
  {\begin{inenv}[#1]{#2}{注}{}}%
  {\end{inenv}}
\newenvironment{example}[2][]%
  {\begin{inenv}[#1]{#2}{例}{}}%
  {\end{inenv}}
\newenvironment{scholium}[2][]%
  {\begin{inenv}[#1]{#2}{评注}{}}%
  {\end{inenv}}

% Keep the upstream book's layout and labels; localize only reader-facing names.
\renewcommand{\proofname}{证明}
\AtEveryBibitem{\phantomsection\label{\thefield{entrykey}}}

\renewcommand{\contentsname}{目录}

\crefname{chapter}{§}{§§}
\crefname{part}{报告}{报告}

% Stable subsection labels; never prepend the current section when reading a stored label.
\renewcommand\thesubsection{\thesection.\alph{subsection}}
\pdfstringdefDisableCommands{%
  \def\({}\def\){}%
  \def\underline#1{#1}%
  \def\mathcal#1{#1}\def\mathscr#1{#1}%
  \def\mathrm#1{#1}\def\mathit#1{#1}\def\operatorname#1{#1}%
}

% Report counters repeat in the original; PDF destinations must remain unique.
\newcommand\fgaAnchorPrefix{front}
\renewcommand\theHpart{\fgaAnchorPrefix.\arabic{part}}
\renewcommand\theHchapter{\fgaAnchorPrefix.\arabic{chapter}}
\renewcommand\theHsection{\theHchapter.\arabic{section}}
\renewcommand\theHsubsection{\theHsection.\arabic{subsection}}
\renewcommand\theHequation{\fgaAnchorPrefix.\arabic{chapter}.\arabic{equation}}

% END preamble.tex

\usepackage{Baskervaldx}
\AtEndPreamble{\setmainfont{Baskervaldx-Reg.otf}[BoldFont=Baskervaldx-Bol.otf,ItalicFont=Baskervaldx-Ita.otf,BoldItalicFont=Baskervaldx-BolIta.otf]}
\special{pdf:mapfile Baskervaldx.map}
\special{pdf:mapfile ESSTIX.map}
\addbibresource{fga-bibliography.bib}
\title{代数几何基础\\[0.5em]\large Fondements de la Géométrie Algébrique\\[0.5em]\normalsize Bourbaki 讨论班 1957--1962 年报告选编\\[1em]中文原生修订稿：八篇报告全稿，原文疑点待核}
\author{Alexander Grothendieck}
\date{1957--1962}
\begin{document}
\frontmatter
\pagestyle{plain}
\maketitle
\begin{center}
法文原汇编出版信息：\\
Secrétariat mathématique\\
11 rue Pierre Curie, Paris 5e\\
1962
\end{center}
\begin{abstract}
本稿以 Tim Hosgood 的英文原生稿为主要译源，并逐项核对已发现的法文差异。现已译入前言及全部八篇报告；原报告日期、局部书目、补记及勘误另据法文恢复。本轮已由独立模型审查覆盖八篇报告的全部源页，并逐项回读所作订正；这不等同于专业人工数学审定，也未重新证明原著全部结果。VI标题的脚注(1)正文仍未定位。原作者姓名、报告与陈述编号、公式和证明予以保留；已证实的英文误录、原本自身的疑误及尚未解决的覆盖问题，详见 SOURCE\_NOTES.json、source/reviews/ 与 PROGRESS.json。原本重号保留，内部引用标签作必要区分。
\end{abstract}
\begin{abstract}
以下保留英译本说明的内容：这是 Alexander Grothendieck 的《Fondements de la Géométrie Algébrique》的英译本。法文原笔记已由 Grothendieck Circle 扫描上传；各篇报告及勘误也可在 \url{https://www.numdam.org/actas/SB/} 找到。

英译者 Tim Hosgood（\url{https://thosgood.com}）对英译中引入的错误负全部责任，并不主张其中数学内容的任何权利。英译者的注释以斜体排印，并冠以“[Trans.]”。

英译本的全部源码可在 \url{https://github.com/thosgood/fga/} 查看、参与修改或提交订正。欢迎订正及意见。英译者诚挚感谢 Steve Hnizdur 发现了许多笔误和错误。
\end{abstract}
\setcounter{tocdepth}{1}
\tableofcontents
\chapter*{法文原汇编篇目}
原汇编目录所列篇幅如下。这里保留原版篇目和页数；本中文稿的排印页码见前面的目录。
\begin{center}
\begin{tabular}{p{0.77\textwidth}r}
篇目 & 原版页数\\\hline
代数几何基础：评注（Bourbaki 讨论班第14卷，1961/62，补编） & 11\\
凝聚代数层的对偶定理（Bourbaki 讨论班第9卷，1956/57，第149号） & 25\\
形式几何与代数几何（Bourbaki 讨论班第11卷，1958/59，第182号） & 28\\
下降技术与代数几何中的存在定理 I：一般理论；忠实平坦态射的下降（Bourbaki 讨论班第12卷，1959/60，第190号） & 29\\
下降技术与代数几何中的存在定理 II：形式模空间理论中的存在定理（Bourbaki 讨论班第12卷，1959/60，第195号） & 22\\
下降技术与代数几何中的存在定理 III：商预概形（Bourbaki 讨论班第13卷，1960/61，第212号） & 20\\
下降技术与代数几何中的存在定理 IV：Hilbert 概形（Bourbaki 讨论班第13卷，1960/61，第221号） & 28\\
下降技术与代数几何中的存在定理 V：Picard 概形；存在定理（Bourbaki 讨论班第14卷，1961/62，第232号） & 19\\
下降技术与代数几何中的存在定理 VI：Picard 概形；一般性质（Bourbaki 讨论班第14卷，1961/62，第236号） & 23
\end{tabular}
\end{center}
\mainmatter
\pagestyle{fancy}
\setcounter{part}{-1}
\renewcommand\thepart{}
\renewcommand\fgaAnchorPrefix{fga-foreword}

% BEGIN chapters/fga-foreword.tex
\part{前言}\label{fga-foreword}\thispagestyle{fancy}\par{}1957 至 1962 年间，我们在 \emph{Séminaire Bourbaki} 上作了八次关于代数几何基础的报告。除第一篇外，这些报告均使用\emph{概形}的语言。报告中陈述的所有结果都将在 Jean Dieudonné 与 Alexander Grothendieck 的《Éléments de Géométrie Algébrique》中得到论述。不过，这些报告的主要内容，目前既不见于 EGA 已经出版或仍在准备中的任何一章，也不见于其他任何书籍或论文；这种情形很可能还会持续几年。这是促使我们将这些报告汇集成册的主要原因：在一份完善的综述问世之前，先让读者能够接触概形理论中的若干思想和关键结果。阅读这些报告，还能使读者迅速熟悉上述结果和思想，而不必先处理系统论述中难以避免的繁复细节；同时，也能为进一步研读这种系统论述提供必要的动机。\par{}为方便读者，我们在这里按各篇报告汇集了一些评注和勘误，特别说明各篇写成以来取得的进展，并补充若干参考文献。\emph{[Trans.] 我们没有在此集中翻译这些评注和勘误，而是将它们插入正文的相应位置；除已直接改正而未另作标记的细小订正外，均以“\emph{[Comp.]}”开头，表示“Complément”。}\par{}这些文章中的许多结果，已在 \emph{Séminaire de géométrie algébrique de l’IHES} 以及后来于 1961 至 1962 年在 Harvard 举行的两个讨论班中得到详细讨论。后两个讨论班的第一个由我主持，第二个由 Mumford–Tate 主持，其笔记目前正由 Lichtenbaum 整理。\emph{[Trans.] 这些报告后来还多次提及 Mumford–Tate 的笔记，但这份笔记实际上并不存在：它从未完成，也没有任何草稿发表过（参见 \href{https://mathoverflow.net/a/452230/73622}{MO:452230/73622}）。}

\section*{原本所用缩写}
\begin{description}
\item[SGA] Séminaire de géométrie algébrique，Institut des Hautes Études Scientifiques，1960--1962。
\item[SHG] Séminaire Grothendieck，Harvard University，1961/62。
\item[SHMT] Séminaire Mumford--Tate，Harvard University，1961/62。
\item[EGA] Grothendieck (A.) et Dieudonné (J.)，\emph{Éléments de géométrie algébrique}，Paris，Presses universitaires de France，1960、1961、\ldots（Institut des Hautes Études Scientifiques，Publications mathématiques，4、8、11、\ldots）。
\item[TDTE] Grothendieck (Alexander)，\emph{Technique de descente et théorèmes d’existence en géométrie algébrique}，I--VI，\emph{Séminaire Bourbaki}，t.~12，1959/60，第190、195次报告；t.~13，1960/61，第212、221次报告；t.~14，1961/62，第232、236次报告。
\end{description}
\begin{flushright}原本评注日期：1962年5月。\end{flushright}

% END chapters/fga-foreword.tex

\clearpage
\renewcommand\thepart{\arabic{part}}
\renewcommand\fgaAnchorPrefix{fga1}

% BEGIN chapters/fga1.tex
\part{凝聚代数层的对偶定理}\label{fga1}\thispagestyle{fancy}\begin{original}{}\label{fga1-original-citation}\skipsinglepar\par{}A. Grothendieck.
“Théorème de dualité pour les faisceaux algébriques cohérents”.
\emph{Séminaire Bourbaki} \textbf{9}（1956–57），第 149 号报告；原报告日期：1957 年 5 月。
\url{http://www.numdam.org/book-part/SB_1956-1958__4__169_0/}\end{original}\par{}下面的结果受到 Serre 的“代数对偶定理”的启发，发现于 1955 年冬季和 1956 年冬季。利用射影空间上同调的一些相当初等的结果 \hyperref[Ser1955]{[3]}，并充分使用 Cartan–Eilenberg 同调代数在 \hyperref[Gro1957]{[2]} 中给出的形式，就能很简便地建立这些结果。\chapter{模层的 \texorpdfstring{\(\operatorname {Ext}\)}{Ext}}\thispagestyle{fancy}\label{fga1-1}\par{}\emph{（参见 [\hyperref[Gro1957]{2}，第 3、4 章]）}\par{}设 \(X\) 是一个拓扑空间，其上给定一个有单位元的环层 \(\mathcal {O}\)，环不必交换。考虑由 \(\mathcal {O}\)-模层组成的 Abel 范畴 \(\mathcal {C}^\mathcal {O}\)，这些模层也简称为 \(\mathcal {O}\)-模。已知这个范畴中的每个对象都有内射分解，因此可以定义具有熟知形式性质的 \(\operatorname {Ext}\) 函子。为避免混淆，更确切地说，我们用 \(\operatorname {Hom}_\mathcal {O}(X;\mathscr {A},\mathscr {B})\)，或简记为 \(\operatorname {Hom}(X;\mathscr {A},\mathscr {B})\)，表示从 \(\mathscr {A}\) 到 \(\mathscr {B}\) 的 \(\mathcal {O}\)-同态所成的 Abel \emph{群}；而 \(\mathscr {H}\kern  -2.5pt\mathit {om}_\mathcal {O}(\mathscr {A},\mathscr {B})\) 表示从 \(\mathscr {A}\) 到 \(\mathscr {B}\) 的同态芽所成的\emph{层}（其中 \(\mathscr {A},\mathscr {B}\in \mathcal {C}^\mathcal {O}\)）。固定 \(\mathscr {A}\in \mathcal {C}^\mathcal {O}\) 后，定义函子 \(h_\mathscr {A}\) 和 \(\underline {h}_\mathscr {A}\)，它们分别取值于 Abel 群范畴 \(\mathcal {C}\) 和 \(X\) 上的 Abel 层范畴 \(\mathcal {C}^Z\)，定义式为：
\begin{eqn}{1.1}\setnumber{1.1}\label{fga1-equation-1.1}\skipsinglepar\[       \begin {aligned}         h_\mathscr {A}(\mathscr {B})         &= \operatorname {Hom}_\mathcal {O}(X;\mathscr {A},\mathscr {B})       \\\underline {h}_\mathscr {A}(\mathscr {B})         &= \mathscr {H}\kern  -2.5pt\mathit {om}_\mathcal {O}(\mathscr {A},\mathscr {B}).       \end {aligned}     \tag{1.1}     \]\end{eqn}
函子 \(h_\mathscr {A}\) 和 \(\underline {h}_\mathscr {A}\) 都是左正合协变函子，因此可以考虑其右导出函子，分别记为 \(\operatorname {Ext}_\mathcal {O}^p(X;\mathscr {A},\mathscr {B})\) 和 \(\mathscr {E}\kern  -1.5pt\mathit {xt}_\mathcal {O}^p(\mathscr {A},\mathscr {B})\)。按照定义，有
\begin{eqn}{1.2}\setnumber{1.2}\label{fga1-equation-1.2}\skipsinglepar\[       \begin {gathered}         \operatorname {Ext}_\mathcal {O}^p(X;\mathscr {A},\mathscr {B})         = (\operatorname {R}^p h_\mathscr {A})(\mathscr {B})         = \operatorname {H}^p(\operatorname {Hom}_\mathcal {O}(X;\mathscr {A},C(\mathscr {B})))       \\\mathscr {E}\kern  -1.5pt\mathit {xt}_\mathcal {O}^p(\mathscr {A},\mathscr {B})         = (\operatorname {R}^p \underline {h}_\mathscr {A})(\mathscr {B})         = \operatorname {H}^p(\mathscr {H}\kern  -2.5pt\mathit {om}_\mathcal {O}(\mathscr {A},C(\mathscr {B})))       \end {gathered}     \tag{1.2}     \]\end{eqn}
其中 \(\operatorname {R}^p\) 表示取右导出函子，\(C(\mathscr {B})\) 表示 \(\mathscr {B}\) 在 \(\mathcal {C}^\mathcal {O}\) 中的任意内射分解。用 \(\Gamma \colon \mathcal {C}^Z\to \mathcal {C}\) 表示“截面”函子；回顾其右导出函子记为 \(B\mapsto \operatorname {H}^p(X,\mathscr {B})\)：
\begin{eqn}{1.3}\setnumber{1.3}\label{fga1-equation-1.3}\skipsinglepar\[       \operatorname {H}^p(X,\mathscr {B})       = (\operatorname {R}^p\Gamma )(\mathscr {B})       = \operatorname {H}^p(\Gamma (C(\mathscr {B}))).     \tag{1.3}     \]\end{eqn}
显然有 \(h_\mathscr {A}=\Gamma \underline {h}_\mathscr {A}\)；此外还可以证明，\(\underline {h}_\mathscr {A}\) 将内射对象映为 \(\Gamma \)-无环对象。由此可按熟知的方法得到：\begin{proposition}{1}\setnumber{1}\label{fga1-proposition-1}\skipsinglepar\par{}对每个 \(\mathcal {O}\)-模 \(\mathscr {A}\)，在 \(\mathcal {C}^\mathcal {O}\) 上都存在一个上同调谱函子，收敛于分次函子 \((\operatorname {Ext}_\mathcal {O}^\bullet (X;\mathscr {A},\mathscr {B}))\)，其初始页为
\begin{eqn}{1.4}\setnumber{1.4}\label{fga1-equation-1.4}\skipsinglepar\[         E_2^{p,q}(\mathscr {A},\mathscr {B})         = \operatorname {H}^p(X,\mathscr {E}\kern  -1.5pt\mathit {xt}_\mathcal {O}^q(\mathscr {A},\mathscr {B})).       \tag{1.4}       \]\end{eqn}\end{proposition}\par{}由此得到“\emph{边缘同态}”，以及一个我们不再写出的五项正合列。\footnote{中译校记：法文原本第 149-02 页为五项正合列；英文稿此处误作短正合列。}\begin{corollary}{1}\setnumber{1}\label{fga1-proposition-1-corollary-1}\skipsinglepar\par{}如果 \(\mathscr {A}\) 局部同构于 \(\mathcal {O}^n\)，则有典范同构
\begin{eqn}{1.5}\setnumber{1.5}\label{fga1-equation-1.5}\skipsinglepar\[         \operatorname {Ext}_\mathcal {O}^p(X;\mathscr {A},\mathscr {B})         \xleftarrow {\sim } \operatorname {H}^p(X,\mathscr {H}\kern  -2.5pt\mathit {om}_\mathcal {O}(\mathscr {A},\mathscr {B}))       \tag{1.5}       \]\end{eqn}
（由谱序列的边缘同态给出）。\footnote{中译校记：式 (1.5) 中的空间依照法文原本记为 $X$；英文稿误为小写 $x$。}特别地，有典范同构
\begin{eqn}{1.6}\setnumber{1.6}\label{fga1-equation-1.6}\skipsinglepar\[         \operatorname {Ext}_\mathcal {O}^p(X;\mathcal {O},\mathscr {B})         = \operatorname {H}^p(X,\mathscr {B}).       \tag{1.6}       \]\end{eqn}\end{corollary}\par{}为使用这些结果，需要知道如何具体描述 \(\operatorname {Ext}_\mathcal {O}^p(\mathscr {A},\mathscr {B})\)。这些函子可以局部计算，也就是说，如果 \(U\) 是 \(X\) 的一个开子集，则
\[     \mathscr {E}\kern  -1.5pt\mathit {xt}_\mathcal {O}^p(\mathscr {A},\mathscr {B})|U     = \mathscr {E}\kern  -1.5pt\mathit {xt}_{\mathcal {O}|U}^p(\mathscr {A}|U,\mathscr {B}|U)   \]
这是因为内射 \(\mathcal {O}\)-模在 \(U\) 上的限制仍是内射 \((\mathcal {O}|U)\)-模。此外，固定 \(x\in  X\) 后，有函子性的同态
\begin{eqn}{1.7}\setnumber{1.7}\label{fga1-equation-1.7}\skipsinglepar\[       \mathscr {H}\kern  -2.5pt\mathit {om}_\mathcal {O}(\mathscr {A},\mathscr {B})_x       \to  \operatorname {Hom}_{\mathcal {O}_x}(\mathscr {A}_x,\mathscr {B}_x)     \tag{1.7}     \]\end{eqn}
它们唯一地延拓为关于 \(\mathscr {B}\) 的上同调函子之间的同态：
\begin{eqn}{1.8}\setnumber{1.8}\label{fga1-equation-1.8}\skipsinglepar\[       \mathscr {E}\kern  -1.5pt\mathit {xt}_\mathcal {O}^p(\mathscr {A},\mathscr {B})_x       \to  \operatorname {Ext}_{\mathcal {O}_x}^p(\mathscr {A}_x,\mathscr {B}_x).     \tag{1.8}     \]\end{eqn}\begin{proposition}{2}\setnumber{2}\label{fga1-proposition-2}\skipsinglepar\par{}如果 \(\mathscr {A}\) 在 \(x\) 的某个邻域内同构于某个同态 \(\mathcal {O}^m\to \mathcal {O}^n\) 的余核，那么对于所有 \(p\)，\hyperref[fga1-equation-1.7]{式~\ref*{fga1-equation-1.7}} 都是同构。特别地，当 \(\mathscr {A}\) 是凝聚 \(\mathcal {O}\)-模 \hyperref[Ser1955]{[3]} 时，情形就是如此。\end{proposition}\begin{proposition}{3}\setnumber{3}\label{fga1-proposition-3}\skipsinglepar\par{}设 \(\mathscr {L}_\bullet =(\mathscr {L}_i)\) 是 \(\mathcal {O}\)-模 \(\mathscr {A}\) 的一个左分解，其各项都是局部同构于某个 \(\mathcal {O}^n\) 的 \(\mathcal {O}\)-模。那么 \(\mathscr {E}\kern  -1.5pt\mathit {xt}_\mathcal {O}(\mathscr {A},\mathscr {B})\) 可与 \(\operatorname {H}^\bullet (\mathscr {H}\kern  -2.5pt\mathit {om}_\mathcal {O}(\mathscr {L}_\bullet ,\mathscr {B}))\) 等同，而 \(\operatorname {Ext}_\mathcal {O}(X;\mathscr {A},\mathscr {B})\) 可与 \(X\) 关于复形 \(\mathscr {H}\kern  -2.5pt\mathit {om}_\mathcal {O}(\mathscr {L}_\bullet ,\mathscr {B})\) 的超上同调等同。\begin{proof}\par{}证明是标准的：考虑双复形 \(\mathscr {H}\kern  -2.5pt\mathit {om}_\mathcal {O}(\mathscr {L}_\bullet ,C(\mathscr {B}))\)，其中 \(C(\mathscr {B})\) 是 \(\mathscr {B}\) 的内射分解，并考虑从 \(\mathscr {H}\kern  -2.5pt\mathit {om}_\mathcal {O}(\mathscr {L}_\bullet ,\mathscr {B})\) 和 \(\mathscr {H}\kern  -2.5pt\mathit {om}_\mathcal {O}(\mathscr {A},C(\mathscr {B}))\) 到这个双复形的自然同态。\end{proof}\end{proposition}\par{}最后指出，\hyperref[fga1-equation-1.2]{式~\ref*{fga1-equation-1.2}} 中引入的两个 \(\operatorname {Ext}\) 函子，不仅是关于 \(\mathscr {B}\) 的上同调函子，而且是\emph{上同调双函子}，关于 \(\mathscr {B}\) 协变，关于 \(\mathscr {A}\) 反变。\chapter{\texorpdfstring{\(\operatorname {Ext}\)}{Ext} 中的复合律}\thispagestyle{fancy}\label{fga1-2}\par{}本节的结果由 Cartier 和 Yoneda 各自独立得到；详情可参见 Cartier 的报告 \hyperref[Car1957]{[1]}。设 \(\mathcal {C}\) 是 Abel 范畴，\(K\) 和 \(L\) 是 \(\mathcal {C}\) 中的两个分次对象。用 \(\operatorname {Hom}(K,L)\) 表示这样一个分次 Abel 群：其 \(n\) 次分量由从 \(K\) 到 \(L\) 的 \(n\) 次齐次同态组成，也就是由同态 \(K^i\to  L^{i+n}\) 组成的族 \((u_i)\)。如果 \(K\) 和 \(L\) 是复形（为固定约定，设微分的次数为 \(+1\)），则在 \(\operatorname {Hom}(K,L)\) 上引入如下微分算子：
\begin{eqn}{2.1}\setnumber{2.1}\label{fga1-equation-2.1}\skipsinglepar\[       \delta (u) = \mathrm {d}u + (-1)^{n+1}u\mathrm {d}       \quad \text {其中 }n=\deg (u)     \tag{2.1}     \]\end{eqn}
这使它成为一个微分次数为 \(+1\) 的复形。\(n\) 次闭链就是作为齐次映射与 \(u\) 反交换的 \(n\) 次映射。于是可以考虑 \(\operatorname {H}^\bullet (\operatorname {Hom}(K,L))\)；它是 \(K\) 与 \(L\) 的同伦型的不变量，有时记为 \(\operatorname {H}^\bullet (K,L)\)。如果还有第三个复形 \(M\)，则同态的复合定义一个与微分算子相容的配对 \(\operatorname {Hom}(K,L)\times \operatorname {Hom}(L,M)\to \operatorname {Hom}(K,M)\)，再取上同调，就得到配对
\begin{eqn}{2.2}\setnumber{2.2}\label{fga1-equation-2.2}\skipsinglepar\[       \operatorname {H}^\bullet (K,L)\times \operatorname {H}^\bullet (L,M) \to  \operatorname {H}^\bullet (K,M)     \tag{2.2}     \]\end{eqn}
记作 \((u,v)\mapsto  vu\)。这些配对显然满足结合律；特别地，\(\operatorname {H}^\bullet (K,K)\) 是一个有单位元的结合分次环，而 \(\operatorname {H}^\bullet (K,L)\)（分别地，\(\operatorname {H}^\bullet (L,K)\)）是这个环上的分次右模（分别地，左模），等等。在 \(0\) 次时，\hyperref[fga1-equation-2.2]{式~\ref*{fga1-equation-2.2}} 就是复形的容许同态的复合。最后，若复形正合列 \(0\to  K'\to  K\to  K''\to 0\) 满足：对每个 \(i\)，\(K'^i\) 可与 \(K^i\) 的一个直和因子等同，则得到群复形 \(\operatorname {Hom}(K'',L)\) 等所成的正合列，从而得到上边缘算子 \(\operatorname {H}^i(K',L)\to \operatorname {H}^{i+1}(K'',L)\)。同样，可以定义与 \(L\) 中的一个正合列相联系的上边缘算子。\hyperref[fga1-equation-2.2]{式~\ref*{fga1-equation-2.2}} 中的配对按通常的意义与这些上边缘算子相容。\par{}现在设 \(\mathcal {C}\) 是这样一个范畴：\(\mathcal {C}\) 中的每个对象 \(A\) 都有一个内射分解 \(C(A)\)。利用双复形定理的许多变体之一，可以看到
\[     \operatorname {H}^\bullet (C(A),C(B)) = \operatorname {H}^\bullet (\operatorname {Hom}(C(A),C(B)))   \]
典范同构于
\[     \operatorname {H}^\bullet (\operatorname {Hom}(A,C(B))) = \operatorname {Ext}^\bullet (A,B).   \]
上面讨论的上边缘算子给出 \(\operatorname {Ext}\) 的上边缘同态。此外，\hyperref[fga1-equation-2.2]{式~\ref*{fga1-equation-2.2}} 中的配对在这里给出结合的配对
\begin{eqn}{2.3}\setnumber{2.3}\label{fga1-equation-2.3}\skipsinglepar\[       \operatorname {Ext}^\bullet (A,B)\times \operatorname {Ext}^\bullet (B,C) \to  \operatorname {Ext}^\bullet (A,C)     \tag{2.3}     \]\end{eqn}
并且与上边缘同态相容。特别地，\(\operatorname {Ext}^\bullet (A,A)\) 是有单位元的结合分次环，等等。（同样可以证明，\(\operatorname {Ext}\) 函子作用于任意函子的导出函子；这里不使用这一事实。）\par{}如果所考虑的范畴是 \(X\) 上的 \(\mathcal {O}\)-模范畴 \(\mathcal {C}^\mathcal {O}\)，便得到配对
\begin{eqn}{2.4}\setnumber{2.4}\label{fga1-equation-2.4}\skipsinglepar\[       \operatorname {Ext}_\mathcal {O}^p(X;\mathscr {A},\mathscr {B})\times \operatorname {Ext}_\mathcal {O}^q(X;\mathscr {B},\mathscr {C}) \to  \operatorname {Ext}_\mathcal {O}^{p+q}(X;\mathscr {A},\mathscr {C})     \tag{2.4}     \]\end{eqn}
并可按照已经说明的方法将它们具体写出。用同样的方法，将 Abel 群范畴换成 \(X\) 上的 Abel 层范畴，并将 \(\operatorname {Hom}\) 函子换成 \(\mathscr {H}\kern  -2.5pt\mathit {om}\) 函子，就又能定义具有相同形式性质的配对，不过这次具有“局部性质”：
\begin{eqn}{2.5}\setnumber{2.5}\label{fga1-equation-2.5}\skipsinglepar\[       \mathscr {E}\kern  -1.5pt\mathit {xt}_\mathcal {O}^p(\mathscr {A},\mathscr {B})\times \mathscr {E}\kern  -1.5pt\mathit {xt}_\mathcal {O}^q(\mathscr {B},\mathscr {C}) \to  \mathscr {E}\kern  -1.5pt\mathit {xt}_\mathcal {O}^{p+q}(\mathscr {A},\mathscr {C}).     \tag{2.5}     \]\end{eqn}
注意到 \hyperref[fga1-equation-1.8]{式~\ref*{fga1-equation-1.8}} 中的同态与 \(\operatorname {Ext}\) 之间的配对相容，就能进一步说明这些配对。\par{}最后，回顾函子 \(\operatorname {H}^p(X,A)\) 之间还存在一种乘法结构，即杯积。于是可以看到，\hyperref[fga1-proposition-1]{命题~\ref*{fga1-proposition-1}} 的谱序列与这些乘法结构相容；更确切地说，谱序列 \(E(A,B)\) 与谱序列 \(E(B,C)\) 之间存在一个取值于谱序列 \(E(A,C)\) 的配对，它收敛于整体 \(\operatorname {Ext}\) 之间的配对，其初始页来自杯积及 \hyperref[fga1-equation-1.4]{式~\ref*{fga1-equation-1.4}} 右端局部 \(\operatorname {Ext}\) 之间的配对。特别地，由此得到“边缘同态”
\begin{eqn}{2.6}\setnumber{2.6}\label{fga1-equation-2.6}\skipsinglepar\[       \operatorname {Ext}_\mathcal {O}^n(X;\mathscr {A},\mathscr {B}) \to  \operatorname {H}^0(X;\mathscr {E}\kern  -1.5pt\mathit {xt}_\mathcal {O}^n(\mathscr {A},\mathscr {B}))     \tag{2.6}     \]\end{eqn}
\begin{eqn}{2.7}\setnumber{2.7}\label{fga1-equation-2.7}\skipsinglepar\[       \operatorname {H}^n(X,\mathscr {H}\kern  -2.5pt\mathit {om}_\mathcal {O}(\mathscr {A},\mathscr {B})) \to  \operatorname {Ext}_\mathcal {O}^n(X;\mathscr {A},\mathscr {B})     \tag{2.7}     \]\end{eqn}
与乘法结构相容。因此，如果只考虑局部同构于某个 \(\mathcal {O}^m\) 的层，结合 \hyperref[fga1-equation-1.5]{式~\ref*{fga1-equation-1.5}} 中的同构，便能用杯积完全说明整体 \(\operatorname {Ext}\) 的复合。\chapter{关于局部上同调的结果}\thispagestyle{fancy}\label{fga1-3}\par{}设 \(A\) 是有单位元的交换环，\(\mathfrak {J}\) 是其中的一个理想。对于任意 \(A\)-模 \(M\)，我们将定义函子性的同态
\begin{eqn}{3.1}\setnumber{3.1}\label{fga1-equation-3.1}\skipsinglepar\[       \begin {aligned}         \operatorname {Ext}_A^p(A/\mathfrak {J},M)         &\to  \operatorname {Hom}_A(\wedge ^p\mathfrak {J}/\mathfrak {J}^2,M\otimes  A/\mathfrak {J})       \\\operatorname {Tor}_p^A(A/\mathfrak {J},M)         &\leftarrow  (\wedge ^p\mathfrak {J}/\mathfrak {J}^2)\otimes \operatorname {Hom}_A(A/\mathfrak {J},M)       \end {aligned}     \tag{3.1}     \]\end{eqn}
其中张量积和外积均在环 \(A\) 上取；还应注意，\(\mathfrak {J}/\mathfrak {J}^2\) 实际上是一个 \(A/\mathfrak {J}\)-模，它作为 \(A\)-模的外幂与作为 \(A/\mathfrak {J}\)-模的外幂相同。\hyperref[fga1-equation-3.1]{式~\ref*{fga1-equation-3.1}} 中同态的定义，来自如下构造：对于 \(\mathfrak {J}\) 中的每一组元素 \(x=(x_1,\ldots ,x_p)\)，定义同态 \(\varphi _x\)
\begin{eqn}{3.2}\setnumber{3.2}\label{fga1-equation-3.2}\skipsinglepar\[       \begin {aligned}         \varphi _x\colon  \operatorname {Ext}_A^p(A/\mathfrak {J},M) &\to  M\otimes  A/\mathfrak {J}       \\\varphi _x\colon  \operatorname {Hom}_A(A/\mathfrak {J},M) &\to  \operatorname {Tor}_p^A(A/\mathfrak {J},M)       \end {aligned}     \tag{3.2}     \]\end{eqn}
使其满足以下条件：
\begin{enumerate}
    
 \item[i.]\(\varphi _{x_1,\ldots ,x_p}\) 关于元素组 \(x_i\in \mathfrak {J}\) 是交替的 \(A\)-多重线性映射；


    
 \item[ii.]只要某个 \(x_i\) 属于 \(\mathfrak {J}^2\)，\(\varphi _{x_1,\ldots ,x_p}\) 就为零。

  \end{enumerate}
事实上，(ii) 可由 (i) 推出，因为当 \(a\in \mathfrak {J}\) 时有 \(a\varphi _x=0\)；注意到 \hyperref[fga1-equation-3.2]{式~\ref*{fga1-equation-3.2}} 中的各个模都被 \(\mathfrak {J}\) 零化，即可看出这一点。\par{}为定义 \(\varphi _x\)，考虑复形 \(K_x\)：它的底层 \(A\)-模为 \(\wedge  A^p\)，微分是与 \(x\) 的内积 \(i_x\)，这里将后者看作 \(A^p\) 上分量为 \(x_1,\ldots ,x_p\) 的线性型。微分的次数为 \(-1\)，复形的次数为正，而它在 \(0\) 维的上同调为 \(A/(x_1A+\ldots +x_pA)\)。由于各个 \(x_i\) 都属于 \(\mathfrak {J}\)，便得到增广映射 \(K_{x,0}\to  A/\mathfrak {J}\)。因此，\(K_x\) 是一个由 \emph{自由}-模组成的增广复形，增广模为 \(A/\mathfrak {J}\)。由此得到熟知的同态
\[     \begin {aligned}       \operatorname {Ext}_A^\bullet (\operatorname {H}_0(K_x),M) &\to  \operatorname {H}^\bullet (\operatorname {Hom}_A(K_x,M))     \\\operatorname {Tor}_\bullet ^A(\operatorname {H}_0(K_x),M) &\leftarrow  \operatorname {H}_\bullet (K_x\otimes  M)     \end {aligned}   \]
再与由增广同态 \(\operatorname {H}_0(K_x)\to  A/\mathfrak {J}\) 诱导的、到 \(\operatorname {Ext}\) 和 \(\operatorname {Tor}\) 的同态复合，便得到同态
\begin{eqn}{3.3}\setnumber{3.3}\label{fga1-equation-3.3}\skipsinglepar\[       \begin {aligned}         \psi _x\colon  \operatorname {Ext}_A^\bullet (A/\mathfrak {J},M)         &\to  \operatorname {H}^\bullet (\operatorname {Hom}_A(K_x,M))       \\\psi _x\colon  \operatorname {Tor}_\bullet ^A(A/\mathfrak {J},M)         &\leftarrow  \operatorname {H}_\bullet (K_x\otimes  M).       \end {aligned}     \tag{3.3}     \]\end{eqn}
立即可以看出，右端在最高维 \(p\) 的上同调为 \(M/(x_1M+\ldots +x_pM)\)（分别为 \(M\) 中被每个 \(x_i\) 零化的元素所成的集合）。由于各个 \(x_i\) 属于 \(\mathfrak {J}\)，因而得到同态
\begin{eqn}{3.4}\setnumber{3.4}\label{fga1-equation-3.4}\skipsinglepar\[       \begin {aligned}         \operatorname {H}^p(\operatorname {Hom}_A(K_x,M)) &\to  M\otimes  A/\mathfrak {J}       \\\operatorname {H}_p(K_x\otimes  M) &\leftarrow  \operatorname {Hom}_A(A/\mathfrak {J},M).       \end {aligned}     \tag{3.4}     \]\end{eqn}
将 \hyperref[fga1-equation-3.3]{式~\ref*{fga1-equation-3.3}} 和 \hyperref[fga1-equation-3.4]{式~\ref*{fga1-equation-3.4}} 中的同态复合，就得到所要定义的 \hyperref[fga1-equation-3.2]{式~\ref*{fga1-equation-3.2}} 中的同态。条件 (i) 的验证较为繁琐，但没有困难。\begin{proposition}{4}\setnumber{4}\label{fga1-proposition-4}\skipsinglepar\par{}设 \(A\) 是有单位元的交换环，\((x_1,\ldots ,x_p)\) 是 \(A\) 中的一个元素序列，并且对于 \(1\leqslant  i\leqslant  p\)，\(x_i\) 在 \(A\) 对由 \((x_1,\ldots ,x_{i-1})\) 生成的理想所作的商环中的像不是零因子。设 \(\mathfrak {J}\) 是由各个 \(x_i\) 生成的理想。那么 \(\mathfrak {J}/\mathfrak {J}^2\) 是自由 \((A/\mathfrak {J})\)-模，以各个 \(x_i\) 的典范像为基；复形 \(K_x\) 是 \(A/\mathfrak {J}\) 的自由分解；对于每个 \(A\)-模 \(M\)，\hyperref[fga1-equation-3.1]{式~\ref*{fga1-equation-3.1}} 中 \(p\) 维的同态都是双射。只要 \(\mathfrak {J}\cdot  M=0\)，任意次数 \(i\) 中相应定义的同态也同样是双射。\end{proposition}\par{}（\hyperref[fga1-proposition-4]{命题~\ref*{fga1-proposition-4}} 中的关键在于 \(K_x\) 的无环性，其他结论都由此推出；在所给条件下，这一事实是熟知的。）\begin{corollary}{1}\setnumber{1}\label{fga1-proposition-4-corollary-1}\skipsinglepar\par{}设 \(A\) 和 \(\mathfrak {J}\) 如 \hyperref[fga1-proposition-4]{命题~\ref*{fga1-proposition-4}} 中所述，并进一步假定 \(A\) 是完美域 \(k\) 上维数为 \(n\) 的正则仿射代数，且 \(A/\mathfrak {J}\) 也是正则仿射代数。用 \(\Omega ^i(A)\) 和 \(\Omega ^i(A/\mathfrak {J})\) 表示 Kähler 微分模。那么存在与局部化相容的典范同构
\begin{eqn}{3.5}\setnumber{3.5}\label{fga1-equation-3.5}\skipsinglepar\[         \operatorname {Ext}_A^p(\Omega ^{n-p}(A/\mathfrak {J}),\Omega ^n(A)) = A/\mathfrak {J}.       \tag{3.5}       \]\end{eqn}\begin{proof}\par{}事实上，\(\Omega ^{n-p}(A/\mathfrak {J})\) 是秩为 \(1\) 的自由 \((A/\mathfrak {J})\)-模；同样，\(\Omega ^n(A)\) 是秩为 \(n\) 的自由 \(A\)-模。因此左端等于
\[         \operatorname {Ext}_A^p(A/\mathfrak {J},A) \otimes  \Omega ^{n-p}(A/\mathfrak {J})' \otimes  \Omega ^n(A)       \]
（其中记号 \('\) 表示对偶 \((A/\mathfrak {J})\)-模）。最后两个因子的张量积可与 \(\wedge ^p(\mathfrak {J}/\mathfrak {J}^2)\) 等同，所以整个式子可与 \(\operatorname {Ext}_A^p(A/\mathfrak {J},\wedge ^p(\mathfrak {J}/\mathfrak {J}^2))\) 等同，再由 \hyperref[fga1-proposition-4]{命题~\ref*{fga1-proposition-4}}，可与
\[         \operatorname {Hom}_A(\wedge ^p \mathfrak {J}/\mathfrak {J}^2,\wedge ^p \mathfrak {J}/\mathfrak {J}^2)       \]
也就是 \(A/\mathfrak {J}\) 等同。\end{proof}\end{corollary}\par{}特别地，\(\operatorname {Ext}_A^p(\Omega ^{n-p}(A/\mathfrak {J}),\Omega ^n(A))\) 中有一个与 \(A/\mathfrak {J}\) 的单位元对应的特定元素，称为理想 \(\mathfrak {J}\) 在 \(A\) 中的\emph{基本类}。（事实上，在更一般的条件下也可以定义它。）\hyperref[fga1-proposition-4-corollary-1]{推论~\ref*{fga1-proposition-4-corollary-1}} 可以写成更具几何性和整体性的形式：\begin{corollary}{2}\setnumber{2}\label{fga1-proposition-4-corollary-2}\skipsinglepar\par{}设 \(X\) 是代数闭域 \(k\) 上的非奇异簇，\(Y\) 是 \(X\) 的非奇异闭子簇，\(\mathcal {O}_X\) 是 \(X\) 的结构层，\(\mathcal {O}_Y\) 是 \(Y\) 的结构层，并将后者看成 \(\mathcal {O}_X\) 的商层。设 \(n\) 是 \(X\) 的维数，\(n-p\) 是 \(Y\) 的维数。用 \(\Omega _X\)（分别地，\(\Omega _Y\)）表示 \(X\)（分别地，\(Y\)）上正则微分形式芽所成的层。那么存在典范同构
\begin{eqn}{3.6}\setnumber{3.6}\label{fga1-equation-3.6}\skipsinglepar\[         \mathscr {E}\kern  -1.5pt\mathit {xt}_{\mathcal {O}_X}^p(\Omega _Y^{n-p},\Omega _X^n) = \mathcal {O}_Y       \tag{3.6}       \]\end{eqn}
以及
\begin{eqn}{3.6 bis}\setnumber{3.6 bis}\label{fga1-equation-3.6bis}\skipsinglepar\[         \mathscr {E}\kern -1.5pt\mathit {xt}_{\mathcal {O}_X}^p(\mathcal {O}_Y,\Omega _X^n) = \Omega _Y^{n-p}.       \tag{3.6 bis}       \]\end{eqn}\end{corollary}\par{}当 \(Y\) 是奇异簇时，\hyperref[fga1-equation-3.6bis]{式~\ref*{fga1-equation-3.6bis}} 可以作为 \(\Omega _Y^{n-p}\) 的\emph{定义}。更确切地说：\begin{proposition}{5}\setnumber{5}\label{fga1-proposition-5}\skipsinglepar\par{}设 \(X\) 是 \(n\) 维非奇异代数簇，\(Y\) 是 \(X\) 中维数为 \(q=n-p\) 的代数子集。设 \(\mathscr {F}\) 是 \(X\) 上支集包含于 \(Y\) 的凝聚代数层，\(\mathscr {L}\) 是 \(X\) 上的局部自由代数层。那么当 \(i<p\) 时，层 \(\mathscr {E}\kern  -1.5pt\mathit {xt}_{\mathcal {O}_X}^i(\mathscr {F},\mathscr {L})\) 为零；当 \(i=p\) 时，存在典范同构
\begin{eqn}{3.7}\setnumber{3.7}\label{fga1-equation-3.7}\skipsinglepar\[         \mathscr {E}\kern  -1.5pt\mathit {xt}_{\mathcal {O}_X}^p(\mathscr {F},\mathscr {L})         = \mathscr {H}\kern  -2.5pt\mathit {om}_{\mathcal {O}_X}(\mathscr {F},\mathscr {E}\kern  -1.5pt\mathit {xt}^p(\mathcal {O}_X/\mathfrak {J},\mathscr {L}))       \tag{3.7}       \]\end{eqn}
其中 \(\mathfrak {J}\) 表示 \(X\) 上任意一个零化 \(\mathscr {F}\) 且以 \(Y\) 为零点集的理想层。特别地，若 \(\mathscr {F}\) 是 \(Y\) 上的凝聚代数层，则
\begin{eqn}{3.7 bis}\setnumber{3.7 bis}\label{fga1-equation-3.7bis}\skipsinglepar\[         \mathscr {E}\kern  -1.5pt\mathit {xt}_{\mathcal {O}_X}^p(\mathscr {F},\mathscr {L})         = \mathscr {H}\kern  -2.5pt\mathit {om}_{\mathcal {O}_Y}(\mathscr {F},\mathscr {E}\kern  -1.5pt\mathit {xt}^p(\mathcal {O}_Y,\mathscr {L})).       \tag{3.7 bis}       \]\end{eqn}
最后，仍设 \(\mathscr {F}\) 是 \(Y\) 上的凝聚代数层，则层 \(\mathscr {E}^i=\mathscr {E}\kern  -1.5pt\mathit {xt}_{\mathcal {O}_X}^{p+i}(\mathscr {F},\Omega _X^n)\) 不依赖于将代数空间 \(Y\) 嵌入非奇异代数簇 \(X\) 的方式。\begin{proof}\par{}由于问题具有局部性，可以假设 \(X\) 是仿射的且 \(\mathscr {L}=\mathcal {O}_X\)。于是归结为一个交换代数问题，更确切地说是局部代数问题：若 \(A\) 是正则局部环，\(M\) 是一个支集维数为 \(\leqslant  q=n-p\) 的 \(A\)-模，则须证明当 \(i<p\) 时有 \(\operatorname {Ext}_A^i(M,A)=0\)，并且有 \(\operatorname {Ext}_A^p(M,A)=\operatorname {Hom}_A(M,\operatorname {Ext}^p(A/\mathfrak {J},A))\)；其中 \(\mathfrak {J}\) 是任意一个零化 \(M\)、且“维数”为 \(\leqslant  q\) 的理想。对于第一个断言，对 \(q\) 作归纳：一个直接的\emph{逐层剥离（dévissage）}将问题归结到 \(M\) 具有 \(A/\mathfrak {J}\) 的形式的情形；再将 \(\mathfrak {J}\) 换成更小的理想，并利用归纳假设和 \(\operatorname {Ext}\) 的正合列，就归结到 \(\mathfrak {J}\) 如 \hyperref[fga1-proposition-4]{命题~\ref*{fga1-proposition-4}} 那样由一个“参数系”生成的情形，此时结论是显然的。上述结果表明，固定一个“维数”为 \(\leqslant  q\) 的理想 \(\mathfrak {J}\) 后，\((A/\mathfrak {J})\)-模范畴上的反变函子 \(E(M)=\operatorname {Ext}_A^p(M,A)\) 是左正合的；而且它把直和变成直积，由此容易得到 \(E(M)=\operatorname {Hom}_A(M,E(A))\)。最后，\hyperref[fga1-proposition-5]{命题~\ref*{fga1-proposition-5}} 的最后一个断言较为微妙，它来自一个局部对偶定理对 \(E^i(F)\) 的内在刻画；这里无法陈述这个定理。\end{proof}\end{proposition}\begin{corollary}{}\label{fga1-proposition-5-corollary}\skipsinglepar\par{}用 \(\omega _Y^q\) 表示层 \(\mathscr {E}\kern -1.5pt\mathit {xt}_{\mathcal {O}_X}^p(\mathcal {O}_Y,\Omega _X^n)\)。那么对于 \(Y\) 上的凝聚代数层 \(\mathscr {F}\)，存在函子性的同构
\begin{eqn}{3.8}\setnumber{3.8}\label{fga1-equation-3.8}\skipsinglepar\[         \mathscr {E}\kern  -1.5pt\mathit {xt}_{\mathcal {O}_X}^p(\mathscr {F},\Omega _X^n)         = \mathscr {H}\kern  -2.5pt\mathit {om}_{\mathcal {O}_X}(\mathscr {F},\omega _Y^q).       \tag{3.8}       \]\end{eqn}\end{corollary}\chapter{与子簇对应的上同调类}\thispagestyle{fancy}\label{fga1-4}\par{}以下各节中，\(X\) 表示定义在域 \(k\) 上的 \(n\) 维代数集；为简单起见，假定该域代数闭。除 \cref{fga1-6} 外，我们都假定 \(X\) 非奇异。用 \(\mathcal {O}_X\) 表示 \(X\) 的结构层，用 \(\Omega _X^\bullet =\bigcup _p\Omega _X^p\) 表示 \(X\) 上微分形式芽所成的层。若 \(Y\) 是 \(X\) 的闭子集，则将 \(Y\) 上的凝聚代数层与 \(X\) 上在 \(Y\) 外为零的凝聚代数层等同；特别地，对 \(\mathcal {O}_Y\) 和 \(\Omega _Y\) 也采用这个约定。\begin{lemma}{1}\setnumber{1}\label{fga1-lemma-1}\skipsinglepar\par{}设 \(\mathscr {F}\) 是 \(X\) 上支集维数为 \(\leqslant  n-p\) 的凝聚代数层，\(\mathscr {L}\) 是 \(X\) 上的局部自由凝聚代数层。那么当 \(i<p\) 时，\(\operatorname {Ext}_{\mathcal {O}_X}^i(X;\mathscr {F},\mathscr {L})\) 为零，并且存在典范同构
\begin{eqn}{4.1}\setnumber{4.1}\label{fga1-equation-4.1}\skipsinglepar\[         \operatorname {Ext}_\mathcal {O}^p(X;\mathscr {F},\mathscr {L})         = \operatorname {H}^0(X,\mathscr {E}\kern  -1.5pt\mathit {xt}_\mathcal {O}^p(\mathscr {F},\mathscr {L})).       \tag{4.1}       \]\end{eqn}
若 \(\mathscr {F}\) 是 \(X\) 的一个维数为 \(\leqslant  n-p\) 的闭子集 \(W\) 上的凝聚代数层，则存在典范同构
\begin{eqn}{4.1 bis}\setnumber{4.1 bis}\label{fga1-equation-4.1bis}\skipsinglepar\[         \mathscr {E}\kern  -1.5pt\mathit {xt}_{\mathcal {O}_X}^p(\mathscr {F},\mathscr {L})         = \mathscr {H}\kern  -2.5pt\mathit {om}_{\mathcal {O}_X}(\mathscr {F}\otimes \mathscr {L}'\otimes \Omega _X^n,\omega _Y^{n-p})       \tag{4.1 bis}       \]\end{eqn}
其中 \(\omega _Y^{n-p}\) 是 \hyperref[fga1-proposition-5-corollary]{命题 5 的推论} 中定义的 \(Y\) 上的层（若 \(Y\) 非奇异，则可将它与 \(\Omega _Y^{n-p}\) 等同）。\begin{proof}\par{}\hyperref[fga1-equation-4.1]{式~\ref*{fga1-equation-4.1}} 中的公式直接来自 \hyperref[fga1-proposition-1]{命题~\ref*{fga1-proposition-1}} 的谱序列以及 \hyperref[fga1-proposition-5]{命题~\ref*{fga1-proposition-5}}。由 \hyperref[fga1-equation-3.8]{式~\ref*{fga1-equation-3.8}} 中的公式可写出
\[         \begin {aligned}           \mathscr {E}\kern  -1.5pt\mathit {xt}_{\mathcal {O}_X}^p(\mathscr {F},\mathscr {L})           &= \mathscr {L}\otimes (\Omega _X^n)'\otimes \mathscr {E}\kern  -1.5pt\mathit {xt}_{\mathcal {O}_X}(\mathscr {F},\Omega _X^n)         \\&= \mathscr {L}\otimes (\Omega _X^n)'\otimes \mathscr {H}\kern  -2.5pt\mathit {om}_{\mathcal {O}_X}(\mathscr {F},\omega _Y^q)         \\&= \mathscr {H}\kern  -2.5pt\mathit {om}_{\mathcal {O}_X}(\mathscr {F}\otimes \mathscr {L}'\otimes \Omega _X^n,\omega _Y^q)         \end {aligned}       \]
其中 \(q=n-p\)，由此便得到 \hyperref[fga1-equation-4.1bis]{式~\ref*{fga1-equation-4.1bis}} 中的公式。\end{proof}\end{lemma}\par{}特别取 \(\mathscr {F}=\mathcal {O}_Y\)、\(\mathscr {L}=\Omega _X^p\)，并注意到 \(\Omega _X^n\otimes (\Omega _X^p)'=\Omega _X^{n-p}\)，便得到典范同构
\begin{eqn}{4.2}\setnumber{4.2}\label{fga1-equation-4.2}\skipsinglepar\[       \operatorname {Ext}_{\mathcal {O}_X}^p(X;\mathcal {O}_Y,\Omega _X^p)       = \operatorname {Hom}_{\mathcal {O}_X}(X;\Omega _X^{n-p},\omega _Y^{n-p}).     \tag{4.2}     \]\end{eqn}
现在为简单起见，假定 \(Y\) \emph{非奇异}，于是 \(\omega _Y^{n-p}=\Omega _Y^{n-p}\)。存在从 \(\Omega _X^{n-p}\) 到 \(\Omega _Y^{n-p}\) 的自然同态，从而得到层 \(\mathscr {E}\kern  -1.5pt\mathit {xt}_{\mathcal {O}_X}^p(\mathcal {O}_Y,\Omega _X^p)\) 的一个典范截面 \(s_Y\)；若 \(Y\) 的各个分支的维数均为 \(n-p\)，便称其为层 \(\mathscr {E}\kern  -1.5pt\mathit {xt}_{\mathcal {O}_X}^p(\mathcal {O}_Y,\Omega _X^p)\) 的\emph{基本截面}。由 \hyperref[fga1-equation-4.1]{式~\ref*{fga1-equation-4.1}}，这个截面定义了 \(\operatorname {Ext}_{\mathcal {O}_X}^p(X;\mathcal {O}_Y,\Omega _X^p)\) 中的一个元素。而自然同态 \(\mathcal {O}_X\to \mathcal {O}_Y\) 定义一个同态
\[     \operatorname {Ext}_{\mathcal {O}_X}^p(X;\mathcal {O}_Y,\Omega _X^p)     \to  \operatorname {Ext}_{\mathcal {O}_X}^p(X;\mathcal {O}_X,\Omega _X^p)     = \operatorname {H}^p(X,\Omega _X^p).   \]
于是得到 \(\operatorname {H}^p(X,\Omega _X^p)\) 中的一个元素，记为 \(P_X(Y)\)，称为\emph{\(Y\) 在 \(X\) 中的上同调类}；它由 \(\mathscr {E}\kern  -1.5pt\mathit {xt}_{\mathcal {O}_X}^p(\mathcal {O}_Y,\Omega _X^p)\) 的截面 \(s_Y\) 经下面的同态图诱导：
\begin{eqn}{4.3}\setnumber{4.3}\label{fga1-equation-4.3}\skipsinglepar\[       \begin {CD}         \operatorname {Ext}_{\mathcal {O}_X}^p(X;\mathcal {O}_Y,\Omega _X^p) @>\sim >> \operatorname {H}^0(X,\mathscr {E}\kern  -1.5pt\mathit {xt}_{\mathcal {O}_X}^p(\mathcal {O}_Y,\Omega _X^p))       \\@VVV       \\\begin {gathered}           \operatorname {Ext}_{\mathcal {O}_X}^p(X;\mathcal {O}_X,\Omega _X^p)           \\=\operatorname {H}^p(X,\Omega _X^p)         \end {gathered}       \end {CD}     \tag{4.3}     \]\end{eqn}\par{}我们将 \(X\) 中 \(n-p\) 维非奇异不可约子簇所生成的自由 Abel 群的任意元素，定义为 \(n-p\) 维\emph{非奇异圈}。这样，映射 \(Y\mapsto  P(Y)\) 就可延拓为从 \(X\) 上 \(n-p\) 维非奇异圈群到群 \(\operatorname {H}^p(X,\Omega _X^p)\) 的同态。\par{}设 \(Z^{n-p}\) 和 \(Z'^{n-p'}\) 分别是 \(n-p\) 维和 \(n-p'\) 维非奇异圈。如果 \(Z\) 的每个分支与 \(Z'\) 的每个分支都横截相交，就称这两个圈\emph{横截相交}。这时圈 \(Z\cdot  Z'\) 有定义，并且是维数为 \(n-p-p'\) 的非奇异圈。于是有：\begin{theorem}{1}\setnumber{1}\label{fga1-theorem-1}\skipsinglepar\par{}若 \(Z^{n-p}\) 和 \(Z'^{n-p'}\) 是横截相交的非奇异圈，则
\begin{eqn}{4.4}\setnumber{4.4}\label{fga1-equation-4.4}\skipsinglepar\[         P_X(Z\cdot  Z') = P_X(Z)\cdot  P_X(Z')       \tag{4.4}       \]\end{eqn}
其中右端的乘法是杯积：
\[       \operatorname {H}^p(X,\Omega _X^p)\times \operatorname {H}^{p'}(X,\Omega _X^{p'})       \to  \operatorname {H}^{p+p'}(X,\Omega _X^{p+p'}).     \]
（这里假定 \(X\) 同构于某个射影空间的局部闭子集。）\par{}\hyperref[fga1-theorem-1]{定理~\ref*{fga1-theorem-1}} 中的最后一个假设，只用于推出 \(X\) 上的每个凝聚代数层都是某个局部自由凝聚代数层的商（Serre），从而都有一个由局部自由层组成的左分解。\begin{proof}\par{}为证明 \hyperref[fga1-theorem-1]{定理~\ref*{fga1-theorem-1}}，可假定 \(Z\) 和 \(Z'\) 是横截相交的不可约非奇异子簇 \(Y\) 和 \(Y'\)。设 \(\mathscr {L}_\bullet \) 是 \(\mathcal {O}_Y\) 的一个由局部自由层组成的左分解；由 \hyperref[fga1-proposition-3]{命题~\ref*{fga1-proposition-3}}，\hyperref[fga1-equation-4.3]{式~\ref*{fga1-equation-4.3}} 中的同态图可与下图等同：
\begingroup\footnotesize\[         \begin{gathered}
(\operatorname{R}^p\Gamma)(\Omega_X^p)
\xleftarrow{\alpha}
(\underline{\operatorname{R}}^p\Gamma)
\bigl(\mathscr{H}\kern-2.5pt\mathit{om}_{\mathcal{O}_X}(\mathscr{L}_\bullet,\Omega_X^p)\bigr)
\overset{\beta}{\xrightarrow{\sim}}
\Gamma\bigl(\operatorname{H}^p\bigl(\mathscr{H}\kern-2.5pt\mathit{om}_{\mathcal{O}_X}(\mathscr{L}_\bullet,\Omega_X^p)\bigr)\bigr)
\end{gathered}       \]\endgroup{}
其中 \(\beta \) 是同构，\(\Gamma \) 是 \(X\) 上 Abel 层范畴上的“截面群”函子，\(\underline {\operatorname {R}}^p\Gamma \) 是它在 \(p\) 维的超上同调，而 \(\operatorname {R}^p\Gamma \) 是它的第 \(p\) 个导出函子。为简单起见，假定 \(\mathscr {L}_0=\mathcal {O}_X\)，且增广映射 \(\mathscr {L}_0\to \mathcal {O}_Y\) 为自然同态（确实可以作此假设而无妨）；那么 \(\alpha \) 由复形同态 \(\mathcal {O}_X\to \mathscr {L}\) 诱导（将 \(\mathcal {O}_X\) 看成集中在 \(0\) 次的复形），这里用到：若 \(\mathscr {K}\) 是集中在 \(0\) 次的层复形，则 \(\underline {\operatorname {R}}^p\Gamma (\mathscr {K})=\operatorname {R}^p\Gamma (\mathscr {K}_0)\)。同态 \(\beta \) 是熟知的“边缘同态”。对 \(\mathcal {O}_{Y'}\) 的一个局部自由分解 \(\mathscr {L}'_\bullet \) 考虑类似的图，并考虑下列配对的交换图：
\par\vspace{1em}\begin{eqn}{4.5}\setnumber{4.5}\label{fga1-equation-4.5}\skipsinglepar\begingroup\scriptsize\[                \begin {CD}             \operatorname {R}^p\Gamma (\Omega _X^p)             @<<<             \underline {\operatorname {R}}^p\Gamma \big (\mathscr {H}\kern  -2.5pt\mathit {om}_{\mathcal {O}_X}(\mathscr {L}_\bullet ,\Omega _X^p)\big )             @>\sim >>             \Gamma \big (\operatorname {H}^p\big (\mathscr {H}\kern  -2.5pt\mathit {om}_{\mathcal {O}_X}(\mathscr {L}_\bullet ,\Omega _X^p)\big )\big )           \\@. @. @.           \\\times  @. \times  @. \times            \\@. @. @.           \\\operatorname {R}^{p'}\Gamma (\Omega _X^{p'})             @<<<             \underline {\operatorname {R}}^{p'}\Gamma \big (\mathscr {H}\kern  -2.5pt\mathit {om}_{\mathcal {O}_X}(\mathscr {L}'_\bullet ,\Omega _X^{p'})\big )             @>\sim >>             \Gamma \big (\operatorname {H}^{p'}\big (\mathscr {H}\kern  -2.5pt\mathit {om}_{\mathcal {O}_X}(\mathscr {L}'_\bullet ,\Omega _X^{p'})\big )\big )           \\@VVV @VVV @VVV           \\\operatorname {R}^{p+p'}\Gamma (\Omega _X^{p+p'})             @<<<             \underline {\operatorname {R}}^{p+p'}\Gamma \big (\mathscr {H}\kern  -2.5pt\mathit {om}_{\mathcal {O}_X}(\mathscr {L}_\bullet \otimes \mathscr {L}'_\bullet ,\Omega _X^{p+p'})\big )             @>\sim >>             \Gamma \big (\operatorname {H}^{p+p'}\big (\mathscr {H}\kern  -2.5pt\mathit {om}_{\mathcal {O}_X}(\mathscr {L}_\bullet \otimes \mathscr {L}'_\bullet ,\Omega _X^{p+p'})\big )\big )           \end {CD}         \tag{4.5}         \]\endgroup\end{eqn}
右边两列中的配对由层复形的配对
\[         \mathscr {H}\kern  -2.5pt\mathit {om}_{\mathcal {O}_X}(\mathscr {L}_\bullet ,\Omega _X^p) \times  \mathscr {H}\kern  -2.5pt\mathit {om}_{\mathcal {O}_X}(\mathscr {L}'_\bullet ,\Omega _X^{p'})         \to  \mathscr {H}\kern  -2.5pt\mathit {om}_{\mathcal {O}_X}(\mathscr {L}_\bullet \otimes \mathscr {L}',\Omega _X^{p+p'})       \]
诱导，后者利用外积 \(\Omega _X^p\times \Omega _X^{p'}\to \Omega _X^{p+p'}\) 定义；第一列中的配对则是相对于外积的杯积。我断言，\hyperref[fga1-equation-4.5]{式~\ref*{fga1-equation-4.5}} 的最后一行可与类似于 \hyperref[fga1-equation-4.3]{式~\ref*{fga1-equation-4.3}} 的同构图等同，只需将 \(Y\) 换成 \(Y\cap  Y'\)，将 \(p\) 换成 \(p+p'\)。为此，只须证明 \(\mathscr {L}\otimes \mathscr {L}'\) 是 \(\mathcal {O}_{Y\cap  Y'}\) 的一个分解（它显然局部自由）。但由 \(Y\) 和 \(Y'\) 横截相交可知，
\[         \operatorname {H}_0(\mathscr {L}\otimes \mathscr {L}')         = \mathcal {O}_Y\otimes \mathcal {O}_{Y'}         = \mathcal {O}_{Y\cap  Y'}       \]
且当 \(i>0\) 时有
\[         \operatorname {H}_i(\mathscr {L}\otimes \mathscr {L}')         = \operatorname {Tor}_i^{\mathcal {O}_X}(\mathcal {O}_Y,\mathcal {O}_{Y'})         = 0       \]
于是 \hyperref[fga1-theorem-1]{定理~\ref*{fga1-theorem-1}} 可由公式
\begin{eqn}{4.6}\setnumber{4.6}\label{fga1-equation-4.6}\skipsinglepar\[           s_Y\cdot  s_{Y'} = s_{Y\cdot  Y'}         \tag{4.6}         \]\end{eqn}
推出（左端的乘法是 \hyperref[fga1-equation-4.5]{式~\ref*{fga1-equation-4.5}} 最后一列中的乘法）。\hyperref[fga1-equation-4.6]{式~\ref*{fga1-equation-4.6}} 中的公式是纯局部的，取 \(\mathscr {L}_\bullet \) 和 \(\mathscr {L}'_\bullet \) 为 \hyperref[fga1-proposition-4]{命题~\ref*{fga1-proposition-4}} 所述的分解，即可容易地证明它。同样还可证明（而且更容易），\(Z\mapsto  P_X(Z)\) 与笛卡尔积相容：
\begin{eqn}{4.7}\setnumber{4.7}\label{fga1-equation-4.7}\skipsinglepar\[           P_{X\times  X'}(Z\times  Z')           = P_X(Z)\otimes  P_{X'}(Z')         \tag{4.7}         \]\end{eqn}
（这个公式在下列情形成立：\(Z\)（分别地，\(Z'\)）是非奇异簇 \(X\)（分别地，\(X'\)）上的非奇异圈，并将 \(Z\times  Z'\) 看成 \(X\times  X'\) 上的非奇异圈。）由 \hyperref[fga1-equation-4.4]{式~\ref*{fga1-equation-4.4}} 和 \hyperref[fga1-equation-4.7]{式~\ref*{fga1-equation-4.7}} 又可推出，\(P_X(Z)\) 与非奇异簇的态射 \(f\colon  X\to  X'\) 下取“逆像”的运算相容：
\begin{eqn}{4.8}\setnumber{4.8}\label{fga1-equation-4.8}\skipsinglepar\[           P_X(f^{-1}(Z')) = f^*(P_{X'}(Z))         \tag{4.8}         \]\end{eqn}
此式在下列情形成立：\(Z\) 是 \(X'\) 上的非奇异圈，且 \(f\) 与 \(Z\) “横截”，也就是说，\(f\) 的图像与 \(X\times  X'\) 中的圈 \(X\times  Z'\) 横截。\end{proof}\end{theorem}\begin{corollary}{1}\setnumber{1}\label{fga1-theorem-1-corollary-1}\skipsinglepar\par{}设 \(X\) 和 \(X'\) 是射影空间中的非奇异局部闭簇，并假定 \(X'\) 完备。设 \(U\) 是 \(X\times  X'\) 上的非奇异圈，\(a\) 和 \(b\) 是 \(X'\) 中的两个点，并且 \(U\) 与圈 \(X\times (a)\) 和 \(X\times (b)\) 横截相交。设 \(Z\) 和 \(Z'\) 是 \(X\) 上的非奇异圈，满足 \(Z\times (a)=(X\times (a))\cdot  U\) 和 \(Z\times (b)=(X\times (b))\cdot  U\)。那么
\[       P_X(Z) = P_X(Z').     \]\begin{proof}\par{}设 \(f_a\colon  X\to  X\times  X'\) 由 \(f_a(x)=(x,a)\) 定义。由 \hyperref[fga1-equation-4.8]{式~\ref*{fga1-equation-4.8}} 有 \(P(Z)=f_a^*(P(U))\)，同样有 \(P(Z')=f_b^*(P(U))\)。利用 Künneth 公式
\[         \operatorname {H}^\bullet (X\times  X',\Omega _{X\times  X'}^\bullet )         = \operatorname {H}^\bullet (X,\Omega _X^\bullet )\otimes \operatorname {H}^\bullet (X',\Omega _{X'}^\bullet )       \]
以及 \(\operatorname {H}^0(X',\Omega _{X'})\) 仅由标量组成这一事实，容易看出 \(f_a^*=f_b^*\)，由此便得所需结论。\end{proof}\end{corollary}\par{}对于任意 \(x\in  X\)，\((x)\) 是 \(X\) 中余维为 \(n\) 的非奇异子簇，从而定义了 \(\operatorname {H}^n(X,\Omega _X^n)\) 中的一个元素 \(\varepsilon _x\)。若 \(X\) 是非奇异射影簇，则由 \hyperref[fga1-theorem-1-corollary-1]{推论~\ref*{fga1-theorem-1-corollary-1}}，\(\varepsilon _x\) 不依赖于所选的点 \(x\)，因而将其记为 \(\varepsilon _X\)，并称为 \(\operatorname {H}^n(X,\Omega _X^n)\) 的\emph{基本类}。\begin{remark}{}\label{fga1-4-remark-i}\skipsinglepar\par{}为得到完善的理论，必须对任意圈 \(Z\) 定义 \(P_X(Z)\)，并对圈的适当相交证明 \hyperref[fga1-theorem-1]{定理~\ref*{fga1-theorem-1}}。（撰写本报告时，这一点尚未在完全一般的情形下做到。）假设已经完成这些工作，那么 \hyperref[fga1-theorem-1-corollary-1]{推论~\ref*{fga1-theorem-1-corollary-1}} 就变成如下断言：
\emph{若 \(Z\) 和 \(Z'\) 是两个代数等价的圈，则 \(P_X(Z)=P_X(Z')\)}
（即使 \(Z\) 和 \(Z'\) 非奇异，这个断言似乎也不能从上述结果推出。）\end{remark}\begin{remark}{}\label{fga1-4-remark-ii}\skipsinglepar\par{}\emph{[Comp.]}
正如我在 1958 年国际数学家大会的报告中指出的（Grothendieck, A. “The cohomology theory of abstract algebraic varieties”，载于 \emph{Proceedings of the international Congress of Mathematicians [1958, Edinburgh]}，Cambridge University Press (1960)，103–118），这里提出的问题现在已经全部解决。\par{}关于凝聚层的对偶性，读者可进一步参阅\emph{同上文献}，第 112–115 页，以及 [\hyperref[GD1960]{GD1960}，III.2] 和 \hyperref[Gro1960b]{[Gro1960b]}。更系统的论述将见于 \hyperref[GD1960]{[GD1960]} 的后续一章（暂定计划中的第 IX 章）。\end{remark}\chapter{对偶定理}\thispagestyle{fancy}\label{fga1-5}\par{}本节中，\(X\) 表示 \(n\) 维非奇异射影簇。\begin{theorem}{2}\setnumber{2}\label{fga1-theorem-2}\skipsinglepar\par{}\(\operatorname {H}^n(X,\Omega _X^n)\) 的基本类 \(\varepsilon _X\) 是这个向量空间的一组基。\begin{proof}\par{}这个断言的证明将在后面给出。\end{proof}\end{theorem}\par{}由上述定理，可将 \(\operatorname {H}^n(X,\Omega _X^n)\) 与域 \(k\) 等同。现在考虑 \cref{fga1-2} 所述的配对，它们特别给出一个配对
\[     \operatorname {Ext}_{\mathcal {O}_X}^p(X;\mathcal {O}_X,\mathscr {F})\times \operatorname {Ext}_{\mathcal {O}_X}^{n-p}(X;\mathscr {F},\Omega _X^n)     \to  \operatorname {Ext}_{\mathcal {O}_X}^n(X;\mathcal {O}_X,\Omega _X^n)   \]
即
\begin{eqn}{5.1}\setnumber{5.1}\label{fga1-equation-5.1}\skipsinglepar\[       \operatorname {H}^p(X,\mathscr {F})\times \operatorname {Ext}_{\mathcal {O}_X}^{n-p}(X;\mathscr {F},\Omega _X^n)       \to  \operatorname {H}^n(X,\Omega _X^n).     \tag{5.1}     \]\end{eqn}
结合 \hyperref[fga1-theorem-2]{定理~\ref*{fga1-theorem-2}}，这个配对定义一个同态
\begin{eqn}{5.2}\setnumber{5.2}\label{fga1-equation-5.2}\skipsinglepar\[       \operatorname {Ext}_{\mathcal {O}_X}^{n-p}(X;\mathscr {F},\Omega _X^n)       \to  (\operatorname {H}^p(X,\mathscr {F}))'.     \tag{5.2}     \]\end{eqn}
这个同态关于 \(\mathscr {F}\) 具有函子性，并与正合列 \(0\to \mathscr {F}'\to \mathscr {F}\to \mathscr {F}''\to 0\) 对应的上边缘映射交换。\begin{theorem}{3}\setnumber{3}\label{fga1-theorem-3}\skipsinglepar\par{}\hyperref[fga1-equation-5.2]{式~\ref*{fga1-equation-5.2}} 中的同态是同构。\end{theorem}\par{}特别地，便重新得到 Serre 的下列结果：\begin{corollary}{}\label{fga1-theorem-3-corollary}\skipsinglepar\par{}设 \(E\) 是 \(X\) 上的代数向量丛，\(\mathcal {O}_X(E)\) 是 \(E\) 的正则截面芽所成的层。那么存在典范同构
（\emph{[Trans.] 原版笔记将这个公式编号为 (5.3)，后面另有同号公式。英文译者认为前者是笔误，并指出对 (5.3) 的引用似乎均指向后者。中译保留原本的重复印号，源码中的引用标签则指向后一个公式。}）
\[       (\operatorname {H}^p(X,\mathcal {O}_X(E)))'       = \operatorname {H}^{n-p}(X,\Omega _X^n\otimes \mathcal {O}_X(E')).     \tag{5.3}\]\begin{proof}\par{}只须应用 \hyperref[fga1-theorem-3]{定理~\ref*{fga1-theorem-3}} 和 \hyperref[fga1-proposition-1-corollary-1]{命题 1 的推论 1} 即可。\end{proof}\end{corollary}\par{}\hyperref[fga1-theorem-2]{定理~\ref*{fga1-theorem-2}} 和 \hyperref[fga1-theorem-3]{定理~\ref*{fga1-theorem-3}} 可由下列断言推出：\begin{statement}{D}\setnumber{D}\label{fga1-d}\skipsinglepar\par{}由 \hyperref[fga1-equation-5.1]{式~\ref*{fga1-equation-5.1}} 中的配对诱导的同态
\begin{eqn}{5.2 bis}\setnumber{5.2 bis}\label{fga1-equation-5.2bis}\skipsinglepar\[         \operatorname {Ext}_{\mathcal {O}_X}^{n-p}(X;\mathscr {F},\Omega _X^n)         \to  (\operatorname {H}^p(X,\mathscr {F}))'\otimes \mathscr {L}       \tag{5.2 bis}       \]\end{eqn}
（其中 \(\mathscr {L}=\operatorname {H}^n(X,\Omega _X^n)\)）是同构。\end{statement}\par{}下面证明 \hyperref[fga1-d]{断言~\ref*{fga1-d}} 蕴含 \hyperref[fga1-theorem-2]{定理~\ref*{fga1-theorem-2}}。设 \(k_x=\mathcal {O}_{(x)}\) 是由单个点 \(x\in  X\) 组成的簇的结构层，考虑典范同态 \(\mathcal {O}_X\to  k_x\) 及其对应的同态
\begin{eqn}{5.3}\setnumber{5.3}\label{fga1-equation-5.3}\skipsinglepar\[       \operatorname {H}^0(X,\mathcal {O}_X) \to  \operatorname {H}^0(X,k_x).     \tag{5.3}     \]\end{eqn}
其转置可与同态
\begin{eqn}{5.4}\setnumber{5.4}\label{fga1-equation-5.4}\skipsinglepar\[       \operatorname {Ext}_{\mathcal {O}_X}^n(X;k_x,\Omega _X^n)\otimes \mathscr {L}' \to  \operatorname {Ext}_{\mathcal {O}_X}^n(X;\mathcal {O}_X,\Omega _X^n)\otimes \mathscr {L}'     \tag{5.4}     \]\end{eqn}
等同，后者由对应于 \(\mathcal {O}_X\to  k_x\) 的 \(\operatorname {Ext}^n\) 之间的同态诱导，也就是说，来自
\begin{eqn}{5.5}\setnumber{5.5}\label{fga1-equation-5.5}\skipsinglepar\[       \operatorname {Ext}_{\mathcal {O}_X}^n(X;k_x,\Omega _X^n)       \to  \operatorname {Ext}_{\mathcal {O}_X}^n(X;\mathcal {O}_X,\Omega _X^n) = \operatorname {H}^n(X,\Omega _X^n)     \tag{5.5}     \]\end{eqn}
由于 \hyperref[fga1-equation-5.3]{式~\ref*{fga1-equation-5.3}} 是同构，\hyperref[fga1-equation-5.4]{式~\ref*{fga1-equation-5.4}} 也是同构，从而 \hyperref[fga1-equation-5.5]{式~\ref*{fga1-equation-5.5}} 也是同构。由 \hyperref[fga1-equation-4.2]{式~\ref*{fga1-equation-4.2}}，\(s_{(x)}\) 是 \(\operatorname {Ext}_{\mathcal {O}_X}^n(X;k_x,\Omega _X^n)\) 的一组基，所以其像 \(\varepsilon _X\) 确实是 \(\operatorname {H}^n(X,\Omega _X^n)\) 的一组基。\par{}只剩下证明 \hyperref[fga1-d]{断言~\ref*{fga1-d}}；它将以纯形式的方式从下列引理所概括的若干初等事实推出。这里假定 \(X\) 是 \(r\) 维射影空间 \(\mathbf {P}\) 的闭子集（可以奇异，也可以非奇异）。用记号 \(\mathcal {O}_\mathbf {P}(m)\) 表示 \(\mathbf {P}\) 上在 \hyperref[Ser1955]{[3]} 中记作 \(\mathcal {O}(m)\) 的层，并用 \(\mathcal {O}_X(m)\) 表示 \(X\) 上相应的层。\begin{lemma}{2}\setnumber{2}\label{fga1-lemma-2}\skipsinglepar\par{}若 \(X=\mathbf {P}\) 且 \(\mathscr {F}=\mathcal {O}_\mathbf {P}(m)\)，则 \hyperref[fga1-d]{断言~\ref*{fga1-d}} 成立。\begin{proof}\par{}这个引理可由直接计算证明。\hyperref[Ser1955]{[3]} 中给出了 \(\operatorname {H}^i(\mathbf {P},\mathcal {O}_\mathbf {P}(m))\) 的具体计算，但也可以用更简单的方法进行。计算杯积 \(\operatorname {H}^i(\mathbf {P},\mathcal {O}_\mathbf {P}(m))\times \operatorname {H}^j(\mathbf {P},\mathcal {O}_\mathbf {P}(m')) \to  \operatorname {H}^{i+j}(\mathbf {P},\mathcal {O}_\mathbf {P}(m+m'))\) 并无困难；而计算 \hyperref[fga1-equation-5.1]{式~\ref*{fga1-equation-5.1}} 中的配对需要用到它。\end{proof}\end{lemma}\begin{lemma}{3}\setnumber{3}\label{fga1-lemma-3}\skipsinglepar\par{}\(X\) 上的每个凝聚代数层 \(\mathscr {F}\) 都同构于 \(\mathcal {O}_X(-m)^k\) 的某个商层，并且 \(m\) 可以取得任意大。\begin{proof}\par{}这是因为当 \(m\) 充分大时，\(\mathscr {F}\otimes \mathcal {O}_X(m)\) “由其截面生成”；参见 \hyperref[Ser1955]{[3]}。\end{proof}\end{lemma}\begin{lemma}{4}\setnumber{4}\label{fga1-lemma-4}\skipsinglepar\par{}设 \(i>0\)。当 \(m\) 充分大时，有 \(\operatorname {H}^{r-i}(\mathbf {P},\mathcal {O}_\mathbf {P}(-m))=0\)；而对于 \(X\) 上的每个凝聚代数层 \(\mathscr {B}\)，当 \(m\) 充分大时，有 \(\operatorname {Ext}_{\mathcal {O}_X}^i(X;\mathcal {O}_X(-m),\mathscr {B})=0\)。\begin{proof}\par{}第一个断言来自上面提到的具体计算；对于第二个断言，注意到有同构
\[         \operatorname {Ext}_{\mathcal {O}_X}^i(X;\mathcal {O}_X(-m),\mathscr {B})         = \operatorname {H}^i(X,\mathscr {B}\otimes \mathcal {O}(m))       \]
（由 \hyperref[fga1-proposition-1-corollary-1]{命题 1 的推论 1}），再利用 \hyperref[Ser1955]{[3]} 中的一个熟知结果，即得结论。\end{proof}\end{lemma}\par{}结合前两个引理，得到：\begin{corollary}{}\label{fga1-lemma-3-and-4-corollary}\skipsinglepar\par{}设 \(i>0\)。那么 \(\mathbf {P}\) 上凝聚代数层范畴上的函子 \(\mathscr {F}\mapsto \operatorname {H}^{r-i}(\mathbf {P},\mathscr {F})\) 是可余抹的，\(X\) 上凝聚代数层范畴上的函子 \(\operatorname {Ext}_{\mathcal {O}_X}^i(X;\mathscr {F},\mathscr {B})\) 也是如此。\end{corollary}\begin{lemma}{5}\setnumber{5}\label{fga1-lemma-5}\skipsinglepar\par{}设 \(\mathscr {A}\) 和 \(\mathscr {B}\) 是 \(X\) 上的凝聚代数层，并令 \(\mathscr {A}(m)=\mathscr {A}\otimes \mathcal {O}_X(m)\)。那么当 \(m\) 充分大时，典范同态
\[       \operatorname {Ext}_{\mathcal {O}_X}^i(X;\mathscr {A}(-m),\mathscr {B})       \to  \operatorname {H}^0(X,\mathscr {E}\kern  -1.5pt\mathit {xt}_{\mathcal {O}_X}^i(\mathscr {A}(-m),\mathscr {B}))       = \operatorname {H}^0(X,\mathscr {E}\kern  -1.5pt\mathit {xt}_{\mathcal {O}_X}^i(\mathscr {A},\mathscr {B})(m))     \]
是同构。\begin{proof}\par{}将 \hyperref[fga1-proposition-1]{命题~\ref*{fga1-proposition-1}} 的谱序列应用于 \(\mathscr {A}(-m)\) 和 \(\mathscr {B}\)，立即可得此结论，因为这时有
\[         E_2^{p,q}(\mathscr {A}(-m),\mathscr {B})         = \operatorname {H}^p(X,\mathscr {E}\kern  -1.5pt\mathit {xt}_{\mathcal {O}_X}^q(\mathscr {A}(-m),\mathscr {B}))         = \operatorname {H}^p(X,\mathscr {E}\kern  -1.5pt\mathit {xt}_{\mathcal {O}_X}^q(\mathscr {A},\mathscr {B})(m))       \]
而当 \(p>0\) 且 \(m\) 充分大时，此式为零。\end{proof}\end{lemma}\par{}现在证明 \(X=\mathbf {P}\) 时的 \hyperref[fga1-d]{断言~\ref*{fga1-d}}。先证明当 \(p=n\) 时，\hyperref[fga1-equation-5.2bis]{式~\ref*{fga1-equation-5.2bis}} 是同构。这时两端都是左正合函子（因为 \(\operatorname {H}^{r+i}(\mathbf {P},\mathscr {F})=0\)），因此由 \hyperref[fga1-lemma-3]{引理~\ref*{fga1-lemma-3}}，只须对 \(\mathscr {F}=\mathcal {O}_\mathbf {P}(-m)\) 证明，而这正是 \hyperref[fga1-lemma-2]{引理~\ref*{fga1-lemma-2}} 所述的情形。由于 \hyperref[fga1-equation-5.2bis]{式~\ref*{fga1-equation-5.2bis}} 中的同态具有函子性，且与上边缘映射相容；又由于当 \(p<n\) 时，\hyperref[fga1-equation-5.2bis]{式~\ref*{fga1-equation-5.2bis}} 的两端都是关于 \(\mathscr {F}\) 的可余抹函子（引理 3 和 4 的推论），由标准论证可知，对所有 \(p\)，\hyperref[fga1-equation-5.2bis]{式~\ref*{fga1-equation-5.2bis}} 都是同构。这样便证明了射影空间的对偶定理。\par{}现在设 \(X\) 任意，但仍非奇异。由 \(\mathbf {P}\) 的对偶定理，有同构
\[     \operatorname {H}^n(X,\mathscr {F})'     = \operatorname {H}^n(\mathbf {P},\mathscr {F})'     = \operatorname {Ext}_{\mathcal {O}_\mathbf {P}}^{r-n}(\mathbf {P};\mathscr {F},\Omega _\mathbf {P}^r).   \]
由 \hyperref[fga1-lemma-1]{引理~\ref*{fga1-lemma-1}}，最右端可与
\[     \operatorname {Hom}_{\mathcal {O}_\mathbf {P}}(\mathbf {P};\mathscr {F},\omega _X^n)     = \operatorname {Hom}_{\mathcal {O}_X}(X;\mathscr {F},\Omega _X^n)     = \operatorname {Ext}_{\mathcal {O}_X}^0(X;\mathscr {F},\Omega _X^n)   \]
等同，从而得到同构
\begin{eqn}{5.6}\setnumber{5.6}\label{fga1-equation-5.6}\skipsinglepar\[       \operatorname {H}^n(X,\mathscr {F})'       = \operatorname {Hom}_{\mathcal {O}_X}(X;\mathscr {F},\Omega _X^n)       = \operatorname {Ext}_{\mathcal {O}_X}^0(X;\mathscr {F},\Omega _X^n).     \tag{5.6}     \]\end{eqn}
取 \(\mathscr {F}=\Omega _X^n\)，就得到同构
\begin{eqn}{5.7}\setnumber{5.7}\label{fga1-equation-5.7}\skipsinglepar\[       \eta \colon  \operatorname {H}^n(X,\Omega _X^n) \xrightarrow {\sim } k.     \tag{5.7}     \]\end{eqn}
利用 \hyperref[fga1-equation-5.7]{式~\ref*{fga1-equation-5.7}}，可以证明 \hyperref[fga1-equation-5.6]{式~\ref*{fga1-equation-5.6}} 中的同构恰好是取 \(p=n\)、\(\mathscr {L}=k\) 时的 \hyperref[fga1-equation-5.2bis]{式~\ref*{fga1-equation-5.2bis}}。因此，当 \(p=n\) 时，\hyperref[fga1-equation-5.2bis]{式~\ref*{fga1-equation-5.2bis}} 是同构。要证明它对所有 \(p\) 都是同构，仍只须证明当 \(p<n\) 时，\hyperref[fga1-equation-5.2bis]{式~\ref*{fga1-equation-5.2bis}} 的两端都是关于 \(\mathscr {F}\) 的可余抹函子；而且结合 \hyperref[fga1-lemma-3]{引理~\ref*{fga1-lemma-3}}，只须证明取 \(\mathscr {F}=\mathcal {O}_X(-m)\)、\(m\) 充分大时两端均为零。左端由 \hyperref[fga1-lemma-4]{引理~\ref*{fga1-lemma-4}} 即知；对于右端，利用 \(\mathbf {P}\) 的对偶定理可写成
\[     \operatorname {H}^p(X,\mathcal {O}_X(-m))'     = \operatorname {Ext}_{\mathcal {O}_\mathbf {P}}^{r-p}(\mathbf {P};\mathcal {O}_X(-m),\Omega _\mathbf {P}^r).   \]
当 \(p<n\)、\(m\) 充分大时，右端为零。这来自 \hyperref[fga1-lemma-5]{引理~\ref*{fga1-lemma-5}}（其中实际上 \(X=\mathbf {P}\)），以及如下事实：\(\mathcal {O}_X\) 作为 \(\mathbf {P}\) 上的凝聚代数层，其上同调维数为 \(\leqslant  r-n\)（因为 \(X\) 非奇异），因而有
\[     \mathscr {E}\kern  -1.5pt\mathit {xt}_{\mathcal {O}_\mathbf {P}}^{r-p}(\mathcal {O}_X,\Omega _\mathbf {P}^r)     = 0     \quad \text {当 }p<n.   \]\chapter{奇异簇的对偶定理}\thispagestyle{fancy}\label{fga1-6}\par{}设 \(X\) 是 \(r\) 维射影空间 \(\mathbf {P}\) 中的 \(n\) 维闭子集。此时 \hyperref[fga1-equation-5.6]{式~\ref*{fga1-equation-5.6}} 可写为
\begin{eqn}{6.1}\setnumber{6.1}\label{fga1-equation-6.1}\skipsinglepar\[       \operatorname {H}^n(X,\mathscr {F})'       \simeq  \operatorname {Hom}_{\mathcal {O}_X}(X;\mathscr {F},\omega _X^n)       = \operatorname {Ext}_{\mathcal {O}_X}^0(X;\mathscr {F},\omega _X^n)     \tag{6.1}     \]\end{eqn}
其中令
（\emph{[Trans.] 下列对偶化层的定义式与后面的同构式，在原本均编号为 (6.2)。英文译者认为前者是笔误，并指出对 (6.2) 的引用似乎均指向后者。中译保留原本的重复印号，源码中的引用标签则指向后一个公式。}）
\[     \omega _X^n     = E^0(\mathcal {O}_X)     = \mathscr {E}\kern  -1.5pt\mathit {xt}_{\mathcal {O}_\mathbf {P}}^{r-n}(\mathcal {O}_X,\Omega _\mathbf {P}^r).   \tag{6.2}\]
如 \hyperref[fga1-proposition-5]{命题~\ref*{fga1-proposition-5}} 中所述，这样定义的层不依赖于将 \(X\) 嵌入非奇异簇 \(\mathbf {P}\) 的方式。在 \hyperref[fga1-equation-6.1]{式~\ref*{fga1-equation-6.1}} 中取 \(\mathscr {F}=\omega _X^n\)，得到
\begin{eqn}{6.2}\setnumber{6.2}\label{fga1-equation-6.2}\skipsinglepar\[       \operatorname {H}^n(X,\omega _X^n)'       \simeq  \operatorname {Hom}_{\mathcal {O}_X}(X;\omega _X^n,\omega _X^n)     \tag{6.2}     \]\end{eqn}
从而 \(\operatorname {H}^n(X,\omega _X^n)'\) 中存在一个与 \(\omega _X^n\) 的恒等态射对应的特定元素：
\begin{eqn}{6.3}\setnumber{6.3}\label{fga1-equation-6.3}\skipsinglepar\[       \eta \colon  \operatorname {H}^n(X,\omega _X^n) \to  k.     \tag{6.3}     \]\end{eqn}
再考虑由 \(\operatorname {Ext}\) 的复合定义的配对：
\begin{eqn}{6.4}\setnumber{6.4}\label{fga1-equation-6.4}\skipsinglepar\[       \operatorname {H}^p(X,\mathscr {F}) \times  \operatorname {Ext}_{\mathcal {O}_X}^{n-p}(X;\mathscr {F},\omega _X^n)       \to  \operatorname {H}^n(X,\omega _X^n)     \tag{6.4}     \]\end{eqn}
并将其与 \hyperref[fga1-equation-6.3]{式~\ref*{fga1-equation-6.3}} 中的同态 \(\eta \) 复合；这样便得到与上边缘映射相容的函子性同态
\begin{eqn}{6.5}\setnumber{6.5}\label{fga1-equation-6.5}\skipsinglepar\[       \operatorname {Ext}_{\mathcal {O}_X}^{n-p}(X;\mathscr {F},\omega _X^n)       \to  \operatorname {H}^p(X,\mathscr {F})'     \tag{6.5}     \]\end{eqn}
（它推广了 \hyperref[fga1-equation-5.2]{式~\ref*{fga1-equation-5.2}}）。可以证明，当 \(p=n\) 时，所得正是 \hyperref[fga1-equation-6.1]{式~\ref*{fga1-equation-6.1}} 中的同构。于是有：\begin{theorem}{3 bis}\setnumber{3 bis}\label{fga1-theorem-3bis}\skipsinglepar\par{}对于任意给定的整数 \(k\geqslant 0\)，关于 \(X\) 的下列四个条件等价：
\begin{enumerate}
      
 \item[i.]当 \(n-k\leqslant  p\leqslant  n\) 时，\hyperref[fga1-equation-6.5]{式~\ref*{fga1-equation-6.5}} 中的函子性同态是同构。


      
 \item[ii.]当 \(n-k\leqslant  p<n\) 且 \(m\) 充分大时，有 \(\operatorname {H}^p(X,\mathcal {O}_X(-m)) = 0\)。


      
 \item[iii.]当 \(n-k\leqslant  p<n\) 时，\(X\) 上凝聚代数层范畴上的函子 \(\operatorname {H}^p(X,\mathscr {F})\) 是可余抹的。


      
 \item[iv.]当 \(0<i\leqslant  k\) 时，有 \(E^i(\mathcal {O}_X) = \mathscr {E}\kern  -1.5pt\mathit {xt}_{\mathcal {O}_\mathbf {P}}^{r-n+i}(\mathcal {O}_X,\omega _\mathbf {P}^r) = 0\)。

    \end{enumerate}\begin{proof}\par{}由 \hyperref[fga1-lemma-4]{引理~\ref*{fga1-lemma-4}}，(i)\(\implies \)(ii)；由 \hyperref[fga1-lemma-3]{引理~\ref*{fga1-lemma-3}}，(ii)\(\implies \)(iii)；利用熟知的标准论证，得到 (iii)\(\implies \)(i)，只须注意到这时当 \(n-k\leqslant  p<n\) 时，\hyperref[fga1-equation-6.5]{式~\ref*{fga1-equation-6.5}} 的两端都是可余抹函子（第一个函子的这一性质由 \hyperref[fga1-lemma-4]{引理~\ref*{fga1-lemma-4}} 保证）；最后，(ii)\(\iff \)(iv) 来自 \hyperref[fga1-proposition-6-corollary]{命题 6 的推论}。\end{proof}\end{theorem}\begin{proposition}{6}\setnumber{6}\label{fga1-proposition-6}\skipsinglepar\par{}设 \(\mathscr {F}\) 是 \(X\) 上的凝聚代数层，\(i\) 是整数。那么当 \(m\) 充分大时，存在同构
\begin{eqn}{6.6}\setnumber{6.6}\label{fga1-equation-6.6}\skipsinglepar\[         \operatorname {H}^i(X,\mathscr {F}(-m))' \simeq  \operatorname {H}^0(X,E^{n-i}(\mathscr {F})(m))       \tag{6.6}       \]\end{eqn}
其中令
\begin{eqn}{6.7}\setnumber{6.7}\label{fga1-equation-6.7}\skipsinglepar\[         E^j(\mathscr {F}) = \mathscr {E}\kern  -1.5pt\mathit {xt}_{\mathcal {O}_\mathbf {P}}^{r-n+j}(\mathscr {F},\Omega _\mathbf {P}^r)       \tag{6.7}       \]\end{eqn}
（与 \cref{fga1-3} 中的 \hyperref[fga1-proposition-5]{命题~\ref*{fga1-proposition-5}} 比较）。\begin{proof}\par{}事实上，由 \(\mathbf {P}\) 的对偶定理，\hyperref[fga1-equation-6.6]{式~\ref*{fga1-equation-6.6}} 的左端同构于 \(\operatorname {Ext}_{\mathcal {O}_\mathbf {P}}^{r-i}(\mathbf {P};\mathscr {F}(-m),\Omega _\mathbf {P}^r)\)，所以由 \hyperref[fga1-lemma-5]{引理~\ref*{fga1-lemma-5}} 即得 \hyperref[fga1-equation-6.6]{式~\ref*{fga1-equation-6.6}}。\end{proof}\end{proposition}\begin{corollary}{}\label{fga1-proposition-6-corollary}\skipsinglepar\par{}当 \(m\) 充分大时 \(\operatorname {H}^i(X,\mathscr {F}(-m))\) 为零的充要条件是 \(E^{n-i}(\mathscr {F})\) 为零。\end{corollary}\par{}回顾，\(E^j(\mathscr {F})\) 不依赖于所考虑的射影嵌入。上述推论中的条件是纯局部的，因此若它对 \(\mathscr {F}\) 成立，则对任何局部同构于某个 \(\mathscr {F}^n\) 的层也成立。特别地，若它对 \(\mathcal {O}_X\) 成立，则对每个局部自由凝聚代数层都成立。例如，若 \(X\) 非奇异，则对所有 \(i<n\) 成立；若 \(X\) 没有由单点组成的分支，则对 \(i=0\) 成立；若 \(S\) 正规且其各个分支的维数均为 \(>1\)，则对 \(i=0,1\) 成立（参见 \hyperref[Ser1955]{[3]}）。按定义，要使它对 \(i<k\) 成立，充要条件是局部环 \(\mathcal {O}_x\)（\(x\in  X\)）具有“同调余维 \(\geqslant  k\)”（这一概念的详情见 \hyperref[Ser1956a]{[4]}）。若 \(k=n\)，则由 定理 3 bis 可知，对偶定理对 \(X\) 成立，即对所有 \(p\) 和所有 \(\mathscr {F}\)，\hyperref[fga1-equation-6.5]{式~\ref*{fga1-equation-6.5}} 都是同构。关于局部环 \(\mathcal {O}_x\)，可以给出许多与此等价的条件（Nagata）；例如满足 Cohen–Macaulay 等维性定理的那些局部环便是如此。又如，若 \(X\) 在非奇异环境簇中局部为“完全交”，也属于这种情形。\chapter{Poincaré 对偶}\thispagestyle{fancy}\label{fga1-7}\par{}设 \(X\) 是 \(n\) 维非奇异射影簇。那么 \(\operatorname {H}^\bullet (X)=\operatorname {H}^\bullet (X,\Omega _X^\bullet )\) 是一个有限维双分次反交换代数；以总次数作为分次，使 \(\operatorname {H}^{p,q}(X)=\operatorname {H}^p(X,\Omega _X^q)\) 的次数为 \(p+q\)，因而 \(\operatorname {H}^\bullet (X)\) 的次数集中在 \(0\) 与 \(2n\) 之间。由 \hyperref[fga1-theorem-2]{定理~\ref*{fga1-theorem-2}} 和 \hyperref[fga1-theorem-3-corollary]{定理 3 的推论}，\(\operatorname {H}^\bullet (X)\) 是一个 \(2n\) 维“Poincaré 代数”，也就是说，\(\operatorname {H}^{2n}(X)\) 带有一个到基域 \(k\) 的同构，而乘法 \(\operatorname {H}^m(X)\times \operatorname {H}^{2n-m}(X)\to \operatorname {H}^{2n}(X)=k\) 给出 \(\operatorname {H}^m(X)\) 与 \(\operatorname {H}^{2n-m}(X)\) 之间的对偶。此外，若 \(Y\) 是另一个非奇异射影簇，则凝聚代数层的 Künneth 公式给出
\begin{eqn}{7.1}\setnumber{7.1}\label{fga1-equation-7.1}\skipsinglepar\[       \operatorname {H}^\bullet (X\times  Y) = \operatorname {H}^\bullet (X)\otimes \operatorname {H}^\bullet (Y)     \tag{7.1}     \]\end{eqn}
这是与 Poincaré 代数结构相容的同构。此外，\(\operatorname {H}^\bullet (X)\) 作为交换代数，是关于 \(X\) 的反变函子，因为态射 \(f\colon  Y\to  X\) 显然定义了分次代数同态
\begin{eqn}{7.2}\setnumber{7.2}\label{fga1-equation-7.2}\skipsinglepar\[       f^*\colon  \operatorname {H}^\bullet (X)\to \operatorname {H}^\bullet (Y).     \tag{7.2}     \]\end{eqn}
由于这里处理的是 Poincaré 代数，通过转置便得到向量空间同态
\begin{eqn}{7.3}\setnumber{7.3}\label{fga1-equation-7.3}\skipsinglepar\[       f_*\colon  \operatorname {H}^\bullet (Y)\to \operatorname {H}^\bullet (X).     \tag{7.3}     \]\end{eqn}
我们在 \cref{fga1-4} 中已经看到，\(f^*\) 在非奇异圈所对应的上同调类上的作用，可从几何上解释为取该圈逆像所对应的上同调类。在目前的研究中，还须说明 \hyperref[fga1-equation-7.3]{式~\ref*{fga1-equation-7.3}} 同样对应于圈的“正像”运算。这一点（至少在适当的非奇异性条件下）可由下列特殊情形推出：\begin{theorem}{4}\setnumber{4}\label{fga1-theorem-4}\skipsinglepar\par{}若 \(f\) 是将 \(X^n\) 的非奇异子簇 \(Y^m\) 恒等地映入 \(X^n\) 的映射，并用 \(1_Y\) 表示 \(\operatorname {H}(Y)\) 的单位元，则
\begin{eqn}{7.4}\setnumber{7.4}\label{fga1-equation-7.4}\skipsinglepar\[         f_*(1_Y) = P_X(Y)       \tag{7.4}       \]\end{eqn}
其中右端为 \(Y\) 在 \(X\) 中对应的上同调类。\par{}\hyperref[fga1-theorem-4]{定理~\ref*{fga1-theorem-4}} 中的 \hyperref[fga1-equation-7.4]{式~\ref*{fga1-equation-7.4}} 等价于
\begin{eqn}{7.4 bis}\setnumber{7.4 bis}\label{fga1-equation-7.4bis}\skipsinglepar\[       \langle  \xi ^{m,m}, P_X(Y) \rangle  \varepsilon _Y       = f_*(\xi ^{m,m})       \quad \text {其中 }\xi ^{m,m}\in \operatorname {H}^m(X,\Omega _X^m)     \tag{7.4 bis}     \]\end{eqn}
其中 \(\varepsilon _Y\) 是 \(\operatorname {H}^m(Y,\Omega _Y^m)\) 的基本元素；对于非奇异射影簇，这给出了与 \(Y\) 对应的上同调类的一个新定义。\begin{proof}\par{}由 \hyperref[fga1-theorem-3]{定理~\ref*{fga1-theorem-3}}，将同态
\[         \operatorname {H}^m(X,\Omega _X^m)         \to  \operatorname {H}^m(Y,\Omega _Y^m)         = \operatorname {H}^m(X,\Omega _Y^m)       \]
的转置视为同态
\begin{eqn}{7.5}\setnumber{7.5}\label{fga1-equation-7.5}\skipsinglepar\[           \begin {CD}             @.             \operatorname {Ext}_{\mathcal {O}_X}^{n-m}(X;\Omega _Y^m,\Omega _X^n)             @>\sim >>             \operatorname {Hom}_{\mathcal {O}_X}(X;\Omega _Y^m,\Omega _Y^m)           \\@. @VVV @.           \\\operatorname {H}^{n-m}(X,\Omega _X^{n-m})             @>\sim >>             \operatorname {Ext}_{\mathcal {O}_X}^{n-m}(X;\Omega _X^m,\Omega _X^n)           \end {CD}         \tag{7.5}         \]\end{eqn}
可以验证，\(\operatorname {H}^m(Y,\Omega _Y^m)\) 的对偶中的元素 \(1_Y\) 可与右端对应于 \(\Omega _Y^m\) 的恒等自同态的元素等同；该元素在 \(\operatorname {H}^{n-m}(X,\Omega _X^{n-m})\) 中的像也确实为 \(P_X(Y)\)。\end{proof}\end{theorem}\par{}这些关系（本来也可以在 \cref{fga1-4} 中给出）可对任意非奇异簇陈述，而且确实成立。例如，第二个关系来自下列典范自同态图的交换性：
\par\vspace{1em}\begin{eqn}{7.6}\setnumber{7.6}\label{fga1-equation-7.6}\skipsinglepar\begingroup\footnotesize\[        \begin {CD}         \operatorname {Ext}_{\mathcal {O}_X}^{n-m}(X;\Omega _X^m,\Omega _X^n)         @<<<         \operatorname {Ext}_{\mathcal {O}_X}^{n-m}(X;\Omega _Y^m,\Omega _X^n)         @>\sim >>         \operatorname {Hom}_{\mathcal {O}_X}(X;\Omega _Y^m,\Omega _Y^m)       \\@VVV @. @VVV       \\\operatorname {Ext}_{\mathcal {O}_X}^{n-m}(X;\mathcal {O}_X,\Omega _X^{n-m})         @<<<         \operatorname {Ext}_{\mathcal {O}_X}^{n-m}(X;\mathcal {O}_Y,\Omega _X^{n-m})         @>\sim >>         \operatorname {Hom}_{\mathcal {O}_X}(X;\Omega _X^m,\Omega _Y^m)       \end {CD}     \tag{7.6}     \]\endgroup\end{eqn}
这样便得到与紧致定向流形的 Poincaré 对偶形式体系完全对应的形式体系。特别地，\hyperref[fga1-theorem-4]{定理~\ref*{fga1-theorem-4}} 可以确定 \(X\times  X\) 的对角线所对应的上同调类。由熟知的论证，便可推出例如下列\emph{Lefschetz 公式}：\begin{theorem}{5}\setnumber{5}\label{fga1-theorem-5}\skipsinglepar\par{}设 \(f\) 是非奇异射影簇 \(X\) 的自态射，且 \(f\) 的各个不动点的重数均为 \(1\)。那么这些不动点的个数，模去 \(k\) 的特征后，等于 \(f\) 在各个 \(\operatorname {H}^i(X)\) 上所定义的自同态的迹的交错和。\end{theorem}\par{}对 \(f\) 不得不加的限制，与 \cref{fga1-4} 中评注所提及的困难有关。但应注意，如果 \(f\) “同伦”于一个各不动点重数均为 \(1\) 的自态射，Lefschetz 公式仍然成立。\chapter{对偶定理的推广}\thispagestyle{fancy}\label{fga1-8}\par{}设 \(X\) 是非奇异代数簇，并且 \(X\) 上的每个凝聚代数层 \(\mathscr {F}\) 都同构于某个局部自由凝聚代数层的商层（若 \(X\) 是某个射影空间中的局部闭子集，就属于这种情形）。那么 \(X\) 上的每个凝聚代数层 \(\mathscr {F}\) 都有一个由局部自由层组成的有限分解 \(\mathscr {L}\)；对于任意两个这样的分解，总能找到第三个，以及从第三个到前两个的、与增广映射相容的同态。同样，若 \(\mathscr {L}\) 是 \(\mathscr {F}\) 的一个有限局部自由分解，并给定同态 \(\mathscr {F}'\to \mathscr {F}\)，则存在 \(\mathscr {F}'\) 的有限局部自由分解 \(\mathscr {L}'\)，以及与 \(\mathscr {F}'\to \mathscr {F}\) 相容的同态 \(\mathscr {L}'\to \mathscr {L}\)；若 \(\mathscr {F}'\to \mathscr {F}\) 满射，还可令该同态为满射。于是，对给定整数 \(r,s\geqslant 0\)，可以定义两个上同调多函子，其变元 \(\mathscr {A}_1,\ldots ,\mathscr {A}_r;\mathscr {B}_1,\ldots ,\mathscr {B}_s\) 取自 \(X\) 上的凝聚代数层范畴；一个取值于 \(X\) 上的凝聚代数层范畴，另一个取值于 \(\operatorname {H}^0(X,\mathcal {O}_X)\)-模范畴。
定义公式如下：
\begin{eqn}{8.1}\setnumber{8.1}\label{fga1-equation-8.1}\skipsinglepar\[       \begin {aligned}         &\underline {T}_r^{s\bullet }(\mathscr {A}_1,\ldots ,\mathscr {A}_r;\mathscr {B}_1,\ldots ,\mathscr {B}_s)       \\&= \operatorname {H}^\bullet (\mathscr {H}\kern  -2.5pt\mathit {om}_{\mathcal {O}_X}(\underline {L}(\mathscr {A}_1)\otimes \ldots \otimes \underline {L}(\mathscr {A}_r), \underline {L}(\mathscr {B}_1)\otimes \ldots \otimes \underline {L}(\mathscr {B}_s))),       \\&T_r^{s\bullet }(\mathscr {A}_1,\ldots ,\mathscr {A}_r;\mathscr {B}_1,\ldots ,\mathscr {B}_s)       \\&= \underline {\operatorname {R}}^\bullet \Gamma (\mathscr {H}\kern  -2.5pt\mathit {om}_{\mathcal {O}_X}(\underline {L}(\mathscr {A}_1)\otimes \ldots \otimes \underline {L}(\mathscr {A}_r), \underline {L}(\mathscr {B}_1)\otimes \ldots \otimes \underline {L}(\mathscr {B}_s)))       \end {aligned}     \tag{8.1}     \]\end{eqn}
其中 \(\underline {L}(\mathscr {F})\) 表示凝聚代数层 \(\mathscr {F}\) 的一个有限局部自由分解，\(\underline {\operatorname {R}}^\bullet \Gamma (\mathscr {K})\) 表示空间 \(X\) 关于层复形 \(\mathscr {K}\) 的超上同调。若 \(r\)（分别地，\(s\)）为零，就将各个 \(\underline {L}(\mathscr {A}_i)\)（分别地，\(\underline {L}(\mathscr {B}_j)\)）的张量积换成 \(\mathcal {O}_X\)。特别地，\(\underline {T}_0^0\) 和 \(T_0^0\) 是没有变元的分次函子：\(\underline {T}_0^0\) 集中在 \(0\) 次，在该次等于层 \(\mathcal {O}_X\)；\(T_0^0\) 则等于 \(\operatorname {H}^\bullet (X,\mathcal {O}_X)\)。\hyperref[fga1-equation-8.1]{式~\ref*{fga1-equation-8.1}} 的右端不依赖于所选分解，这对于 \(\underline {T}\) 是显然的（因为问题是局部的）；对于 \(T\)，则来自前面的一般论述，还须利用层复形 \(\mathscr {K}=\mathscr {H}\kern  -2.5pt\mathit {om}_{\mathcal {O}_X}(\underline {L}(\mathscr {A}_1)\otimes \ldots ,\underline {L}(\mathscr {B}_1)\otimes \ldots )\) 的超上同调谱序列：它收敛于 \(X\) 关于 \(\mathscr {K}\) 的超上同调，初始页为 \(\operatorname {H}^p(X,\operatorname {H}^q(\mathscr {K}))\)，即
\begin{eqn}{8.2}\setnumber{8.2}\label{fga1-equation-8.2}\skipsinglepar\[       E_2^{p,q}       = \operatorname {H}^p(X,(\underline {T}_r^s)^{(q)}(\mathscr {A}_1,\ldots ,\mathscr {A}_r;\mathscr {B}_1,\ldots ,\mathscr {B}_s)).     \tag{8.2}     \]\end{eqn}
由此还可看出，这个谱序列本身不依赖于所选分解，其终项就是 \hyperref[fga1-equation-8.1]{式~\ref*{fga1-equation-8.1}} 的左端。注意到每个正合列 \(0\to \mathscr {F}'\to \mathscr {F}\to \mathscr {F}''\to 0\) 都可由一个有限局部自由复形的正合列来分解，就容易定义关于各个变元 \(\mathscr {A}_i,\mathscr {B}_j\) 的上边缘映射。\par{}在函子系 \(\underline {T}_r^{s\bullet }\)（分别地，\(T_r^{s\bullet }\)）上，定义与张量演算类似的运算；其定义可直接由 \hyperref[fga1-equation-8.1]{式~\ref*{fga1-equation-8.1}} 的定义公式给出。于是有一个\emph{复合}运算（它推广了 \cref{fga1-2} 所述的运算）：
\begin{eqn}{8.3}\setnumber{8.3}\label{fga1-equation-8.3}\skipsinglepar\[       \begin {gathered}         T_r^{s\bullet }(\mathscr {A}_1,\ldots ,\mathscr {A}_r;\mathscr {B}_1,\ldots ,\mathscr {B}_s)         \times  T_{r'}^{s'\bullet }(\mathscr {A}'_1,\ldots ,\mathscr {A}'_{r'};\mathscr {B}'_1,\ldots ,\mathscr {B}'_{s'})       \\\to  T_{r+r'}^{(s+s')\bullet }(\mathscr {A}_1,\ldots ,\mathscr {A}_r,\mathscr {A}'_1,\ldots ,\mathscr {A}'_{r'};\mathscr {B}_1,\ldots ,\mathscr {B}_s,\mathscr {B}'_1,\ldots ,\mathscr {B}'_{s'})       \end {gathered}     \tag{8.3}     \]\end{eqn}
它显然满足结合律，并与函子性同态、上边缘同态及谱序列相容。同样，还存在对称运算，其具体描述留给读者。
每当一个变元 \(\mathscr {A}_i\) 等于某个变元 \(\mathscr {B}_j\) 时，还存在一个\emph{缩并}运算：
\begin{eqn}{8.4}\setnumber{8.4}\label{fga1-equation-8.4}\skipsinglepar\[       \begin {gathered}         T_r^{s\bullet }(\mathscr {A}_1,\ldots ,\mathscr {A}_{i-1},\mathscr {C},\mathscr {A}_{i+1},\ldots ,\mathscr {A}_r;\mathscr {B}_1,\ldots ,\mathscr {B}_{j-1},\mathscr {C},\mathscr {B}_{j+1},\ldots ,\mathscr {B}_s)       \\\to  -T_{r-1}^{(s-1)\bullet }(\mathscr {A}_1,\ldots ,\widehat {\mathscr {A}_i},\ldots ,\mathscr {A}_r;\mathscr {B}_1,\ldots ,\widehat {\mathscr {B}_j},\ldots ,\mathscr {B}_s).       \end {gathered}     \tag{8.4}     \]\end{eqn}
此外，若某个变元 \(\mathscr {A}_i\) 是局部自由层，就可删去它，同时将某个 \(\mathscr {A}_j\)（\(j\neq  i\)）换成 \(\mathscr {A}_j\otimes \mathscr {A}_i\)，或将某个 \(\mathscr {B}_k\) 换成 \(\mathscr {B}_k\otimes \mathscr {A}'_i\)（其中 \(\mathscr {A}'_i=\mathscr {H}\kern  -2.5pt\mathit {om}_{\mathcal {O}_X}(\mathscr {A}_i,\mathcal {O}_X)\)）；若某个变元 \(\mathscr {B}_j\) 局部自由，也有类似规则。特别地，总可以删去等于 \(\mathcal {O}_X\) 的变元。若除至多一个变元 \(\mathscr {B}_i\) 外，其余变元都局部自由，则上述规则给出函子性同构
\begin{eqn}{8.5}\setnumber{8.5}\label{fga1-equation-8.5}\skipsinglepar\[       T_r^{s\bullet }(\mathscr {A}_1,\ldots ,\mathscr {A}_r;\mathscr {B}_1,\ldots ,\mathscr {B}_s)       = \operatorname {H}^\bullet (X,\mathscr {A}'_1\otimes \ldots \otimes \mathscr {A}'_{r}\otimes \mathscr {B}_1\otimes \ldots \otimes \mathscr {B}_s)     \tag{8.5}     \]\end{eqn}
（因为可以归结到 \(r=0\)、\(s=1\) 的情形，此时结论显然；也可直接利用初始项为 \hyperref[fga1-equation-8.2]{式~\ref*{fga1-equation-8.2}} 的谱序列）。以上各种运算也可相应地对 \(\underline {T}_r^s\) 定义。这样引入的各种运算之间的关系，与张量演算中相应运算之间的关系相同。\par{}现在设 \(n\) 是 \(X\) 的维数。依次应用张量复合 \hyperref[fga1-equation-8.3]{式~\ref*{fga1-equation-8.3}}，再对重复出现的变元进行缩并 \hyperref[fga1-equation-8.4]{式~\ref*{fga1-equation-8.4}}，便得到配对
\begin{eqn}{8.6}\setnumber{8.6}\label{fga1-equation-8.6}\skipsinglepar\[       \begin {gathered}         (T_r^s)^p(\mathscr {A}_1,\ldots ;\mathscr {B}_1,\ldots )         \times  (T_s^r)^{n-p}(\mathscr {B}_1,\ldots ;\mathscr {A}_1,\ldots ,\mathscr {A}_r\otimes \Omega _X^n)       \\\to \operatorname {H}^n(X,\Omega _X^n).       \end {gathered}     \tag{8.6}     \]\end{eqn}\begin{theorem}{6}\setnumber{6}\label{fga1-theorem-6}\skipsinglepar\par{}若 \(X\) 是非奇异射影簇，则 \hyperref[fga1-equation-8.6]{式~\ref*{fga1-equation-8.6}} 中的配对都是对偶。\begin{proof}\par{}这以纯形式的方式来自 \hyperref[fga1-theorem-3-corollary]{定理 3 的推论}。事实上，容易由该推论得知，若 \(\mathscr {K}\) 是由\emph{局部自由}凝聚代数层组成的复形，则 \(X\) 关于 \(\mathscr {K}\) 的超上同调与 \(X\) 关于 \(\mathscr {K}'\otimes \Omega _X^n\) 的超上同调，通过自然配对
\begin{eqn}{8.7}\setnumber{8.7}\label{fga1-equation-8.7}\skipsinglepar\[           \underline {\operatorname {R}}^p\Gamma (\mathscr {K})           \times  \underline {\operatorname {R}}^{n-p}\Gamma (\mathscr {K}'\otimes \Omega _X^n)           \to  \underline {\operatorname {R}}^n\Gamma (\Omega _X^n)           = \operatorname {H}^n(X,\Omega _X^n).         \tag{8.7}         \]\end{eqn}
互为对偶。利用初始页为 \(\operatorname {H}^p(\operatorname {H}^q(X,\mathscr {K}))\) 的谱序列，以及对 \(\mathscr {K}'\otimes \Omega _X^n\) 的类似谱序列，即可看出这一点。由上述结果和定义 \hyperref[fga1-equation-8.1]{式~\ref*{fga1-equation-8.1}}，便可推出 \hyperref[fga1-theorem-6]{定理~\ref*{fga1-theorem-6}}。\end{proof}\end{theorem}\begin{remarks}{}\label{fga1-8-remarks}\skipsinglepar\begin{enumerate}\item{}对于 \hyperref[fga1-theorem-6]{定理~\ref*{fga1-theorem-6}} 之前的定义，并不需要 \(X\) 非奇异，因为没有必要只考虑\emph{有限}分解。但是，若 \(X\) 奇异，就不能再先验地保证各个 \((\underline {T}_r^s)^p(\mathscr {A}_1,\ldots ;\mathscr {B}_1,\ldots )\) 是\emph{凝聚}层，因为在层复形
\[         \mathscr {H}\kern  -2.5pt\mathit {om}_{\mathcal {O}_X}(\underline {L}(\mathscr {A}_1)\otimes \ldots ,\underline {L}(\mathscr {B}_1)\otimes \ldots )       \]
中，任意给定总次数的分量都可能有无限多个。

    \item{}容易验证，在 \hyperref[fga1-equation-8.1]{式~\ref*{fga1-equation-8.1}} 的公式中，可将\emph{一个} \(\underline {L}(\mathscr {B}_i)\) 换成 \(\mathscr {B}_i\)。结合 \hyperref[fga1-proposition-3]{命题~\ref*{fga1-proposition-3}}，便得到
\begin{eqn}{8.8}\setnumber{8.8}\label{fga1-equation-8.8}\skipsinglepar\[           \begin {aligned}             \underline {T}_1^{1\bullet }(\mathscr {A};\mathscr {B})             &= \mathscr {E}\kern  -1.5pt\mathit {xt}_{\mathcal {O}_X}^\bullet (\mathscr {A},\mathscr {B})           \\T_1^{1\bullet }(\mathscr {A};\mathscr {B})             &= \operatorname {Ext}_{\mathcal {O}_X}^\bullet (X;\mathscr {A},\mathscr {B}).           \end {aligned}         \tag{8.8}         \]\end{eqn}
特别地，在 \hyperref[fga1-equation-8.6]{式~\ref*{fga1-equation-8.6}} 中取 \(r=s=1\)、\(\mathscr {A}_1=\mathcal {O}_X\)，便重新得到 \hyperref[fga1-theorem-3]{定理~\ref*{fga1-theorem-3}}。\hyperref[fga1-equation-8.8]{式~\ref*{fga1-equation-8.8}} 还蕴含 \(T_0^{1\bullet }(\mathscr {B})=\operatorname {H}^\bullet (X,\mathscr {B})\) 和 \(T_1^{0\bullet }(\mathscr {A})=\operatorname {Ext}_{\mathcal {O}_X}^\bullet (X;\mathscr {A},\mathcal {O}_X)\)。

    \item{}由 \hyperref[fga1-equation-8.1]{式~\ref*{fga1-equation-8.1}} 可见，函子 \(\underline {T}_r^{s\bullet }\) 和 \(T_r^{s\bullet }\) 一般在正次数\emph{以及}负次数中都有分量。利用上一个评注可知，若 \(X\) 的维数为 \(n\)，则当 \(s>0\) 时，\(\underline {T}_r^{s\bullet }\) 的非零分量集中在 \(-(s-1)n\) 与 \(rn\) 次之间；当 \(s=0\) 时，则集中在 \(0\) 与 \(rn\) 次之间。当 \(s>0\) 时，\(T_r^{s\bullet }\) 的非零分量集中在 \(-(s-1)n\) 与 \((r+1)n\) 次之间；当 \(s=0\) 时，则集中在 \(0\) 与 \((r+1)n\) 次之间（而且，如果我没有弄错，当 \(r>0\) 时，分别集中在 \(-(s-1)n\) 与 \(rn\) 次之间，以及 \(0\) 与 \(rn\) 次之间）。\end{enumerate}\end{remarks}
\begingroup\emergencystretch=2em\section*{原报告参考文献}
\begin{enumerate}
\item[1.]\setnumber{1}\label{fga1-bib-1} CARTIER (Pierre). — Les groupes $\operatorname{Ext}^s(A,B)$，Séminaire A. Grothendieck : Algèbre homologique，t. 1，1957，n\textsuperscript{o} 3.
\item[2.]\setnumber{2}\label{fga1-bib-2} GROTHENDIECK (Alexandre). — Sur quelques points d’algèbre homologique，Tohoku math. J.，t. 9，1957，p. 119–183.
\item[3.]\setnumber{3}\label{fga1-bib-3} SERRE (Jean-Pierre). — Faisceaux algébriques cohérents，Annals of Math.，t. 61，1955，p. 197–278.
\item[4.]\setnumber{4}\label{fga1-bib-4} SERRE (Jean-Pierre). — Sur la dimension homologique des anneaux et des modules noethériens，Proc. Intern. Symp. on alg. number Theory [1955, Tokyo et Nikko]，Tokyo，Science Council of Japan，1956，p. 175–189.
\end{enumerate}\endgroup
% 原文：FGA 法文汇编 PDF 物理第 39 页，报告页码 149-25。
% 1959 年 4 月原报告补记；与英文读者版 §4 的后出 [Comp.] 说明分列。
% 中译工作稿；原补记完整恢复。
\section*{补记}

第 13 页评注中指出的困难，现在已经完全解决；其方法是将对偶定理推广到任意簇（允许任意奇点）。这一理论可以在“概形理论”的一般框架中表述。关于这一理论，我们请读者参阅 J. DIEUDONNE 与 A. GROTHENDIECK 正在编写的《Éléments de géométrie algébrique》。若干说明可见：

\begin{quote}
GROTHENDIECK (A.). — \emph{Cohomology theory of algebraic varieties}，1958 年 Edinburgh 国际数学家大会（待刊）。
\end{quote}

\begin{flushright}
[1959 年 4 月]
\end{flushright}

% END chapters/fga1.tex

\clearpage
\setcounter{chapter}{0}
\renewcommand\fgaAnchorPrefix{fga2}

% BEGIN chapters/fga2.tex
\part{形式几何与代数几何}\label{fga2}\thispagestyle{fancy}
\begin{original}{}\label{fga2-original-citation}\skipsinglepar
\par{}A. Grothendieck。"Géométrie formelle et géométrie algébrique"。\emph{Séminaire Bourbaki} \textbf{11}（1958–59），第 182 次报告，1959 年 5 月。\url{http://www.numdam.org/book-part/SB_1958-1960__5__193_0/}\end{original}
\par{}\cref{fga2-1}至\cref{fga2-5}的主要内容已收入 [\hyperref[GD1960]{EGA}, III] 的已出版部分；\cref{fga2-6}与\cref{fga2-7}的内容见 [\hyperref[Gro1960b]{SGA}, III]。关于基本群的研究，见 [\hyperref[Gro1960b]{SGA}, V, IX, X, XI]；关于 Lefschetz 型定理以及许多尚未解决的问题，另见 [\emph{SGA 1962}, X, XII, XIII]。只有温和分歧覆盖理论（参见\hyperref[fga2-theorem-14]{定理~\ref*{fga2-theorem-14}}）尚未有正式成文的论述。\hyperref[fga2-theorem-14-corollary-1]{定理 14 的推论}完全确定了代数闭域上代数曲线的、次数与特征互素的 Galois 覆盖，已有三项工作实质性地使用了这一推论：
\begin{enumerate}\item{}Igusa 对任意特征下非奇异射影曲面的 Picard 不等式的证明；

    \item{}Ogg 与 Šafarevič 各自独立进行的研究：在任意特征下，对一变元函数域上定义的 Abel 簇，研究其主齐性丛所成的群；以及

    \item{}Artin 最近对代数簇的“Weil 上同调”的若干关键定理所作的证明。\end{enumerate}
\chapter{概形}\thispagestyle{fancy}\label{fga2-1}
\par{}我们知道，域 \(k\) 上定义的仿射代数空间本质上由它的仿射代数 \(A\)（定义在 \(k\) 上的正则函数所成的环）决定，而代数空间的态射 \(X\to  Y\) 与 \(k\)-代数同态 \(A(Y)\to  A(X)\) 一一对应。与代数空间对应的仿射代数是有限型 \(k\)-代数，并且在“经典”观点下不含幂零元；反过来，每个这样的代数都可实现为 \(k\) 上某个代数空间的仿射代数。因此，已有一套熟悉的对应，可以把仿射代数空间的问题解释为交换代数的问题。人们早已注意到，这样所得的陈述更为一般，因为通常不必要求所考虑的环具有刚才所述的形式，大多数情况下只需 Noether 条件。尤其是，无论是否给定了基域，都不必排除环中含有幂零元的情形。可是，直到现在，几何学家仍不愿接受这一信息，坚持只考虑不含幂零元的仿射代数，也就是结构层中不含幂零元的代数空间，甚至大多数时候只考虑“绝对不可约”的代数空间。报告人认为，这种态度严重阻碍了代数几何中真正自然的方法的发展。
\par{}设 \(A\) 为交换环。众所周知，\(A\) 的素理想集合 \(X=\operatorname {Spec}(A)\) 带有一个自然的拓扑，即“\emph{Zariski 拓扑}”，也称谱拓扑。此外，\(X\) 上有交换环层 \(\mathcal {O}_X\)，其在 \(\mathfrak {p}\in  X\) 处的茎为局部化环 \(A_\mathfrak {p}\)，其截面环可与 \(A\) 等同。这样，\(X\) 就成为\emph{赋环空间}，称为 \(A\) 的\emph{素谱}。环同态 \(f\colon  A\to  B\) 定义一个赋环空间的态射 \(f'\colon \operatorname {Spec}(B)\to \operatorname {Spec}(A)\)，其底层集合映射恰为 \(\mathfrak {p}\mapsto  f^{-1}(\mathfrak {p})\)。这样得到的赋环空间态射 \(\operatorname {Spec}(B)\to \operatorname {Spec}(A)\)，恰好是那些使同态 \(\mathcal {O}_x\to \mathcal {O}_y\)（其中 \(x=f'(y)\)）为局部同态的态射，也就是说，极大理想的原像是极大理想。
\par{}我们把同构于某个 \(\operatorname {Spec}(A)\) 的赋环空间称为\emph{仿射概形}；把局部仿射的赋环空间称为\emph{预概形}，也就是说，每一点都有一个开邻域，在诱导的结构下是仿射概形。预概形的\emph{态射}可按显然的方式定义；在局部，它们对应于环同态。
\par{}固定一个预概形 \(S\)，并考虑预概形态射 \(X\to  S\) 时，\(S\) 起着基域或基环的作用，或者更恰当地说，起着纤维化中的底空间的作用。此时称 \(X\) 为\emph{\(S\)-预概形}；若 \(S=\operatorname {Spec}(A)\)，则还意味着 \(\mathcal {O}_X\) 是一个\emph{\(A\)-代数}层。因此，每个预概形都可唯一地视为 \(\mathbb {Z}\)-预概形。当然，\(S\)-预概形组成一个范畴；还可证明，在这个范畴中，任意两个对象 \(X\) 与 \(Y\) 的积总是存在，记为 \(X\times _S Y\)。这一积的概念使我们能够对态射 \(S'\to  S\) 定义 \(S\)-预概形的\emph{基变换}，因为 \(X\times _S S'\) 可视为 \(S'\)-预概形。
\par{}若 \(X\times _S X\) 的对角线是闭的，就称 \(X\) 在 \(S\) 上是\emph{分离}的。我们把在 \(\mathbb {Z}\) 上分离的预概形称为\emph{概形}；此时，它在任意底上都分离。为简便起见，以下只谈概形，并进一步假定它们是\emph{Noether}的，也就是说，可写成有限个仿射开集的并，而这些仿射开集是 Noether 环的谱。如果对 \(S\) 的每个仿射开子集 \(U\)，它在 \(X\) 中的原像都是有限个仿射开集的并，且各自的环是 \(U\) 的环上的有限型代数，就称 \(X\) 在 \(S\) 上\emph{为有限型}。正是这样的 \(S\)-概形适于作真正的几何研究。特别地，对每个 \(s\in  S\)，\(X\) 在 \(s\) 上的纤维 \(f^{-1}(s)\) 是局部环 \(\mathcal {O}_s\)（即 \(S\) 中点 \(s\) 的局部环）的剩余域 \(k(s)\) 上的代数概形。因此，在一定程度上，可把 \(X\) 看成一族“代数空间”\(f^{-1}(s)\)，其参数 \(s\) 遍历 \(S\)，也就是在局部观点下遍历某个给定环的素理想集合。当然，各个 \(k(s)\) 可以具有不同的特征。若 \(S=\operatorname {Spec}(k)\)，其中 \(k\) 为域，就基本恢复了通常的“代数空间”概念，唯一的区别是现在允许结构层含有幂零元。
\par{}受熟知的思想启发，我们可以定义\emph{射影态射}，并进一步定义\emph{固有态射}。这样的态射是有限型的，把闭子集映为闭子集，并且在任意基变换后仍保持这一性质。
\par{}设 \(X\) 为概形（仍一律假定为 Noether 概形），则层 \(\mathcal {O}_X\) 是 \hyperref[Gro1960]{[2]} 意义下的\emph{凝聚环层}。因此，\(X\) 上的凝聚模层也恰好是那些局部同构于某个态射 \(\mathcal {O}_X^m\to \mathcal {O}_X^n\) 的余核的层。
\chapter{形式概形}\thispagestyle{fancy}\label{fga2-2}
\par{}设 \(X\) 为概形，\(X'\) 为 \(X\) 的闭子集。则存在 \(\mathcal {O}_X\) 的凝聚子层 \(\mathscr {J}\)，使得 \(X'=\operatorname {supp}\mathcal {O}_X/\mathscr {J}\)，甚至存在最大的这种子层。给 \(X'\) 配备层 \(\mathcal {O}_X/\mathscr {J}\)，便使 \(X'\) 成为一个概形，记为 \(X_0\)；这样的概形称为\emph{\(X\) 的闭子概形}。此外，对任意 \(n\)，可考虑配备了 \(\mathcal {O}_X/\mathscr {J}^{n+1}\) 的 \(X'\)，记为 \(X_n\)。这是 \(X\) 的闭子预概形，其底层集合仍为 \(X'\)，但结构层不同，为 \(\mathcal {O}_{X_n}=\mathcal {O}_X/\mathscr {J}^{n+1}\)。显然，各个 \(\mathcal {O}_{X_n}\) 构成 \(X\) 上环层的射影系，其射影极限 \(\overline {\mathcal {O}_X}\) 称为\emph{\(\mathcal {O}_X\) 沿 \(X'\) 的形式完备化}。赋予这一环层后，\(X'\) 称为\emph{\(X\) 沿 \(X'\) 的形式完备化}；它因而是赋环空间，但通常不是概形。对 \(X\) 上任意凝聚层 \(\mathscr {F}\)，同样可考虑 \(\mathscr {F}\) 沿 \(X'\) 的形式完备化 \(\overline {\mathscr {F}}=\varprojlim _n\mathscr {F}_n\)（其中 \(\mathscr {F}_n=\mathscr {F}\otimes _{\mathcal {O}_X}\mathcal {O}_X/\mathscr {J}^{n+1}\)）；这是 \(\overline {X}\) 上的模层。它的截面称为\emph{\(\mathscr {F}\) 沿 \(X'\) 的形式截面}，可与 \(\varprojlim _n\Gamma (X',\mathscr {F}_n)\) 的元素等同。取 \(\mathscr {F}=\mathcal {O}_X\) 时，便得到 Zariski 所谓 \(X\) 沿 \(X'\) 的“全纯函数”；由于这一用语与经典术语相混，我们不沿用它。
\par{}所谓\emph{形式概形}（默认满足 Noether 条件），是指一个拓扑空间 \(\mathfrak {X}\)，配备拓扑环层 \(\mathcal {O}_{\mathfrak {X}}\)，并满足下述条件：有拓扑环层的同构 \(\mathcal {O}_{\mathfrak {X}}=\varprojlim _n\mathcal {O}_n\)，其中各个 \(\mathcal {O}_n\) 构成 \(\mathfrak {X}\) 上环层的射影系，每一个都使 \(\mathfrak {X}\) 成为概形 \(\mathfrak {X}_n\)；对 \(m\geqslant  n\)，同态 \(\mathcal {O}_m\to \mathcal {O}_n\) 是满射，核为 \(\mathscr {J}_m^{n+1}\)，这里 \(\mathscr {J}_m\) 是 \(\mathcal {O}_m\to \mathcal {O}_0\) 的核。我们将证明，\(\mathcal {O}_{\mathfrak {X}}\) 是\emph{Noether 局部环}的\emph{凝聚}层。
\par{}依定义，上述形式完备化 \(\overline {X}\) 是形式概形；反之，每个形式概形都\emph{局部地}具有这种形式。事实上，给定一个\emph{仿射}形式概形（即 \(\mathfrak {X}_0\) 为仿射概形，从而各个 \(\mathfrak {X}_n\) 都仿射），等价于给定一个分离、完备的 \(\mathscr {J}\)-进 Noether 拓扑环。
\par{}普通概形的通常定义，如态射、有限型态射、固有态射等，都可以毫无困难地推广到形式概形。
\chapter{三个基本定理}\thispagestyle{fancy}\label{fga2-3}
\par{}设 \(f\colon  X\to  Y\) 为概形之间的固有态射（仍一律假定为 Noether 概形），\(Y'\) 为 \(Y\) 的闭子集，\(X'\) 为它在 \(X\) 中的原像。考虑相应的形式完备化 \(\overline {Y}\) 和 \(\overline {X}\)。则 \(f\) 诱导形式概形的态射 \(\overline {f}\colon \overline {X}\to \overline {Y}\)，它也是固有的。设 \(\mathscr {F}\) 为 \(X\) 上的凝聚层，则 \(\overline {\mathscr {F}}\) 为 \(\overline {X}\) 上的凝聚层。在\hyperref[fga2-theorem-1]{定理~\ref*{fga2-theorem-1}}中，我们忘掉 \(X\)、\(Y\) 和 \(\mathscr {F}\)，只考虑形式概形之间的固有态射 \(\overline {f}\)，以及 \(\overline {X}\) 上的凝聚层 \(\overline {\mathscr {F}}\)。（不过，报告人只在从某些 \(X\)、\(Y\)、\(f\) 和 \(F\) 出发的情形下写出了完整证明。）
\begin{theorem}[{有限性定理}]{1}\setnumber{1}\label{fga2-theorem-1}\skipsinglepar
\begin{enumerate}
    
 \item[i.]各个 \(\operatorname {R}^q\overline {f}_*(\overline {\mathscr {F}})\) 都是 \(\overline {Y}\) 上的凝聚层。


    
 \item[ii.]自然同态 \[         \operatorname {R}^q\overline {f}_*(\overline {\mathscr {F}})         \to  \varprojlim _n\operatorname {R}^q (f_n)_*(\mathscr {F}_n)       \] 都是同构。

  \end{enumerate}\end{theorem}
\par{}在\hyperref[fga2-theorem-1]{定理~\ref*{fga2-theorem-1}}的陈述中，假定已经给定 \(\mathcal {O}_Y\) 的凝聚子层 \(\mathscr {J}\)，它定义 \(Y'\)；再取原像，便得到定义 \(X'\) 的 \(\mathcal {O}_X\) 的凝聚子层。由此，按定义得到 \(\mathscr {F}_n\)、\(X_n\)、\(Y_n\) 和 \(f_n\colon  X_n\to  Y_n\)，如\cref{fga2-2}所述。如果从两个形式概形之间任意的固有态射出发，上述说明中需要对记号作哪些小改动是显然的。
\par{}\hyperref[fga2-theorem-1]{定理~\ref*{fga2-theorem-1}}仅涉及“形式上同调”；下面的定理把它与“代数上同调”联系起来，类似于 Serre 关于代数上同调与解析上同调比较的著名定理 \hyperref[Ser1956]{[4]}。
\begin{theorem}[{第一比较定理}]{2}\setnumber{2}\label{fga2-theorem-2}\skipsinglepar
\par{}各个 \(\operatorname {R}^q f_*(\mathscr {F})\) 都是 \(Y\) 上的凝聚层（这是\hyperref[fga2-theorem-1]{定理~\ref*{fga2-theorem-1}}的特例），且自然同态 \[       \overline {\operatorname {R}^q f_*(\mathscr {F})}       \to  \varprojlim _n \operatorname {R}^q (f_n)_*(\mathscr {F}_n)     \] 都是同构。\end{theorem}
\begin{corollary}{1}\setnumber{1}\label{fga2-theorem-2-corollary-1}\skipsinglepar
\par{}有典范同构 \(\overline {\operatorname {R}^q f_*(\mathscr {F})} = \operatorname {R}^q\overline {f}_*(\overline {\mathscr {F}})\)。\end{corollary}
\par{}当 \(q=0\) 时，这一推论推广了 Zariski 的“全纯函数基本定理”，我们将由此导出 Zariski 的“连通性定理”的推广。还应注意，\hyperref[fga2-theorem-1]{定理~\ref*{fga2-theorem-1}}中的 (ii) 在 \(q=0\) 时显然成立，但\hyperref[fga2-theorem-2]{定理~\ref*{fga2-theorem-2}}及其等价表述（\hyperref[fga2-theorem-2-corollary-1]{推论~\ref*{fga2-theorem-2-corollary-1}}）在此情形下却绝非显然。事实上，证明对 \(q\) 作下降归纳：当 \(q\) 充分大时，等式两边均为零，结论显然；因此，\(q=0\) 的情形恰是最后一步归纳，也就可以称为“最困难”的情形。
\begin{corollary}{2}\setnumber{2}\label{fga2-theorem-2-corollary-2}\skipsinglepar
\par{}设 \(Y=\operatorname {Spec}(A)\)，且 \(Y'\) 由 \(A\) 的理想 \(\mathscr {J}\) 定义。则对 \(X\) 上任意凝聚层 \(\mathscr {F}\)，各个 \(\operatorname {H}^q(X,\mathscr {F})\) 都是有限型 \(A\)-模，其 \(\mathscr {J}\)-进完备化为 \(\operatorname {H}^q(\overline {X},\overline {\mathscr {F}})\)。\end{corollary}
\par{}最后，把这一推论应用于 \(H=\mathscr {H}\kern  -2.5pt\mathit {om}_{\mathcal {O}_X}(\mathscr {F},\mathscr {G})\)，得到：
\begin{corollary}{3}\setnumber{3}\label{fga2-theorem-2-corollary-3}\skipsinglepar
\par{}设 \(Y=\operatorname {Spec}(A)\)，且 \(Y'\) 由 \(A\) 的理想 \(\mathscr {J}\) 定义。设 \(\mathscr {F}\) 和 \(\mathscr {G}\) 为 \(X\) 上的凝聚层。则 \(\operatorname {Hom}(\mathscr {F},\mathscr {G})\) 是有限型 \(A\)-模，其 \(\mathscr {J}\)-进完备化可与 \(\operatorname {Hom}(\overline {\mathscr {F}},\overline {\mathscr {G}})\) 等同。\end{corollary}
\par{}当然，自然映射 \(\operatorname {Hom}(\mathscr {F},\mathscr {G})\to \operatorname {Hom}(\overline {\mathscr {F}},\overline {\mathscr {G}})\) 把同态 \(u\colon \mathscr {F}\to \mathscr {G}\) 送到其“按连续性”的延拓 \(\overline {u}\colon \overline {\mathscr {F}}\to \overline {\mathscr {G}}\)；因此，\(\overline {\mathscr {F}}\) 成为关于 \(\mathscr {F}\) 的函子。
\par{}现在假定 \(A\) 关于其 \(\mathscr {J}\)-进拓扑分离且完备。则上述\hyperref[fga2-theorem-2-corollary-2]{推论~\ref*{fga2-theorem-2-corollary-2}}和\hyperref[fga2-theorem-2-corollary-3]{推论~\ref*{fga2-theorem-2-corollary-3}}给出：\[     \begin {aligned}       \operatorname {H}^q(X,\mathscr {F}) &= \operatorname {H}^q(\overline {X},\overline {\mathscr {F}}),     \\\operatorname {Hom}(\mathscr {F},\mathscr {G}) &= \operatorname {Hom}(\overline {\mathscr {F}},\overline {\mathscr {G}}).     \end {aligned}   \] 后一个恒等式说明，\(X\) 上凝聚层的范畴可与 \(\overline {X}\) 上凝聚层范畴的一个\emph{子范畴}等同，其中态射取为诱导的态射。事实上，甚至有：
\begin{theorem}{3}\setnumber{3}\label{fga2-theorem-3}\skipsinglepar
\par{}\(\overline {X}\) 上的模层凝聚，当且仅当它同构于形如 \(\overline {\mathscr {F}}\) 的层，其中 \(\mathscr {F}\) 为 \(X\) 上的凝聚层；由上述\hyperref[fga2-theorem-2-corollary-3]{推论~\ref*{fga2-theorem-2-corollary-3}}，这个凝聚层在典范同构的意义下唯一确定。（回顾一下，此时 \(Y=\operatorname {Spec}(A)\)，其中 \(A\) 为分离、完备的 \(\mathscr {J}\)-进 Noether 拓扑环。）\end{theorem}
\begin{corollary}{1}\setnumber{1}\label{fga2-theorem-3-corollary-1}\skipsinglepar
\par{}\(X\) 的闭子概形与 \(\overline {X}\) 的闭形式子概形一一对应。\end{corollary}
\par{}事实上，两者分别对应于 \(\mathcal {O}_X\) 与 \(\mathcal {O}_{\overline {X}}\) 的凝聚子层。把态射的图像视为闭子概形，则上述\hyperref[fga2-theorem-3-corollary-1]{推论~\ref*{fga2-theorem-3-corollary-1}}给出：
\begin{corollary}{2}\setnumber{2}\label{fga2-theorem-3-corollary-2}\skipsinglepar
\par{}设 \(X\) 和 \(Z\) 为 \(A\) 上的固有概形，其中后者为分离、完备的 \(\mathscr {J}\)-进 Noether 环。则映射 \(g\mapsto \overline {g}\) 在从 \(X\) 到 \(Z\) 的 \(Y\)-态射与从 \(\overline {X}\) 到 \(\overline {Z}\) 的 \(\overline {Y}\)-态射之间给出一一对应。\end{corollary}
\par{}换言之，\(A\) 上的固有代数概形在 \(\overline {Y}\) 上固有形式概形的范畴中给出一个子范畴，其中态射取诱导的态射。然而应注意，\emph{存在不能“代数化”的 \(\overline {Y}\) 上固有形式概形}，也就是说，它们并不同构于某个 \(\overline {X}\)，其中 \(X\) 在 \(A\) 上固有；正如存在并不来自复数域上代数簇的紧复解析簇一样。这样的形式概形自然出现于“模空间理论”中。不过，下面有一个值得注意的情形，其形式概形可以代数化：
\begin{theorem}{4}\setnumber{4}\label{fga2-theorem-4}\skipsinglepar
\par{}设 \(A\) 为完备 Noether 局部环，剩余域为 \(k\)；设 \(\mathfrak {X}\) 为 \(A\) 上的固有形式概形，其中底环带有其 \(\mathfrak {r}(A)\)-进拓扑。假定 
\begin{enumerate}
      
 \item[i.]\(\mathcal {O}_{\mathfrak {X}}\) 的各局部环都是\emph{平坦} \(A\)-模；等价地，用 \(A\) 的极大理想的幂给 \(\mathcal {O}_{\mathfrak {X}}\) 和 \(A\) 赋予滤过，则相伴分次代数满足 \[           \operatorname {gr}(\mathcal {O}_{\mathfrak {X}})           \simeq  \operatorname {gr}^0(\mathcal {O}_{\mathfrak {X}})\otimes _k\operatorname {gr}(A);         \]


      
 \item[ii.]\(\operatorname {H}^2(\mathfrak {X}_0,\mathcal {O}_{\mathfrak {X}_0})=0\)，这里将 \(\mathfrak {X}_0=\mathfrak {X}\otimes _Ak\) 视为 \(k\) 上的代数概形；


      
 \item[iii.]\(\mathfrak {X}_0\) 是射影的。

    \end{enumerate} 则在这些条件下，\(\mathfrak {X}\) 可以代数化；更确切地，它同构于 \(\overline {X}\)，其中 \(X\) 为某个射影 \(A\)-概形。\end{theorem}
\par{}特别地，若 \(\mathfrak {X}_0\) 为 \(k\) 上的\emph{光滑曲线}，则\hyperref[fga2-theorem-4]{定理~\ref*{fga2-theorem-4}}中的条件 (ii)、(iii) 成立；因此，\hyperref[fga2-theorem-4]{定理~\ref*{fga2-theorem-4}} 尤其可应用于给定亏格曲线的“模空间理论”……这里略述\hyperref[fga2-theorem-4]{定理~\ref*{fga2-theorem-4}}的证明方法：可以证明（参见下面的\hyperref[fga2-proposition-3]{命题~\ref*{fga2-proposition-3}}），(i)、(ii) 蕴含，\(\mathfrak {X}_0\) 上每个局部同构于结构层的凝聚层，都可由 \(\mathfrak {X}\) 上同类的层约化得到。因此，从 \(\mathfrak {X}_0\) 上一个“丰沛”层出发（由 (iii)，这样的层存在），把它提升为 \(\mathfrak {X}\) 上的可逆层；再利用\hyperref[fga2-theorem-1]{定理~\ref*{fga2-theorem-1}}，证明这一可逆层的某个倍数定义了 \(X\) 到概形 \(\mathbb {P}_A^r\) 的形式完备化中的一个浸入，后者是 \(A\) 上维数为 \(r\) 的“标准射影空间”。
\par{}\hyperref[fga2-theorem-1]{定理~\ref*{fga2-theorem-1}}、\hyperref[fga2-theorem-2]{定理~\ref*{fga2-theorem-2}}和\hyperref[fga2-theorem-3]{定理~\ref*{fga2-theorem-3}}的证明，参见 [\hyperref[GD1960]{1}]。
\chapter{对 Zariski 连通性定理与“主定理”的应用}\thispagestyle{fancy}\label{fga2-4}
\par{}设 \(f\colon  X\to  Y\) 为概形间的固有态射。由有限性定理（\hyperref[fga2-theorem-1]{定理~\ref*{fga2-theorem-1}}），\(f_*(\mathcal {O}_X)=\underline {A}\) 是 \(Y\) 上的凝聚层，且为交换代数层，因此对应于一个 \(Y\)-概形 \(g\colon  Y'\to  Y\)，它在 \(Y\) 上有限；其定义条件是在 \(Y\) 上仿射，即仿射开集的原像仿射，并且 \(g_*(\mathcal {O}_{Y'})=\underline {A}\)。显然，此时 \(f\) 可典范地分解为 \(f=gf'\)，其中 \(f'\colon  X\to  Y'\) 为从 \(X\) 到 \(Y\) 的态射，现在满足 \(f'_*(\mathcal {O}_X)=\mathcal {O}_{Y'}\)。\(f\) 的这一分解称为 \(f\) 的\emph{Stein 分解}。把第一比较定理（\hyperref[fga2-theorem-2]{定理~\ref*{fga2-theorem-2}}）及其\hyperref[fga2-theorem-2-corollary-1]{推论 1}应用于 \(f'\) 与只含一个点 \(y'\) 的子集 \(Y'\)，便得到 \((f')^{-1}(y')=X'\) 连通：否则，\(X\) 沿 \(X'\) 的形式截面不会构成局部环，然而完备化 \(f'_*(\mathcal {O}_X)_{y'}=\mathcal {O}_{y'}\) 却是局部环！我们证明了：
\begin{theorem}[{Zariski 的“连通性定理”}]{5}\setnumber{5}\label{fga2-theorem-5}\skipsinglepar
\par{}设 \(f\colon  X\to  Y\) 为固有态射。则 \(f\) 唯一地（在同构意义下）分解为 \(f=gf'\)，其中 \(g\colon  Y'\to  Y\) 有限，而 \(f'\colon  X\to  Y'\) 满足 \(f'_*(\mathcal {O}_X)=\mathcal {O}_{Y'}\)，从而 \(g_*(\mathcal {O}_{Y'})=f_*(\mathcal {O}_X)\)。此外，\(f'\) 的各纤维连通；也就是说，\(f\) 的纤维 \(f^{-1}(y)\) 的连通分支集合，与 \(Y'\) 中位于 \(y\) 上方的点的集合，亦即 \(f_*(\mathcal {O}_X)_y\) 中极大理想的集合，一一对应。\end{theorem}
\par{}由此立即得到连通性定理的通常变体。这里只提出下面的推论：
\begin{corollary}{1}\setnumber{1}\label{fga2-theorem-5-corollary-1}\skipsinglepar
\par{}\(X\) 的点 \(x\) 在其纤维 \(f^{-1}(y)\) 中孤立，当且仅当纤维 \((f')^{-1}(y')\)（其中 \(y'=f'(x)\)）只含点 \(x\)，也等价于 \(f'\) 诱导 \(x\) 的某个邻域到 \(y'\) 的某个邻域的同构。这些点组成开子集 \(U\)，且 \(f'\) 诱导 \(U\) 到 \(Y'\) 的某个开子集的同构。\end{corollary}
\par{}为证明 \(f'\) 在 \(x\) 处为局部同构，注意由 \(f'_*(\mathcal {O}_X)=\mathcal {O}_{Y'}\) 可知，\(f'\) 诱导\emph{同构} \(\mathcal {O}_{y'}\to \mathcal {O}_x\)；又注意，当 \(V\) 遍历 \(y'\) 的一个基本邻域系时，各个 \((f')^{-1}(V)\) 构成 \(x\) 的一个基本邻域系，因为 \(f'\) 为闭映射，其在 \(y'\) 处的纤维仅含点 \(x\)。这样立即得到下面的结论；在“几何”情形下，它由 Chevalley 给出：
\begin{corollary}{2}\setnumber{2}\label{fga2-theorem-5-corollary-2}\skipsinglepar
\par{}\(f\) 为有限态射，当且仅当它固有且各纤维有限。\end{corollary}
\par{}事实上，在\hyperref[fga2-theorem-5-corollary-2]{推论~\ref*{fga2-theorem-5-corollary-2}}的假设下，由上述讨论可知 \(f'\) 为同构。
\par{}设 \(f\colon  X\to  Y\) 为未必固有的态射，但假定 \(X\) 可作为开子集包含于某个固有 \(Y\)-概形 \(\overline {f}\colon \overline {X}\to  Y\) 中；特别地，当 \(\overline {f}\) 为拟射影态射时就是如此。应用\hyperref[fga2-theorem-5-corollary-1]{推论 1}可知，\(\overline {f'}\) 诱导一个同构，把 \(X\) 中在各自纤维内孤立的点所成的集合 \(U\) 映到 \(Y'\) 的一个开子集，并且 \(U\) 确为开子集。于是得到 Zariski“主定理”的下述整体形式：
\begin{theorem}{6}\setnumber{6}\label{fga2-theorem-6}\skipsinglepar
\par{}设 \(f\colon  X\to  Y\) 为有限型态射。则 \(X\) 中在各自纤维内孤立的点所成的集合 \(U\) 是开集；若 \(f\) 为拟射影态射，（\emph{[补注]} 此假设可换成较弱的假设“\(f\) 分离”，依据是以下结果，见 [\hyperref[Gro1960b]{SGA}, VIII, 6.2]：每个拟有限且分离的态射 \(f\colon  X\to  Y\) 都是拟射影态射。）则 \(U\) 在 \(Y\) 上同构于某个概形 \(Y'\) 的开子集，而该概形在 \(Y\) 上有限。\end{theorem}
\par{}由于有限型态射局部仿射，\emph{因而更是}局部拟射影，故立即从\hyperref[fga2-theorem-6]{定理~\ref*{fga2-theorem-6}}得到主定理的通常局部形式。
\chapter{对固有平坦态射的上同调研究的应用}\thispagestyle{fancy}\label{fga2-5}
\par{}设 \(f\colon  X\to  Y\) 为固有态射，\(\mathscr {F}\) 为 \(X\) 上的凝聚层，并假定 \(\mathscr {F}\) 是 \(Y\)-平坦的，即各个 \(\mathscr {F}_x\) 都是环 \(\mathcal {O}_y\) 上的平坦模，其中 \(y=f(x)\)。这也意味着：对任意 \(y\in  Y\)，若沿纤维 \(f^{-1}(y)\) 用 \(\mathfrak {m}_y^n\mathscr {F}\) 给 \(\mathscr {F}\) 赋予滤过，其中 \(\mathfrak {m}_y\) 为 \(\mathcal {O}_y\) 的极大理想，则相伴分次对象同构于 \((\mathscr {F}/\mathfrak {m}_y\mathscr {F})\otimes _{k(y)}\operatorname {gr}(\mathcal {O}_y)\)；换言之，对任意整数 \(n\)，有 \[     \mathfrak {m}_y^n\mathscr {F}/\mathfrak {m}_y^{n+1}\mathscr {F}     = \mathscr {F}_y\otimes _{k(y)}(\mathfrak {m}_y^n/\mathfrak {m}_y^{n+1})   \]。这里 \(\mathscr {F}_y\) 表示 \(\mathscr {F}\) 在 \(X_y\) 上诱导的层 \(\mathscr {F}/\mathfrak {m}_y\mathscr {F}\)，其中 \(X_y\) 表示视为 \(y\) 的剩余域 \(k(y)\) 上固有概形的纤维 \(f^{-1}(y)\)。利用这一同构以及\hyperref[fga2-theorem-2]{定理~\ref*{fga2-theorem-2}}，由 \(X_y\) 以 \(\mathscr {F}_y\) 为系数的上同调，可以得到 \(\operatorname {R}^q f_*(\mathscr {F})\) 在 \(y\) 邻域中的估计，有时还能进行计算。这里，\hyperref[fga2-theorem-2]{定理~\ref*{fga2-theorem-2}}具有形式 \[     \overline {\operatorname {R}^q f_*(\mathscr {F})_y}     = \varprojlim _n \operatorname {H}^q(\mathscr {F}_y,\mathscr {F}/\mathfrak {m}_y^n\mathscr {F}).   \] 我们只指出下面的推论：
\begin{proposition}{1}\setnumber{1}\label{fga2-proposition-1}\skipsinglepar
\par{}设 \(f\colon  X\to  Y\) 为固有态射，\(\mathscr {F}\) 为 \(X\) 上的凝聚 \(Y\)-平坦层。设 \(y\in  Y\)，\(q\) 为整数，并假定 \(\operatorname {H}^q(X_y,\mathscr {F}_y)=0\)。则 \(\operatorname {R}^q f_*(\mathscr {F})\) 在 \(y\) 的某个邻域上为零，并且对任意 \(n\)，自然同态 \[       \operatorname {R}^{q-1}f_*(\mathscr {F})_y       \to  \operatorname {H}^{q-1}(X_y,\mathscr {F}_y/\mathfrak {m}_y^n\mathscr {F}_y)     \] 为满射。\end{proposition}
\par{}特别地，若 \(f\) 为平坦态射，即 \(\mathcal {O}_X\) 是 \(Y\)-平坦的，则 \(X\) 上每个局部自由凝聚层 \(\mathscr {F}\) 都是 \(Y\)-平坦的。设 \(\mathscr {F}\) 和 \(\mathscr {G}\) 为两个这样的层，把\hyperref[fga2-proposition-1]{命题~\ref*{fga2-proposition-1}}应用于 \(\mathscr {H}\kern  -2.5pt\mathit {om}_{\mathcal {O}_X}(\mathscr {F},\mathscr {G})\) 和 \(q=1\)，得到：
\begin{theorem}{7}\setnumber{7}\label{fga2-theorem-7}\skipsinglepar
\par{}设 \(f\) 为固有平坦态射，\(\mathscr {F}\)、\(\mathscr {G}\) 为 \(X\) 上的局部自由凝聚层，且 \(y\in  Y\)；假定 \[\operatorname {H}^1(X_y,\mathscr {H}\kern  -2.5pt\mathit {om}_{\mathcal {O}_X}(\mathscr {F}_y,\mathscr {G}_y))=0\]则每个同态 \(u_0\colon \mathscr {F}_y\to \mathscr {G}_y\) 都由同态 \(u\colon \mathscr {F}|V\to \mathscr {G}|V\) 诱导，其中 \(V=f^{-1}(U)\) 为 \(y\) 的某个邻域 \(U\) 的原像。\end{theorem}
\begin{corollary}{1}\setnumber{1}\label{fga2-theorem-7-corollary-1}\skipsinglepar
\par{}若 \(u_0\) 为同构（分别为单态射、满态射），则当 \(U\) 充分小时，\(u\) 也为同构（分别为单态射、满态射）。\end{corollary}
\par{}特别地：
\begin{corollary}{2}\setnumber{2}\label{fga2-theorem-7-corollary-2}\skipsinglepar
\par{}设 \(E_0\) 为 \(X_y\) 上的局部自由凝聚层，且 \(\operatorname {H}^1(X_y;\mathscr {H}\kern  -2.5pt\mathit {om}_{\mathcal {O}_X}(E_0,E_0))=0\)。则任意两个局部自由层，只要它们在 \(X_y\) 上的限制都同构于 \(E_0\)，便在 \(X_y\) 的某个邻域上彼此同构。\end{corollary}
\par{}因此：
\begin{corollary}{3}\setnumber{3}\label{fga2-theorem-7-corollary-3}\skipsinglepar
\par{}假定 \(\operatorname {H}^1(X_y,\mathcal {O}_{X_y})=0\)。则 \(X\) 上任意两个可逆层，即局部同构于 \(\mathcal {O}_X\) 的层，只要它们在 \(X_y\) 上的限制同构，它们自身便彼此同构。\footnote{译注：法文原文此处也省略了“在该纤维的某个邻域上”的限制；由上面的推论 2 只能直接得到邻域内的同构。这里保留原文陈述，不将其作为整体同构的断言使用。}\end{corollary}
\par{}由此得到：
\begin{proposition}{2}\setnumber{2}\label{fga2-proposition-2}\skipsinglepar
\par{}设 \(Y\) 为连通概形，\(E\) 为 \(Y\) 上的局部自由凝聚层。考虑与 \(E\) 关联的射影空间丛 \(X=\mathbb {P}(E)\)，并赋予它通常的可逆层 \(\mathcal {O}_X(1)\)。则 \(X\) 上每个可逆层 \(\mathscr {L}\) 都同构于形如 \(f^*(\mathscr {L}')\otimes \mathcal {O}_X(n)\) 的层，其中 \(\mathscr {L}'\) 为 \(Y\) 上的可逆层，\(n\) 为整数。此外，\(n\) 唯一确定，\(\mathscr {L}'\) 在同构意义下确定。
\begin{proof}
\par{}上述\hyperref[fga2-theorem-7-corollary-3]{推论~\ref*{fga2-theorem-7-corollary-3}}证明，在每条纤维的某个邻域上，\(\mathscr {L}\) 同构于某个 \(\mathcal {O}_X(n)\)。其余部分大体上是形式的。\end{proof}\end{proposition}
\par{}\hyperref[fga2-proposition-2]{命题~\ref*{fga2-proposition-2}}使我们能够确定从 \(X=\mathbb {P}(E)\) 到另一个射影丛的 \(Y\)-态射。特别地，可得：
\begin{corollary}{1}\setnumber{1}\label{fga2-proposition-2-corollary-1}\skipsinglepar
\par{}设 \(u\) 为 \(X=\mathbb {P}(E)\) 的自同构。则存在 \(Y\) 上的可逆层 \(\mathscr {L}'\)，以及从 \(E\) 到 \(E\otimes \mathscr {L}'\) 的同构 \(v\)，使得 \(u\) 为与 \(\mathbb {P}(E)\xrightarrow {\sim }\mathbb {P}(E\otimes \mathscr {L}')=\mathbb {P}(E)\) 相应的同构；二元组 \((v,\mathscr {L}')\) 在同构意义下确定。\end{corollary}
\par{}设 \(\Gamma \) 为 \(Y\) 上满足 \(E\otimes \mathscr {L}'\) 同构于 \(E\) 的可逆丛 \(\mathscr {L}'\) 的同构类集合。其元素都是挠元，因为若 \(n\) 为 \(E\) 的秩，则取 \(n\) 次外幂可知，必有 \((\mathscr {L}')^{\otimes  n}\xrightarrow {\sim }\mathcal {O}_Y\)。因此，上述推论可表述为以下群的正合列：\[     e \to  \operatorname {Aut}(E)/\Gamma (Y,\mathcal {O}_Y^*) \to  \operatorname {Aut}_Y(X) \to  \Gamma  \to  e   \] （这也可由如下群\emph{层}的正合列所诱导的上同调正合列推出：\[     e \to  \mathcal {O}_X^\times  \to  \mathscr {A}\kern  -1.75pt\mathit {ut}(E) \to  \mathscr {A}\kern  -1.75pt\mathit {ut}_Y(X) \to  e,   \] 其中 \(\mathcal {O}_X^\times \) 为 \(\mathcal {O}_X\) 的“单位元”层，并与 \(\mathscr {A}\kern  -1.75pt\mathit {ut}(E)\) 的中心等同。）
\chapter{对完备 \(\mathscr {J}\)-进环上层与概形的存在和唯一性定理的应用}\thispagestyle{fancy}\label{fga2-6}
\par{}\hyperref[fga2-theorem-7]{定理~\ref*{fga2-theorem-7}}利用\hyperref[fga2-theorem-1]{定理~\ref*{fga2-theorem-1}}和\hyperref[fga2-theorem-2]{定理~\ref*{fga2-theorem-2}}，给出了局部自由凝聚层的唯一性结果。现在利用\hyperref[fga2-theorem-3]{定理~\ref*{fga2-theorem-3}}，可得到层、概形态射以及概形的\emph{存在性}定理。以下，\(A\) 表示一个分离且\emph{完备}的 Noether 局部环。一般的方法仍然是作\emph{形式}构造；这实质上是在作\emph{Artin 环上的代数几何}，然后利用三个基本定理，从中导出“代数”性质的结论。
\begin{proposition}{3}\setnumber{3}\label{fga2-proposition-3}\skipsinglepar
\par{}设 \(\mathfrak {X}\) 为 \(A\) 上固有且平坦的形式概形，\(\mathscr {F}_0\) 为 \(X_0\) 上的局部自由层，满足 \[\operatorname {H}^2(X_0,\mathscr {H}\kern  -2.5pt\mathit {om}_{\mathcal {O}_{X_0}}(\mathscr {F}_0,\mathscr {F}_0))=0\]则存在 \(\mathfrak {X}\) 上的局部自由层 \(\mathscr {F}\)，其在 \(X_0\) 上诱导的层同构于 \(\mathscr {F}_0\)。（若 \(\operatorname {H}^1(X_0,\mathscr {H}\kern  -2.5pt\mathit {om}_{\mathcal {O}_{X_0}}(\mathscr {F}_0,\mathscr {F}_0))=0\)，则这个 \(\mathscr {F}\) 在同构意义下还唯一。）\end{proposition}
\par{}逐步在各个 \(X_n\) 上构造彼此相容的局部自由层 \(\mathscr {F}_n\)。构造 \(\mathscr {F}_n\) 时会遇到位于 \[\operatorname {H}^2(X_0,\mathscr {H}\kern  -2.5pt\mathit {om}_{\mathcal {O}_{X_0}}(\mathscr {F}_0,\mathscr {F}_0))\otimes _{A/\mathscr {J}}(\mathscr {J}^n/\mathscr {J}^{n+1})\] 中的障碍，但依假设，该群为零。现在利用\hyperref[fga2-theorem-3]{定理~\ref*{fga2-theorem-3}}，得到：
\begin{corollary}{1}\setnumber{1}\label{fga2-proposition-3-corollary-1}\skipsinglepar
\par{}设 \(X\) 为 \(A\) 上固有且平坦的概形，\(\mathscr {F}_0\) 如上。则存在 \(X\) 上的局部自由层 \(\mathscr {F}\)，其在 \(X_0\) 上诱导的层同构于 \(\mathscr {F}_0\)。若 \(\operatorname {H}^1(X_0,\mathscr {H}\kern  -2.5pt\mathit {om}_{\mathcal {O}_{X_0}}(\mathscr {F}_0,\mathscr {F}_0))=0\)，则这个 \(\mathscr {F}\) 在同构意义下还唯一。\end{corollary}
\par{}设 \(X_0\) 为域 \(k\) 上的有限型概形，并假定 \(X_0\) 在 \(k\) 上\emph{光滑}，这里指\emph{绝对}光滑，但不必在 \(k\) 上固有。设 \(A\) 为剩余域为 \(k\) 的 Artin 局部环。我们要寻找 \(A\) 上平坦且满足 \(X\otimes _A k=X_0\) 的概形 \(X\)；这是“\emph{模空间理论}”，或 \(X_0\) 的“结构变形”研究的出发点。给定这样的 \(X\)，等价于在拓扑空间 \(X_0\) 上给定层 \(\mathcal {O}_X\)，并配备下述结构：
\begin{enumerate}
    
 \item[i.]\(\mathcal {O}_X\) 是 \(A\)-代数层；

    
 \item[ii.]\(\mathcal {O}_X\) 配备一个与 \(A\)-代数结构相容的增广同态 \(\mathcal {O}_X\to \mathcal {O}_{X_0}\)；

  \end{enumerate} 同时要求这些数据满足以下条件：增广诱导同构 \(\mathcal {O}_X\otimes _A k\xrightarrow {\sim }\mathcal {O}_{X_0}\)；\(\mathcal {O}_X\) 在 \(A\) 上平坦，也就是说，用 \(A\) 的极大理想 \(\mathfrak {m}\) 的各次幂给 \(\mathcal {O}_X\) 赋予滤过，则相伴分次代数同构于 \(\operatorname {gr}^0(\mathcal {O}_X)\otimes _k\operatorname {gr}(A)\)，即有同构 \(\mathfrak {m}^n\mathcal {O}_X/\mathfrak {m}^{n+1}\mathcal {O}_X = \mathcal {O}_{X_0}\otimes _k(\mathfrak {m}^n/\mathfrak {m}^{n+1})\)。基本事实如下：
\begin{theorem}{8}\setnumber{8}\label{fga2-theorem-8}\skipsinglepar
\par{}设 \(X_0\) 为域 \(k\) 上的光滑有限型概形，并假定 \(X_0\) 仿射。设 \(A\) 为剩余域为 \(k\) 的 Artin 局部环。则存在 \(A\)-概形 \(X\)，它在 \(A\) 上平坦，且满足 \(X\otimes _A k=X_0\)。此外，任意两个这样的概形必同构。\end{theorem}
\par{}注意，这里的同构不是典范的，因为一般而言，\(X\) 有非平凡的 \(A\)-自同构，在 \(X_0\) 上诱导恒等态射。此外，通常也无法“典范地”选择满足这些条件的 \(X\)；例外是 \(A\) 为 \(k\)-代数的情形，即\emph{等特征}的情形，此时可取 \(X=X_0\otimes _k A\)，也就是 \(\mathcal {O}_X=\mathcal {O}_{X_0}\otimes _k A\)；事实上，这不要求 \(X_0\) 仿射。在不等特征的情形，若 \(X_0\) 不仿射，我一般不知道是否能把 \(X_0\)“提升”为一个定义在 \(A\) 上的 \(X\)。不过，对任意整数 \(n>0\)，置 \(A_{n-1}=A/\mathfrak {m}^n\)，并假定已经把 \(X_0\) 提升为平坦的 \(A_{n-1}\)-概形 \(X_{n-1}\)；我们打算将 \(X_{n-1}\) 提升为平坦的 \(A_n\)-概形 \(X_n\)。由\hyperref[fga2-theorem-8]{定理~\ref*{fga2-theorem-8}}已知，这在局部是可能的；也容易验证，若 \(U_n\) 提升 \(X_{n-1}\) 的开子集 \(U_{n-1}\)，则 \(U_n\) 的、在 \(U_{n-1}\) 上诱导恒等态射的自同构所成的群层，典范同构于 \(\mathfrak {G}_{X_0/k}\otimes _k\mathfrak {m}^n/\mathfrak {m}^{n+1}\) 在 \(U_{n-1}\) 上的限制，因而特别是交换的；这里 \(\mathfrak {G}_{X_0/k}\) 表示 \(X_0\) 上 \(k\)-导子的芽所成的层。由此容易得出，构造提升 \(X_{n-1}\) 的 \(X_n\) 时有一个障碍，位于 \(\operatorname {H}^2(X_0,\mathfrak {G}_{X_0/k})\otimes \mathfrak {m}^n/\mathfrak {m}^{n+1}\) 中。因此：
\begin{corollary}{1}\setnumber{1}\label{fga2-theorem-8-corollary-1}\skipsinglepar
\par{}设 \(X_0\) 为 \(k\) 上的光滑有限型概形，并假定 \[       \operatorname {H}^2(X_0,\mathfrak {G}_{X_0/k})       = 0.     \] 则对每个剩余域为 \(k\) 的 Artin 局部环 \(A\)，都存在平坦的 \(A\)-概形 \(X\)，使得 \(X\otimes _A k=X_0\)。\end{corollary}
\par{}此外，若能找到\emph{一个}在 \(A\) 上平坦且提升 \(X_0\) 的 \(X\)，则由\hyperref[fga2-theorem-8]{定理~\ref*{fga2-theorem-8}}可知，提升 \(X_0\) 的平坦 \(A\)-概形的同构类集合可与 \(\operatorname {H}^1(X_0,\mathscr {A}\kern  -1.75pt\mathit {ut}(X))\) 等同，其中 \(\mathscr {A}\kern  -1.75pt\mathit {ut}(X)\) 表示 \(A\)-代数层 \(\mathscr {O_X}\) 的\emph{与增广相容}的自同构的芽所成的层。\(\mathcal {O}_X\) 的滤过定义 \(\mathscr {A}\kern  -1.75pt\mathit {ut}(X)\) 的一个滤过，用此层除以滤过的第 \(n\) 个子群，所得商为 \(\mathscr {A}\kern  -1.75pt\mathit {ut}(X_n)\)；相伴分次对象是交换的，可与 \(\mathfrak {G}_{X_0/k}\otimes _k\operatorname {gr}(A)\) 等同。特别地，若 \(\mathfrak {m}^{n+1}\) 是 \(\mathfrak {m}\) 的第一个为零的幂，则 \(F^n(\mathscr {A}\kern  -1.75pt\mathit {ut}(X))\)，也就是滤过的最后一级，包含在 \(\mathscr {A}\kern  -1.75pt\mathit {ut}(X)\) 的中心内，并同构于 \(\mathfrak {G}_{X_0/k}\otimes _k\mathfrak {m}^n\)；它也是 \(X\) 的、在 \(X_{n-1}=X\otimes _A A/\mathfrak {m}^n\) 上诱导恒等态射的自同构的芽所成的层。利用这些结果，立即得到下述陈述：
\begin{corollary}{2}\setnumber{2}\label{fga2-theorem-8-corollary-2}\skipsinglepar
\par{}设 \(X_0\) 为 \(k\) 上的光滑有限型概形，\(A\) 为 Artin 局部环，剩余域为 \(k\)，极大理想为 \(\mathfrak {m}\)。假定 \(\mathfrak {m}^{n+1}=0\)。置 \(A_{n-1}=A/\mathfrak {m}^n\)，并设 \(X_{n-1}\) 为满足 \(X_{n-1}\otimes _Ak=X_0\) 的平坦 \(A_{n-1}\)-概形。则满足 \(X\otimes _AA_{n-1}=X_{n-1}\) 的平坦 \(A\)-概形 \(X_n\)，模去在 \(X_{n-1}\) 上诱导恒等态射的同构后，其类的集合或者为空，或者是 \(\operatorname {H}^1(X_0,\mathfrak {G}_{X_0/k})\otimes _k\mathfrak {m}^n\) 下的主齐性空间。\end{corollary}
\par{}（注意，一般而言，后一个主齐性空间并无优先选定的原点，因为没有优先选定的、将 \(X_{n-1}\) 提升到 \(X_n\) 的方法。）
\begin{corollary}{3}\setnumber{3}\label{fga2-theorem-8-corollary-3}\skipsinglepar
\par{}设 \(X_0\) 为 \(k\) 上的光滑有限型概形，并假定 \(\operatorname {H}^1(X_0,\mathfrak {G}_{X_0/k})=0\)。则对每个剩余域为 \(k\) 的 Artin 局部环 \(A\)，在同构意义下，满足 \(X\otimes _Ak=X_0\) 的\emph{平坦 \(A\)-概形 \(X\) 至多存在一个}。\end{corollary}
\par{}\hyperref[fga2-theorem-8]{定理~\ref*{fga2-theorem-8}}的\hyperref[fga2-theorem-8-corollary-1]{推论~\ref*{fga2-theorem-8-corollary-1}}与\hyperref[fga2-theorem-8-corollary-3]{推论~\ref*{fga2-theorem-8-corollary-3}}立即蕴含一些表面上更一般的陈述：只需假定 \(A\) 为\emph{剩余域为 \(k\) 的完备 Noether 局部环}，同时把 \(X\) 理解为 \(A\) 上的形式概形：
\begin{theorem}{9}\setnumber{9}\label{fga2-theorem-9}\skipsinglepar
\par{}设 \(k\) 为域，\(X_0\) 为 \(k\) 上的\emph{光滑}有限型概形。对每个剩余域为 \(k\) 的完备 Noether 局部环 \(A\)，以 \(F(A)\) 表示 \(A\) 上有限型、在 \(A\) 上平坦并满足 \(X\otimes _Ak=X_0\) 的形式概形 \(\mathfrak {X}\) 的类所成的集合，其中模去的是在 \(X_0\) 上诱导恒等态射的同构。采用这些记号，对所有 \(A\)，有：
\begin{enumerate}
      
 \item[i.]若 \(\operatorname {H}^1(X_0,\mathfrak {G}_{X_0/k})=0\)，则 \(F(A)\) 至多有一个元素；

      
 \item[ii.]若 \(\operatorname {H}^2(X_0,\mathfrak {G}_{X_0/k})=0\)，则 \(F(A)\) 至少有一个元素。

    \end{enumerate}\end{theorem}
\begin{corollary}{1}\setnumber{1}\label{fga2-theorem-9-corollary-1}\skipsinglepar
\par{}假定 \(X_0\) 在 \(k\) 上固有。在\hyperref[fga2-theorem-9]{定理~\ref*{fga2-theorem-9}}的条件 (i) 下，对所有 \(A\)，在模去于 \(X_0\) 上诱导恒等态射的同构后，至多存在一个在 \(A\) 上固有、平坦且满足 \(X\otimes _Ak=X_0\) 的概形 \(X\)。\end{corollary}
\par{}可利用\hyperref[fga2-theorem-3-corollary-2]{定理 3 的推论 2}。例如：
\begin{corollary}{2}\setnumber{2}\label{fga2-theorem-9-corollary-2}\skipsinglepar
\par{}若 \(X\) 为固有平坦的 \(A\)-概形，且 \(X\otimes _A k\) 同构于 \(k\) 上维数为 \(r\) 的标准射影概形 \(\mathfrak {P}_k^r\)，则 \(X\) 同构于 \(\mathfrak {P}_A^r\)。\end{corollary}
\par{}（此结果也可由\hyperref[fga2-proposition-3-corollary-1]{命题 3 的推论 1}推出。）
\begin{corollary}{3}\setnumber{3}\label{fga2-theorem-9-corollary-3}\skipsinglepar
\par{}设 \(X_0\) 为 \(k\) 上的光滑射影概形，并假定 \[       H^2(X_0,\mathcal {O}_{X_0})       = H^2(X_0,\mathfrak {G}_{X_0/k})       = 0.     \] 则对所有 \(A\)，都存在平坦射影的 \(A\)-概形 \(X\)，使得 \(X\otimes _Ak=X_0\)。\end{corollary}
\par{}将\hyperref[fga2-theorem-9]{定理~\ref*{fga2-theorem-9}}的 (ii) 与\hyperref[fga2-theorem-4]{定理~\ref*{fga2-theorem-4}}结合即可。特别地：
\begin{corollary}{4}\setnumber{4}\label{fga2-theorem-9-corollary-4}\skipsinglepar
\par{}设 \(X_0\) 为 \(k\) 上一条完备光滑代数曲线的概形。则对每个剩余域为 \(k\) 的完备 Noether 局部环 \(A\)，都存在 \(A\) 上的“光滑曲线概形”\(X\)，使得 \(X\otimes _Ak=X_0\)。\end{corollary}
\begin{remarks}{}\label{fga2-6-remarks}\skipsinglepar
\begin{enumerate}\item{}\hyperref[fga2-theorem-9-corollary-3]{推论~\ref*{fga2-theorem-9-corollary-3}}与\hyperref[fga2-theorem-9-corollary-4]{推论~\ref*{fga2-theorem-9-corollary-4}}在 \(k\) 的特征为 \(p\neq 0\) 时尤其有意义：可取 \(A\) 为剩余域为 \(k\)、\emph{特征为 \(0\)}的离散赋值环，例如“尽可能小的 \(A\)”，即 \(p\) 生成极大理想的情形。（事实上，由 Cohen 的定理，只需对这样的环 \(A\) 知道\hyperref[fga2-theorem-9-corollary-3]{推论~\ref*{fga2-theorem-9-corollary-3}}和\hyperref[fga2-theorem-9-corollary-4]{推论~\ref*{fga2-theorem-9-corollary-4}}即可。）关于这一点，据专家们说，目前尚不知是否存在域 \(k\) 上的概形，不能由这样的环 \(A\) 上定义的平坦概形\(\mod  p\) 约化得到。至少，本节的结果给出了系统研究这一问题的途径。首先应考察我们在 \(H^2(X_0,\mathfrak {G}_{X_0/k})\) 中遇到的第一个障碍是否必然为零。

    \item{}注意，\hyperref[fga2-theorem-3]{定理~\ref*{fga2-theorem-3}}及其相应的方法只适用于\emph{完备}的底环；为明确起见，还可假定底环为局部环。要从关于局部环的完备化的已知结果过渡到该局部环本身的相应结果，还需要第四个“基本定理”，而它的准确陈述尚待确定。

    \item{}应将本节，尤其是上述\hyperref[fga2-theorem-9-corollary-1]{推论~\ref*{fga2-theorem-9-corollary-1}}、\hyperref[fga2-theorem-9-corollary-2]{推论~\ref*{fga2-theorem-9-corollary-2}}，以及下一节的结果，与 Kodaira–Spencer 关于复结构变形的结果相比较。利用刚才提及的猜想性定理，我们应能在\hyperref[fga2-theorem-9-corollary-1]{推论~\ref*{fga2-theorem-9-corollary-1}}的条件下，但不再假定 \(A\) 完备时，推出存在包含 \(A\) 的环 \(A'\)，它在 \(A\) 上有限且无分歧，使得 \(X\otimes _AA'\) 与 \(X'\otimes _AA'\) 在 \(A'\) 上同构；这里 \(X\)、\(X'\) 是给定的固有平坦 \(A\)-概形，满足 \(X\otimes _Ak=X'\otimes _Ak=X_0\)。至少，当 \(X_0=\mathfrak {P}_k^r\) 时，可利用\hyperref[fga2-proposition-2-corollary-1]{命题 2 的推论 1}证明这一点。无论如何，若 \(H^1(X_0,\mathfrak {G}_{X_0/k})=0\)，就可证明 \(X\) 和 \(X'\) 在 \(Y=\operatorname {Spec}(A)\) 的任意点 \(y\) 上的纤维同构，至少在转到剩余域 \(k(y)\) 的代数闭包后如此。（不必假定 \(A\) 为局部环时，还有一个表面上更强的局部结果。）关于射影空间的“结构变形”，再提出以下问题，它来自 Kodaira–Spencer 的一个相应问题。设 \(X\) 为局部整环 \(A\) 上的固有平坦概形，后者的分式域为 \(K\)，剩余域为 \(k\)；假定 \(X\otimes _AK\) 同构于 \(\mathfrak {P}_K^r\)。那么，\(X\otimes _Ak=X_0\) 是否同构于 \(\mathfrak {P}_k^R\)，至少在 \(k\) 的代数闭包上如此？在此问题中，可假定 \(A\) 为完备离散赋值环。当 \(X_0\) 为 Abel 簇时，也有类似的问题。\end{enumerate}\end{remarks}
\begin{remark}{}\label{fga2-6-remark-i}\skipsinglepar
\par{}\emph{[补注]} \emph{（关于上面的注 1。）} J.-P. Serre 已在 \hyperref[Ser1961]{[Ser1961]} 中构造了一个特征为 \(p>0\) 的代数闭域 \(k\) 上维数为 \(3\) 的非奇异射影簇；它不能由剩余域为 \(k\)、分式域特征为 \(0\) 的局部整环上的固有概形约化得到。据说 Mumford 对非奇异射影\emph{曲面}也得到了类似结果。\end{remark}
\begin{remark}{}\label{fga2-6-remark-ii}\skipsinglepar
\par{}\emph{[补注]} \emph{（关于上面的注 2 和注 3。）} 对这里猜测的结果，我现在已不如以前乐观。不过，上面注 3 末尾提出的射影空间结构变形问题，已由 Hironaka 肯定地解决；Abel 簇的类似问题也已由 Koizumi 解决。\end{remark}
\chapter{对“模空间理论”的应用}\thispagestyle{fancy}\label{fga2-7}
\par{}报告人自己接触这一理论也为时不久，因此只能限于作一些简略的说明。为简便起见，我们在\emph{域} \(k\) 上工作，也就是说，处于\emph{等特征}的情形；尽管看来借助\hyperref[fga2-theorem-8]{定理~\ref*{fga2-theorem-8}}，无需作本质性的改动，也能讨论更一般的情形。目前我们尚未越过“形式”阶段，但报告人仍希望由此在某些情形下构造出真正的模概形，尤其希望对每个整数 \(g\)，构造一个整数上的概形，起到亏格为 \(g\) 的光滑曲线的泛模概形的作用。
\begin{remark}{}\label{fga2-7-remark}\skipsinglepar
\par{}\emph{[补注]} Mumford 最近已经构造了亏格为 \(g\) 的曲线的模概形，参见 \emph{Mumford–Tate seminar}，Harvard University，1961–62。\hyperref[fga2-theorem-10]{定理~\ref*{fga2-theorem-10}}还证明，模空间理论中的“层级为 \(n\) 的 Jacobi 概形”是非奇异的，甚至在 \(\underline {\mathbb {Z}}\) 上光滑。\end{remark}
\par{}继续采用\hyperref[fga2-theorem-9]{定理~\ref*{fga2-theorem-9}}的框架和记号，并假定 \(A\) 是 \(k\) 上有限秩的局部代数；为简便起见，假定该域代数闭。此时，\(F(A)\) 可视为关于 \(A\) 的、取值于集合范畴的协变函子：\(k\)-代数同态 \(A\to  B\) 定义映射 \(F(A)\to  F(B)\)，因为每个满足 \(X\otimes _Ak=X_0\) 的平坦 \(A\)-概形 \(X\) 都给出具有相同性质的 \(B\)-概形 \(X\otimes _AB\)。\emph{假定}我们能找到一个完备 Noether 局部 \(k\)-代数 \(\mathcal {O}\)，以及一个函子性同构 \[     \operatorname {Hom}(\mathcal {O},A) \xrightarrow {\sim } F(A)     \tag{*}   \]，其中左边表示 \(k\)-代数同态。容易看出，这样的 \(\mathcal {O}\) 在典范同构意义下确定。因此，把 \(\mathcal {O}\) 的形式谱 \(\mathfrak {Y}\)，即只含一个点、且其拓扑环层恰为 \(\mathcal {O}\) 的拓扑空间，称为\emph{\(X_0\) 的形式模概形}。（注意，它未必存在。）设 \(\mathfrak {r}\) 为 \(\mathcal {O}\) 的极大理想，并对所有 \(n\) 置 \(\mathcal {O}_n=\mathcal {O} /\mathfrak {r}^{n+1}\)，从而 \(\mathcal {O}_0=k\)。则典范同态 \(\mathcal {O}\to \mathcal {O}_n\) 是 \(\operatorname {Hom}(\mathcal {O},\mathcal {O}_n)\) 的一个元素，因而定义 \(F(\mathcal {O}_n)\) 的一个元素，也就是一个平坦 \(\mathcal {O}_n\)-概形 \(X_n\)，其模 \(\mathfrak {r}\) 的约化为 \(X_0\)。这些 \(X_n\) 经由标量扩张，在这里即约化，彼此诱导；因此它们来自一个形式概形 \(\mathfrak {X}\)，后者由形式模概形 \(\mathfrak {Y}\) 唯一确定；此外，\(\mathfrak {X}\) 在 \(\mathfrak {Y}\) 上平坦，且 \(\mathfrak {X}_0=X_0\)。立即可见，同构 \((*)\) 可作如下描述：把每个 \(k\)-代数同态 \(\mathcal {O}\to  A\) 对应到 \(A\)-概形 \(\mathfrak {X}\otimes _{\mathcal {O}} A\) 的类；也就是说，对每个 \(k\)-概形态射 \(\mathfrak {Y}'=\operatorname {Spec}(A)\to \mathfrak {Y}\)，对应由基变换给出的 \(\mathfrak {Y}'\)-概形 \(\mathfrak {X}\otimes _{\mathfrak {Y}}\mathfrak {Y}'\)。进一步可见，即使只假定 \(A\) 为完备 Noether 局部 \(k\)-代数而不必为 Artin 代数，同构 \((*)\) 及其上述描述仍然成立。当然，和以往一样，\(\mathcal {O}\) \emph{先验地}完全可以含有幂零元；而且很可能存在这样的情形：\(\mathcal {O}\) 本身为 Artin 环，却不等于 \(k\)。这说明 Kodaira–Spencer 的观点，即仅考虑为\emph{正则}环的 \(A\)，在一般情形下\emph{先验地}就不充分。
\par{}还须给出 \(X_0\) 的形式模概形存在的充分条件，这里假定它在 \(k\) 上固有。一般而言，容易对函子 \(A\to  F(A)\)，即从有限秩局部 \(k\)-代数到集合的函子，给出简单的充要条件，使它具有 \(\operatorname {Hom}(\mathcal {O},A)\) 的形式，其中 \(\mathcal {O}\) 为适当的环。这里不作详述。只指出，在我们研究的情形中，这些条件对 \(X_0\) 施加了非平凡的上同调条件；尽管报告人尚未构造出反例，但这些条件似乎不大可能总是满足。不过，条件 \(H^0(X,\mathfrak {G}_{X_0/k})=0\) 很可能是保证形式模概形存在的\emph{充分}条件，尽管它远非必要条件。这里只陈述一个处理特别简单情形的定理；它在解析空间理论中的对应结果是熟知的，参见 Kodaira–Spencer。利用上一节的结果，容易证明它：
\begin{theorem}{10}\setnumber{10}\label{fga2-theorem-10}\skipsinglepar
\par{}设 \(X_0\) 为域 \(k\) 上的光滑固有概形，且满足 \[       H^0(X_0,\mathfrak {G}_{X_0/k})       = H^2(X_0,\mathfrak {G}_{X_0/k})       = 0.     \] 则 \(X_0\) 的形式模概形存在，对应于一个正则局部环 \(\mathcal {O}\)，也就是 \(k\) 上的形式幂级数代数。\end{theorem}
\par{}如前所述，\(\mathcal {O}\) 上的形式概形 \(\mathfrak {X}\) 一般未必能够代数化；不过，已知当 \(X_0\) 射影且 \(H^2(X_0,\mathcal {O}_{X_0})=0\) 时，这一点成立（\hyperref[fga2-theorem-4]{定理~\ref*{fga2-theorem-4}}），例如 \(X_0\) 的维数为 \(1\) 时就是如此。这使我们有希望在整数上构造给定亏格曲线的模概形……
\par{}还应指出，本节所述的这些方法既可用于 Picard 簇的构造与研究，也可用于许多其他构造。不久我们将再讨论这一点。
\chapter{对基本群的应用}\thispagestyle{fancy}\label{fga2-8}
\par{}以上方法使我们能够仿照拓扑理论，着手系统地研究基本群。本节陈述的前两个定理推广了 Lang–Serre 最近一篇工作的结果。
\par{}设 \(X\) 为概形。满足以下条件的 \(X\)-概形 \(X'\)，称为\emph{\(X\) 的无分歧覆盖}：
\begin{enumerate}
    
 \item[i.]\(X'\) 在 \(X\) 上有限，即由 \(X\) 上的凝聚代数层 \(\mathscr {A}=\mathscr {A}(X')\) 定义；

    
 \item[ii.]\(\mathscr {A}\) 是 \(X\) 上的局部自由层；

    
 \item[iii.]对所有 \(x\in  X\)，商 \(\mathscr {A}_x/\mathfrak {m}_x\mathscr {A}_x = \mathscr {A}_x\otimes _{\mathcal {O}_X}k(x)\) 是 \(k(x)\) 上的可分代数。

  \end{enumerate} 这一无分歧覆盖的概念，由 Serre 和报告人提出，具有一切合理预期的初等性质，这里不一一列举。只指出，它导出了一套\emph{Galois 理论}，仿照经典 Galois 理论及拓扑覆盖的 Galois 理论而建立，并包含经典 Galois 理论；总体而言，这里的证明反而比通常的经典证明更简单。更确切地说，概形 \(X\) 的一个\emph{几何点}，是指从代数闭域 \(\Omega \) 的谱 \(\xi \) 到 \(X\) 的态射 \(a\)；也就是给定 \(X\) 的某个点 \(x=|a|\) 的剩余域 \(k(x)\) 的一个代数闭扩域，该点称为几何点 \(a\) 的\emph{所在地}。若 \(X'\) 为 \(X\) 的无分歧覆盖，可给它配上“位于 \(a\) 上方的 \(X'\) 的几何点”的集合 \(E_a(X')\)；这就是由位于 \(x\) 上方的点 \(x'\in  X'\)，以及一个到 \(\Omega \) 的 \(k(x)\)-同态组成的二元组的集合。于是，对固定的 \((X,a)\)，得到一个函子 \(F(X,a)\)，从 \(X\) 的无分歧覆盖 \(X'\) 所成的范畴 \(R(X)\) 到有限集合的范畴。如果 \(X\) 连通，则由 \(R(X)\) 和 \(F(X,a)\) 组成的二元组具有所有必要的形式性质，可与某个适当的全不连通紧拓扑群 \(\pi \)，即有限群的射影极限，所定义的类似二元组同构：取 \(\pi \) 连续作用的有限集合所成的范畴 \(\mathcal {C}(\pi )\)，以及从此范畴到有限集合范畴的恒等函子 \(F(\pi )(E)=E\)。要求 \((\mathcal {C}(\pi ),F(\pi ))\) 同构于给定的二元组，也就在典范同构意义下确定了群 \(\pi \)。具体而言，称 \(\pi \) 为\emph{连通概形 \(X\) 在几何点 \(a\) 处的基本群}，记为 \(\pi _1(X,a)\)。若 \(X\) 不连通，可将其换成包含 \(x=|a|\) 的连通分支。若 \(X\) 连通，则对 \(X\) 的任意两个几何点 \(a'\)、\(a''\)，群 \(\pi _1(X,a')\) 与 \(\pi _1(X,a'')\) 同构，且同构模内自同构而确定。因此，照通常做法，可为我们的目的选择最合适的 \(a\)，例如在假定 \(X\) 不可约时选其泛点。当然，\(\pi _1(X,a)\) 是\emph{关于带基点概形 \((X,a)\) 的协变函子}。于是，按照熟悉的对应，关于不可分覆盖分类的每个陈述\footnote{译注：法文原文为“revêtements inséparables”，英文亦照录为“inseparable coverings”。本段的范畴与基本群实际围绕无分歧覆盖构造。此处保留原文用词，并标出这一疑误。}都可翻译成群论的语言，只须注意这里涉及\emph{拓扑群}。
\par{}我们的目标是对固有态射 \(f\colon  X\to  Y\) 建立纤维丛同伦正合列的对应物。显然，由于我们还不知道高阶同伦群应当是什么，所得结果必然不完整。为了能应用\cref{fga2-3}中的基本定理，必须先说明关于 Artin 环或域上概形的几个初等引理；这仍遵循一般的处理步骤！
\begin{lemma}{1}\setnumber{1}\label{fga2-lemma-1}\skipsinglepar
\par{}设 \((X',a')\) 为与 \(\pi _1(X,a)\) 在带标记点 \(e\) 的有限集合 \(E\) 中的带基点表示相应的带基点无分歧覆盖。则典范态射 \(\pi _1(X',a')\to \pi _1(X,a)\) 把定义域与 \(\pi _1(X,a)\) 中点 \(e\) 的稳定子等同，因而为单射。\end{lemma}
\begin{lemma}{2}\setnumber{2}\label{fga2-lemma-2}\skipsinglepar
\par{}设 \(X\) 为域 \(k\) 上的代数概形，\(k'\) 为 \(k\) 的纯不可分扩张。则 \(X\otimes _kk'\) 的每个无分歧覆盖，都是 \(X\) 的某个无分歧覆盖的原像，即标量扩张；后者在典范同构意义下确定。\end{lemma}
\par{}特别地，由这两个引理可得：对 \(k\) 的\emph{任意}代数扩张 \(K\)，以及 \(X'=X\otimes _kK\) 的每个投影到 \(X\) 的几何点 \(a\) 的几何点 \(a'\)，函子性同态 \(\pi _1(X',a')\to \pi _1(X,a)\) 都是单射。
\begin{lemma}{3}\setnumber{3}\label{fga2-lemma-3}\skipsinglepar
\par{}设 \(X\) 为 Artin 局部环 \(A\) 上的完备概形，满足 \(H^0(X,\mathcal {O}_X)=A\)。设 \(X'\) 为 \(X\) 的无分歧覆盖，并置 \(A'=H^0(X',\mathcal {O}_{X'})\)；因而它是 \(A\) 上的有限环扩张，但先验地可能在 \(A\) 上分歧。设 \(X_0\) 与 \(X'_0\) 分别为 \(X\) 与 \(X'\) 对应的约化子概形，即分别对 \(\mathcal {O}_X\) 与 \(\mathcal {O}_{X'}\) 中的幂零元层取商而得。设 \(k\) 为 \(A/\mathfrak {r}(A)\) 的子域，且 \(A/\mathfrak {r}(A)\) 在它上面有限；于是 \(X_0\) 是 \(k\) 上的完备代数概形，\(X'_0\) 是其无分歧覆盖。最后，设 \(\Omega \) 为 \(k\) 的代数闭扩域，考虑 \(X_0\otimes _k\Omega \) 的无分歧覆盖 \(X'_0\otimes _k\Omega \)。则以下两个条件等价：
\begin{enumerate}
      
 \item[i.]\(X'_0\otimes _k\Omega \) 在 \(X_0\otimes _k\Omega \) 上完全分解；

      
 \item[ii.]自然态射 \(X'\to  X\otimes _AA'\) 是同构。

    \end{enumerate} 在这些条件下，\(A'\) 是 \(A\) 的\emph{无分歧}扩张。最后，若 \(X'\) 连通，则条件 (i) 等价于下面这个表面上较弱的条件：
\begin{enumerate}
      
 \item[i bis.]\(X'_0\otimes _k\Omega \) 在 \(X_0\otimes _k\Omega \) 上有正则截面。

    \end{enumerate}\end{lemma}
\par{}满足\hyperref[fga2-lemma-3]{引理~\ref*{fga2-lemma-3}}的条件 (ii) 时，称 \(X\) 的无分歧覆盖 \(X'\) 为\emph{几何平凡}的。
\begin{lemma}{4}\setnumber{4}\label{fga2-lemma-4}\skipsinglepar
\par{}设 \(f\colon  X\to  Y\) 为满足 \(f_*(\mathcal {O}_X)=\mathcal {O}_Y\) 的固有态射。设 \(a\) 为 \(X\) 的几何点，\(b\) 为它在 \(Y\) 上的投影。则 \(\pi _1(X,a)\to \pi _1(Y,b)\) 为\emph{满射}。\end{lemma}
\par{}实际上，只须证明：若 \(Y\) 的一个无分歧覆盖 \(Y'\)，对应于局部自由代数层 \(\mathscr {A}\)，使 \(X\otimes _YY'\) 不连通，则 \(Y'\) 也不连通。事实上，此时 \(f^*(\mathscr {A})\) 为两个非零环层的直和，从而其直接像亦然，而该直接像恰为 \(\mathscr {A}\otimes  f_*(\mathcal {O}_X)=\mathscr {A}\)。
\begin{lemma}{5}\setnumber{5}\label{fga2-lemma-5}\skipsinglepar
\par{}设 \(X\) 为域 \(k\) 上的完备概形，并假定 \(H^0(X,\mathcal {O}_X)\) 是局部环 \(A\)，且 \(A/\mathfrak {r}(A)\) 在 \(k\) 上纯不可分。设 \(\Omega \) 为 \(k\) 的代数闭包，置 \(\overline {X}=X\otimes _k\Omega \)，它是连通的。选取 \(\overline {X}\) 的几何点 \(\overline {a}\)，其投影为 \(X\) 的几何点 \(a\)。则有正合列 \[       e       \to  \pi _1(\overline {X},\overline {a})       \to  \pi _1(X,a)       \to  \pi _1(k,b)       \to  e     \] 其中 \(\pi _1(k,b)\) 为 \(\Omega \) 在 \(k\) 上的 Galois 群。
\begin{proof}
\par{}第一个同态的单射性已由\hyperref[fga2-lemma-1]{引理~\ref*{fga2-lemma-1}}与\hyperref[fga2-lemma-2]{引理~\ref*{fga2-lemma-2}}证明；中间的正合性由\hyperref[fga2-lemma-3]{引理~\ref*{fga2-lemma-3}}得到；最后一个同态的满射性由\hyperref[fga2-lemma-4]{引理~\ref*{fga2-lemma-4}}得到，只有这一步用到了 \(A/\mathfrak {r}(A)\) 为纯不可分扩张。\end{proof}\end{lemma}
\begin{proposition}{4}\setnumber{4}\label{fga2-proposition-4}\skipsinglepar
\par{}设 \(f\colon  X\to  Y\) 为固有平坦态射，且对所有 \(y\in  Y\)，代数 \(H^0(f^{-1}(y),\mathcal {O}_{f^{-1}(y)})\) 在剩余域 \(k(y)\) 上可分；例如，若 \(f^{-1}(y)\) 是 \(k(y)\) 上的\emph{可分概形}，即约化且其各不可约分支对应的域都是 \(k(y)\) 的可分扩张，便满足此条件。则与 \(f_*(\mathcal {O}_X)\) 关联的 \(Y\) 的覆盖 \(Y'\) 无分歧。
\begin{proof}
\par{}借助\hyperref[fga2-theorem-2]{定理~\ref*{fga2-theorem-2}}，证明很容易。\end{proof}\end{proposition}
\par{}把\hyperref[fga2-proposition-4]{命题~\ref*{fga2-proposition-4}}与\hyperref[fga2-lemma-1]{引理~\ref*{fga2-lemma-1}}结合，实际上可把固有平坦态射（具有可分纤维）的同伦研究归结为 \(f_*(\mathcal {O}_X)=\mathcal {O}_Y\) 的情形，因为利用 Stein 分解，可将 \(Y\) 换成 \(Y'\)。
\begin{remark}{}\label{fga2-8-remark}\skipsinglepar
\par{}纤维为可分概形（分别为光滑概形）的有限型平坦态射，称为\emph{可分}态射（分别为\emph{光滑}态射）。可以证明，若 \(f\) 平坦且 \(f^{-1}(y)\) 可分（分别为光滑），则存在 \(f^{-1}(y)\) 的一个邻域，使 \(f\) 在其上可分（分别为光滑）。对“绝对正规”也有相同结果，这就是\emph{Bertini 定理}。\end{remark}
\par{}设 \(f\colon  X\to  Y\) 为满足 
\begin{enumerate}
    
 \item[i.]
      \(f_*(\mathcal {O}_X) = \mathcal {O}_Y\)
    

  \end{enumerate} 的固有态射，并设 \(X'\) 为 \(X\) 上的有限概形。设 \(Y'\) 为与 \(X'\to  Y\) 的 Stein 分解相应的 \(Y\) 的覆盖，参见\hyperref[fga2-theorem-5]{定理~\ref*{fga2-theorem-5}}。设 \(y\in  Y\)，则 \(X'\) 在 \(y\) 上的纤维 \(F'\) 的连通分支集合，可与位于 \(y\) 上方的点 \(y'\in  Y'\) 的集合等同（\hyperref[fga2-theorem-5]{定理~\ref*{fga2-theorem-5}}）。考虑显然的态射 
\begin{eqn}{*}\setnumber{*}\label{fga2-equation-star}\skipsinglepar\[       X'\to  X\times _Y Y'     \tag{*}     \]\end{eqn}，它由自然态射 \(X'\to  X\) 与 \(X\to  Y'\) 诱导。如果 \(X'\) 具有 \(X\times _Y Y''\) 的形式，其中 \(Y''\) 是 \(Y\) 的无分歧覆盖，则此态射为同构；此时 \(Y'\) 恰为 \(Y''\)，且\hyperref[fga2-equation-star]{式~\ref*{fga2-equation-star}}为恒等态射。我们正要给出条件，使 \(X'\) 具有刚才所述的形式，也就是说，使 \(Y'\) 无分歧且\hyperref[fga2-equation-star]{式~\ref*{fga2-equation-star}}为同构。为此，引入 \(X\) 在 \(y\) 处的纤维 \(F\)，这是 \(k(y)\) 上的固有概形，\(F'\) 是它的覆盖；若 \(X\) 无分歧，此覆盖亦无分歧。设 \(F'_1\) 为 \(F'\) 的一个连通分支，对应于 \(Y'\) 中位于 \(y\) 上方的点 \(y'_1\)。再假定：
\begin{enumerate}
    
 \item[ii.]\(X'\) 在 \(F'_1\) 的各点处相对于 \(X\) 无分歧，因而 \(F'_1\) 是 \(F\) 的无分歧覆盖；并且

    
 \item[iii.]\(F'_1\) 是 \(F\) 的几何平凡覆盖，参见\hyperref[fga2-lemma-3]{引理~\ref*{fga2-lemma-3}}。

  \end{enumerate}
\begin{theorem}{11}\setnumber{11}\label{fga2-theorem-11}\skipsinglepar
\par{}在上述条件下，存在 \(Y'\) 中 \(y'_1\) 的一个开邻域 \(U'\)，使\hyperref[fga2-equation-star]{式~\ref*{fga2-equation-star}}在 \(U'\) 上为同构。此外，\(Y'\) 在 \(y'_1\) 处相对于 \(Y\) 无分歧，但在 \(Y'\) 中位于 \(y\) 上方的其他点 \(y'\) 处可能分歧。
\begin{proof}
\par{}当然，条件 (ii)、(iii) 对定理的结论也是必要的。借助\hyperref[fga2-lemma-3]{引理~\ref*{fga2-lemma-3}}和\hyperref[fga2-theorem-2]{定理~\ref*{fga2-theorem-2}}，定理的证明很容易。\end{proof}\end{theorem}
\begin{corollary}{1}\setnumber{1}\label{fga2-theorem-11-corollary-1}\skipsinglepar
\par{}仍假定上述 (i) 成立。\(X\) 上的一个无分歧覆盖同构于 \(Y\) 的某个无分歧覆盖 \(Y'\) 的原像，当且仅当 \(X'\) 在每条纤维 \(f^{-1}(y)\) 上诱导的覆盖几何平凡。
\begin{proof}
\par{}由\hyperref[fga2-theorem-11]{定理~\ref*{fga2-theorem-11}}，\(Y\) 中满足此条件的点组成开集，故只需在闭点 \(y\) 处验证即可……\end{proof}\end{corollary}
\par{}注意，下面的陈述与\hyperref[fga2-theorem-11-corollary-1]{推论~\ref*{fga2-theorem-11-corollary-1}}等价：
\par{}\emph{同态 \(\pi _1(X)\to \pi _1(Y)\) 的核，是各个 \(\pi _1(\overline {f^{-1}(y)})\) 在 \(\pi _1(X)\) 中的像所生成的闭正规子群；由\hyperref[fga2-lemma-4]{引理~\ref*{fga2-lemma-4}}，该同态为满射。这里 \(\overline {f^{-1}(y)}\) 表示概形 \(f^{-1}(y)\otimes _{k(y)}\overline {k(y)}\)，其中 \(\overline {k(y)}\) 为 \(k(y)\) 的一个代数闭包。}
\par{}注意，因为无法为所有纤维\emph{选择}同一个基点，所以在为 \(X\)，从而为 \(Y\) 选定基点后，同态 \(\pi _1(\overline {f^{-1}(y)})\to \pi _1(X)\) 仍只在模去与 \(\pi _1(X)\) 的内自同构复合的意义下确定。
\begin{corollary}{2}\setnumber{2}\label{fga2-theorem-11-corollary-2}\skipsinglepar
\par{}在\hyperref[fga2-theorem-11]{定理~\ref*{fga2-theorem-11}}的一般条件下，进一步假定 \(Y\)、\(X\)、\(X'\) 都整，并分别以 \(K\)、\(L\)、\(L'\) 表示其函数域。则存在 \(L'\) 中 \(K\) 的可分子扩张 \(K'\)，它与 \(L\) 线性无交，且 \(L'=LK'\)，从而 \(L'=L\otimes _KK'\)。
\begin{proof}
\par{}（把\hyperref[fga2-lemma-3]{引理~\ref*{fga2-lemma-3}}的最后一部分应用于 \(X\) 的泛纤维。）\end{proof}\end{corollary}
\par{}\hyperref[fga2-theorem-11]{定理~\ref*{fga2-theorem-11}}最有意义的应用情形是 \(f\) 为\emph{可分}态射时。此时，\(X'\) 在 \(Y\) 上也可分，故由\hyperref[fga2-proposition-4]{命题~\ref*{fga2-proposition-4}}，\(Y'\) 在 \(Y\) 上无分歧，从而\hyperref[fga2-equation-star]{式~\ref*{fga2-equation-star}}的右端 \(X\times _YY'\) 在 \(X\) 上无分歧。由此容易得到：
\begin{corollary}{3}\setnumber{3}\label{fga2-theorem-11-corollary-3}\skipsinglepar
\par{}除 (i) 外，再假定 \(f\) 可分。设 \(X'\) 为 \(X\) 的连通无分歧覆盖。\(X\) 是 \(Y\) 的某个无分歧覆盖 \(Y'\) 的原像\footnote{译注：法文原文与英译此处均写作 \(X\)，但本推论讨论的对象是覆盖 \(X'\)。正文保留原记号，此处标出疑误。}，当且仅当在某条几何纤维 \(\overline {F}=\overline {f^{-1}(y)}\) 上诱导的覆盖 \(\overline {F}'\) 有一个正则截面。\end{corollary}
\par{}注意，不必假定 \(\overline {F}'\) 在 \(\overline {F}\) 上\emph{几何}平凡；尽管这一条件\emph{先验地}强得多，但它\emph{事后}仍将成立。\hyperref[fga2-theorem-11-corollary-3]{推论~\ref*{fga2-theorem-11-corollary-3}}实际上等价于下面的陈述：
\begin{corollary}{4}\setnumber{4}\label{fga2-theorem-11-corollary-4}\skipsinglepar
\par{}设 \(f\colon  X\to  Y\) 为满足 \(f_*(\mathcal {O}_X)=\mathcal {O}_Y\) 的固有可分态射。设 \(\overline {F}\) 为点 \(y\in  Y\) 的几何纤维，在 \(\overline {F}\) 中选一个几何点，经态射 \(\overline {F}\to  X\to  Y\)，得到 \(X\)、\(Y\) 中的几何点；分别用这三个点作 \(\overline {F}\)、\(X\)、\(Y\) 的基本群的基点。在这些条件下，有正合列 \[       \boxed {\pi _1(\overline {F}) \to  \pi _1(X) \to  \pi _1(Y) \to  e.}     \]\end{corollary}
\par{}由此容易得到下面两个 Serre–Lang 的陈述，这里去掉了所有正规性假设：
\begin{corollary}{5}\setnumber{5}\label{fga2-theorem-11-corollary-5}\skipsinglepar
\par{}设 \(X\)、\(Y\) 为域 \(k\) 上的连通概形，\(X\) 或 \(Y\) 在 \(k\) 上固有，并假定约化概形 \(X_\mathrm {red}\) 在 \(k\) 上可分且完备；当 \(k\) 为完美域时，可分性自动成立。在 \(X\)（分别在 \(Y\)）中选一个几何点 \(a\)（分别为 \(b\)），从而得到 \(X\times _kY\) 中的几何点 \(c=(a,b)\) 及自然态射 \[       \pi _1(X\times _kY,c) \to  \pi _1(X,a)\times \pi _1(Y,b)     \]，它由从 \(\pi _1(X\times  Y,c)\) 到 \(\pi _1(X,a)\) 与 \(\pi _1(Y,b)\) 的函子性态射诱导。此态射为单射；若 \(k\) 代数闭，则还是双射。\end{corollary}
\par{}（\hyperref[fga2-theorem-11-corollary-5]{推论~\ref*{fga2-theorem-11-corollary-5}}中的满射性几乎显然。）于是与 Serre–Lang 一样，可得：
\begin{corollary}{6}\setnumber{6}\label{fga2-theorem-11-corollary-6}\skipsinglepar
\par{}设 \(X\) 为代数闭域 \(k\) 上的连通代数概形，\(K\) 为 \(k\) 的代数闭扩域。则 \(X\) 与 \(X\times _kK\) 的基本群相同；也就是说，后一个概形的每个无分歧覆盖，都由 \(X\) 的某个无分歧覆盖经标量扩张得到，后者在同构意义下唯一。\footnote{译注：法文原文此推论也未重述上文的完备性条件。此处按原文保留，不据此宣称任意非固有概形在任意特征下均具有这种基本群不变性。}\end{corollary}
\begin{remarks}{}\label{fga2-8-remarks-i}\skipsinglepar
\begin{enumerate}\item{}利用\hyperref[fga2-proposition-4]{命题~\ref*{fga2-proposition-4}}可见，\hyperref[fga2-theorem-11-corollary-4]{推论~\ref*{fga2-theorem-11-corollary-4}}中的假设 \(f_*(\mathcal {O}_X)=\mathcal {O}_Y\) 并非必要。在一般情形，应当在 \(\pi _1(Y)\) 之后用 \(\pi _0(\overline {F})\to \pi _0(X)\to \pi _0(Y)\to  e\) 继续，而不是接上平凡群 \(e\)，与代数拓扑中的情形相同。

    \item{}一般而言，目前还无法说明 \(\pi _1(\overline {F})\to \pi _1(X)\) 的核，尽管它理应涉及某个 \(\pi _2(Y)\)。不过，似乎应能借助下面的\hyperref[fga2-theorem-12]{定理~\ref*{fga2-theorem-12}}，证明当 \(Y\) 为局部环 \(A\) 的谱时，\(\pi _1(\overline {F})\to \pi _1(X)\) 为\emph{单射}；该定理说明，当 \(A\) \emph{完备}时确是如此。\end{enumerate}\end{remarks}
\par{}\hyperref[fga2-theorem-11]{定理~\ref*{fga2-theorem-11}}只用了\hyperref[fga2-theorem-1]{定理~\ref*{fga2-theorem-1}}与\hyperref[fga2-theorem-2]{定理~\ref*{fga2-theorem-2}}；现在我们将应用\hyperref[fga2-theorem-3]{定理~\ref*{fga2-theorem-3}}，并借助下面这个初等引理：
\begin{lemma}{6}\setnumber{6}\label{fga2-lemma-6}\skipsinglepar
\par{}设 \(X\) 为概形，\(X_0\) 为对应的约化概形，即将幂零元全部消去后得到的概形。则 \(X_0\) 的每个无分歧覆盖 \(X'_0\)，都由 \(X\) 的某个无分歧覆盖 \(X'\) 诱导，后者在典范同构意义下确定。\end{lemma}
\par{}这个纯粹局部性质的引理，在这里起的作用，与\hyperref[fga2-theorem-8]{定理~\ref*{fga2-theorem-8}}在模空间理论中的作用类似。把它与存在性定理，即\hyperref[fga2-theorem-3]{定理~\ref*{fga2-theorem-3}}结合，得到：
\begin{theorem}{12}\setnumber{12}\label{fga2-theorem-12}\skipsinglepar
\par{}设 \(A\) 为剩余域为 \(k\) 的完备 Noether 局部环，\(X\) 为 \(A\) 上的固有概形。则 \(X_0=X\otimes _Ak\) 的每个无分歧覆盖 \(X'_0\) 都由 \(X\) 的某个无分歧覆盖 \(X'\) 诱导，后者在同构意义下唯一。\end{theorem}
\par{}换言之：
\begin{corollary}{1}\setnumber{1}\label{fga2-theorem-12-corollary-1}\skipsinglepar
\par{}在 \(X_0\) 中选一个几何点，作为 \(X_0\) 与 \(X\) 的基本群的基点。则典范同态 \(\pi _1(X_0)\to \pi _1(X)\) 是\emph{同构}。\end{corollary}
\par{}现在把\hyperref[fga2-lemma-5]{引理~\ref*{fga2-lemma-5}}应用于 \(X_0\)，为简便起见，假定 \(H^0(X_0,\mathcal {O}_{X_0})=k\)；注意，由于 \(A\) 完备，\(A\) 的无分歧扩张与其剩余域的无分歧扩张对应，即 \(\pi _1(Y)=\pi _1(k)\)，其中 \(Y=\operatorname {Spec}(A)\)。得到正合列：\[     e \to  \pi _1(\overline {X_0}) \to  \pi _1(X) \to  \pi _1(Y) \to  e.   \]
\begin{corollary}{2}\setnumber{2}\label{fga2-theorem-12-corollary-2}\skipsinglepar
\par{}设 \(f\colon  X\to  Y\) 为固有平坦态射，\(y_1\) 为 \(Y\) 的一点，\(y_0\) 为 \(y_1\) 的一个特化。考虑相应的“几何”纤维 \(\overline {X_1}\) 与 \(\overline {X_0}\)，并假定 \(\overline {X_0}\) 可分且连通，这也蕴含 \(\overline {X_1}\) 满足相同条件。则可找到一个群同态 \(\pi _1(X_1)\to \pi _1(X_0)\)，模内自同构而确定。此外，此同态为\emph{满射}。\end{corollary}
\par{}我们或许希望\hyperref[fga2-theorem-12-corollary-2]{推论~\ref*{fga2-theorem-12-corollary-2}}中的同态总是\emph{双射}。可惜，若 \(k(y_0)\) 的特征 \(>0\)，一般并非如此。不过，下面将得到此同态的核的一个上界，至少在 \(\overline {X_0}\) 光滑时如此；由此可知，若 \(k(y_0)\) 的特征为 \(0\)，则上述同态为双射，这也可通过超越方法证明。至少，无论如何，我们已经可以用泛纤维的基本群，给出特殊纤维的 \(\pi _1\) 的上界。例如，利用特征 \(p\) 的代数曲线可以提升为特征 \(0\) 的曲线这一事实（\hyperref[fga2-theorem-9-corollary-4]{定理 9 的推论 4}），通过超越方法便得到：
\begin{corollary}{3}\setnumber{3}\label{fga2-theorem-12-corollary-3}\skipsinglepar
\par{}设 \(X_0\) 为任意特征的代数闭域上一条完备光滑曲线的概形，\(g\) 为 \(X_0\) 的亏格。则 \(\pi _1(X_0)\) 有 \(2g\) 个拓扑生成元，满足熟知的关系。\end{corollary}
\par{}利用熟知的超平面截面方法，由此得到：
\begin{corollary}{4}\setnumber{4}\label{fga2-theorem-12-corollary-4}\skipsinglepar
\par{}设 \(X\) 为任意特征的代数闭域上的光滑射影概形。则 \(\pi _1(X)\) 有有限个拓扑生成元。\end{corollary}
\par{}我们来描述同态 \(\pi _1(\overline {X_1})\to \pi _1(\overline {X_0})\) 的核。为此，可假定 \(Y\) 为完备离散赋值环 \(V=\mathcal {O}_y\) 的谱，其中 \(y=y_0\)。问题等价于：给定 \(\overline {X_1}\) 的一个无分歧覆盖 \(\overline {X'_1}\)，必要时可假定它为 Galois 覆盖，在什么条件下，它来自 \(X_0\) 的一个无分歧覆盖？先验地，给定的覆盖由 \(X_1\otimes _KK'\) 的某个无分歧覆盖 \(X'_1\) 经标量扩张得到，其中 \(K'\) 是 \(V\) 的分式域 \(K\) 的代数闭包 \(\overline {K}\) 中的一个有限子扩张；若 \(\overline {X'_1}\) 为群 \(G\) 的 Galois 覆盖，则也可选 \(X'_1\) 为群 \(G\) 的 Galois 覆盖。因此：\emph{\(X'_1\otimes _{K'}\overline {K}=\overline {X'_1}\) 来自 \(X_0\) 的一个无分歧覆盖，当且仅当在 \(\overline {K}\) 中存在 \(K'\) 的有限扩张 \(K''\)，使 \(X''_1=X'_1\otimes _{K'}K''\) 具有 \(X''\otimes _{V''}K''\) 的形式，其中 \(V''\) 为 \(V\) 在 \(K''\) 中的整闭包，\(X''\) 为 \(X\otimes _V V''\) 的一个无分歧覆盖}。例如，假定 \(X_0\) 绝对正规，则 \(X\otimes _V V''\) 正规，因为它在 \(V''\) 上平坦，且特殊纤维正规；其函数域即为 \(K''(X_1)\)，也就是下列概形的函数域：\[     X_1\otimes _K K''     = (X\otimes _VK)\otimes _KK''     = X\otimes _VK''     = (X\otimes _VV'')\otimes _{V''}K''.   \] 设 \(L''=K''(X'_1)\) 为 \(X'_1\otimes _{K'}K''\) 的函数域，这是 \(K''(X_1)\) 的有限可分扩张，上述条件也意味着\emph{\(L''\) 是 \(X\otimes _VV''\) 的函数域的无分歧扩张}；即 \(X\otimes _V V''\) 在 \(L''\) 中的正规化，在 \(X\otimes _VV''\) 上无分歧。只须验证 \(L''\) 在 \(X\otimes _VV''\) 的特殊纤维各点处无分歧，因为它在泛纤维 \(X_1\otimes _KK''\) 上已无分歧。现在若 \(X_0\) \emph{光滑}，则由 Nagata–Zariski 的“\emph{纯性定理}”，甚至只须验证 \(L''\) 在 \(X\otimes _VV''\) 的特殊纤维的泛点的局部环 \(\mathcal {O}''\) 上无分歧；这是一个离散赋值环，等于 \(X\) 的特殊纤维的泛点的局部环 \(\mathcal {O}\subset  K(X_1)\) 在 \(K''(X_1)\) 中的正规化。于是得到：
\begin{corollary}{5}\setnumber{5}\label{fga2-theorem-12-corollary-5}\skipsinglepar
\par{}在上述条件与记号下，\(\overline {X_1}=X_1\otimes _K\overline {K}\) 的无分歧覆盖 \(X'_1\otimes _{K'}\overline {K}\) 来自 \(\overline {X_0}\) 的无分歧覆盖，当且仅当存在 \(\overline {K}/K'\) 的有限子扩张 \(K''\)，使 \(K''(X'_1)\) 在离散赋值环 \(\mathcal {O}''\subset  K''(X_1)\) 上无分歧。\end{corollary}
\par{}现在注意，\(\mathcal {O}''\) 是离散赋值环 \(\mathcal {O}'\subset  K'(X_1)\) 在 \(K''(X_1)\) 中的正规化，而该环是 \(\mathcal {O}\) 在 \(K'(X_1)\) 中的正规化；此外，\(\mathcal {O}'\) 包含 \(V\) 在 \(K'\) 中的正规化 \(V'\)，\(V'\) 的一致化参数 \(u\) 也是 \(\mathcal {O}'\) 的一致化参数。现在假定 \(X'_1\) 为 Galois 覆盖，其 Galois 群 \(G\) \emph{的阶为 \(n\)，且与 \(k(y_0)\) 的特征 \(p\) 互素}；这里的特征也是 \(\mathcal {O}'\) 的剩余域的特征。则 \(K'(X_1')\) 在 \(\mathcal {O}'\) 上“温和分歧”。由“\emph{Abhyankar 引理}”容易推出，若添入 \(\mathcal {O}'\) 的一个一致化参数的 \(n\) 次方根 \(v\)，那么它在 \(\mathcal {O}'\) 于 \(K'(X_1)(v)\) 中的正规化上便无分歧。可是，\(v\) 可以取为 \(V'\) 的一个一致化参数的 \(n\) 次方根，这表明\hyperref[fga2-theorem-12-corollary-5]{推论~\ref*{fga2-theorem-12-corollary-5}}的条件成立。（利用 Abhyankar 引理与纯性定理的这个想法，是 Serre 告诉我的。）为表述由此所得的结果，对每个全不连通紧群 \(\pi \)，引入商群 \(\overline {\pi }\)：它是 \(\pi \) 模去其各 Sylow \(p\)-子群生成的闭子群所得的商，也就是 \(\pi \) 的所有阶与 \(p\) 互素的离散商群的射影极限。采用此记号，可得：
\begin{theorem}{13}\setnumber{13}\label{fga2-theorem-13}\skipsinglepar
\par{}设 \(f\colon  X\to  Y\) 为固有平坦态射，\(y_1\) 为 \(Y\) 的一点，\(y_0\) 为 \(y_1\) 的特化。假定 \(\overline {X_0}\) 连通且光滑。则由\hyperref[fga2-theorem-12-corollary-2]{定理 12 的推论 2}中的满同态诱导的同态 \(\overline {\pi _1}(\overline {X_1})\to \overline {\pi _1}(\overline {X_0})\) 是\emph{同构}。\end{theorem}
\par{}换言之：
\begin{corollary}{1}\setnumber{1}\label{fga2-theorem-13-corollary-1}\skipsinglepar
\par{}Galois 群的阶与 \(k(y_0)\) 的特征 \(p\) 互素的无分歧 Galois 覆盖，在 \(\overline {X_0}\) 与 \(\overline {X_1}\) 上具有相同的分类。\end{corollary}
\par{}特别地，若 \(k(y_0)\) 的特征为 \(0\)，则可用代数方法看出，\(\pi _1(\overline {X_1})\to \pi _1(\overline {X_0})\) 是双射。
\par{}最后指出，以上方法还给出下面这个比\hyperref[fga2-theorem-13]{定理~\ref*{fga2-theorem-13}}更一般的结果：
\begin{theorem}{14}\setnumber{14}\label{fga2-theorem-14}\skipsinglepar
\par{}设 \(f\colon  X\to  Y\) 为固有光滑态射，\(D\) 为 \(X\) 的闭子概形，它在 \(Y\) 上光滑，且在每一点处的余维均为 \(1\)。给定 \(f\) 的一条纤维 \(Z=f^{-1}(z)\)，置 \(Z'=Z\setminus  Z\cap  D\)，并以 \(\pi _1^\mathrm {t}(\overline {Z'})\) 表示基本群 \(\pi _1(Z')\) 的商群，它分类 \(\overline {Z'}\) 的那些在 \(\overline {Z\cap  D}\) 处“温和分歧”的无分歧覆盖。设 \(y_0\)、\(y_1\) 如\hyperref[fga2-theorem-13]{定理~\ref*{fga2-theorem-13}}中所述。则存在\emph{满}同态 \(\pi _1^\mathrm {t}(\overline {X'_1})\to \pi _1^\mathrm {t}(\overline {X'_0})\)，模内自同构而确定，且相应的同态 \(\overline {\pi _1^\mathrm {t}}(\overline {X'_1})\to \overline {\pi _1^\mathrm {t}}(\overline {X'_0})\) 是同构。\end{theorem}
\par{}由此得到\hyperref[fga2-theorem-13]{定理~\ref*{fga2-theorem-13}}的各推论及\hyperref[fga2-theorem-12-corollary-4]{定理 12 的推论 4}的相应变体。同样，利用\hyperref[fga2-theorem-9-corollary-3]{定理 9 的推论 3}，通过超越方法得到：
\begin{corollary}{1}\setnumber{1}\label{fga2-theorem-14-corollary-1}\skipsinglepar
\par{}设 \(X_0\) 为任意特征的代数闭域上一条完备光滑曲线的概形，\(S=(s_i)_{1\leqslant  i\leqslant  n}\) 为 \(X_0\) 中含有 \(n\) 个元素的有限子集。则 \(\pi _1^\mathrm {t}(X_0\setminus  S)\) 有 \(2g+n\) 个拓扑生成元 \(x_i,y_i,\sigma _j\)，其中 \(1\leqslant  i\leqslant  g\) 且 \(1\leqslant  j\leqslant  n\)，满足关系 \[       \left (\prod _i x_iy_ix_i^{-1}y_i^{-1}\right )\sigma _1\ldots \sigma _n = 1,     \]，这里的 \(\sigma _j\) 是各点 \(s_j\) 对应的惯性群的生成元。对任意\emph{阶与特征互素}的有限群 \(G\)，若它由满足上述关系的元素 \(\overline {x_i},\overline {y_i},\overline {\sigma _j}\) 生成，则存在 \(X_0\setminus  S\) 的一个群为 \(G\) 的无分歧 Galois 覆盖，其在各点 \(s_j\) 的惯性群由 \(\overline {\sigma _j}\) 生成。\end{corollary}
\par{}若 \(X_0\) 的亏格为 \(0\) 且 \(n=3\)，就得到了“\emph{三点问题}”的解，至少对次数与特征互素的 Galois 覆盖如此。（这里，\hyperref[fga2-theorem-9]{定理~\ref*{fga2-theorem-9}}实际上用不上；看来可从三点问题所涉及的特例推出上述推论。）
\begin{remarks}{}\label{fga2-8-remarks-ii}\skipsinglepar
\begin{enumerate}\item{}更完整的研究，很可能涉及 \(X\)、\(X_0\)、\(X_1\) 的广义 Galois 覆盖，其 Galois 群最终可能是无穷小的，应能使我们确定\hyperref[fga2-theorem-12-corollary-2]{定理 12 的推论 2}中的核。不过，对允许非“温和”分歧的覆盖作研究，似乎困难得多。

    \item{}把\hyperref[fga2-lemma-6]{引理~\ref*{fga2-lemma-6}}与 Grauert 关于非奇异射影概形沿超平面截面的形式完备化的结果结合，或者与\hyperref[fga2-theorem-11]{定理~\ref*{fga2-theorem-11}}后注 2 中提到而尚未证明的定理结合，也能在“抽象”代数几何中证明关于基本群的经典\emph{Lefschetz 定理}。\end{enumerate}\end{remarks}

\section*{原报告参考文献}
% Restored from the original French report, printed 182-28, physical PDF page 67.
% Numbered bibliography is distinct from the coordinator-owned consolidated bibliography.
\begin{enumerate}
\item[{[1]}]\setnumber{1}\label{fga2-bib-1} DIEUDONNÉ (J.) 与 GROTHENDIECK (A.). \emph{Eléments de géométrie algébrique}. Publications mathématiques de l'Institut des Hautes Etudes Scientifiques（待出版）。
\item[{[2]}]\setnumber{2}\label{fga2-bib-2} GROTHENDIECK (Alexander). \emph{The cohomology theory of abstract algebraic varieties}. International Congress of Mathematicians [1958. Edinburgh]（待出版）。
\item[{[3]}]\setnumber{3}\label{fga2-bib-3} SERRE (Jean-Pierre). \emph{Faisceaux algébriques cohérents}. Annals of Math., 第 2 辑，第 61 卷，1955，197--278 页。
\item[{[4]}]\setnumber{4}\label{fga2-bib-4} SERRE (Jean-Pierre). \emph{Géométrie algébrique et géométrie analytique}. Ann. Institut Fourier Grenoble, 第 6 卷，1955--1956，1--42 页。
\end{enumerate}

% END chapters/fga2.tex

\clearpage
\renewcommand\fgaAnchorPrefix{fga3}
\part{下降技术与代数几何中的存在定理}
\setcounter{part}{0}
\renewcommand\thepart{3.\Roman{part}}
\setcounter{chapter}{0}
\renewcommand\thechapter{\Alph{chapter}}
\renewcommand\thesubsection{\thesection.\alph{subsection}}
\renewcommand\fgaAnchorPrefix{fga3.i}

% BEGIN chapters/fga3.i.tex
\part{一般理论与沿忠实平坦态射的下降}\label{fga3.i}\thispagestyle{fancy}\begin{original}{}\label{fga3.i-original-citation}\skipsinglepar\par{}A. Grothendieck.
    "Technique de descente et théorèmes d'existence en géométrie algébrique, I: Généralités. Descente par morphismes fidèlement plats".
    \emph{Séminaire Bourbaki} \textbf{12} (1959–60), 第 190 次报告，1959 年 12 月。
    \url{http://www.numdam.org/book-part/SB_1958-1960__5__299_0/}\end{original}\par{}\emph{［编者补注］} 关于下降理论的各种细节，另见［\hyperref[Gro1960b]{Gro1960b}，VI、VII 和 VIII］。\par{}从技术角度看，本文及后续文章可以看作 Hilbert 著名的“定理 90”的变式。将下降方法引入代数几何似乎应归功于 A. Weil，当时称为“基域的下降”。Weil 只考虑有限可分域扩张。随后 P. Cartier 研究了高度为 1 的根式扩张情形。由于缺乏概形语言，尤其是所考虑的环中缺乏幂零元，Cartier 未能表述这两种情形的本质一致性。\par{}目前看来，以下将解释的一般下降技术，必要时结合“形式几何”的基本定理（参见 \hyperref[fga2]{[3]}），是代数几何中大多数存在定理的基础。（\emph{［编者补注］ 如今看来，说下降技术是“代数几何中大多数存在定理的基础”有些过强。对本系列，即“代数几何中的下降技术与存在定理”，前两讲所研究的非射影技术，这在很大程度上成立；但对射影技术，即第 \hyperref[fga3.iv]{IV}、\hyperref[fga3.v]{V} 和 \hyperref[fga3.vi]{VI} 讲，则并非如此。}）还应指出，上述下降技术当然能移植到“解析几何”中，我们可以希望在不太遥远的将来，专家们能够证明第 \hyperref[fga3.ii]{II} 讲将给出的形式几何存在定理的“解析”对应物。无论如何，Kodaira–Spencer 最近的工作已相当清楚地说明，需要更接近概形理论的方法，自然应与超越技术互补，因为他们的方法似乎不适于在“模簇”的奇点邻域内定义并研究这些模簇。\par{}本讲，即第 I 讲，将讨论下降的最初等情形，也就是标题所指的情形。不过，下文 \cref{fga3.i-b.1} 中的 \hyperref[fga3.i-b.1-theorem-1]{定理~\ref*{fga3.i-b.1-theorem-1}}、\hyperref[fga3.i-b.1-theorem-2]{定理~\ref*{fga3.i-b.1-theorem-2}} 与 \hyperref[fga3.i-b.1-theorem-3]{定理~\ref*{fga3.i-b.1-theorem-3}} 的应用已非常多。这里只举其中若干例子，不追求可能的最大一般性。\par{}我们将自由使用概形语言，读者可参阅前引文章及 \hyperref[GR1958]{[2]}。不过须明确指出，本文考虑的预概形不必为 Noether 预概形，态射也不必有限型。例如，若 \(A\) 为完备化为 \(\overline {A}\) 的 Noether 局部环，则需要考虑非 Noether 环 \(\overline {\overline {A}}\otimes _A\overline {A}\)，以及对应于相关环之间的包含关系的仿射概形态射。\chapter{范畴论预备知识}\thispagestyle{fancy}\label{fga3.i-a}\section{纤维范畴、下降资料与 \(\mathcal {F}\) 下降态射}\label{fga3.i-a.1}\subsection{}\label{fga3.i-a.1.a}\begin{definition}{1.1}\setnumber{1.1}\label{fga3.i-a.1-definition-1.1}\skipsinglepar\par{}一个 \emph{以 \(\mathcal {C}\) 为基底的纤维范畴 \(\mathcal {F}\)}（或称 \emph{在 \(\mathcal {C}\) 上} 的纤维范畴）由以下资料组成：\begin{itemize}\item{}一个范畴 \(\mathcal {C}\)；

      \item{}对每个 \(X\in \mathcal {C}\)，一个范畴 \(\mathcal {F}_X\)；

      \item{}对每个 \(\mathcal {C}\) 态射 \(f\colon  X\to  Y\)，一个函子 \(f^*\colon \mathcal {F}_Y\to \mathcal {F}_X\)；对 \(\xi \in \mathcal {F}_Y\)，也将其写作 \[           f^*(\xi ) = \xi  \times _Y X,\]这里把 \(X\) 视为“\(Y\) 上的 \(\mathcal {C}\) 的对象”，即赋有态射 \(f\) 的对象；

      \item{}对任意两个可合成态射 \(X\xrightarrow {f}Y\xrightarrow {g}Z\)，一个函子同构 \[           c_{f,g}\colon  (gf)^* \to  f^*g^*         \]；\end{itemize} 这些资料须满足下列条件：\begin{enumerate}
      
 \item[i.]
        \(\operatorname {id}^*=\operatorname {id}\)
      


      
 \item[ii.]若 \(f\) 或 \(g\) 为恒等同构，则 \(c_{f,g}\) 为恒等同构；


      
 \item[iii.]对任意三个可合成态射 \(X\xrightarrow {f}Y\xrightarrow {g}Z\xrightarrow {h}T\)，利用形如 \(c_{u,v}\) 的同构得到的下图交换：\[           \begin {CD}             (h(gf))^* @= ((hg)f)^*           \\@VVV @VVV           \\(gf)^*h^* @. f^*(hg)^*           \\@VVV @VVV           \\(f^*g^*)h^* @= f^*(g^*h^*)           \end {CD}         \]

    \end{enumerate}\end{definition}\begin{example}{1}\setnumber{1}\label{fga3.i-a.1-example-1}\skipsinglepar\par{}设 \(\mathcal {C}\) 为一切纤维积均存在的范畴。可定义以 \(\mathcal {C}\) 为基底的纤维范畴 \(\mathcal {F}\) 如下：取 \(\mathcal {F}_X\) 为 \(X\) 上的 \(\mathcal {C}\) 的对象范畴，对应于态射 \(f\colon  X\to  Y\) 的函子 \(f^*\colon \mathcal {F}_Y\to \mathcal {F}_X\) 由 \emph{纤维积} \(Z\mapsto  Z\times _Y X\) 定义。\end{example}\begin{example}{2}\setnumber{2}\label{fga3.i-a.1-example-2}\skipsinglepar\par{}设 \(\mathcal {C}\) 为预概形范畴，对 \(X\in \mathcal {C}\)，令 \(\mathcal {F}_X\) 为 \(X\) 上的拟凝聚模层范畴。若 \(f\colon  X\to  Y\) 为预概形态射，则 \(f^*\colon \mathcal {F}_Y\to \mathcal {F}_X\) 为 \emph{模层逆像} 函子。这样便得到 \(\mathcal {C}\) 上的一个纤维范畴。\end{example}\subsection{}\label{fga3.i-a.1.b}\begin{definition}{1.2}\setnumber{1.2}\label{fga3.i-a.1-definition-1.2}\skipsinglepar\par{}若 \(u\) 给出从 \(E\) 到 \(E'\) 中满足 \(v_1(x')=v_2(x')\) 的元素 \(x'\in  E'\) 所成子集的双射，则称集合映射图式 \[       E \xrightarrow {u}       E' \overset {v_1}{\underset {v_2}{\rightrightarrows }}       E''     \] 是 \emph{正合} 的。\end{definition}\begin{definition}{1.3}\setnumber{1.3}\label{fga3.i-a.1-definition-1.3}\skipsinglepar\par{}设 \(\mathcal {F}\) 为以 \(\mathcal {C}\) 为基底的纤维范畴，并考虑 \(\mathcal {C}\) 中满足 \(\alpha \beta _1=\alpha \beta _2\) 的态射图式 \[       S \xleftarrow {\alpha }       S' \overset {\beta _1}{\underset {\beta _2}{\leftleftarrows }}       S''     \]。若对 \(\mathcal {F}_S\) 中的每一对元素 \((\xi ,\eta )\)，集合图式 \begin{eqn}{+}\setnumber{+}\label{fga3.i-a.1-definition-1.3-equation}\skipsinglepar\[         \operatorname {Hom}(\xi ,\eta ) \xrightarrow {\alpha ^*}         \operatorname {Hom}(\alpha ^*(\xi ),\alpha ^*(\eta )) \overset {\beta _1^*}{\underset {\beta _2^*}{\rightrightarrows }}         \operatorname {Hom}(\gamma ^*(\xi ),\gamma ^*(\eta ))       \tag{+}       \]\end{eqn} 都正合，其中 \(\gamma =\alpha \beta _1=\alpha \beta _2\)，则称此图式为 \emph{\(\mathcal {F}\) 正合} 的。\par{}为简便起见，上图利用 \(c_{\beta _i,\alpha }\) 将 \(\beta _i^*\alpha ^*\) 与 \((\alpha \beta _i)^*=\gamma ^*\) 等同。\end{definition}\begin{definition}{1.4}\setnumber{1.4}\label{fga3.i-a.1-definition-1.4}\skipsinglepar\par{}设 \(\mathcal {F}\) 为以 \(\mathcal {C}\) 为基底的纤维范畴，考虑 \(\mathcal {C}\) 中的态射 \(\beta _1,\beta _2\colon  S''\to  S'\)。设 \(\xi '\in \mathcal {F}_{S'}\)。\(\xi '\) 上关于二元组 \((\beta _1,\beta _2)\) 的一个 \emph{粘合资料}，定义为从 \(\beta _1^*(\xi ')\) 到 \(\beta _2^*(\xi ')\) 的一个同构。若 \(\xi ',\eta '\in \mathcal {F}_{S'}\) 均赋有粘合资料，则当下图交换时，称 \(\mathcal {F}_{S'}\) 中的态射 \(u\colon \xi '\to \eta '\) \emph{与粘合资料相容}：\[       \begin {CD}         \beta _1^*(\xi ') @>>> \beta _2^*(\xi ')       \\@VVV @VVV       \\\beta _1^*(\eta ') @>>> \beta _2^*(\eta ').       \end {CD}     \]\end{definition}\par{}由此定义，\(\mathcal {F}_{S'}\) 中赋有关于 \(\beta _1\) 和 \(\beta _2\) 的粘合资料的对象构成一个 \emph{范畴}。若 \(\alpha \colon  S'\to  S\) 是满足 \(\alpha \beta _1=\alpha \beta _2\) 的态射，则对每个 \(\xi \in \mathcal {F}_{S'}\)，\(\mathcal {F}_{S'}\) 的对象 \(\xi '=\alpha ^*(\xi )\) 赋有一个典范粘合资料，因为 \[     \beta _i^*\alpha ^*(\xi )     \simeq  (\alpha \beta _i)^*(\xi )     = \gamma ^*(\xi ),   \]，这里仍置 \(\gamma =\alpha \beta _1=\alpha \beta _2\)；此外，若 \(u\colon \xi \to \eta \) 为 \(\mathcal {F}_S\) 中的态射，则 \[     \alpha ^*(u)\colon  \alpha ^*(\xi ) \to  \alpha ^*(\eta )   \] 是 \(\mathcal {F}_{S'}\) 中与典范粘合资料相容的态射。于是得到从范畴 \(\mathcal {F}_S\) 到 \(\mathcal {F}_{S'}\) 中赋有关于二元组 \((\beta _1,\beta _2)\) 的粘合资料的对象范畴的 \emph{典范函子}。因此，\hyperref[fga3.i-a.1-definition-1.3]{定义~\ref*{fga3.i-a.1-definition-1.3}} 也可重新表述为：当上述函子为 \emph{全忠实} 时，称 \hyperref[fga3.i-a.1-definition-1.3-equation]{式~\ref*{fga3.i-a.1-definition-1.3-equation}} 为 \emph{\(\mathcal {F}\) 正合} 的；即上述函子给出范畴 \(\mathcal {F}_S\) 与 \(\mathcal {F}_{S'}\) 中赋有关于 \((\beta _1,\beta _2)\) 的粘合资料的对象范畴的一个子范畴之间的等价。\begin{definition}{1.5}\setnumber{1.5}\label{fga3.i-a.1-definition-1.5}\skipsinglepar\par{}若赋有此资料的 \(\xi '\) 同构于某个 \(\xi \in \mathcal {F}_S\) 所对应的 \(\alpha ^*(\xi )\)，则称 \(\xi '\in \mathcal {F}_{S'}\) 上的粘合资料关于 \(\alpha \) 是 \emph{有效} 的。\end{definition}\par{}若 \hyperref[fga3.i-a.1-definition-1.3-equation]{式~\ref*{fga3.i-a.1-definition-1.3-equation}} 为 \(\mathcal {F}\) 正合，则 \hyperref[fga3.i-a.1-definition-1.5]{定义~\ref*{fga3.i-a.1-definition-1.5}} 中的对象 \(\xi \) 在唯一同构的意义下确定，并且 \emph{范畴 \(\mathcal {F}_S\) 等价于 \(\mathcal {F}_{S'}\) 中赋有有效粘合资料的对象范畴}。\subsection{}\label{fga3.i-a.1.c}\par{}最重要的情形是 \[     S'' = S' \times _S S',   \]，其中各 \(\beta _i\) 是从 \(S'\times _S S'\) 到其两个因子的投影 \(p_1\) 和 \(p_2\)；这里假设 \(\mathcal {C}\) 中一切纤维积存在。某个 \(\xi '\in \mathcal {F}_{S'}\) 上的粘合资料 \(\varphi \colon  p_1^*(\xi ')\to  p_2^*(\xi ')\) 若有效，立即得到两个必要条件：\begin{enumerate}
    
 \item[i.]\(\Delta ^*(\varphi ) = \operatorname {id}_{\xi '} \)，其中 \(\Delta \colon  S'\to  S'\times _S S'\) 表示对角态射，并将 \(\Delta ^* p_i^*(\xi ')\) 与 \((p_i\Delta )^*(\xi ')=\xi '\) 等同；


    
 \item[ii.]\(p_{31}^*(\varphi ) = p_{32}^*(\varphi )p_{21}^*(\varphi )\)，其中 \(p_{ij}\) 表示从 \(S'\times _S S'\times _S S'\) 到其第 \(i\) 和第 \(j\) 个因子的部分乘积的典范投影。

  \end{enumerate}\begin{definition}{1.6}\setnumber{1.6}\label{fga3.i-a.1-definition-1.6}\skipsinglepar\par{}\(\xi '\in \mathcal {F}_{S'}\) 上关于态射 \(\alpha \colon  S'\to  S\) 的 \emph{下降资料}，定义为 \(\xi '\) 上关于两个典范投影 \(S'\times _S S'\to  S'\) 所成的二元组 \((p_1,p_2)\) 的粘合资料，并要求满足上述条件（i）和（ii）。\end{definition}\begin{definition}{1.7}\setnumber{1.7}\label{fga3.i-a.1-definition-1.7}\skipsinglepar\par{}若态射图式 \[       S \xleftarrow {\alpha }       S' \overset {p_1}{\underset {p_2}{\leftleftarrows }}       S'\times _S S'     \] 为 \(\mathcal {F}\) 正合（\hyperref[fga3.i-a.1-definition-1.3]{定义~\ref*{fga3.i-a.1-definition-1.3}}），则称态射 \(\alpha \colon  S'\to  S\) 为 \emph{\(\mathcal {F}\) 下降态射}；若此外 \(\mathcal {F}_{S'}\) 的任意对象上的每个下降资料（\hyperref[fga3.i-a.1-definition-1.6]{定义~\ref*{fga3.i-a.1-definition-1.6}}）均有效，则称 \(\alpha \) 为 \emph{严格 \(\mathcal {F}\) 下降态射}。后一个严格性条件也可更直观地表述为：“给出 \(\mathcal {F}_S\) 的一个对象，等价于给出 \(\mathcal {F}_{S'}\) 的一个赋有下降资料的对象。”\end{definition}\par{}注意，若 \(\mathcal {F}\) 下降态射（\emph{［编者补注］} 这里不必假设 \(\alpha \) 是 \(\mathcal {F}\) 下降态射）\(\alpha \colon  S'\to  S\) 有一个 \emph{截面} \(s\colon  S\to  S'\)，即满足 \(\alpha  s=\operatorname {id}_S\) 的态射 \(s\)，则它是严格 \(\mathcal {F}\) 下降态射：若 \(\xi '\in \mathcal {F}_{S'}\) 赋有关于 \(\alpha \) 的下降资料，那么“它来自” \(\xi =s^*(\xi ')\)。\subsection{}\label{fga3.i-a.1.d}\par{}可以更直观地表述上述概念：对 \(S\) 上的 \(\mathcal {C}\) 的对象 \(T\)，引入集合 \[     \operatorname {Hom}_S(T,S') = S'(T),   \]，其元素记为 \(t\)、\(t'\) 等。给定对象 \(\xi '\in \mathcal {F}_{S'}\)，每个 \(t\in  S'(T)\) 都对应于 \(\mathcal {F}_T\) 中的对象 \(t^*(\xi ')\)，亦记为 \(\xi '_t\)。于是，\(\xi '\) 上关于 \((p_1,p_2)\) 的粘合资料由如下资料定义：对 \(S\) 上的每个 \(T\) 及每一对点 \(t,t'\in  S'(T)\)，给出同构 \[     \varphi _{t',t}\colon  \xi '_t \to  \xi '_{t'}   \]，并满足关于 \(T\) 的显然的函子性条件。此时 \cref{fga3.i-a.1.c} 的条件（i）和（ii）可写为 \begin{enumerate}
    
 \item[i bis.]对所有 \(T\) 及所有 \(t\in  S'(T)\)，均有 \(\varphi _{t,t}=\operatorname {id}_{\xi '_t}\)；


    
 \item[ii bis.]对所有 \(T\) 及所有 \(t,t',t''\in  S'(T)\)，均有 \(\varphi _{t,t''}=\varphi _{t,t'}\varphi _{t',t''}\)。

  \end{enumerate}。令 \(t=t'=t''\)，可见（ii bis）蕴含 \(\varphi _{t,t}^2=\varphi _{t,t}\)；所以，由于假设 \(\varphi _{t,t}\) 是同构，它蕴含（i bis）。因而（i bis）是（ii bis）的推论，（i）也是（ii）的推论。但若不预先假设各个 \(\varphi _{t,t'}\) 都是同构，即不假设 \(\varphi \colon  p_1^*(\xi ')\to  p_2^*(\xi ')\) 是同构，则（ii bis）未必再蕴含（i bis）；不过，（ii bis）与（i bis）合起来确实蕴含各个 \(\varphi _{t,t'}\) 都是同构，因为此时有 \(\varphi _{t,t'}\varphi _{t',t}=\varphi _{t,t}=\operatorname {id}_{\xi '_t}\)。\section{正合图式、严格满态射、下降态射及例子}\label{fga3.i-a.2}\subsection{}\label{fga3.i-a.2.a}\begin{definition}{2.1}\setnumber{2.1}\label{fga3.i-a.2-definition-2.1}\skipsinglepar\par{}设 \(\mathcal {C}\) 为一个范畴。若对每个 \(Z\in \mathcal {C}\)，相应的集合图式 \[       \operatorname {Hom}(Z,T) \to        \operatorname {Hom}(Z,T') \rightrightarrows        \operatorname {Hom}(Z,T'')     \] 都正合（\hyperref[fga3.i-a.1-definition-1.2]{定义~\ref*{fga3.i-a.1-definition-1.2}}），则称态射图式 \[       T \xrightarrow {\alpha }       T' \overset {\beta _1}{\underset {\beta _2}{\rightrightarrows }}       T''     \] \emph{正合}。此时称 \((T,\alpha )\)，或滥用语言称 \(T\)，为态射对 \((\beta _1,\beta _2)\) 的一个 \emph{核}。\end{definition}\par{}显然，此核在唯一同构的意义下确定。若 \(\mathcal {C}\) 为集合范畴，则上述定义与 \hyperref[fga3.i-a.1-definition-1.2]{定义~\ref*{fga3.i-a.1-definition-1.2}} 相容。对偶地，可以定义 \(\mathcal {C}\) 中态射图式 \[     S \xleftarrow {\alpha }     S' \overset {\beta _1}{\underset {\beta _2}{\leftleftarrows }}     S''   \] 的正合性；此时称 \((S,\alpha )\) 为态射对 \((\beta _1,\beta _2)\) 的一个 \emph{余核}。\begin{definition}{2.2}\setnumber{2.2}\label{fga3.i-a.2-definition-2.2}\skipsinglepar\par{}若态射 \(\alpha \colon  S'\to  S\) 是满态射，并且对每个态射 \(u\colon  S'\to  Z\)，下列必要条件也是 \(u\) 可分解为 \(S'\to  S\to  Z\) 的充分条件，则称其为 \emph{严格满态射}：对每个 \(S''\in \mathcal {C}\) 及每对满足 \(\alpha \beta _1=\alpha \beta _2\) 的态射 \(\beta _1,\beta _2\colon  S''\to  S'\)，均有 \(u\beta _1=u\beta _2\)。\end{definition}\par{}若纤维积 \(S'\times _S S'\) 存在，则这等价于图式 \[     S \xleftarrow {\alpha }     S' \overset {p_1}{\underset {p_2}{\leftleftarrows }}     S'\times _S S'   \] 正合，即 \(S\) 为二元组 \((p_1,p_2)\) 的余核。无论如何，余核态射都是严格满态射。还要指出，若严格满态射也是单态射，则它是同构。将对偶概念 \emph{严格单态射} 的展开留给读者。\par{}为了更精确地说明 \(\mathcal {F}\) 下降态射与严格满态射之间的关系，引入如下定义：\begin{definition}{2.3}\setnumber{2.3}\label{fga3.i-a.2-definition-2.3}\skipsinglepar\par{}若对 \(S\) 上的每个 \(T\)，纤维积 \(T'=S'\times _S T\) 存在且投影 \(T'\to  T\) 是满态射，相应地是严格满态射，则称态射 \(\alpha \colon  S'\to  S\) 为 \emph{普遍满态射}，相应地称为 \emph{普遍严格满态射}。\end{definition}\par{}在性质很好的范畴中，例如集合范畴、拓扑空间上的层范畴、abelian 范畴等等，上述四种“满性”概念都一致；但在预概形范畴，或给定非空预概形 \(S\) 上的预概形范畴中，它们各不相同，即使只考虑在 \(S\) 上有限的 \(S\) 概形也是如此。\begin{definition}{2.4}\setnumber{2.4}\label{fga3.i-a.2-definition-2.4}\skipsinglepar\par{}若态射 \(\alpha \colon  S'\to  S\) 是 \(\mathcal {F}\) 下降态射，相应地是严格 \(\mathcal {F}\) 下降态射（参见 \hyperref[fga3.i-a.1-definition-1.7]{定义~\ref*{fga3.i-a.1-definition-1.7}}），则称其为 \emph{下降态射}，相应地称为 \emph{严格下降态射}；其中 \(\mathcal {F}\) 是 \(\mathcal {C}\) 上的纤维范畴，其纤维由 \(\mathcal {C}\) 的对象上的 \(\mathcal {C}\) 的对象组成（参见 \hyperref[fga3.i-a.1-example-1]{例~\ref*{fga3.i-a.1-example-1}}）。\end{definition}\begin{proposition}{2.1}\setnumber{2.1}\label{fga3.i-a.2-proposition-2.1}\skipsinglepar\par{}若 \(\mathcal {C}\) 中一切有限乘积和有限纤维积均存在，则 \(\mathcal {C}\) 中的下降态射与 \(\mathcal {C}\) 中的普遍严格满态射一致。\end{proposition}\subsection{}\label{fga3.i-a.2.b}\begin{example}{}\skipsinglepar\par{}设 \(\mathcal {C}\) 为预概形范畴，\(S\in \mathcal {C}\)，且 \(S'\)、\(S''\) 是 \(S\) 上的 \emph{有限} 预概形；即它们对应于 \(S\) 上的代数层 \(\mathscr {A}'\)、\(\mathscr {A}''\)，作为模层拟凝聚且有限型，若 \(S\) 局部 Noether 则为凝聚层。设 \(\alpha \colon  S'\to  S\) 为 \(S'\) 的结构态射，\(\beta _1\) 和 \(\beta _2\) 为从 \(S''\) 到 \(S'\) 的 \(S\) 态射，由代数同态 \(\mathscr {A}'\to \mathscr {A}''\) 定义，后者也记为 \(\beta _1\) 和 \(\beta _2\)。利用有限态射是闭映射这一事实，即第一 Cohen–Seidenberg 定理，容易证明 \(\mathcal {C}\) 中的图式 \begin{eqn}{+}\setnumber{+}\label{fga3.i-a.2.b-equation}\skipsinglepar\[         S \xleftarrow {\alpha }         S' \overset {\beta _1}{\underset {\beta _2}{\leftleftarrows }}         S''       \tag{+}       \]\end{eqn} 正合，当且仅当 \(S\) 上的层图式 \[       \mathscr {O}_S = \mathscr {A} \xrightarrow {\alpha }       \mathscr {A}' \overset {\beta _1}{\underset {\beta _2}{\rightrightarrows }}       \mathscr {A}''     \] 正合。特别地，若 \(\alpha \colon  S'\to  S\) 为对应于 \(S\) 上代数层 \(\mathscr {A}'\) 的有限态射，则 \(\alpha \) 是严格满态射，当且仅当层图式 \[       \mathscr {O}_S = \mathscr {A} \to        \mathscr {A}' \overset {p_1}{\underset {p_2}{\rightrightarrows }}       \mathscr {A}'\otimes _{\mathscr {A}}\mathscr {A}'     \] 正合；它是满态射，当且仅当 \(\mathscr {A}\to \mathscr {A}'\) 为单射。若 \(S\) 是环为 \(A\) 的仿射概形，则 \(S'\) 是环为 \(A'\) 的仿射概形，\(A'\) 在 \(A\) 上有限；所以 \(S'\to  S\) 是严格满态射，当且仅当 \(A\to  A'\) 给出从 \(A\) 到 \(A'\) 中满足 \[       1_{A'}\otimes _A x' - x'\otimes _A 1_{A'}       = 0     \] 的元素 \(x'\in  A'\) 所成子环的同构；它是满态射，当且仅当 \(A\to  A'\) 为单射。如前所述，即使 \(S\) 是 Artin 局部环的概形，一个为满态射的有限态射 \(S'\to  S\) 也未必为严格满态射。不过，可以证明，\emph{若 \(S\) 为 Noether 预概形，则每个为满态射的有限态射 \(S'\to  S\) 都是有限个严格满态射的复合}，其中各态射也有限。这也表明，两个严格满态射的复合未必是严格满态射。\end{example}\subsection{}\label{fga3.i-a.2.c}\par{}若 \hyperref[fga3.i-a.2.b-equation]{式~\ref*{fga3.i-a.2.b-equation}} 是有限态射的正合图式，则对每个预概形的 \emph{平坦} 态射 \(T\to  S\)，由基底变换 \(T\to  S\) 从 \hyperref[fga3.i-a.2.b-equation]{式~\ref*{fga3.i-a.2.b-equation}} 诱导的图式仍正合。因此，若 \(X\)、\(Y\) 是 \(S\) 预概形，且 \(X\) 在 \(S\) 上 \emph{平坦}，则如下映射图式正合：\[     \operatorname {Hom}_S(X,Y) \to      \operatorname {Hom}_{S'}(X',Y') \rightrightarrows      \operatorname {Hom}_{S''}(X'',Y'').   \] 其中 \(X'\)、\(Y'\) 是 \(X\)、\(Y\) 在 \(S'\) 上的逆像，而 \(X''\)、\(Y''\) 是它们在 \(S''\) 上的逆像。特别地，若 \(\mathcal {F}\) 表示预概形范畴 \(\mathcal {C}\) 上的纤维范畴，且对 \(X\in \mathcal {C}\)，\(\mathcal {F}_X\) 为 \emph{平坦} \(X\) 预概形的范畴，则图式 \hyperref[fga3.i-a.2.b-equation]{式~\ref*{fga3.i-a.2.b-equation}} 是 \(\mathcal {F}\) 正合的。若不加平坦性假设，此结论不成立；特别地，有限严格满态射未必为下降态射。类似地，若 \(\mathcal {F}\) 表示使 \(\mathcal {F}_X\) 为预概形 \(X\) 上 \emph{平坦} 拟凝聚层范畴的纤维范畴，则 \hyperref[fga3.i-a.2.b-equation]{式~\ref*{fga3.i-a.2.b-equation}} 是 \(\mathcal {F}\) 正合的；这里平坦性假设同样是必要的。在这两种情形下，\(S'\) 上的平坦对象的粘合资料是否 \emph{有效}，特别是在 \(S''=S'\times _S S'\) 时下降资料是否有效，都是精细的问题；本系列文章的主要目标之一就是在许多特殊情形下回答它。报告人不知道：对每个有限严格满态射 \(S'\to  S\)，\(S'\) 上平坦拟凝聚层的每个下降资料是否都有效；即使假设 \(S\) 是 Artin 局部环的谱，并只考虑秩为 \(1\) 的局部自由层，也不知道。更一般地，设 \(A\) 为环，\(A'\) 为 \(A\) 代数，均交换，使图式 \[     A \to      A' \rightrightarrows      A'\otimes _A A'   \] 正合；这等价于说，相应的 \(A'\) 与 \(A\) 的谱之间的态射 \(S'\to  S\) 是 \(\mathcal {F}\) 下降态射，其中 \(\mathcal {F}\) 为平坦拟凝聚层的纤维范畴。设 \(M'\) 为赋有向 \(A\) 的下降资料的平坦 \(A'\) 模，即赋有 \((A'\otimes _A A')\) 模同构 \[     \varphi \colon  M'\otimes _A A'     \xrightarrow {\sim } A'\otimes _A M'   \]，且满足 \cref{fga3.i-a.1.c} 的条件（i）和（ii）；将它们写成模的形式留给读者。此资料相对于平坦拟凝聚层的纤维范畴是否有效？设 \(M\) 为 \(M'\) 中满足 \[     \varphi (x'\otimes _A 1_{A'})     = 1_{A'}\otimes _A x',   \] 的元素 \(x'\in  M'\) 所成子集，它是 \(M'\) 的 \(A\) 子模。典范单射 \(M\to  M'\) 定义 \(A'\) 模同态 \(M\otimes _A A'\to  M'\)。\emph{于是 \(\varphi \) 的有效性等价于如下事实：\(M\) 为平坦 \(A\) 模，且上述同态为同构。}\begin{remark}{}\label{fga3.i-a.2.c-remark}\skipsinglepar\par{}上述讨论并未对图式 \hyperref[fga3.i-a.2.b-equation]{式~\ref*{fga3.i-a.2.b-equation}} 的态射作任何平坦性假设，因此，为建立下降技术，必须对所考虑的 \(S\) 和 \(S'\) 上的对象作平坦性假设。在 \cref{fga3.i-b.2} 中，将对 \(\alpha \colon  S'\to  S\) 作平坦性假设，从而可对 \(S\) 和 \(S'\) 上不再要求平坦的对象建立下降技术。无论如何，总有某个地方涉及平坦性假设。这是平坦性概念在代数几何中如此重要的主要原因之一；只考虑以 \emph{域} 为基底时，这种作用不会显现，因为域上一切对象实际上都平坦！\end{remark}\section{对 étale 态射的应用}\label{fga3.i-a.3}\par{}设 \(A\) 为局部环，\(B\) 为 \(A\) 上的局部代数，其极大理想诱导 \(A\) 的极大理想。若满足以下条件，则称 \(B\) 在 \(A\) 上 \emph{étale}，而不用其他地方所说的“非分歧”：\begin{enumerate}
    
 \item[i.]\(B\) 在 \(A\) 上平坦；

    
 \item[ii.]\(B/\mathfrak {m}B\) 为 \(A/\mathfrak {m}=k\) 的有限可分扩张，其中 \(\mathfrak {m}\) 表示 \(A\) 的极大理想。

  \end{enumerate} 若 \(A\)、\(B\) 均为 Noether 环，且 \(k\) 代数闭，则这等价于延拓 \(A\to  B\) 的完备化同态 \(\overline {A}\to \overline {B}\) 是同构。若 \(\mathscr {O}_x\) 在 \(\mathscr {O}_{f(x)}\) 上 étale，则称有限型态射 \(f\colon  T\to  S\) \emph{在 \(x\in  T\) 处 étale}，或称 \(T\) \emph{在 \(x\) 处相对于 \(S\) 为 étale}；若 \(f\) 在每个 \(x\in  T\) 处 étale，则称 \(f\) 为 \emph{étale} 态射，或 \emph{étale 态射}，也称 \(T\) \emph{在 \(S\) 上 étale}。注意，若 \(S\) 局部 Noether，则 \(T\) 中使 \(f\) 为 étale 的点构成开集，Zariski“主定理”使我们能够精确描述 \(T/S\) 在这种点的邻域内的结构，给出一个众所周知类型的方程。\par{}若 \(S\) 为复数域上的有限型概形，则对应于一个解析空间 \(\overline {S}\)，其意义与 Serre \hyperref[Ser1956]{[5]} 相同，只是 \(\overline {S}\) 的结构层可以有幂零元，而这并不实质性改变 \hyperref[Ser1956]{[5]} 的内容。容易看出，\(f\) 是 étale 态射，当且仅当 \(\overline {f}\colon \overline {T}\to \overline {S}\) 是 étale 态射，即 \(\overline {T}\) 的每一点都有一个邻域，使 \(\overline {f}\) 在该邻域上诱导到 \(\overline {S}\) 的开子集的同构。特别地，\(S\) 的每个 \emph{étale 覆盖} \(T\)，即每个有限 étale 态射 \(f\colon  T\to  S\)，都对应于 \(\overline {S}\) 的一个 étale 覆盖 \(\overline {T}\)；它连通，当且仅当 \(T\) 连通 \hyperref[Ser1956]{[5]}。还容易看出，若 \(T\)、\(T'\) 为 \(S\) 上的 étale 概形，则自然映射 \[     \operatorname {Hom}_S(T,T') \to  \operatorname {Hom}_{\overline {S}}(\overline {T},\overline {T}')   \] 是双射；即从 \(S\) 上的 étale 概形范畴到 \(\overline {S}\) 上的 étale 解析空间范畴的函子 \(T\mapsto \overline {T}\) 是“全忠实”的，因此给出前一个范畴与后一个范畴的一个子范畴之间的等价。Grauert–Remmert 的定理 \hyperref[GR1958]{[2]} 蕴含：若 \(S\) 正规，则由此得到 \(S\) 的 \emph{étale 覆盖} 范畴与 \(\overline {S}\) 的 \emph{有限} étale 覆盖范畴之间的等价；即 \(\overline {S}\) 的每个 étale 覆盖 \(\mathscr {C}\) 均与某个 \(\overline {T}\) 在 \(\overline {S}\) 上同构，其中 \(T\) 为 \(S\) 的 étale 覆盖。我们将证明 \emph{不对 \(S\) 作任何正规性假设，Grauert–Remmert 定理仍成立}。先设 \(S'\to  S\) 为有限严格满态射，并假设定理对 \(S'\) 已成立；要证明它对 \(S\) 成立。设 \(\mathscr {C}\) 为 \(\overline {S}\) 的 étale 覆盖，考虑它在 \(S'\) 上的逆像 \(\mathscr {C}'\)；后者对应于 \(S'\) 上的一个凝聚解析代数层 \(\mathfrak {A}'\)，该层为 \(\overline {S}\) 上定义 \(\mathscr {C}\) 的代数层 \(\mathfrak {A}\) 的逆像。由关于 \(S'\) 的假设，\(\mathscr {C}'\) 来自 \(S'\) 的一个 étale 覆盖 \(T'\)，即 \(\mathfrak {A}'\) 来自 \(S'\) 上的凝聚代数层 \(\mathscr {A}'\)。此外，\(\mathfrak {A}'\) 赋有关于 \(\overline {S}'\to \overline {S}\) 的典范下降资料，即其在 \(\overline {S}'\times _{\overline {S}}\overline {S}'=\overline {S'\times _SS'}\) 上的两个逆像之间的同构，满足条件（i）和（ii）；由前述讨论，此同构来自相应代数层之间的同构，即来自 \(\mathscr {A}'\) 上关于 \(S'\to  S\) 的下降资料。容易证明后者有效，因为它在 \(\mathfrak {A}'\) 上有效，而上一节已明确描述的下降资料的有效性，可以在所涉及的模的 \emph{完备化} 上局部检验。由此得到 \(S\) 上的凝聚代数层 \(\mathscr {A}\)，定义了 \(S\) 的一个覆盖 \(T\)，这就是所求覆盖。显然，当 \(S'\to  S\) 仅为有限多个有限严格满态射的复合，即由 \cref{fga3.i-a.2} 的分解结果，仅为任意有限满态射时，上述结果仍成立。因此，若 \(S\) 为 \emph{既约} 概形，即 \(\mathscr {O}_S\) 没有幂零元，则引入其正规化 \(S'\)，便知 Grauert–Remmert 定理仍成立。容易过渡到一般情形。\par{}一个完全类似的论证，同样利用有限严格满态射的分解结果以及下降资料有效性的“形式”性质，能证明如下结果：设 \(S\) 为局部 Noether 预概形，\(S'\to  S\) 为有限、满射且根式的态射；等价地，它是一个有限型态射，使得对 \(S\) 上的每个 \(T\)，态射 \(T'=S'\times _S T\to  T\) 都是同胚，也可说 \(S'\to  S\) 是 \emph{普遍同胚}。对 \(S\) 上的每个 étale 预概形 \(T\)，考虑其逆像 \(T'=T\times _S S'\)，它在 \(S'\) 上 étale。\emph{于是函子 \(T\mapsto  T'\) 是 \(S\) 上的 étale 预概形 \(T\) 的范畴与 \(S'\) 上的 étale 预概形 \(T'\) 的范畴之间的等价。} 这里使用了：对 \(S\) 上的 étale 预概形 \(T_1\) 和 \(T_2\)，映射 \[     \operatorname {Hom}_S(T_1,T_2) \to  \operatorname {Hom}_{S'}(T'_1,T'_2)   \] 为双射；这可毫无困难地直接证明。还使用了：若 \(S'=(S,\mathscr {O}_S/\mathscr {J})\)，其中 \(\mathscr {J}\) 为 \(\mathscr {O}_S\) 的幂零凝聚理想层，则所述定理成立，参见［\hyperref[Mur1958]{4}，引理 6］。还要注意，这里不假设所涉及的 \(T\) 在 \(S\) 上有限。特别地，此结果蕴含态射 \(S'\to  S\) 诱导 \(S'\) 的基本群与 \(S\) 的基本群之间的同构，即“\emph{预概形基本群的拓扑不变性}”。\section{与一次上同调的关系}\label{fga3.i-a.4}\subsection{}\label{fga3.i-a.4.a}\par{}设 \(\mathcal {C}\) 为任意两个对象的乘积总存在的范畴，且 \(T\in \mathcal {C}\)。对每个有限集 \(I\neq \varnothing \)，可以考虑 \(T^I\)，从而得到从非空有限集范畴到范畴 \(\mathcal {C}\) 的协变函子，即可称为 \(\mathcal {C}\) 的一个 \emph{单纯对象}，记为 \(K_T\)。此对象协变地依赖于 \(T\)；此外，\emph{若 \(u\) 和 \(v\) 是态射 \(T\to  T'\)，则相应态射 \(K_T\to  K_{T'}\) 同伦}。若 \(\operatorname {Hom}(T,T')\neq \varnothing \)，则称 \(T\) \emph{支配} \(T'\)，这在 \(\mathcal {C}\) 上给出一个向上有向的预序。由前述讨论，若 \(T\) 支配 \(T'\)，则存在单纯对象同态 \(K_T\to  K_{T'}\) 的一个典范同伦类；特别地，若 \(K_T\)、\(K_{T'}\) 所对应的 \(T\) 与 \(T'\) 相互支配，则 \(K_T\) 与 \(K_{T'}\) 同伦等价。现设 \(F\) 为从 \(\mathcal {C}\) 到 \emph{abelian} 范畴 \(\mathcal {C}'\) 的函子，明确起见，设其反变。于是 \[     C^\bullet (T,F) = F(K_T)   \] 是 \(\mathcal {C}'\) 的余单纯对象，从而以熟知的方式定义 \(\mathcal {C}'\) 中的上链复形，可取其上同调：\[     \operatorname {H}^\bullet (T,F)     = \operatorname {H}^\bullet (C^\bullet (T,F))     = \operatorname {H}^\bullet (F(K_T))   \] 若可能混淆，可在 \(\operatorname {H}^\bullet \) 上添下标“\(\mathcal {C}\)”。这是关于 \(F\) 的上同调函子；当 \(T\) 变化时的变性，由先前关于 \(K_T\) 的讨论得出。更确切地说，固定 \(F\)，让 \(T\) 在以支配关系预序的 \(\mathcal {C}\) 中变化，则各个 \(\operatorname {H}^\bullet (T,F)\) 构成 \(\mathcal {C}'\) 中分次对象的归纳系；特别地，若 \(T\)、\(T'\) 相互支配，则 \(\operatorname {H}^\bullet (T,F)\) 与 \(\operatorname {H}^\bullet (T',F)\) 典范同构。\par{}假设 \(\mathcal {C}\) 中一切纤维积均存在。固定 \(S\in \mathcal {C}\)，即可将上述讨论用于 \(S\) 上的 \(\mathcal {C}\) 的对象范畴 \(\mathcal {C}_S\)；若要明确指出是在范畴 \(\mathcal {C}_S\) 中工作，就用 \(C^\bullet (T/S,F)\) 与 \(\operatorname {H}^\bullet (T/S,F)\) 代替 \(C^\bullet (T,F)\) 与 \(\operatorname {H}^\bullet (T,F)\)。于是 \(C^\bullet (T/S,F)\) 是 \(\mathcal {C}'\) 中的上链复形，其 \(n\) 次项为 \(F(T\times _S T\times _S\ldots \times _S T)\)，其中共有 \(n+1\) 个因子 \(T\)。\par{}注意，照常可以不假设范畴 \(\mathcal {C}'\) 为 abelian 范畴就定义 \(\operatorname {H}^0(T/S,F)\)：它是对应于两个投影 \(p_1,p_2\colon  T\times _S T\to  T\) 的态射对 \((F(p_1),F(p_2))\)、\[     F(T) \to  F(T\times _S T)   \] 的核（\hyperref[fga3.i-a.2-definition-2.1]{定义~\ref*{fga3.i-a.2-definition-2.1}}），只要此核存在。特别地，有自然态射 \[     F(S) \to  \operatorname {H}^0(T/S,F)   \]，称为 \emph{增广}；在良好情形下它是同构，特别是当 \(T\to  S\) 为严格满态射且 \(F\)“左正合”时如此。类似地，若 \(F\) 取值于某范畴 \(\mathcal {C}''\) 中的群范畴，也可定义 \(\operatorname {H}^1(T/S,F)\)；当 \(\mathcal {C}''\) 为集合范畴，即 \(F\) 取值于不必交换的群范畴时，\(\operatorname {H}^1(T,F)\) 是 \(C^1(T/S,F) = F(T\times _S T)\) 中满足 \[     F(p_{31})(g)     = F(p_{32})(g) F(p_{21})(g)   \] 的元素 \(g\) 所成子群 \(Z^1(T/S,F)\)，对算子群 \(F(T)\) 的作用取商。该群通过 \[     \rho (g')\cdot  g     = F(p_2)(g') g F(p_1)(g')^{-1}.   \] 作用于 \(C^1(T/S,F)\)，因而特别作用于子集 \(Z^1(T/S,F)\)。\subsection{}\label{fga3.i-a.4.b}\par{}例如，设 \(\mathcal {F}\) 为以 \(\mathcal {C}\) 为基底的纤维范畴，\(\xi ,\eta \in \mathcal {F}_S\)；对 \(S\) 上的所有 \(S'\)，置 \[     F_{\xi ,\eta }(S')     = \operatorname {Hom}(\xi \times _S S', \eta \times _S S').   \]。于是 \(F_{\xi ,\eta }\) 为从 \(\mathcal {C}_S\) 到集合范畴的反变函子。在此设定下，\emph{断言增广态射} \[     F_{\xi ,\eta }(S) \to  \operatorname {H}^0(S'/S,F_{\xi ,\eta })   \] \emph{对每对元素 \(\xi ,\eta \in \mathcal {F}_S\) 都是同构，就等价于断言 \(\alpha \colon  S'\to  S\) 是 \(\mathcal {F}\) 下降态射}（\hyperref[fga3.i-a.1-definition-1.7]{定义~\ref*{fga3.i-a.1-definition-1.7}}）。\subsection{}\label{fga3.i-a.4.c}\par{}类似地，对 \(\xi \in \mathcal {F}_S\)，对 \(S\) 上的 \(\mathcal {C}\) 的任意对象 \(S'\)，用 \[     G_\xi (S')     = \operatorname {Aut}(\xi \times _S S'),   \] 定义从 \(\mathcal {C}_S\) 到群范畴的反变函子 \(G_\xi \)。在此设定下，断言 \emph{\(Z^1(S'/S,G)\) 典范地等同于 \(\xi '=\xi \times _S S'\) 上关于 \(S'\to  S\) 的下降资料的集合}（\hyperref[fga3.i-a.1-definition-1.6]{定义~\ref*{fga3.i-a.1-definition-1.6}}），并且 \emph{\(\operatorname {H}^1(S'/S,G)\) 可等同于如下对象的同构类集：\(\mathcal {F}_{S'}\) 的对象赋有关于 \(\alpha \colon  S'\to  S\) 的下降资料，且作为 \(\mathcal {F}_{S'}\) 的对象同构于 \(\xi '=\xi \times _S S'\)}。于是，\emph{若 \(\alpha \colon  S'\to  S\) 为 \(\mathcal {F}\) 下降态射}（参见 \cref{fga3.i-a.4.b}），\emph{则 \(\operatorname {H}^1(S'/S,G)\) 包含如下对象的同构类集为子集：\(\mathcal {F}_S\) 的对象 \(\eta \) 满足 \(\eta \times _S S'\) 在 \(\mathcal {F}_{S'}\) 中同构于 \(\xi \times _S S'\)}；此外，\emph{此包含为恒等，当且仅当 \(\xi '=\xi \times _S S'\) 上关于 \(\alpha \colon  S'\to  S\) 的每个下降资料都有效}。特别地，若 \(\alpha \colon  S'\to  S\) 是严格 \(S\) 下降态射，就会如此。\begin{remark}{}\label{fga3.i-a.4.c-remark}\skipsinglepar\par{}形如 \(C^\bullet (T/S,F)\) 的上链复形包含大多数已知的标准复形为特例，例如 Čech 上同调、群上同调的复形等，并在代数几何中，特别是在预概形的“Weil 上同调”中发挥重要作用。\end{remark}\subsection{}\label{fga3.i-a.4.d}\begin{example}{1}\setnumber{1}\label{fga3.i-a.4-example-1}\skipsinglepar\par{}设 \(S'\) 为 \(S\in \mathcal {C}\) 上的对象，\(\Gamma \) 为 \(S'\) 的一个自同构群，使得 \(S'\)“在 \(S\) 上形式地为 \(\Gamma \) 主对象”，即自然态射 \[       \Gamma \times  S' \to  S'\times _S S'     \] 是同构，其中 \(\Gamma \times  S'\) 表示 \(\Gamma \) 个 \(S'\) 的直和；假设 \(\mathcal {C}\) 中所需直和均存在。设 \(F\) 为从 \(\mathcal {C}\) 到 abelian 群范畴的反变函子。则 \emph{\(C^\bullet (S'/S,F)\) 典范同构于标准齐次上链的单纯群 \(C^\bullet (\Gamma ,F(S'))\)，所以 \(\operatorname {H}^\bullet (S'/S,F)\) 典范同构于 \(\operatorname {H}^\bullet (\Gamma ,F(S'))\)}。\end{example}\subsection{}\label{fga3.i-a.4.e}\begin{example}{2}\setnumber{2}\label{fga3.i-a.4-example-2}\skipsinglepar\par{}设 \(\mathcal {C}\) 为预概形范畴。以 \(\operatorname {G_a}\) 表示从 \(\mathcal {C}\) 到 abelian 群范畴的反变函子，由 \[       \operatorname {G_a}(X)       = \operatorname {H}^0(X,\mathscr {O}_X).     \] 定义，其记号意为“加法群”。类似地，以 \[       \operatorname {G_m}(X)       = \operatorname {H}^0(X,\mathscr {O}_X)^\times      \] 定义函子 \(\operatorname {G_m}\)，意为“乘法群”，即环 \(\operatorname {H}^0(X,\mathscr {O}_X)\) 的可逆元群；更一般地，以 \[       \operatorname {GL}(n)(X)       = \operatorname {GL}(n,\operatorname {H}^0(X,\mathscr {O}_X)),     \] 定义函子 \(\operatorname {GL}(n)\)，意为“\(n\) 阶线性群”，这是从 \(\mathcal {C}\) 到群范畴的函子，当 \(n>1\) 时群不必交换；取 \(n=1\) 就恢复 \(\operatorname {G_m}\)。还可以将 \(\operatorname {GL}(n)\) 看作自同构函子（参见 \cref{fga3.i-a.4.c}）：考虑以 \(\mathcal {C}\) 为基底的纤维范畴 \(\mathcal {F}\)，其中对 \(X\in \mathcal {C}\)，\(\mathcal {F}_X\) 为 \(X\) 上的局部自由层范畴，因为此时 \(\operatorname {GL}(n)(X)=\operatorname {Aut}_{\mathcal {F}_X}(\mathscr {O}_X^n)\)。由 \cref{fga3.i-a.4.b}，若 \(\alpha \colon  S'\to  S\) 为 \(\mathcal {F}\) 下降态射（参见 \cref{fga3.i-a.2.c}），则 \(\operatorname {H}^1(S'/S,\operatorname {GL}(n))\) 包含 \(S\) 上逆像到 \(S'\) 后同构于 \(\mathscr {O}_{S'}^n\) 的局部自由层的同构类集；此包含为等号，当且仅当 \(\mathscr {O}_{S'}^n\) 上关于 \(\alpha \colon  S'\to  S\) 的每个下降资料均有效。若 \(S\) 为局部环的谱，则这蕴含 \(\operatorname {H}^1(S'/S,\operatorname {GL}(n))=(e)\)，因为此时 \(S\) 上的每个局部自由层均平凡。\par{}注意，关于态射 \(\alpha \colon  S'\to  S\) 的下列条件等价：\begin{enumerate}
      
 \item[i.]增广同态 \(\operatorname {H}^0(S,\mathscr {O}_S) = \operatorname {G_a}(S)\to \operatorname {H}^0(S'/S,\operatorname {G_a})\) 是同构；

      
 \item[ii.]\(\alpha \colon  S'\to  S\) 是 \(\mathcal {F}\) 下降态射，其中 \(\mathcal {F}\) 是上述 \(\mathcal {C}\) 上的纤维范畴；

      
 \item[iii.]若 \(S'\to S\) 有限，则上述条件还等价于：\(\alpha \colon  S'\to  S\) 是严格满态射（参见 \cref{fga3.i-a.2.c}）。

    \end{enumerate}\par{}现在假设 \(S=\operatorname {Spec}(A)\) 和 \(S'=\operatorname {Spec}(A')\)；则 \[       C^n(S'/S,\operatorname {G_a})       = C^n(A'/A,\operatorname {G_a})       = \underbrace {A'\otimes _A A'\otimes _A\ldots \otimes _A A'}_{n+1\text { 个因子 }A'}     \]，其上边缘算子 \(C^n(A'/A,\operatorname {G_a})\to  C^{n+1}(A'/A,\operatorname {G_a})\) 是面算子 \[       \partial _i(x_0\otimes  x_1\otimes \ldots \otimes  x_n)       = x_0\otimes \ldots \otimes  x_{i-1}\otimes 1_{A'}\otimes  x_i\otimes \ldots \otimes  x_n.     \] 的交错和。类似地，\(C^n(S'/S,\operatorname {G_m})=C^n(A'/A,\operatorname {G_m})\) 可等同于 \((\bigotimes _A^{n+1}A')^\times \)，而 \(C^\bullet (A'/A,\operatorname {G_m})\) 的单纯运算由 \(C^\bullet (S'/S,\operatorname {G_a})\) 中的运算诱导。也可同样明确地写出 \(C^\bullet (A'/A,\operatorname {GL}(n))\) 的单纯运算。\emph{在报告人所知道的一切情形下，对 \(i>0\) 有 \(\operatorname {H}^i(A'/A,\operatorname {G_a})=0\)；若 \(A\) 为局部环，则 \(\operatorname {H}^1(A'/A,\operatorname {G_m})=0\)，更一般地，\(\operatorname {H}^1(A'/A,\operatorname {GL}(n))=(e)\)} 若 \(S'\to  S\) 为 \(\mathcal {F}\) 下降态射，即图式 \(A\to  A'\rightrightarrows  A'\otimes _A A'\) 正合，则可与 \cref{fga3.i-a.2.c} 比较。注意，\emph{Hilbert 的“定理 90”恰好就是说：若 \(A\) 为域，\(A'\) 为 \(A\) 的有限 Galois 扩张，则 \(\operatorname {H}^1(S'/S,\operatorname {G_m})=0\)}（参见 \hyperref[fga3.i-a.4-example-1]{例~\ref*{fga3.i-a.4-example-1}}），\emph{也可表述为，在所考虑的情形下，\(S'\to  S\) 是秩为 \(1\) 的局部自由层的纤维范畴的严格下降态射}。最后这种表述最适于推广 Hilbert 定理：既改变对态射 \(S'\to  S\) 的假设，也改变所涉及的拟凝聚层。\par{}最后指出，对极大理想为 \(\mathfrak {m}\) 的局部 \emph{Artin} 环 \(A\) 及 \(A\) 代数 \(A'\)，下列性质等价；这里对每个整数 \(k>0\)，将环 \(A/\mathfrak {m}^{k+1}\)、相应地 \(A'/\mathfrak {m}^{k+1}A'\) 记为 \(A_k\)、相应地 \(A'_k\)：\begin{enumerate}
      
 \item[i.]对所有 \(k\)，有 \(\operatorname {H}^1(A'_k/A_k,\operatorname {G_a})=0\)；

      
 \item[ii.]对所有 \(k\)，有 \(\operatorname {H}^1(A'_k/A_k,\operatorname {G_m})=0\)；

      
 \item[iii.]对所有 \(k\) 和所有 \(n\)，有 \(\operatorname {H}^1(A'_k/A_k,\operatorname {GL}(n))=(e)\)。

    \end{enumerate} 若 \(S'\to  S\) 为严格满态射，则上述条件蕴含它对 \(A'\) 上不必有限型的自由模是 \emph{严格} 下降态射。\end{example}\begin{remark}{}\label{fga3.i-a.4.e-remark}\skipsinglepar\par{}当 \(S\)、相应地 \(S'\) 为域 \(A\)、相应地 \(A'\) 的谱时，群 \(\operatorname {H}^i(S'/S,\operatorname {G_m})\) 的定义归功于 Amitsur。群 \(\operatorname {H}^2(S'/S,\operatorname {G_m})\) 作为 Brauer 群的“整体”变式尤其有趣，参见［\hyperref[GD1960]{1}，VII］。\end{remark}\chapter{沿忠实平坦态射的下降}\thispagestyle{fancy}\label{fga3.i-b}\section{下降定理的陈述}\label{fga3.i-b.1}\begin{definition}{1.1}\setnumber{1.1}\label{fga3.i-b.1-definition-1.1}\skipsinglepar\par{}若对所有 \(x'\in  S'\)，\(\mathscr {O}_{x'}\) 都是环 \(\mathscr {O}_{\alpha (x')}\) 上的平坦模，即 \(\mathscr {O}_{x'}\otimes _{\mathscr {O}_{\alpha (x')}}M\) 是关于 \(\mathscr {O}_{\alpha (x')}\) 模 \(M\) 的正合函子，则称预概形态射 \(\alpha \colon  S'\to  S\) 为 \emph{平坦} 的。若一个态射平坦且满射，则称其为 \emph{忠实平坦} 的。\end{definition}\par{}例如，若 \(S=\operatorname {Spec}(A)\)、\(S'=\operatorname {Spec}(A')\)，则 \(S'\) 在 \(S\) 上平坦，当且仅当 \(A'\) 为平坦 \(A\) 模；\(S'\) 在 \(S\) 上忠实平坦，当且仅当 \(A'\) 为忠实平坦 \(A\) 模，即 \(A'\otimes _A M\) 是关于 \(A\) 模 \(M\) 的 \emph{正合} 且 \emph{忠实} 的函子。这也就是说，按 Serre \hyperref[Ser1956]{[5]} 的术语，二元组 \((A,A')\) 平坦。若 \(S'\) 在 \(S\) 上忠实平坦，则 \(S\) 上拟凝聚层的逆像函子正合且忠实；换言之，\(S\) 上一个拟凝聚层同态序列正合的充要条件是其在 \(S'\) 上的逆像正合。特别地，\(S\) 上一个拟凝聚层同态是单态射、满态射或同构的充要条件，是其在 \(S'\) 上的逆像具有相应性质。限制到 \(S\) 的任意开子集上方时，此性质仍成立，并以此形式刻画忠实平坦态射。\begin{definition}{1.2}\setnumber{1.2}\label{fga3.i-b.1-definition-1.2}\skipsinglepar\par{}若 \(S\) 的每个拟紧开子集 \(U\) 的逆像拟紧，即为仿射开子集的 \emph{有限} 并，则称态射 \(\alpha \colon  S'\to  S\) 为 \emph{拟紧} 的。\end{definition}\par{}显然，只需对 \(S\) 的 \emph{仿射} 开子集验证此性质。例如，仿射态射，即仿射开子集的逆像仍仿射的态射，是拟紧的。\par{}平坦态射、忠实平坦态射以及拟紧态射这三类态射，各自对合成与“基底扩张”稳定，当然也都包含一切同构。\begin{theorem}{1}\setnumber{1}\label{fga3.i-b.1-theorem-1}\skipsinglepar\par{}设 \(\alpha \colon  S'\to  S\) 为 \emph{忠实平坦} 且 \emph{拟紧} 的预概形态射。则 \(\alpha \) 是拟凝聚层的纤维范畴 \(\mathcal {F}\)（参见 \hyperref[fga3.i-a.1-example-2]{§A, 例 2}）的一个 \emph{严格下降态射}（参见 \hyperref[fga3.i-a.1-definition-1.7]{定义~\ref*{fga3.i-a.1-definition-1.7}}）。\end{theorem}\par{}此陈述包含两件事：\begin{enumerate}
    
 \item[i.]若 \(\mathcal {F}\) 与 \(\mathscr {G}\) 为 \(S\) 上的拟凝聚层，\(\mathcal {F}'\) 与 \(\mathscr {G}'\) 为它们在 \(S'\) 上的逆像，则自然同态 \[         \operatorname {Hom}(\mathcal {F},\mathscr {G}) \to  \operatorname {Hom}(\mathcal {F}',\mathscr {G}')       \] 是从左边到右边的如下子群的双射：它由与这些层上的典范下降资料相容的同态 \(\mathcal {F}'\to \mathscr {G}'\) 组成，即沿 \(S''=S'\times _S S'\) 到 \(S'\) 的两个投影取逆像后，得到相同的同态 \(\mathcal {F}''\to \mathscr {G}''\)；

    
 \item[ii.]\(S'\) 上每个赋有关于态射 \(\alpha \colon  S'\to  S\) 的下降资料（参见 \hyperref[fga3.i-a.1-definition-1.6]{定义~\ref*{fga3.i-a.1-definition-1.6}}）的拟凝聚层 \(\mathcal {F}'\)，连同此资料，同构于 \(S\) 上的某个拟凝聚层 \(\mathcal {F}\) 的逆像。

  \end{enumerate} 在（i）中置 \(\mathcal {F}=\mathscr {O}_S\)，便得到：\begin{corollary}{1}\setnumber{1}\label{fga3.i-b.1-corollary-1}\skipsinglepar\par{}设 \(\mathscr {G}\) 为 \(S\) 上的拟凝聚层，分别以 \(\mathscr {G}'\) 和 \(\mathscr {G}''\) 表示其在 \(S'\) 和 \(S''=S'\times _S S'\) 上的逆像，\(p_1\) 和 \(p_2\) 为 \(S''\) 到 \(S'\) 的两个投影。则集合映射图式 \[       \Gamma (\mathscr {G}) \xrightarrow {\alpha ^*}       \Gamma (\mathscr {G}') \overset {p_1^*}{\underset {p_2^*}{\rightrightarrows }}       \Gamma (\mathscr {G}'')     \] 为 \emph{正合} 的（参见 \hyperref[fga3.i-a.1-definition-1.1]{定义~\ref*{fga3.i-a.1-definition-1.1}}）。\end{corollary}\par{}此外，将（i）、（ii）与 \hyperref[fga3.i-a.1-definition-1.1]{定义~\ref*{fga3.i-a.1-definition-1.1}} 结合，得到：\begin{corollary}{2}\setnumber{2}\label{fga3.i-b.1-corollary-2}\skipsinglepar\par{}设 \(\mathscr {G}\) 如 \hyperref[fga3.i-b.1-corollary-1]{推论~\ref*{fga3.i-b.1-corollary-1}} 所述。则 \(\mathscr {G}\) 的拟凝聚子层与 \(\mathscr {G}'\) 中沿两个投影 \(p_1\)、\(p_2\) 到 \(S''\) 的逆像给出 \(\mathscr {G}''\) 的同一子层的拟凝聚子层之间，存在双射对应。\end{corollary}\par{}当然，存在用商层表述的等价陈述。正如 \cref{fga3.i-a.4.e} 所见，应将 \hyperref[fga3.i-b.1-theorem-1]{定理~\ref*{fga3.i-b.1-theorem-1}} 视为 Hilbert“定理 90”的推广，它包含各种用 \(1\) 次上同调表述的形式为特例。为证明它，容易归约到 \(S=\operatorname {Spec}(A)\)、\(S'=\operatorname {Spec}(A')\) 的情形；对（i），又容易只需证明 \hyperref[fga3.i-b.1-corollary-1]{推论~\ref*{fga3.i-b.1-corollary-1}}，即对每个 \(A\) 模 \(M\)，图式 \[     M = A\otimes _A M \to      A'\otimes _A M \rightrightarrows      A'\otimes _A A'\otimes _A M   \] 正合。这由下列更一般的引理得出：\begin{lemma}{1.1}\setnumber{1.1}\label{fga3.i-b.1-lemma-1.1}\skipsinglepar\par{}设 \(A'\) 为忠实平坦 \(A\) 代数。则对每个 \(A\) 模 \(M\)，由 \(M\) 增广的复形 \(C^\bullet (A'/A,\operatorname {G_a})\otimes _A M\)（参见 \cref{fga3.i-a.4.e}）是 \(M\) 的一个 \emph{消解}。\begin{proof}\par{}只需证明，将上面的增广复形沿 \(A\) 到 \(A'\) 作基底扩张后，所得复形满足同样结论。这就归结为把 \(A\) 换成 \(A'\)，把 \(A'\) 换成 \(A'\otimes _A A'\) 后证明此陈述，因此可只考虑存在 \(A\) 代数同态 \(A'\to  A\) 的情形；几何地说，就是 \(S'\) 在 \(S\) 上有截面的情形。此时断言由 \cref{fga3.i-a.4.a} 的一般讨论得出。\end{proof}\end{lemma}\par{}顺便指出下列推论，它推广了 Galois 上同调中的熟知陈述，可与 \cref{fga3.i-a.4.e} 比较：\begin{corollary}{}\label{fga3.i-b.1-lemma-1.1-corollary}\skipsinglepar\par{}若 \(A'\) 在 \(A\) 上忠实平坦，则 \(\operatorname {H}^0(A'/A,\operatorname {G_a})=A\)，并且对 \(i\geqslant 1\)，有 \(\operatorname {H}^i(A'/A,\operatorname {G_a})=0\)。\end{corollary}\par{}为证明 \hyperref[fga3.i-b.1-theorem-1]{定理~\ref*{fga3.i-b.1-theorem-1}} 的（ii），与（i）一样，先归约到 \(S'\) 在 \(S\) 上有截面的情形，此时结果由（i）得出，参见 \cref{fga3.i-a.1.c}。\par{}显然，可以在所考虑的拟凝聚层或层系上引入各种附加结构，\emph{任意地} 改变 \hyperref[fga3.i-b.1-theorem-1]{定理~\ref*{fga3.i-b.1-theorem-1}} 及其推论。例如，在 \(S\) 上给出一个拟凝聚交换代数层，“等价于”在 \(S'\) 上给出一个这种层，并赋有关于 \(\alpha \colon  S'\to  S\) 的下降资料。考虑到 \(S\) 上此类拟凝聚层与 \(S\) 上仿射预概形之间的函子性对应，就得到下列定理的第二个断言：\begin{theorem}{2}\setnumber{2}\label{fga3.i-b.1-theorem-2}\skipsinglepar\par{}设 \(\alpha \colon  S'\to  S\) 如 \hyperref[fga3.i-b.1-theorem-1]{定理~\ref*{fga3.i-b.1-theorem-1}} 所述。则 \(\alpha \) 是一个一般不严格的 \emph{下降态射}（参见 \hyperref[fga3.i-a.2-definition-2.4]{§A, 定义 2.4}），且对预概形上的仿射概形的纤维范畴，它是 \emph{严格下降态射}（参见 \hyperref[fga3.i-a.1-definition-1.7]{§A, 定义 1.7}）。\end{theorem}\par{}定理的第一个断言意味着：设 \(X\)、\(Y\) 为 \(S\) 上的预概形，\(X'\)、\(Y'\) 为它们在 \(S'\) 上的逆像，\(X''\)、\(Y''\) 为它们在 \(S''=S'\times _S S'\) 上的逆像；则自然映射图式 \[     \operatorname {Hom}_S(X,Y) \xrightarrow {\alpha ^*}     \operatorname {Hom}_{S'}(X',Y') \overset {p_1^*}{\underset {p_2^*}{\rightrightarrows }}     \operatorname {Hom}_{S''}(X'',Y'')   \] 为 \emph{正合} 的。也就是说，\(\alpha ^*\) 给出从 \(\operatorname {Hom}_S(X,Y)\) 到 \(\operatorname {Hom}_{S'}(X',Y')\) 中与 \(X'\) 和 \(Y'\) 上的典范下降资料相容的同态所成子集的双射；即这些同态沿 \(S''\) 到 \(S'\) 的两个投影取逆像后相同。若 \(Y\) 在 \(S\) 上仿射，则这容易由 \hyperref[fga3.i-b.1-theorem-1]{定理~\ref*{fga3.i-b.1-theorem-1}} 与 \hyperref[fga3.i-b.1-corollary-1]{推论~\ref*{fga3.i-b.1-corollary-1}} 得出；一般情形需将 \hyperref[fga3.i-b.1-theorem-1]{定理~\ref*{fga3.i-b.1-theorem-1}} 与下列结果结合：\begin{lemma}{1.2}\setnumber{1.2}\label{fga3.i-b.1-lemma-1.2}\skipsinglepar\par{}设 \(\alpha \colon  S'\to  S\) 为忠实平坦拟紧态射。则 \(S\) 可等同于 \emph{\(S'\) 的拓扑商空间}；也就是说，\(S\) 的每个满足 \(\alpha ^{-1}(U)\) 为开集的子集 \(U\)，本身也是开集。\end{lemma}\par{}为补充 \hyperref[fga3.i-b.1-theorem-2]{定理~\ref*{fga3.i-b.1-theorem-2}}，还须给出 \(S'\) 预概形 \(X'\) 上下降资料的有效性判别法，此时不假设 \(X'\) 在 \(S'\) 上仿射。首先指出，\emph{这样的下降资料未必有效}；即使 \(S\) 为域 \(k\) 的谱，\(S'\) 为 \(k\) 的二次扩张 \(k'\) 的谱，而 \(X'\) 为 \(S'\) 上 \(2\) 维固有代数概形，也可能如此；Serre 利用 Nagata 的非射影曲面说明了这一点。\emph{\(X'/S'\) 上关于忠实平坦拟紧态射 \(\alpha \colon  S'\to  S\) 的下降资料有效的充要条件是，\(X'\) 是一些在 \(S'\) 上仿射且在 \(X'\) 上的下降资料下“稳定”的开子集 \(X'_i\) 的并} 若态射 \(\alpha \colon  S'\to  S\) 是 \emph{根式} 的，即单射且剩余扩张为根式扩张，则对任意 \(X'/S'\) 及 \(X'\) 上的任意下降资料，这当然成立。还可证明，若 \(\alpha \colon  S'\to  S\) \emph{有限}，并且 \(X'\) 的每个包含于 \(X'\) 在 \(S\) 上的某个纤维的有限子集，也包含于 \(X'\) 的某个在 \(S\) 上仿射的开子集，则同样成立；这就是 \emph{Weil 判别法}。特别地，当 \(X'/S'\) 拟射影时如此；此时可证明“下降所得”的预概形 \(X/S\) 也拟射影，并且若 \(X'/S'\) 射影，它也射影。归纳如下：\begin{theorem}{3}\setnumber{3}\label{fga3.i-b.1-theorem-3}\skipsinglepar\par{}设 \(\alpha \colon  S'\to  S\) 为预概形的忠实平坦拟紧态射。若 \(\alpha \) \emph{根式}，则它是 \emph{严格下降态射}。若 \(\alpha \) 有限，则相对于预概形上的拟射影、或射影预概形的纤维范畴，它是严格下降态射。\end{theorem}\begin{remark}{}\label{fga3.i-b.1-remark}\skipsinglepar\par{}我不知道在上述第二个断言中，\(\alpha \) \emph{有限} 的假设是否确实必要；可以证明，无论如何，可以“形式地”将其替换为下列看起来较弱的假设：\emph{\(S\) 的每个点都有一个开邻域 \(U\)，并存在 \(U\) 上有限且忠实平坦的 \(U'\)，以及从 \(U'\) 到 \(S'\) 的 \(S\) 态射}。上述结果未涵盖的一类情形是 \(S=\operatorname {Spec}(A)\)、\(S'=\operatorname {Spec}(\overline {A})\)，其中 \(A\) 为 Noether 局部环，\(\overline {A}\) 为其完备化；又如 \(S'\) 在 \(S\) 上拟有限，即局部同构于某个有限 \(S\) 概形的开子集，但并不有限。在这两种情形下，报告人也不知道如下问题的答案：设 \(X\) 为 \(S\) 概形，使得 \(X'=X\times _S S'\) 在 \(S'\) 上射影，是否 \(X\) 也在 \(S\) 上射影？\par{}\emph{［编者补注］} 一个拟有限、étale 且满射的态射 \(S'\to  S\)，或形如 \(\operatorname {Spec}(\overline {A})\to \operatorname {Spec}(A)\) 的态射，并不总是严格下降态射；即使 \(A\) 为代数闭域 \(k\) 上代数曲线的局部环且 \(S=\operatorname {Spec}(A)\)，也可能如此。因此，可以找到固有光滑态射 \(f\colon  X\to  S\)，使 \(X\) 成为 \(S\) 上的主 \(E\) 丛，其中 \(E\) 为椭圆曲线，并使 \(f'\colon  X'\to  S'\) 射影，而 \(f\) 不射影。这也给出 abelian 概形下的一个 \emph{非等平凡} 主齐性丛的例子。\end{remark}\section{对态射某些性质的下降的应用}\label{fga3.i-b.2}\par{}设 \(P\) 为预概形态射的一个类，\(\alpha \colon  S'\to  S\) 为预概形态射，\(f\colon  X\to  Y\) 为 \(S\) 预概形的态射，\(f'\colon  X'\to  Y'\) 为 \(f\) 沿 \(\alpha \) 的逆像。可以问，关系“\(f'\in  P\)”是否蕴含“\(f\in  P\)”？若假设 \(\alpha \) \emph{忠实平坦} 且 \emph{拟紧}，则许多重要情形下答案为肯定；后一个假设有时过强。若 \(P\) 为满射态射类或根式态射类，容易由 \(\alpha \) 的满射性直接看出；若为平坦态射类、忠实平坦态射类或光滑态射类，则由 \(\alpha \) 的忠实平坦性得出；有限型态射类的情形也容易直接看出。利用 \hyperref[fga3.i-b.1-theorem-1]{定理~\ref*{fga3.i-b.1-theorem-1}}、\hyperref[fga3.i-b.1-theorem-2]{定理~\ref*{fga3.i-b.1-theorem-2}} 和 \hyperref[fga3.i-b.1-lemma-1.2]{引理~\ref*{fga3.i-b.1-lemma-1.2}}，可见当 \(P\) 为以下各类之一时也成立：同构、开浸入、闭浸入、浸入（假设 \(f\) 有限型且 \(Y\) 局部 Noether）、仿射态射、有限态射、拟有限态射、开态射、闭态射、同胚、分离态射或固有态射。这里未涵盖的唯一重要情形是射影或拟射影态射，这已在 \cref{fga3.i-b.1} 的注中讨论。\section{沿有限忠实平坦态射的下降}\label{fga3.i-b.3}\par{}设 \(\alpha \colon  S'\to  S\) 为 \emph{有限} 态射，对应于 \(S\) 上作为模层 \emph{局部自由} 且有限型、处处非零的代数层 \(\mathscr {A}'\)。于是 \(\alpha \) 为忠实平坦拟紧态射，可以应用上述结果。在 \(S'\) 上给出拟凝聚层 \(\mathcal {F}'\)，等价于在 \(S\) 上给出拟凝聚层 \(\alpha _*(\mathcal {F}')\) 及其 \(\mathscr {A}'\) 模结构；注意 \(\mathscr {A}'=\alpha _*(\mathscr {O}_{S'})\)。为简便起见，也将 \(S\) 上的此层记为 \(\mathcal {F}'\)。类似地，\(\mathcal {F}'\) 在 \(S'\times _S S'\) 上的两个逆像 \(p_i^*(\mathcal {F}')\)，对应于拟凝聚 \((\mathscr {A}'\otimes _{\mathscr {O}_S}\mathscr {A}')\) 模层 \(\mathcal {F}'\otimes _{\mathscr {O}_S}\mathscr {A}'\) 和 \(\mathscr {A}'\otimes _{\mathscr {O}_S}\mathcal {F}'\)。给出从前者到后者的 \((S'\times _S S')\) 同态，等价于给出一个 \((\mathscr {A}'\otimes \mathscr {A}')\) 模同态；又因 \(\mathscr {A}'\) 局部自由，这等价于给出如下 \((\mathscr {A}'\otimes \mathscr {A}')\) 模同态：\[     \mathscr {U}     = \mathscr {H}\kern  -2.5pt\mathit {om}_{\mathscr {O}_S}(\mathscr {A}',\mathscr {A}')     = \mathscr {A}'\otimes \check {\mathscr {A}}'     \to  \mathscr {H}\kern  -2.5pt\mathit {om}_{\mathscr {O}_S}(\mathcal {F}',\mathcal {F}')   \] 也就是说，对 \(\mathscr {U}\) 在开子集 \(V\) 上的每个截面 \(\xi \)，给出一个 \(\mathscr {O}_S\) 模同态 \(T_\xi \colon \mathcal {F}'|V\to \mathcal {F}'|V\)，并满足条件 \begin{eqn}{3.1}\setnumber{3.1}\label{fga3.i-b.3-equation-3.1}\skipsinglepar\[       \begin {aligned}         T_{f\xi }(x)         &= fT_\xi (x),       \\T_{\xi  f}(x)         &= T_\xi (fx),       \end {aligned}     \tag{3.1}     \]\end{eqn}，其中 \(f\) 和 \(x\) 分别为 \(\mathscr {A}'\) 与 \(\mathcal {F}'\) 在包含于 \(V\) 的 \(S\) 的某个开子集上的截面。下降资料的条件（i）与（ii）（参见 \cref{fga3.i-a.1.c}）于是分别写为 \begin{eqn}{3.2}\setnumber{3.2}\label{fga3.i-b.3-equation-3.2}\skipsinglepar\[       T_{1_U}(x)       = x       \qquad \text {即 }T_{1_U}=\operatorname {id}_{\mathcal {F}'}     \tag{3.2}     \]\end{eqn} \begin{eqn}{3.3}\setnumber{3.3}\label{fga3.i-b.3-equation-3.3}\skipsinglepar\[       T_{\xi \eta }       = T_\xi  T_\eta .     \tag{3.3}     \]\end{eqn}。换言之，\emph{\(\mathcal {F}'\) 上的下降资料等价于 \(\mathscr {O}_S\) 代数层 \(\mathscr {U}=\mathscr {H}\kern  -2.5pt\mathit {om}_{\mathscr {O}_S}(\mathscr {A}',\mathscr {A}')\) 在 \(\mathscr {O}_S\) 代数层 \(\mathscr {H}\kern  -2.5pt\mathit {om}_{\mathscr {O}_S}(\mathcal {F}',\mathcal {F}')\) 中的一个表示，并满足两个线性条件（\hyperref[fga3.i-b.3-equation-3.1]{式~\ref*{fga3.i-b.3-equation-3.1}}）}。若给定 \(S'\) 上拟凝聚层的一个配对 \[     \mathcal {F}'_1\times \mathcal {F}'_2 \to  \mathcal {F}'_3   \]，可视为 \(S\) 上 \(\mathscr {A}'\) 模层的配对，以及各 \(\mathcal {F}'_i\) 上由同态 \(T_i\colon \mathscr {U}\to \mathscr {H}\kern  -2.5pt\mathit {om}_{\mathscr {O}_S}(\mathcal {F}'_i,\mathcal {F}'_i)\)（\(i=1,2,3\)）定义的粘合资料，则这些资料在显然的意义下 \emph{与给定配对相容}，当且仅当满足下列条件：\par{}对 \(\mathscr {U}\) 在某个开子集上的每个截面 \(\xi \)，以 \(\Delta \xi =\sum \xi '_i\otimes _{\mathscr {A}'}\xi ''_i\) 表示 \(\mathscr {U}\otimes _{\mathscr {A}'}\mathscr {U}\) 的如下截面：这里把 \(\mathscr {U}\) 按其左结构视为 \(\mathscr {A}'\) 模，并以公式 \[     \xi \cdot (fg)     = \sum _i\xi '_i(f)\xi ''_i(g)   \] 定义该截面，其中 \(f\)、\(g\) 为 \(\mathscr {A}'\) 在较小开子集上的截面。那么，对 \(\mathscr {A}'\) 在较小子集上的每对截面 \((x,y)\)，有公式 \begin{eqn}{3.4}\setnumber{3.4}\label{fga3.i-b.3-equation-3.4}\skipsinglepar\[       T_\xi ^{(3)}(x\cdot  y)       = \sum _i T_{\xi '_i}^{(1)}x\cdot  T_{\xi ''_i}^{(2)}y     \tag{3.4}     \]\end{eqn}。此性质也可表述为：同态 \(T^{(i)}\) 相对于给定配对 \emph{与 \(\mathscr {U}\) 的对角映射相容}。特别地，\hyperref[fga3.i-b.3-equation-3.1]{式~\ref*{fga3.i-b.3-equation-3.1}}、\hyperref[fga3.i-b.3-equation-3.2]{式~\ref*{fga3.i-b.3-equation-3.2}}、\hyperref[fga3.i-b.3-equation-3.3]{式~\ref*{fga3.i-b.3-equation-3.3}} 与 \hyperref[fga3.i-b.3-equation-3.4]{式~\ref*{fga3.i-b.3-equation-3.4}} 使我们可以用带对角映射的代数的表示来理解 \(S'\) 上拟凝聚 \emph{代数} 层上的下降资料，因此限制到交换代数，也可理解仿射 \(S'\) 概形上的下降资料。\par{}由此可得，任意 \(S'\) 预概形 \(X'\) 上的下降资料也有类似解释：给出这样的 \(X'\)，等价于给出 \emph{在 \(S\) 上} 的预概形 \(X'\)，并赋有一个 \(\mathscr {O}_S\) 代数同态 \[     \mathscr {A}'\to \mathscr {O}_{X'},   \]；而 \(X'\) 上的下降资料，等价于一个与态射 \(h\colon  X'\to  S'\) 相容的层同态 \[     \mathscr {U}     \to  \mathscr {H}\kern  -2.5pt\mathit {om}_{h^{-1}(\mathscr {O}_S)}(\mathscr {O}_{X'},\mathscr {O}_{X'})   \]，并满足类似于上面 \hyperref[fga3.i-b.3-equation-3.1]{式~\ref*{fga3.i-b.3-equation-3.1}}、\hyperref[fga3.i-b.3-equation-3.2]{式~\ref*{fga3.i-b.3-equation-3.2}}、\hyperref[fga3.i-b.3-equation-3.3]{式~\ref*{fga3.i-b.3-equation-3.3}} 和 \hyperref[fga3.i-b.3-equation-3.4]{式~\ref*{fga3.i-b.3-equation-3.4}} 的条件。\begin{example}[{Weil 的例子}]{1}\setnumber{1}\label{fga3.i-b.3-example-1}\skipsinglepar\par{}假设 \(S'/S\) 为 Galois 群为 \(\Gamma \) 的一个 \emph{Galois étale 覆盖}（参见 \cref{fga3.i-a.3} 与 \cref{fga3.i-a.4}）。则 \(S'\) 上的拟凝聚层 \(\mathcal {F}'\)，相应地 \(S'\) 预概形 \(X'\)，上的下降资料，等价于一个将 \(\Gamma \) 表示为 \((S',\mathcal {F}')\)，相应地 \((S',X')\)，的自同构的表示，并与 \(\Gamma \) 在 \(S'\) 上的作用相容。此结果是“形式的”，即可以用范畴语言证明；但从本节的观点，还得到 \(\mathscr {U}\) 的具体结构，包括其环结构、环同态 \(\mathscr {A}'\to \mathscr {U}\) 及对角映射。下列结果完全确定了此结构：\emph{作为左 \(A'\) 模，\(\mathscr {U}\) 有一组基，由对应于 \(\Gamma \) 的元素的 \(\mathscr {U}\) 的截面给出}。\end{example}\begin{example}[{Cartier 的例子}]{2}\setnumber{2}\label{fga3.i-b.3-example-2}\skipsinglepar\par{}设 \(p\) 为素数，假设 \(p\mathscr {O}_S=0\)，即 \(\mathscr {O}_S\) 的 \emph{特征为 \(p\)}；假设 \((\mathscr {A}')^p\subset \mathscr {O}_S=\mathscr {A}\)，即 \(S'/S\) \emph{为高度 \(1\) 的根式态射}；并假设 \(\mathscr {A}\) 上的代数层 \(\mathscr {A}'\) \emph{局部有一个 \(p\) 基}，即一个截面族 \((x_i)\)，使 \(\mathscr {A}'\) 作为代数由各 \(x_i\) 生成，且仅受条件 \(x_i^p=0\) 约束。假设指标 \(i\) 的集合有限，基数为 \(n\)。令 \(\mathfrak {C}\) 为 \(A'\) 的 \(A\) 导子层；它是秩为 \(n\) 的局部自由 \(A'\) 模层，此外还是 \(\mathscr {A}\) 上而非 \(\mathscr {A}'\) 上的 Lie \(p\) 代数层，并满足下列条件：\begin{eqn}{3.5}\setnumber{3.5}\label{fga3.i-b.3-equation-3.5}\skipsinglepar\[         [X,fY] = X(f)Y + f[X,Y].       \tag{3.5}       \]\end{eqn}\end{example}\begin{lemma}{}\label{fga3.i-b.3-lemma}\skipsinglepar\par{}作为赋有代数同态 \(\mathscr {A}'\to \mathscr {U}\) 的 \(\mathscr {O}_S\) 代数，\(\mathscr {U}=\mathscr {H}\kern  -2.5pt\mathit {om}_{\mathscr {O}_S}(\mathscr {A}',\mathscr {A}')\) 由左 \(A'\) 子模 \(\mathfrak {C}\) 生成，并满足下列附加关系：\begin{eqn}{3.6}\setnumber{3.6}\label{fga3.i-b.3-equation-3.6}\skipsinglepar\[         \begin {cases}           Xf-fX &= X(f)         \\XY-YX &= [X,Y]         \\X^p &= X^{(p)}.         \end {cases}       \tag{3.6}       \]\end{eqn}\end{lemma}\par{}由上述引理，\(S'\) 上的拟凝聚层 \(\mathcal {F}'\) 的下降资料，等价于对每个 \(X\in \mathfrak {C}\) 给出 \(\mathcal {F}'\) 的一个 \(\mathscr {O}_S\) 自同态 \(\overline {X}\)，满足以下条件：\begin{eqn}{3.7}\setnumber{3.7}\label{fga3.i-b.3-equation-3.7}\skipsinglepar\[       \overline {fX} = f\overline {X}     \tag{3.7}     \]\end{eqn} \begin{eqn}{3.8}\setnumber{3.8}\label{fga3.i-b.3-equation-3.8}\skipsinglepar\[       \overline {X}(fx) = X(f)x + f\overline {X}(x)     \tag{3.8}     \]\end{eqn} \begin{eqn}{3.9}\setnumber{3.9}\label{fga3.i-b.3-equation-3.9}\skipsinglepar\[       \overline {[X,Y]} = [\overline {X},\overline {Y}]     \tag{3.9}     \]\end{eqn} \begin{eqn}{3.10}\setnumber{3.10}\label{fga3.i-b.3-equation-3.10}\skipsinglepar\[       \overline {X^{(p)}} = \overline {X}^p.     \tag{3.10}     \]\end{eqn} 这可称为一个 \emph{\(\mathcal {F}'\) 上的线性联络}，它还是 \emph{平坦} 的，并且 \emph{与 \(p\) 次幂相容}。类似地可以写出 \(S'\) 预概形 \(X'\) 上下降资料的概念，只需将 \hyperref[fga3.i-b.3-equation-3.4]{式~\ref*{fga3.i-b.3-equation-3.4}} 换为各个 \(\overline {X}\) 都是 \(\mathscr {O}_{X'}\) 的 \emph{导子} 的条件。由于态射 \(S'\to  S\) 为根式态射，\hyperref[fga3.i-b.1-theorem-3]{定理~\ref*{fga3.i-b.1-theorem-3}} 保证每个这样的下降资料均有效，因而定义一个 \(S\) 预概形 \(X\)。\par{}注意，对 \(\mathcal {F}'\) 或 \(X'\) 无须作任何平坦性、非奇异性或有限性假设。\section{对有理性判别法的应用}\label{fga3.i-b.4}\par{}设 \(X\) 为 \(S\) 预概形，使得 \(\mathscr {O}_X\) 在 \(S\) 上的直像为 \(\mathscr {O}_S\)；对任意平坦基底扩张 \(S'\to  S\)，此性质仍成立。若 \(\mathcal {F}\) 为 \(X\) 上的 \emph{可逆层}，即秩为 \(1\) 的局部自由层，则 \(\mathcal {F}\) 的自同构与 \(\mathscr {O}_S\) 的可逆截面之间存在双射对应，前者可等同于 \(\mathscr {O}_X\) 的可逆截面。现设 \(s\) 为 \(X\) 在 \(S\) 上的截面，将 \(S\) 上可逆层 \(s^*(\mathcal {F})\) 的截面定义为 \emph{\(\mathcal {F}\) 在 \(s\) 上方的截面}。由前述讨论，设 \(\mathcal {F}_i\)（\(i=1,2\)）是 \(X\) 上的可逆层，各自赋有 \(s\) 上方的一个截面；那么 \emph{若 \(\mathcal {F}_1\) 与 \(\mathcal {F}_2\) 同构，则恰好存在一个从 \(\mathcal {F}_1\) 到 \(\mathcal {F}_2\) 的同构，与所给截面相容}，即将前一个截面送到后一个截面。独立于截面 \(s\)，若 \(S\) 的每个点都有开邻域 \(U\)，使 \(\mathcal {F}_1\) 与 \(\mathcal {F}_2\) 在 \(X|U\) 上的限制同构，则称 \(X\) 上的两个可逆层 \(\mathcal {F}_1\) 和 \(\mathcal {F}_2\) \emph{等价}。于是 \emph{\(X\) 上的每个可逆层 \(\mathcal {F}\) 都等价于一个在 \(s\) 上方赋有标定截面的可逆层 \(\mathcal {F}_1\)}，取 \(\mathcal {F}_1=Fs^*(\mathcal {F})^{-1}\) 即可，\emph{并且 \(\mathcal {F}_1\) 在同构意义下确定}。换言之，\(X\) 上可逆层的 \emph{模等价} 分类，与赋有一个标定截面的可逆层的 \emph{模同构} 分类相同。\par{}这些性质在平坦基底扩张 \(\alpha \colon  S'\to  S\) 下仍成立，只需将截面 \(s\) 换为其沿 \(\alpha \) 的逆像 \(s'\)；因此结合 \hyperref[fga3.i-b.1-theorem-1]{定理~\ref*{fga3.i-b.1-theorem-1}}，得到：\par{}\emph{设预概形 \(X/S\) 如上，并有截面 \(s\)；设 \(\alpha \colon  S'\to  S\) 为忠实平坦拟紧态射，\(\mathcal {F}'\) 为 \(X'=X\times _S S'\) 上的可逆层。\(\mathcal {F}'\) 等价于 \(X\) 上某个可逆层 \(\mathcal {F}\) 在 \(X'\) 上的逆像的充要条件是，其在 \(X'\times _X X'=X\times _S(S'\times _S S')\) 上的两个逆像 \(p_1^*(\mathcal {F}')\) 和 \(p_2^*(\mathcal {F}')\) 等价。此时，\(\mathcal {F}\) 在等价意义下确定。} 此时称 \(\mathcal {F}'\) 在 \(S\) 上 \emph{有理}。\par{}当 \(\alpha \colon  S'\to  S\) 如上一节的 \hyperref[fga3.i-b.3-example-1]{例~\ref*{fga3.i-b.3-example-1}} 与 \hyperref[fga3.i-b.3-example-2]{例~\ref*{fga3.i-b.3-example-2}} 时考虑此原则，就恢复了 \emph{Weil 和 Cartier 的有理性判别法}。注意，那些作者只考虑 \(S\)、\(S'\) 为域的谱的情形；此时 \(S\) 更是局部环的谱，上述等价关系恰好就是同构关系。在第一种情形，\(\mathcal {F}'\) 在 \(S\) 上有理，当且仅当其在 \(\Gamma \) 下的像均等价于 \(\mathcal {F}'\)。为表达第二种情形的有理性判别法，更一般地考虑 \(X'/X\) 的对角态射 \(X'\to  X''=X'\times _X X'\)、相应的 \(X'\times _X X'\) 上的理想层 \(\mathscr {I}\) 及层 \(\mathscr {I}/\mathscr {I}^2\)；后者可等同于其在 \(X'\) 上的逆像 \(\Omega _{X'/X}^1\)，即 \emph{\(X'\) 相对于 \(X\) 的 \(1\) 次微分层}。由于各 \(\mathcal {F}''_i=p_i(\mathcal {F}')\)（\(i=1,2\)）在对角线上的限制同构，二者均同构于 \(\mathcal {F}'\)，也就是说 \(\mathcal {F}''_1(\mathcal {F}''_2)^{-1}=\mathcal {F}''\) 在对角线上的限制平凡，因此 \(\mathcal {F}''\) 在 \((X'',\mathscr {O}_{X''}/\mathscr {I}^2)\) 上的限制，在同构意义下，由 \[     \operatorname {H}^1(X'',\mathscr {I}/\mathscr {I}^2)     = \operatorname {H}^1(X',\Omega _{X'/X}^1).   \] 的一个确定元素 \(\xi \) 给出。此外，精确地说，有 \(\Omega _{X'/X}^1=\Omega _{S'/S}^1\otimes _{\mathscr {O}_S}\mathscr {O}_X\)，所以 \emph{若 \(\Omega _{S'/S}^1\) 在 \(S\) 上局部自由}，如 Cartier 情形那样，\emph{则 \(\xi \) 定义 \(S'\) 上 \(\operatorname {R}^1f'(\mathscr {O}_{X'})\otimes \Omega _{S'/S}^1\) 的一个截面}，称为 \emph{\(X'/S\) 上可逆层 \(\mathcal {F}'\) 的 Atiyah–Cartier 类}；\emph{此类消失，当且仅当 \(\mathcal {F}'\) 沿} \[     (X'',\mathscr {O}_{X''}/\mathscr {I}^2)     = X\times _S(S'',\mathscr {O}_{S''}/\mathscr {J}^2)   \] \emph{的两个投影到 \(X'\) 的逆像等价}，其中 \(\mathscr {J}\) 是 \(S''=S'\times _S S'\) 上由对角态射 \(S'\to  S'\times _S S'\) 定义的理想层。因此，此类的消失显然是 \(\mathcal {F}'\) 在 \(X''=X\times _S S''\) 本身上的两个逆像等价的 \emph{必要} 条件，从而也是 \(\mathcal {F}'\) 等价于 \(X\) 上某个可逆层 \(\mathcal {F}\) 的逆像的必要条件。Atiyah–Cartier 类也可理解为在 \(S'\) 上局部存在 \(\mathcal {F}'\) 的关于 \(X'/X\) 的导子的一个 \emph{联络} 的障碍；只要知道对应于 \(S'/S\) 的导子到 \(X'\) 的自然延拓的 \(\mathcal {F}'\) 的导子，这个联络就确定了。由此及上一节结果，容易得出：在上述 \hyperref[fga3.i-b.3-example-2]{例~\ref*{fga3.i-b.3-example-2}} 的情形，且 \(X/S\) 有截面时，Atiyah–Cartier 类消失也是 \(\mathcal {F}'\) 在 \(S\) 上有理的充分条件。\section{对 abelian 概形的基底限制的应用}\label{fga3.i-b.5}\par{}设 \(S\) 为预概形。\(S\) 上的 \emph{abelian 概形}，定义为 \(S\) 上的一个光滑固有概形 \(X\)，使其在各点 \(x\in  S\) 处的纤维为 \(k(x)\) 上 abelian 簇的概形。假设 \(S\) 为 Noether 且 \emph{正则} 的，即其局部环均正则，则利用 Murre \hyperref[Mur1958]{[4]} 的 \emph{连通性定理}，可以证明 \emph{\(X\) 在 \(S\) 上的每个有理截面都处处有定义}，即它是一个截面，这推广了 Weil 的经典定理；至少在前引定理目前已证实的“等特征”情形下如此。由此更一般地得到：若 \(X'\) 为 \(S\) 上的光滑概形，则从 \(X'\) 到 \(X\) 的每个有理 \(S\) 映射都处处有定义。由此得到如下结果，它推广了 Chow–Lang 的一个结果：\emph{设 \(S\) 为 Noether 正则预概形，\(K\) 表示其有理函数环}，它是若干域的直和，\emph{设 \(X\) 为 \(K\) 上的 abelian 概形；若 \(X\) 同构于形如 \(X_0\times _S\operatorname {Spec}(K)\) 的 \(K\) 概形，其中 \(X_0\) 为 \(S\) 上的 abelian 概形，则 \(X_0\) 在唯一同构的意义下确定。}\par{}利用上述唯一性结果，可见将 \(X\) 的基底限制下去的问题在 \(S\) 上是局部的，因此只需知道如何限制到 \(\operatorname {Spec}(\mathscr {O}_x)\)，其中 \(x\in  S\)。同样，若 \(S'\to  S\) 为有限型 \emph{光滑} 态射，\(Y'\) 为 \(S'\) 的有理函数环，且 \(X\otimes _K K'\) 形如 \(X'_0\times _{S'}\operatorname {Spec}(K')\)，则 \emph{\(X'_0\) 赋有关于 \(\alpha \) 的典范下降资料}。考虑到 \hyperref[fga3.i-b.1-theorem-3]{定理~\ref*{fga3.i-b.1-theorem-3}}，得到：\begin{proposition}{5.1}\setnumber{5.1}\label{fga3.i-b.5-proposition-5.1}\skipsinglepar\par{}设 \(S\) 为不可约、正则且 Noether 的预概形，有理函数域为 \(Y\)；设 \(K'\) 为 \(K\) 的一个 \emph{在 \(S\) 上非分歧} 的有限扩张，\(S'\) 为 \(S\) 在 \(K'\) 中的正规化，因此为 \(S\) 的 étale 覆盖；设 \(X\) 为 \(K\) 上的 abelian 概形，使 \(X\otimes _K K'\) 形如 \(X'_0\times _{S'}\operatorname {Spec}(K')\)，其中 \(X'_0\) 为 \(S'\) 上射影的 abelian 概形。则 \(X\) 形如 \(X_0\times _S\operatorname {Spec}(K)\)，其中 \(X_0\) 为 \(S\) 上射影的 abelian 概形。\end{proposition}\begin{remark}{}\label{fga3.i-b.5-remark}\skipsinglepar\par{}报告人不知道，能否将 \(S'\to  S\) 为满射 étale 覆盖的假设——它使我们能应用 \hyperref[fga3.i-b.1-theorem-3]{定理~\ref*{fga3.i-b.1-theorem-3}}——替换为有限型 \emph{光滑} 且 \emph{满射} 态射的假设，即使假设它为 étale 态射也不知道；也不知道，若不假设 \(X'_0\) 在 \(S'\) 上射影，命题是否仍成立，而这个条件也许自动满足。\end{remark}\section{对局部平凡性与等平凡性判别法的应用}\label{fga3.i-b.6}\par{}设 \(S\) 为预概形，\(G\) 为 \(S\) 上的一个“\emph{群预概形}”，\(P\) 为 \(S\) 上的一个预概形，其上“\emph{\(G\) 作用}”，设为右作用。若由 \(G\) 在 \(P\) 上的作用诱导的熟知态射 \[     G\times _S P \to  P\times _S P   \] 是 \emph{同构}，则称 \(P\) 关于 \(G\) 是 \emph{形式主齐性} 的。以下假设 \(G\) 在 \(S\) 上 \emph{平坦}，从而在 \(S\) 上忠实平坦；并只将 \(S\) 上 \emph{忠实平坦} 且 \emph{拟紧} 的形式主齐性丛 \(P\) 称为 \(G\) 的 \emph{主齐性丛}。显然，这等价于能找到基底 \(S\) 的一个 \emph{忠实平坦} 且 \emph{拟紧扩张} \(S'\to  S\)，使 \(G'=G\times _S S'\) 的形式主齐性丛 \(P'=P\times _S S'\) \emph{平凡}，即同构于 \(G'\)，也就是有一个截面；特别地，可以取 \(S'=P\)。还要指出，若 \(S\) 局部 Noether，则对 \(P\) 的忠实平坦性假设，等价于对每个 \(s\in  S\)，\(\overline {P}_s=P\times _S\operatorname {Spec}(\overline {\mathscr {O}}_s)\) 在 \(\overline {\mathscr {O}}_s\) 上忠实平坦；其中 \(\overline {\mathscr {O}}_s\) 表示局部环 \(\mathscr {O}_s\) 的完备化，这由 \(\overline {\mathscr {O}}_s\) 在 \(\mathscr {O}_s\) 上忠实平坦这一事实得出。此外，若 \(P\) 在 \(S\) 上有限型且 \(S\) 局部 Noether，则满足上述条件的点 \(s\) 的集合可构造；所以若 \(S\) 为“Jacobson 预概形”，例如域上，或更一般地 Jacobson 环上的有限型概形，则只需对 \(S\) 的 \emph{闭} 点验证所述条件。这使我们归结到基底为完备局部环 \(A\) 的谱的情形。若 \(S=\operatorname {Spec}(A)\)，其中 \(A\) 为完备 Noether 局部环，且 \(P\) 在 \(S\) 上有限型，则 \(P/S\) 的忠实平坦性还等价于存在一个在 \(S\) 上 \emph{有限且平坦} 的 \(S'\)，使 \(P'\) 平凡；若此外 \(G\) 在 \(S\) 上 \emph{光滑}，则可假设 \(S'\) 在 \(S\) 上 \emph{étale}。此时，若再假设 \(A\) 的剩余域代数闭，即“\emph{几何情形}”，则 \(P\) 在 \(A\) 上忠实平坦，当且仅当它平凡。因此，若 \(S\) 为代数闭域上的代数预概形，且 \(G\) 在 \(S\) 上光滑且有限型，则 \(S\) 上的忠实平坦性条件等价于 Serre［\hyperref[Ser1958a]{6}，第 1—12 页］的解析平凡性条件（SLF）。\par{}还可以对 \(P\) 考虑其他更强的、具有“局部平凡性”性质的条件。特别地，若对每个 \(s\in  S\)，存在 \(s\) 的一个开邻域 \(U\)，以及一个 \emph{有限且忠实平坦} 态射，相应地一个 \emph{满射 étale 覆盖}，\(U'\to  U\)，使 \(P'=P\times _S U'\) 平凡，则称 \(P\) 为 \emph{等平凡} 的，相应地为 \emph{严格等平凡} 的。这里与 Serre \hyperref[GD1960]{[1]}\footnote{译校注：原刊此处写作“Serre [1]”，但本报告原书目 [1] 为 Dieudonné–Grothendieck 的 EGA；疑指 Serre 的 [6]。保留原编号，不据猜测改换引文。} 的术语不同，他用“局部等平凡”指我们所谓的“严格等平凡”。严格等平凡性主要在 \(G\) 在 \(S\) 上光滑时有用，但在其他情形是不适当的概念。\par{}若 \(G\) 在 \(S\) 上 \emph{仿射}，则由 \cref{fga3.i-b.2}，\(G\) 的每个主齐性丛 \(P\) 均仿射；因此借助 \hyperref[fga3.i-b.1-theorem-2]{定理~\ref*{fga3.i-b.1-theorem-2}}，能够沿忠实平坦拟紧态射将这种丛“下降”。特别地，取 \(G=\operatorname {GL}(n)_S\)，它由如下条件定义：\(S\) 预概形上的取值于群范畴的函子 \(S'\mapsto \operatorname {Hom}(S',G)\) 可等同于 \cref{fga3.i-a.4} 所述函子 \(\operatorname {GL}(n)(S')=\operatorname {GL}(n,\operatorname {H}^0(S',\mathscr {O}_{S'}))\)。利用以下事实：\begin{enumerate}
    
 \item[i.]\(G\) 的每个主齐性丛，相应地 \(S\) 上的每个秩为 \(n\) 的局部自由层，在适当的 \(S\) 的忠实平坦拟紧扩张后，都同构于“平凡”对象 \(G\)，相应地 \(\mathscr {O}_S^n\)；


    
 \item[ii.]可以沿这种态射下降所考虑类型的对象，即 \(G\) 的主齐性丛，相应地秩为 \(n\) 的局部自由层；最后，


    
 \item[iii.]任意 \(S'/S\) 上的平凡丛的自同构群函子性地同构于 \(S'\) 上秩为 \(n\) 的平凡局部自由层的自同构群，

  \end{enumerate} 可以“形式地”得出，在 \(S\) 上，或在某个 \(S'/S\) 上，给出群 \(G\) 的一个主齐性丛，“等价于”给出一个秩为 \(n\) 的局部自由层。更精确地说，有一个 \emph{纤维范畴的等价}。特别地，得到：\begin{proposition}{6.1}\setnumber{6.1}\label{fga3.i-b.6-proposition-6.1}\skipsinglepar\par{}群 \(\operatorname {GL}(n)_S\) 的每个主齐性丛都局部平凡。\end{proposition}\par{}由熟知的论证，对其他结构群，例如 \(\operatorname {SL}(n)_S\)、\(\operatorname {Sp}(n)_S\) 及这些群的乘积，也得到相同结果。因此也得到：若 \(F\) 为 \(G=\operatorname {GL}(n)_S\) 的一个在 \(S\) 上平坦的闭子群，使商 \(G/F\) 存在，且 \(G\) 是 \(G/F\) 上结构群为 \(F\times _S(G/F)\) 的等平凡、相应地严格等平凡的主齐性丛，则结构群为 \(F\) 的 \emph{每个} 主齐性丛都等平凡、相应地严格等平凡。这适用于迄今使用过的 \(S\) 上的所有“线性群”；特别适用于 \(G=S\times _k\Gamma \) 的情形，其中 \(S\) 为域 \(k\) 上的预概形，\(\Gamma \) 为 \(k\) 上经典意义下的线性群，特别地它光滑。这就为这种群回答了 Serre 的一个问题（\emph{同前引处}）。\par{}还要指出，对我们所知道的大多数 \(S\) 上的单纯群，无论是否线性，并且当然对上述形如 \(S\times _k\Gamma \) 的所有群，都可证明每个等平凡主齐性丛都严格等平凡。这特别回答了 Serre 的另一个问题（\emph{同前引处}，第 1—14 页）；这里考虑到，沿 Cartier \emph{式方法} 下降（参见 \hyperref[fga3.i-b.3-example-2]{例~\ref*{fga3.i-b.3-example-2}}）所得的主齐性丛显然等平凡。\begin{remark}{}\label{fga3.i-b.6-remark}\skipsinglepar\par{}这些问题的一个本质困难，撇开商概形 \(G/F\) 的存在性问题，就是缺少沿忠实平坦 \emph{非有限} 态射的下降资料的有效性判别法。\end{remark}
\section*{原报告参考文献}
\begin{enumerate}
\item[1.]\setnumber{1}\label{fga3.i-bib-1} DIEUDONNÉ (J.) et GROTHENDIECK (Alexander). \emph{Éléments de géométrie algébrique}，即将发表于 \emph{Publications mathématiques de l'Institut des Hautes Études Scientifiques}。
\item[2.]\setnumber{2}\label{fga3.i-bib-2} GRAUERT (H.) und REMMERT (R.). \emph{Komplexe Räume}，\emph{Math. Annalen}，138 (1958)，245--318。
\item[3.]\setnumber{3}\label{fga3.i-bib-3} GROTHENDIECK (Alexander). \emph{Géométrie formelle et géométrie algébrique}，\emph{Séminaire Bourbaki}，11 (1958/59)，第 182 次报告。
\item[4.]\setnumber{4}\label{fga3.i-bib-4} MURRE (J. P.). \emph{On a connectedness theorem for a birational transformation at a simple point}，\emph{Amer. J. Math.}，80 (1958)，3--15。
\item[5.]\setnumber{5}\label{fga3.i-bib-5} SERRE (Jean-Pierre). \emph{Géométrie algébrique et géométrie analytique}，\emph{Ann. Inst. Fourier Grenoble}，6 (1955--56)，1--42。
\item[6.]\setnumber{6}\label{fga3.i-bib-6} SERRE (Jean-Pierre). \emph{Espaces fibrés algébriques}，\emph{Séminaire Chevalley}，2，1958：\emph{Anneaux de Chow et applications}，第 1 次报告。
\end{enumerate}

% END chapters/fga3.i.tex

\clearpage
\setcounter{chapter}{0}
\renewcommand\fgaAnchorPrefix{fga3.ii}

% BEGIN chapters/fga3.ii.tex
\part{存在定理与模空间的形式理论}\label{fga3.ii}\thispagestyle{fancy}\begin{original}{}\label{fga3.ii-original-citation}\skipsinglepar\par{}A. Grothendieck.
    "Technique de descente et théorèmes d'existence en géométrie algébrique, II: Le théorème d'existence et théorie formelle des modules".
    \emph{Séminaire Bourbaki} \textbf{12} (1959–60), 第 195 次报告，1960 年 2 月。
    \url{http://www.numdam.org/book-part/SB_1958-1960__5__369_0/}\end{original}\chapter{可表与 pro 可表函子}\thispagestyle{fancy}\label{fga3.ii-a}\section{可表函子}\label{fga3.ii-a.1}\par{}设 \(\mathcal {C}\) 为一个范畴。对每个 \(X\in \mathcal {C}\)，令 \(h_X\) 为由下式给出的从 \(\mathcal {C}\) 到集合范畴 \(\mathtt {Set}\) 的反变函子：\[     \begin {aligned}       h_X\colon  \mathcal {C} &\to  \mathtt {Set}     \\Y&\mapsto  \operatorname {Hom}(Y,X).     \end {aligned}   \] 显然，\(\mathcal {C}\) 中的一个态射 \(X\to  X'\) 诱导一个函子态射 \(h_X\to  h_{X'}\)；\(h_X\) 关于 \(X\) 协变，也就是说，我们定义了一个 \emph{典范协变函子} \[     h\colon  \mathcal {C} \to  \operatorname {Hom}(\mathcal {C}^\circ ,\mathtt {Set})   \]，其起点为 \(\mathcal {C}\)，终点为从 \(\mathcal {C}\) 的对偶范畴 \(\mathcal {C}^\circ \) 到集合范畴的协变函子范畴。现回顾如下事实：\begin{proposition}{1.1}\setnumber{1.1}\label{fga3.ii-a.1-proposition-1.1}\skipsinglepar\par{}这个函子 \(h\) 是 \emph{全忠实} 的；换言之，对 \(\mathcal {C}\) 的任意两个对象 \(X,X'\)，自然映射 \[       \operatorname {Hom}(X,X') \to  \operatorname {Hom}(h_X,h_{X'})     \] 是 \emph{双射}。\end{proposition}\par{}特别地，如果函子 \(F\in \operatorname {Hom}(\mathcal {C}^\circ ,\mathtt {Set})\) 同构于形如 \(h_X\) 的函子，那么 \emph{\(X\) 在唯一同构的意义下确定}。此时称函子 \(F\) 是 \emph{可表} 的。上述命题于是表明，典范函子 \(h\) 在范畴 \(\mathcal {C}\) 与 \(\operatorname {Hom}(\mathcal {C}^\circ ,\mathtt {Set})\) 中由可表函子组成的满子范畴之间定义了一个 \emph{等价}。这一事实是 \emph{“泛问题的解”这一观念} 的基础；这种问题总是归结为判定一个给定的从 \(\mathcal {C}\) 到 \(\mathtt {Set}\) 的函子是否可表（此处为反变函子，对偶情形则为协变函子）。还要指出，仅由范畴中乘积的定义 \hyperref[Gro1957]{[1]} 就可知，函子 \(h\colon  X\mapsto  h_X\) 与乘积可交换，只要这些乘积存在；更一般地，它与有限或无限射影极限可交换，特别是与纤维积、取“核”\hyperref[fga2]{[2]}等可交换，只要这些对象存在：当 \(X\times  X'\) 存在时，有函子同构 \[     h_{X\times  X'} \xrightarrow {\sim } h_X\times  h_{X'}   \]，即有关于 \(Y\) 函子性的双射 \[     h_{X\times  X'}(Y) \xrightarrow {\sim } h_X(Y)\times  h_{X'}(Y).   \]。特别地，给定 \(\mathcal {C}\) 中的态射 \[     X\times  X' \to  X''   \]（即 \(\mathcal {C}\) 中 \(X\)、\(X'\) 与 \(X''\) 之间的一个“\emph{合成律}”），等价于给定一个态射 \(h_{X\times  X'}=h_X\times  h_{X'}\to  h_{X''}\)，也就等价于对每个 \(Y\in \mathcal {C}\) 给定一个 \emph{集合} 的合成律 \[     h_X(Y)\times  h_{X'}(Y) \to  h_{X''}(Y)   \]，使得对 \(\mathcal {C}\) 中的每个态射 \(Y\to  Y'\)，集合映射族 \[     h_{X^{(i)}}(Y) \to  h_{X^{(i)}}(Y')     \qquad \text {其中 }i=0,1,2   \] 是关于 \(Y\) 和 \(Y'\) 的两个合成律之间的态射。由此可见，在 \(\mathcal {C}\) 的对象 \(X\) 上给出“\(\mathcal {C}\) 群”结构、“\(\mathcal {C}\) 环”结构等等，无论在理论上还是在实践中，最便于处理的表达方式都是：对每个 \(Y\in \mathcal {C}\)，集合 \(h_X(Y)\) 上具有通常意义下的群运算（相应地，环运算等等），而对应于态射 \(Y\to  Y'\) 的映射 \(h_X(Y)\to  h_X(Y')\) 应为群同态（相应地，环同态等等）。例如，在任意基底上的预概形 \(S\) 上定义各种经典群 \(\operatorname {G_a}\)、\(\operatorname {G_m}\)、\(\operatorname {GL}(n)\) 等，写出这些群之间的经典关系，或者在由拟凝聚层 \(\mathscr {F}\) 定义的 \(S\) 上仿射概形 \(V(\mathscr {F})\) 上赋予“向量丛”结构，以及定义、研究许多相关的旗簇（Grassmann 簇、射影丛）等等，这都是最直观且最便于处理的方式；\emph{一般的处理原则，就是用典范函子 \(h\)，直接把 \(\mathcal {C}\) 的对象视为从 \(\mathcal {C}\) 到集合范畴的某些特定反变函子，即可表函子}。\par{}例如，对于仿射概形，为从几何语言转到交换代数语言而必须进行的通常的反转箭头操作，引导我们将上述讨论对偶化，特别是还要引入 \emph{协变可表函子 \(\mathcal {C}\to \mathtt {Set}\)}，即形如 \(Y\mapsto \operatorname {Hom}(X,Y)=h'_X(Y)\) 的函子。\section{pro 可表函子与 pro 对象}\label{fga3.ii-a.2}\par{}设 \(\mathcal {X}=(X_i)_{i\in  I}\) 为 \(\mathcal {C}\) 中对象的一个射影系；它对应于协变函子 \[     h'_{\mathcal {X}}     = \varinjlim _i h'_{X_i}   \]，更明确地可写为 \[     h'_{\mathcal {X}}(Y)     = \varinjlim _i h'_{X_i}(Y)     = \varinjlim _i\operatorname {Hom}(X_i,Y)   \]，这是从 \(\mathcal {C}\) 到 \(\mathtt {Set}\) 的函子。从 \(\mathcal {C}\) 到 \(\mathtt {Set}\) 的函子，若同构于一个这种形式且 \emph{\(I\) 为滤过集} 的函子，就称为 \emph{pro 可表} 的。由上一节可知，它们恰好是同构于 \emph{可表函子的滤过归纳极限} 的函子。设 \(\mathcal {X}'=(X_j)_{j\in  J}\) 为 \(\mathcal {C}\) 中另一个滤过射影系，以另一个滤过预序指标集 \(J\) 为指标。容易证明有典范双射 \[     \operatorname {Hom}(h_{\mathcal {X}'},h_{\mathcal {X}})     = \varprojlim _j\varinjlim _i\operatorname {Hom}(X_i,X'_j)   \]，这推广了 \hyperref[fga3.ii-a.1-proposition-1.1]{命题~\ref*{fga3.ii-a.1-proposition-1.1}}。这就引导我们引入 \emph{\(\mathcal {C}\) 的 pro 对象范畴 \(\operatorname {Pro}(\mathcal {C})\)}：其对象是 \(\mathcal {C}\) 中对象的射影系，以任意滤过预序集为指标；对象 \(\mathcal {X}=(X_i)_{i\in  I}\) 与 \(\mathcal {X}'=(X_j)_{j\in  J}\) 之间的态射由 \[     \operatorname {Pro}\operatorname {Hom}(\mathcal {X},\mathcal {X}')     = \varprojlim _j\varinjlim _i\operatorname {Hom}(X_i,X'_j),   \] 给出，pro 同态的合成是显然的。由构造本身，\(\mathcal {X}\mapsto  h'_{\mathcal {X}}\) 可以视为关于 \(\mathcal {X}\) 的反变函子，从而建立 \emph{\(\mathcal {C}\) 的 pro 对象范畴 \(\operatorname {Pro}(\mathcal {C})\) 的对偶范畴，与从 \(\mathcal {C}\) 到 \(\mathtt {Set}\) 的 pro 可表协变函子范畴之间的等价}。当然，\(\mathcal {C}\) 的对象 \(X\) 典范地定义一个 pro 对象，仍记为 \(X\)，所以 \emph{\(\mathcal {C}\) 等价于 \(\operatorname {Pro}(\mathcal {C})\) 的一个满子范畴}。于是，若 \(\mathcal {X}=(X_i)_{i\in  I}\) 是 \(\mathcal {C}\) 的任意 pro 对象，在上述等同下便有 \[     \mathcal {X}     = \varprojlim _i X_i   \]，其中射影极限是 \emph{在 \(\operatorname {Pro}(\mathcal {C})\) 中取} 的（因为 \(h_{\mathcal {X}}=\varinjlim _i h_{X_i}\)）。\par{}需要注意，即使 \(X_i\) 的射影极限 \emph{在 \(\mathcal {C}\) 中存在}，它通常也 \emph{不} 同构于 \(\operatorname {Pro}(\mathcal {C})\) 中的射影极限 \(\mathcal {X}\)；在 \(\mathcal {C}\) 为集合范畴时这一点已经很明显。由定义本身，\(\varprojlim {}_{\mathcal {C}}X_i=L\) 由如下条件确定：以 \(Y\in \mathcal {C}\) 为变量、取值于 \(\mathtt {Set}\) 的函子 \[     \varprojlim _i\operatorname {Hom}_{\mathcal {C}}(Y,X_i)     = \operatorname {Hom}_{\operatorname {Pro}(\mathcal {C})}(Y,\mathcal {X})   \] 可由 \(\mathcal {L}\) 表示，即同构于 \(\operatorname {Hom}_{\mathcal {C}}(Y,\mathcal {L})\)；因此 \emph{\(\lim {}_{\mathcal {C}}X_i\) 已经由 \emph{pro 对象} \(\mathcal {X}\) 定义}，并且，只要有定义，它便以明确的函子性方式依赖于 pro 对象 \(\mathcal {X}\)；所以将其记为 \(\lim {}_{\mathcal {C}}(\mathcal {X})\) 并无问题。若 \(\mathcal {C}\) 中射影极限总存在，则 \(\lim {}_{\mathcal {C}}(\mathcal {X})\) 是从 \(\operatorname {Pro}(\mathcal {C})\) 到 \(\mathcal {C}\) 的函子，并且有一个典范的函子同态 \(\lim _\mathcal {C}(\mathcal {X})\to \mathcal {X}\)。由于每个函子 \[     F\colon  \mathcal {C} \to  \mathcal {C}'   \]（为简便起见，设其协变）都能显然地延拓为函子 \[     \operatorname {Pro}(F)\colon  \operatorname {Pro}(\mathcal {C}) \to  \operatorname {Pro}(\mathcal {C}'),   \]，所以若 \(\mathcal {C}'\) 中射影极限总存在，则 \(F\) 还典范地定义一个复合函子 \[     \overline {F}     = \varprojlim {}_{\mathcal {C}'}\circ\operatorname {Pro}(F)\colon  \operatorname {Pro}(\mathcal {C}) \to  \mathcal {C}'   \]，将 \(\mathcal {X}=(X_i)_{i\in  I}\) 送到 \(\varprojlim {}_{\mathcal {C}'}F(X_i)\)。\par{}若 pro 对象 \(\mathcal {X}\) 同构于一个过渡态射 \(X_i\to  X_j\) 均为 \emph{满态射} 的 pro 对象 \((X_i)_{i\in  I}\)，就称其为 \emph{严格 pro 对象}；由这种对象定义的函子称为 \emph{严格 pro 可表} 的。此时还可要求 \(I\) 是滤过的 \emph{有序} 集，并要求每个满态射 \(X_i\to  X'\) 都等价于某个适当的 \(j\in  I\) 所对应的满态射 \(X_i\to  X_j\)，而此条件唯一确定这个指标。在这些条件下，射影系 \((X_i)_{i\in  I}\) 在 \emph{唯一同构} 的意义下确定（这里采用射影系同构的通常意义）。由此可得 \emph{严格 pro 对象的射影系 \(\mathcal {X}^{(\alpha )}\) 的射影极限在 \(\operatorname {Pro}(\mathcal {C})\) 中总存在}，并且沿用上述 \(F\) 和 \(\overline {F}\) 的记号，有 \[     \overline {F}\varprojlim _\alpha \mathcal {X}^{(\alpha )}     = \varprojlim _\alpha {}_{\mathcal {C}'}F(X^{(\alpha )}).   \]。特别地，若 \(\mathcal {C}\) 的每个 pro 对象都严格（参见下一节），则延拓函子 \(\overline {F}\) 与射影极限可交换。\section{pro 可表函子的刻画}\label{fga3.ii-a.3}\par{}设 \(\mathcal {C}\) 和 \(\mathcal {C}'\) 为一切有限射影极限均存在的范畴（即以有限但不必滤过的预序集为指标的极限存在）；等价地，有限乘积与有限纤维积均存在。特别地，这蕴含存在一个“右单位对象” \(e\)，使得对每个 \(X\)，\(\operatorname {Hom}(X,e)\) 都只有一个元素。设 \(F\) 为从 \(\mathcal {C}\) 到 \(\mathcal {C}'\) 的协变函子。则下列条件等价：\begin{enumerate}
    
 \item[i.]\(F\) 与有限射影极限可交换；


    
 \item[ii.]\(F\) 与有限乘积及有限纤维积可交换；


    
 \item[iii.]\(F\) 与有限乘积可交换，并且对 \(\mathcal {C}\) 中的任意正合图式 \[         X\to  X'\rightrightarrows  X''       \]（参见 \hyperref[fga3.i-a.2-definition-2.1]{[3], A, 定义 2.1}），此图式在 \(F\) 下的像 \[         F(X)\to  F(X')\rightrightarrows  F(X'')       \] 是正合的。

  \end{enumerate} 此时称 \(F\) 是 \emph{左正合} 的。\par{}以下假设 \(\mathcal {C}\) 中有限射影极限总存在。由定义立即可知，可表函子是左正合的；再取极限可得 \emph{pro 可表函子是左正合的}。\par{}为得到逆命题，设 \[     F\colon  \mathcal {C} \to  \mathtt {Set}   \] 为协变函子，并设 \(X\in \mathcal {C}\)、\(\xi \in  F(X)\)。若对所有 \(X'\in \mathcal {C}\)、\(\xi '\in  F(X')\)，以及每个满足 \(\xi =F(u)(\xi ')\) 的严格单态射 \(u\colon  X'\to  X\)（参见 \hyperref[fga3.i-a.2]{[3], §A.2}），\(u\) 都是同构，则称 \(\xi \)（或二元组 \((X,\xi )\)）是 \emph{极小} 的。若存在态射 \(v\colon  X\to  X''\) 使得 \(\xi ''=F(v)(\xi )\)，则又称二元组 \((X,\xi )\) \emph{支配} \((X'',\xi '')\)，其中 \(\xi \in  F(X)\)、\(\xi ''\in  F(X'')\)；\emph{若 \(\xi \) 极小且 \(F\) 左正合，则这个态射 \(v\) 唯一}；\emph{若 \(\xi ''\) 极小，则 \(v\) 为满射}。由此容易推出如下命题：\begin{proposition}{3.1}\setnumber{3.1}\label{fga3.ii-a.3-proposition-3.1}\skipsinglepar\par{}\(F\) 严格 pro 可表的充要条件是满足以下两个条件：\begin{enumerate}
      
 \item[i.]\(F\) 左正合；

      
 \item[ii.]每个满足 \(\xi \in  F(X)\) 的二元组 \((X,\xi )\) 都被某个 \emph{极小} 二元组支配。

    \end{enumerate}\end{proposition}\par{}若 \(\mathcal {C}\) 的每个对象都是 Artin 的，则第二个条件显然成立：在那些存在 \(\xi '\in  F(X')\) 使得 \(\xi \) 是 \(\xi '\) 的像的子对象中，取 \(X\) 的一个极小子对象 \(X'\) 即可。因此：\begin{corollary}{}\label{fga3.ii-a.3-proposition-3.1-corollary}\skipsinglepar\par{}设 \(\mathcal {C}\) 为一个所有对象均为 Artin 对象且一切有限射影极限均存在的范畴。那么从 \(\mathcal {C}\) 到 \(\mathtt {Set}\) 的 pro 可表函子恰好是左正合函子，并且它们实际上都严格 pro 可表。\end{corollary}\par{}最后这个事实还蕴含 \emph{此时 \(\mathcal {C}\) 的每个 pro 对象都严格}。\section{例：Galois 型群与 pro 代数群}\label{fga3.ii-a.4}\par{}若 \(\mathcal {C}\) 为通常的有限群范畴，则 \(\operatorname {Pro}(\mathcal {C})\) 等价于全不连通紧拓扑群范畴。（\emph{［英译注］此处隐含 Hausdorff 条件。}）这种群及其推广——将通常的有限群换为给定基预概形上的有限群概形，例如域 \(k\) 上的有限代数群——充当预概形的基本群、同伦群以及绝对和相对同调群。在所有这些例子中，\hyperref[fga3.ii-a.3-proposition-3.1]{命题~\ref*{fga3.ii-a.3-proposition-3.1}} 的推论均适用，而 \(\pi _1\) 确实应通过其伴随函子来定义\hyperref[fga2]{[2]}。若从域上（更一般地，从 Noether 预概形上）的代数群或拟代数群范畴出发，也同样如此：我们得到 Serre \hyperref[Ser1958]{[4]} 的“\emph{pro 代数群}”。\section{例：“形式簇”}\label{fga3.ii-a.5}\par{}设 \(\Lambda \) 为 Noether 环，\(\mathcal {C}\) 为那些作为 \(\Lambda \) 上有限型模且为 Artin 模的 \(\Lambda \) 代数所组成的范畴，简称为 \emph{Artin \(\Lambda \) 代数}。于是 \hyperref[fga3.ii-a.3-proposition-3.1]{命题~\ref*{fga3.ii-a.3-proposition-3.1}} 的推论的条件成立。此时，范畴 \(\operatorname {Pro}(\mathcal {C})\) 等价于 \(\Lambda \) 上那些同构于代数 \(O_i\in \mathcal {C}\) 的拓扑射影极限 \[     O     = \varprojlim  O_i   \] 的 \emph{拓扑代数} \(O\) 所组成的范畴；即其拓扑 \emph{线性}、\emph{分离} 且 \emph{完备}，并且对 \(O\) 的每个开理想 \(\mathfrak {J}_i\)，代数 \(O/\mathfrak {J}_i\) 都是 \(\Lambda \) 上的 \emph{Artin} 代数。与这种代数相伴的函子 \(\mathcal {C}\to \mathtt {Set}\) 恰为 \[     \begin {aligned}       F(A)       &= h'_{O}(A)     \\&= \{\text {拓扑 }\Lambda \text { 代数的连续同态 }O\to  A\}     \\&= \varinjlim _i \operatorname {Hom}_{\Lambda \text { 代数}}(O_i,A).     \end {aligned}   \]。还要指出，范畴 \(\mathcal {C}\) 本质上是一些类似范畴的乘积，这些范畴对应于 \(\Lambda \) 的极大理想 \(\mathfrak {m}\) 处的局部环 \(\Lambda _{\mathfrak {m}}\) 的完备化；所以如有需要，可以只考虑 \(\Lambda \) 为这种完备局部环的情形。无论如何，\(O\) 典范地分解为其 \emph{局部分支} 的拓扑乘积，这些分支对应于 \(O\) 所定义的 \emph{形式概形}\hyperref[fga2]{[2]}的“点”。这样的点由某个 \(F(K)\) 的对象 \(\xi \) 定义，其中 \(K\in \mathcal {C}\) 为一个 \emph{域}，例如相应局部分支的剩余域；两个二元组 \((\xi ,K)\) 与 \((\xi ',K')\) 定义同一点，当且仅当它们均被同一个 \((\xi '',K'')\) 支配，或它们均支配同一个 \((\xi ''',K''')\)。（若各个 \({\Lambda }/\mathfrak {m}\) 均代数闭，则只需取各 \(F({\Lambda }/\mathfrak {m})\) 的不交并所得的集合。）\par{}重要的是给出条件，保证 \(O\) 中对应于某个 \(\xi \in  F(K)\) 的局部分支 \(O_\xi \) 是 \emph{Noether} 环。若 \(\Lambda \) 为完备局部环（仍假设其为 Noether 环），则这等价于 \(O_\xi \) 同构于 \emph{形式幂级数环 \(\Lambda [{[t_1,\ldots ,t_n]}]\) 的商环}。为给出这种判别法，对每个环 \(A\) 引入其“双数” \(A\) 代数 \(I_A\)，这里的双数相对于 \(A\) 定义为 \[     I_A     = A[t]/t^2A[t].   \]。令 \(\varepsilon \colon  I_A\to  A\) 为增广同态；若 \(A\in \mathcal {C}\)，它定义一个映射 \[     F(\varepsilon )\colon  F(I_A) \to  F(A).   \]。利用 \(F\) 的左正合性，可以在由“模 \(\xi \) 约化”的 \(\xi '\in  F(I_A)\) 组成的子集 \[     F(I_A,\xi )     = F(\xi )^{-1}(\xi ) \subset  F(I_A)   \] 上内在地定义 \(A\) 模结构；借助 \(F\) 用 \(O\) 表示的具体形式，可知这个 \(A\) 模可等同于 \(\operatorname {Hom}_\Lambda (\mathfrak {m}_{\xi }/\mathfrak {m}_\xi ^2,A)\)，其中 \(m_\xi \) 是同态 \(\xi \colon  O\to  A\) 的核；也就是说，若 \(A\) 是域，则它是 \(O\) 的局部分支 \(O_\xi \) 的极大理想。由此立即得到如下命题：\begin{proposition}{5.1}\setnumber{5.1}\label{fga3.ii-a.5-proposition-5.1}\skipsinglepar\par{}设 \(\xi \in  F(K)\)，其中 \(K\in \mathcal {C}\) 是域。\(O\) 中相应的局部分支 \(O_\xi \) 为 \emph{Noether} 环的充要条件是：\(F(I_K)\) 中模 \(\xi \) 约化的元素所组成的集合 \(F(I_K,\xi )\) 是 \(K\) 上的 \emph{有限维} 向量空间。在这些条件下，有典范同构 \[       F(I_K,\xi )       = \operatorname {Hom}(\mathfrak {m}_{\xi }/\mathfrak {m}_\xi ^2+\mathfrak {n}_\xi \mathscr {O}_\xi , K)     \]，其中 \(\mathfrak {n}_\xi \) 是 \(\Lambda \) 的极大理想，即同态 \(\Lambda \to  K\) 的核。因此，特别地，\(K\) 向量空间 \(F(I_K,\xi )\) 的维数等于向量空间 \(\mathfrak {m}_{\xi }/\mathfrak {m}_\xi ^2\) 在域 \(O_{\xi }/\mathfrak {m}_\xi =K(\xi )\) 上的维数。\par{}\emph{［编者补注］} 上面的公式仅在 \(\Lambda \) 为域时正确；一般情形下，须将 \(\mathfrak {m}_{\xi }/\mathfrak {m}_\xi ^2\) 换为此空间对 \(\mathfrak {n}_{\xi }/\mathfrak {n}_\xi ^2\) 的像的商，其中 \(\mathfrak {n}\) 是 \(\Lambda \) 的极大理想。\end{proposition}\par{}假设 \(O_\xi \) 为 Noether 环，并为简化记号，设 \(\Lambda \) 完备且局部，并且 \(O=O_\xi \)。（\emph{［编者补注］} 以下定义仅在剩余扩张 \(k'/k\) \emph{可分} 时正确；一般情形见［\hyperref[Gro1960b]{Gro1960b}，III，1.1］。）若 \(O\) 是如下完备化代数上的有限 étale 代数，则称 \emph{\(O\) 在 \(\Lambda \) 上光滑}：先取 \(\Lambda [t_1,\ldots ,t_n]\) 的一个诱导 \(\Lambda \) 的极大理想的极大理想，在该处局部化，再作完备化。若 \(O\) 在 \(\Lambda \) 上的剩余扩张平凡，例如 \(\Lambda \) 的剩余域代数闭，则这等价于 \(O\) 本身同构于这样一个形式幂级数代数。最后，若不再假设 \(O\) 必为 Noether 环，只要 \(O\) 同构于一些 Noether 且在上述意义下 \(\Lambda \) 光滑的商 \(\Lambda \) 代数的拓扑射影极限，仍称 \emph{\(O\) 在 \(\Lambda \) 上光滑}。这立即推广到不再假设 \(\Lambda \) 和 \(O\) 局部的情形。于是有如下命题：\begin{proposition}{5.2}\setnumber{5.2}\label{fga3.ii-a.5-proposition-5.2}\skipsinglepar\par{}\(O\) 在 \(\Lambda \) 上光滑的充要条件是伴随函子 \(F\) 将满态射送到满态射。\end{proposition}\par{}这就是说，对于 \(\mathcal {C}\) 中的每个 \emph{满射} 同态 \(A\to  A'\)，态射 \(F(A)\to  F(A')\) 也是 \emph{满射}。当然，只需在 \(A\) 为 \emph{局部环} 的情形下检验此条件，并且还可逐步归约到同态 \(A\to  A'\) 的核所给出的 \(A\) 的理想被 \(A\) 的极大理想零化的情形。在实践中，这就归结为验证某个障碍为零；该障碍与给出函子 \(F\) 的具体情形的 \emph{无穷小} 不变量有关，这是一个 \emph{上同调} 性质的问题。\par{}最后，在上述框架中略谈 \emph{定义环}。仍设 \(F\) 为从 \(\mathcal {C}\) 到 \(\mathtt {Set}\) 的函子，并假设其由拓扑 \(\Lambda \) 代数 \(O\) pro 表示。那么对每个 \(A\in \mathcal {C}\) 和每个 \(\xi \in  F(A)\)，\(A\) 中存在一个 \emph{最小} 子环 \(A'\)，使得 \(\xi \) 是 \(F(A')\) 中某个元素 \(\xi '\) 的像，且该元素唯一确定。事实上，只需将 \(\xi \) 看作从 \(O\) 到 \(A\) 的同态，并取 \(A'\) 为 \(O\) 在此 \(\xi \) 下的像。于是称 \(A'\) 为 \emph{对象 \(\xi \in  F(A)\) 的定义环}。若 \(u\colon  A\to  B\) 为代数同态，且 \(\eta =F(u)(\xi )\)，则 \(\eta \) 的定义环是 \(\xi \) 的定义环在 \(u\) 下的像。若从一个从 \(\mathcal {C}\) 到 \(\mathtt {Set}\) 的函子 \(F\) 出发，则定义环的存在性以及刚才所述的性质，差不多与 \(F\) pro 可表的条件 \emph{等价}；也就是说，这些性质通常远非显然。\chapter{两个存在定理}\thispagestyle{fancy}\label{fga3.ii-b}\par{}沿用 \cref{fga3.ii-a.5} 的记号，给定协变函子 \[     F\colon  \mathcal {C} \to  \mathtt {Set},   \]，我们要寻找便于运用的判别法，判定 \(F\) pro 可表，即可通过上述形式的 \(\Lambda \) 代数 \(O\) 表达。根据 \hyperref[fga3.ii-a.3-proposition-3.1]{§A, 命题 3.1} 的推论，这成立的充要条件是 \(F\) 为 \emph{左正合}。在下降技术的目前发展阶段（参见 \hyperref[fga3.i-a.2.c]{[3], §A.2.c} 所提出的问题），对最重要的情形，这个形式的判别法尚不能直接验证，所以需要看起来要求较弱的判别法。\begin{theorem}{1}\setnumber{1}\label{fga3.ii-b-theorem-1}\skipsinglepar\par{}函子 \(F\) pro 可表的充要条件是满足以下两个条件：\begin{enumerate}
      
 \item[i.]\(F\) 与有限乘积可交换；

      
 \item[ii.]对每个代数 \(A\in \mathcal {C}\) 以及 \(\mathcal {C}\) 中使图式 \[           A \to  A' \rightrightarrows  A'\otimes _A A'         \] 正合的同态 \(A\to  A'\)（参见 \hyperref[fga3.i-a.1-definition-1.2]{[3], §A, 定义 1.2}），图式 \[           F(A) \to  F(A') \rightrightarrows  F(A'\otimes _A A')         \] 也正合。

    \end{enumerate} 此外，条件（ii）只需在 \(A\) 为局部环、并且还满足下列两种情形之一时检验：\begin{enumerate}
      
 \item[a.]\(A'\) 是 \(A\) 上的 \emph{自由} 模；

      
 \item[b.]商模 \(A'/A\) 是 \emph{长为 \(1\)} 的 \(A\) 模。

    \end{enumerate}\begin{proof}\par{}此定理的证明颇为精细，无法在此勾勒。这里只指出，它主要依赖于对 \emph{Artin 代数的谱} 中的 \emph{等价关系}（范畴意义下）的研究；这种研究还提出了更多问题，它们的解决似乎是理论进一步发展所必需的。\end{proof}\end{theorem}\par{}在应用中，条件（i）的验证总是显然的。条件（ii）的验证分为两种情形：（a）当 \(A'\) 为自由 \(A\) 模时，可以用 \emph{沿平坦态射的下降技术}（参见 \hyperref[fga3.i-b.1]{[3]，B.1 的定理 1、2 和 3}）处理，这没有困难；为了处理情形（b），将使用如下结果：\begin{theorem}{2}\setnumber{2}\label{fga3.ii-b-theorem-2}\skipsinglepar\par{}设 \(A\) 为极大理想为 \(\mathfrak {m}\) 的 Artin 局部环，\(A'\) 为包含 \(A\) 的 \(A\) 代数，满足 \(\mathfrak {m}A'\subset  A\)，且 \(A\to  A'\rightrightarrows  A'\otimes _A A'\) \emph{正合}；特别地，当 \(A'/A\) 为长为 \(1\) 的 \(A\) 模时即如此。设 \(\mathcal {F}\) 为随预概形变化、由在其上平坦的拟凝聚层组成的纤维范畴（参见 \hyperref[fga3.i-a.1-definition-1.1]{[3], §A, 定义 1.1}）。则态射 \(\operatorname {Spec}(A')\to \operatorname {Spec}(A)\) 为 \emph{严格 \(\mathcal {F}\) 下降态射}（参见 \hyperref[fga3.i-a.1-definition-1.7]{[3], §A, 定义 1.7}）。\begin{proof}\par{}证明时先证明 \[         \operatorname {H}^i(A'/A,\operatorname {G_a}) = 0         \qquad \text {其中 }i\geqslant 1       \]（参见 \hyperref[fga3.i-a.4.e]{[3], §A.4.e}）；假设 \(\mathfrak {m}A'\subset  A\) 使我们容易归约到 \(A\) 为域（即 \(A/\mathfrak {m}\)）的情形。随后即可应用 \hyperref[fga3.i-a.4.e]{[3], §A.4.e} 所述的等价。\end{proof}\end{theorem}\par{}换言之，给出平坦 \(A\) 模 \(M\)，完全等价于给出平坦 \(A'\) 模 \(M'\)，连同从 \(M'\otimes _A A'\) 到 \(A'\otimes _A M'\) 的 \((A'\otimes _A A')\) 同构，并且此同构满足下降资料的通常传递性条件（\emph{同前引处}）。\chapter{对若干特殊情形的应用}\thispagestyle{fancy}\label{fga3.ii-c}\section{关于由预概形表示的函子的一般说明}\label{fga3.ii-c.1}\par{}设 \(S\) 为局部 Noether 预概形。若对每个投到 \(y\in  Y\) 的 \(x\in  X\)，均存在 \(y\) 的一个环为 \(A\) 的仿射邻域，以及 \(x\) 的一个环为 \(B\) 的仿射邻域，后者位于上述 \(y\) 的仿射邻域之上，使得 \(B\) 是有限型 \(A\) 代数，则称 \(S\) 上的预概形 \(X\) 在 \(S\) 上为 \emph{局部有限型} 的。经典泛问题的解给出了许多 \(S\) 上局部有限型但在 \(S\) 上并非有限型的预概形的重要例子；因此必须能把曲线的 Picard 概形看作无穷多个连通分支的并，而不能将其与单位元的连通分支，即“Picard 簇”混淆。所以有时在 \(S\) 上局部有限型预概形范畴 \(\mathcal {C}\) 中研究反变函子 \(F\) 的可表性问题是有用的。\emph{这些文章的主要目标，是发展一种一般技术，用以判别这样的函子 \(F\) 何时可表，并通过 \(F\) 的性质研究相应 \(S\) 预概形 \(X\) 的性质。} 顺便指出，在此研究中，会遇到 \(S\) 上并不相对于 \(S\) 分离的预概形的非病态例子，特别是很好的 \(S\) 概形的“Picard 预概形”；因此，不应把那些不是概形的预概形逐出代数几何。\par{}设 \(X\) 为 \(S\) 上局部有限型预概形，\[     F\colon  Y \mapsto  \operatorname {Hom}_S(Y,X)   \] 为其伴随反变函子。考虑 \(F\) 在 \(\mathcal {C}\) 的子范畴 \(\mathcal {C}_0\) 上的限制 \(F_0\)；该子范畴由 \(S\) 上既为 Artin 预概形又在 \(S\) 上有限的预概形 \(Y\) 组成。若 \(S=\operatorname {Spec}(\Lambda )\)，则 \(\mathcal {C}_0\) 是 \cref{fga3.ii-b} 中所考虑的 Artin \(\Lambda\) 代数范畴的对偶范畴。若 \(Y=\operatorname {Spec}(A)\)，其中 \(A\) 为 \emph{局部} Artin 环，则 \(Y\) 仅含一个位于 \(S\) 的某个闭点 \(s\) 上方的点 \(y\)；从 \(Y\) 到 \(X\) 的 \(S\) 同态，即 \(F(Y)\) 的元素，由一个位于 \(s\) 上方的点 \(x\in  X\)，以及一个从 \(\mathscr {O}_x\) 到 \(A\) 的 \(\mathscr {O}_s\) 同态给定。若这样的同态存在，则 \(x\) 必为 \(X\) 的闭点，因为其剩余域在 \(s\) 的剩余域上代数。因此可知，\emph{\(F\) 在“Artin \(Y\) 代数”上的限制 \(F_0\) pro 可表，且由一个拓扑 \(Y\) 代数表示；其局部分支是 \(X\) 在 \(X\) 的那些位于 \(Y\) 的闭点上方的闭点 \(x\) 处的局部环的完备化 \(\widehat {\mathscr {O}_x}\)}。这表明，仅知道 \(F_0\) 就能得到关于 \(X\) 的结构的精确信息，即其在上述各点处的局部环的完备化的结构。要指出，即使 \(S\) 是代数闭域的谱，也只有系统地考虑那些 \(\mathscr {O}_Y\) 可以含有幂零元的“簇” \(Y\)，特别是使用 Artin 局部环的谱，才能得到经典泛问题的“恰当表述”，并理解其“无穷小”方面。\par{}若从给定函子 \(F\) 出发，想知道它是否可表，则用 \hyperref[fga3.ii-b-theorem-1]{定理~\ref*{fga3.ii-b-theorem-1}} 与 \hyperref[fga3.ii-b-theorem-2]{定理~\ref*{fga3.ii-b-theorem-2}} 研究函子 \(F_0\)，便能提供近乎完整的线索。其一，正如经常发生的那样，仅检验集合 \(F(I_K,\xi )\) 的性质及其函子性行为（参见 \cref{fga3.ii-a}），就能发现 \(F_0\) 已经不 pro 可表；这解释了迄今为对秩 \(>1\) 的向量丛分类而以足够自然的方式定义模簇的尝试何以失败。其二，也许能够证明 \(F_0\) 确实可表，但向量空间 \(F(I_K,\xi )\) 并非有限维，此时只能满足于“形式”解。其三，\(F_0\) 确实可由完备 Noether 局部环的乘积表示，这就强烈提示 \(F\) 本身也可表；结合以后将要发展的类似但更具整体性的性质，很可能就足以推出它确实可表。最后，还会遇到有趣的几何问题（见下面 \cref{fga3.ii-c.4} 和 \cref{fga3.ii-c.5}），其中我们只有函子 \(F_0\)，它并不来自任何“整体”函子 \(F\)；此时，能为它配备一个“形式模概形”就足够了。\par{}为结束这些一般讨论，指出概形理论如何解释某些表面上的反常现象，例如 Igusa 曲面 \(V\) 的“Picard 簇” \(P\) 只有一个点，而 \(\operatorname {H}^1(V,\mathscr {O}_V)\neq 0\)。此时 \(P\) 是一个非平凡的“纯无穷小”群，即由一个在基域 \(k\) 上秩有限的局部代数 \(\mathscr {O}\) 定义，并赋有对应于 \(P\) 的乘法结构的对角映射；若 \(\mathfrak {m}\) 为 \(\mathscr {O}\) 的极大理想，则 \(\mathfrak {m}/\mathfrak {m}^2\) 的对偶典范同构于 \(\operatorname {H}^1(V,\mathscr {O}_V)\)（参见下面 \cref{fga3.ii-c.3}）。只有当 Picard 群是经典意义下的代数群，即在基域 \(k\) 上光滑时，\(\operatorname {H}^1(V,\mathscr {O}_V)\) 的维数——它总等于 \(\mathfrak {m}/\mathfrak {m}^2\) 的维数——才等于 Picard 群的维数。\section{\texorpdfstring{概形 \(\underline {\operatorname {Hom}}_S(X,Y)\)、\(\prod _{X/S}Z\)、\(\underline {\operatorname {Aut}}(X)\) 等}{概形 Hom-S(X,Y)、Π-X/S Z、Aut(X) 等}}\label{fga3.ii-c.2}\par{}设 \(X\) 与 \(Y\) 为 \(S\) 上的预概形；对 \(S\) 上的每个预概形 \(T\)，令 \(X_T=X\times _S T\)、\(Y_T=Y\times _S T\)，并将集合 \[     F(T)     = \operatorname {Hom}_T(X_T,Y_T)     = \operatorname {Hom}_S(X_T,Y)     = \operatorname {Hom}_S(X\times _S T,Y)   \] 看作关于 \(T\) 的反变函子。若它可表，则用 \(\underline {\operatorname {Hom}}_S(X,Y)\) 表示代表它的 \(S\) 上预概形，此时有函子性同构 \[     \operatorname {Hom}_S(T,\underline {\operatorname {Hom}}_S(X,Y)) \xrightarrow {\sim } \operatorname {Hom}_S(T\times _S X,Y).   \]。这个泛问题有一些变式，其解可概括如下：\(S\) 预概形 \(X\) 的 \emph{\(S\) 自同构} 的预概形，它是 \emph{群} 预概形；从一个 \(S\) 上 \emph{群} 预概形到另一个的 \emph{\(S\) 同态} 的预概形，若后一个群概形交换，则它将是交换群预概形；等等。还可以推广 \(\underline {\operatorname {Hom}}_S(X,Y)\) 的定义：考虑 \(S\) 上的预概形 \(X\) 上的一个预概形 \(Z\)，以及函子 \[     F(T)     = \operatorname {Hom}_{X_T}(X_T,Z_T)   \]，即 \(Z_T\) 在 \(X_T\) 上的“截面”集。若该函子可表，将表示它的 \(S\) 预概形记为 \(\Pi _{X/S}Z\)，于是按定义有函子性同构 \[     \operatorname {Hom}_S(T,\Pi _{X/S}Z)     = \operatorname {Hom}_{X_T}(X_T,Z_T).   \]。置 \(Z=Y\times _S X\)，便恢复 \(\underline {\operatorname {Hom}}_S(X,Y)\)。由这些定义可得关于如此引入的新预概形的一套公式，它既显然又有用，兹不列出，因为它在一切乘积与纤维积存在的范畴中均成立。更实质的问题是 \(\underline {\operatorname {Hom}}_S(X,Y)\) 型概形的 \emph{存在性}。首先指出，对固定的 \(X\)，只有当 \(X\) 在 \(S\) 上 \emph{平坦} 时，\(\underline {\operatorname {Hom}}_S(X,Y)\)（相应地，\(\Pi _{X/S}Z\)）才可能对 \(S\) 上的所有 \(Y\)（相应地，对 \(X\) 上的所有 \(Z\)）都存在。此外，容易相信，要对足够一般的 \(Y\) 合理地期待解的存在，还须 \(X\) 在 \(S\) 上 \emph{固有}。不过，这些条件似乎足以保证 \(\underline {\operatorname {Hom}}_S(X,Y)\) 与 \(\Pi _{X/S}Z\) 存在，必要时对 \(Y/S\)（相应地，\(Z/X\)）再作某种“拟射影”假设。至少在许多情形下可验证这一点，例如 \(Y\) 在 \(S\) 上 \emph{仿射} 时；又如 \(X\) 在 \(S\) 上 \emph{有限} 时，可用直接的初等构造验证。于是 \hyperref[fga3.ii-b-theorem-1]{定理~\ref*{fga3.ii-b-theorem-1}} 与 \hyperref[fga3.ii-b-theorem-2]{定理~\ref*{fga3.ii-b-theorem-2}} 给出：\begin{proposition}{2.1}\setnumber{2.1}\label{fga3.ii-c.2-proposition-2.1}\skipsinglepar\par{}设 \(\Lambda \) 为 Noether 环，\(X\)、\(Y\) 为 \(\Lambda \) 上的任意预概形。在 Artin \(\Lambda \) 代数范畴 \(\mathcal {C}_0\) 上考虑函子 \[       F(A)       = \operatorname {Hom}_A(X_A,Y_A)     \]。若 \(X\) 在 \(\Lambda \) 上平坦，则此函子 pro 可表。\end{proposition}\par{}此外，可以证明，对所有 \(A\in \mathcal {C}_0\) 及所有 \(\xi \in  F(A)\)，都有典范同构 \[     F(I_A,\xi )     = \operatorname {H}^1\Big (X_A,\underline {\operatorname {Hom}}_{\mathscr {O}_{X_A}}\big (\xi ^*(\Omega _{Y_A/A}^1),\mathscr {O}_{X_A}\big )\Big )   \]，其中 \(\Omega _{Y_A/A}^1\) 是 \(Y_A\) 相对于 \(A\) 的 Kähler \(1\) 次微分层。取 \(A\) 为域，由 \hyperref[fga3.ii-a.5-proposition-5.1]{§A, 命题 5.1} 以及 \hyperref[fga2]{[2]} 的有限性定理得到如下推论：\begin{corollary}{}\label{fga3.ii-c.2-proposition-2.1-corollary}\skipsinglepar\par{}假设 \(X\) 在 \(S\) 上平坦且固有，并且 \(Y\) 在 \(S\) 上局部有限型。则 \(F\) pro 可表，且相应拓扑 \(\Lambda \) 代数的局部分支均为 \emph{Noether} 环。\end{corollary}\begin{remarks}{}\label{fga3.ii-c.2-remarks}\skipsinglepar\par{}本节考虑的问题及其他许多问题，在经典代数几何框架内，通常通过射影空间中闭链的“Chow 坐标”研究，从而可将这些闭链看作适当射影簇的点。这种方法，更一般地，Chow 坐标的使用，从概形的角度看似乎有不可弥补的不足，因为它消除了参数化簇中的幂零元，特别是不能令人满意地研究闭链的无穷小变动；这还没有计及它依赖于射影空间、因而并非内在的性质。尽管 Chow 坐标确有某种技术价值，很可能只是暂时的，但不幸的是，许多代数几何学家研究簇族或闭链族时只使用这种语言，这似乎严重妨碍了对这些概念的理解。若想在概形理论中得到 Chow 簇的对应物，就会导向如下泛问题：设 \(X\) 为 \(S\) 上的预概形，对 \(S\) 上的每个预概形 \(T\)，考虑 \(X_T=X\times _S T\) 中那些在 \(T\) 上 \emph{平坦} 的闭子预概形组成的集合 \(F(T)\)；我们要用 \(S\) 上的某个预概形表示这个关于 \(T\) 的函子。更一般地，可以从 \(X\) 上的拟凝聚层 \(\mathscr {G}\) 出发，将 \(F(T)\) 取为 \(\mathscr {G}_T\) 的那些在 \(T\) 上平坦的商层组成的集合。若 \(X\) 在 \(S\) 上固有，\(S\) 局部 Noether，且 \(F\) 还凝聚，则此问题似乎有解，且给出在 \(S\) 上局部有限型的概形 \(C\)。无论如何，只假设 \(S\) 局部 Noether，\(F\) 在“Artin \(S\) 代数”上的限制就 pro 可表；若还假设 \(X\) 在 \(S\) 上固有且 \(F\) 凝聚，则相应拓扑环 \(\mathscr {O}\) 的局部分支为 Noether 环。当然，即使证明了 \(X\) 在 \(S\) 上的“Chow 概形”的存在性，仍须将其分解为 \(S\) 上互不相交的开子集 \(C_i\)，这对应于固定所变动闭链的次数、维数等连续不变量，并要求这些开子集在 \(S\) 上为有限型，而不只是局部有限型；还须阐明此概形与经典 Chow 簇的关系，并精确说明何时 \(C_i\) 在 \(S\) 上 \emph{射影}，或至少拟射影。\end{remarks}\begin{remark}{}\label{fga3.ii-c.2-remark}\skipsinglepar\par{}\emph{［编者补注］} 射影情形下，这里所述的问题已由“Hilbert 概形”彻底解决（参见 \hyperref[fga3.iv]{FGA~\ref*{fga3.iv}}）。但 Nagata 与 Hironaka 的例子表明，若不作射影性假设，所述函子未必可表，即使只对一个 \(3\) 维完备非奇异簇的 \(0\) 维子簇进行分类也是如此；这与这样的簇的对称平方未必存在有关。\end{remark}\section{Picard 概形}\label{fga3.ii-c.3}\par{}\emph{［编者补注］} 更完整的研究见 \hyperref[fga3.v]{FGA~\ref*{fga3.v}}。\par{}设 \(f\colon  X\to  S\) 为 \(S\) 预概形，考虑 \(X\) 的结构层的单位所组成的乘法层 \(\mathscr {O}_X^\times \)，以及群 \[     P(X/S)     = \operatorname {H}^0(S,\operatorname {R}^1f_*(\mathscr {O}_X^\times )),   \]，称为 \(X/S\) 的 \emph{相对 Picard 群}。此群的一个元素由如下资料定义：\(S\) 的一个开覆盖 \((U_i)\)，以及每个 \(f^{-1}(U_i)\) 上的一个可逆层 \(\mathscr {L}_i\)，使得对所有 \(i,j\)，\(\mathscr {L}_i|f^{-1}(U_i\cap  U_j)\) 同构于 \(\mathscr {L}_j|f^{-1}(U_i\cap  U_j)\)，或至少在 \(U_i\cap  U_j\) 上局部同构；也就是说，这两个层在 \hyperref[fga3.i-b.4]{[3], §B.4} 的意义下“等价”。若 \(X/S\) 有一个截面，则 \(P(X/S)\) 正是 \(X/S\) 上的可逆层模去“等价”后的类集（\emph{同前引处}）。现在对 \(S\) 上的所有 \(T\)，置 \[     F(T)     = P(X_T/T)   \]，这给出关于 \(T\) 的协变函子，称为 \(X/S\) 的 \emph{Picard 函子}；若此函子可表，则将表示它的 \(S\) 上的预概形称为 \(X/S\) 的 \emph{Picard 预概形}，记为 \(\mathscr {P}(X/S)\)。此时有函子同构 \[     \operatorname {Hom}_S(T,\mathscr {P}(X/S)) \xrightarrow {\sim } P(X_T/T).   \]。取 Picard 预概形与基底扩张相容；特别地，\(X\) 在 \(S\) 上各纤维的 Picard 预概形，即各个 \(s\in  S\) 的剩余域 \(K(s)\) 上的预概形，就是 \(\mathscr {P}(X/S)\) 的纤维。当然，由于 \(P(X_T/T)=F(T)\) 是交换群，Picard 预概形是群预概形。还应指出，Rosenlicht 的 \emph{广义 Jacobian} 正是完备曲线（可能带有奇点）的 Picard 概形中单位元的连通分支；一旦存在性得到证明，它们的大部分性质也就应当清楚了。\begin{remark}{}\label{fga3.ii-c.3-remark}\skipsinglepar\par{}此处采用的定义，只有当 \(Y\) 的每个点都有一个开邻域 \(U\)，使 \(X\) 在其上有截面时才合理。一般情形下，须略微修改 Picard 函子的定义，才能仍然得到存在定理。\par{}这里，Picard 预概形存在的合理条件如下：\(X\) 在 \(S\) 上 \emph{固有} 且 \emph{平坦}；\(f_*(\mathscr {O}_X)=\mathscr {O}_S\)；以及 \(X\) 在 \(S\) 上局部有截面。应用下降技术时，此条件自然出现，\emph{因为可以在 \(X\) 上的可逆层 \(\mathscr {L}\) 上赋予截面 \(s\) 上方的一个标定截面，以消除它们的自同构}（\hyperref[fga3.i-b.4]{[3], §B.4}）。特别地，可得如下结果：\end{remark}\begin{proposition}{3.1}\setnumber{3.1}\label{fga3.ii-c.3-proposition-3.1}\skipsinglepar\par{}假设 \(X\) 在 \(S=\operatorname {Spec}(\Lambda )\) 上平坦，其中 \(\Lambda \) 为 Noether 环，并假设对 \(S\) 上的每个有限型 \(T\)，均有 \({f_T}_*(\mathscr {O}_{X_T})=\mathscr {O}_T\)。若 \(f\) 固有、可分且各纤维可分，或若 \(S\) 是域的谱，则由 Künneth 公式可知，最后一个条件等价于 \(f_*(\mathscr {O}_X)=\mathscr {O}_S\)。那么 \(X/S\) 在 Artin \(\Lambda \) 代数范畴上的 Picard 函子 pro 可表。\par{}此外，此时有 \[       F(I_A,\xi )       = \operatorname {H}^1(X_A,\mathscr {O}_{X_A}),     \]，特别地：\end{proposition}\begin{corollary}{}\label{fga3.ii-c.3-proposition-3.1-corollary}\skipsinglepar\par{}若 \(X\) 在 \(S\) 上固有，则与 Picard 函子相应的拓扑 \(\Lambda \) 代数的局部分支为 Noether 环。\end{corollary}\begin{remarks}{}\label{fga3.ii-c.3-remarks-i}\skipsinglepar\par{}本节的定义和结果可推广到 \(X\) 上主丛的分类，其中结构群 \(G\) 为 \(S\) 上的群概形，且 \emph{仿射}、\emph{在 \(S\) 上平坦} 并且 \emph{交换}。若 \(G\) 不交换，因而一个主丛的伴随群丛——其截面是该主丛的自同构——不再平凡，则 \hyperref[fga3.ii-c.3-proposition-3.1]{命题~\ref*{fga3.ii-c.3-proposition-3.1}} 不再按原陈述成立。不过，可以修改泛问题，使之重新有解，至少目前在形式几何中可以做到。\emph{在本节和以下各节中，每当要为仅在同构意义下定义的对象类寻找“模概形”时，都应牢记同一条准则：消除待分类对象可能具有的自同构；必要时，引入辅助结构，例如标定的点、截面的标定元素、指定微分形式等，但要使这些结构足够轻微，以免实质性改变原问题。}\end{remarks}\section{一个簇的形式模空间}\label{fga3.ii-c.4}\par{}设 \(\Lambda \) 为剩余域为 \(k\) 的 \emph{局部} Noether 环；通常 \(\Lambda \) 取为 \(k\)，或 Cohen \(p\) 环。设 \(X_0\) 为 \(k\) 上的预概形。对每个局部 Artin \(\Lambda \) 代数 \(A\)，考虑这样的 \(A\) 预概形 \(X\) 的同构类集 \(F(A)\)：它在 \(A\) 上 \emph{平坦}，并赋有同构 \begin{eqn}{*}\setnumber{*}\label{fga3.ii-c.4-equation-star}\skipsinglepar\[       X\otimes _A k(A) \xleftarrow {\sim } X_0\otimes _k k(A)     \tag{*}     \]\end{eqn}，其中 \(k(A)\) 是 \(A\) 的剩余域。当然，这种平坦 \(A\) 预概形之间的同构应保持结构中给定的上述同构。若 \(A\) 为不必局部的 Artin \(\Lambda \) 代数，局部分支为 \(A_i\)，则取 \(F(A)\) 为各个 \(F(A_i)\) 的乘积。于是 \(F\) 成为关于 \(A\) 的乘法函子，称为 \(X_0\)（及 \(\Lambda \)）的 \emph{模函子}。若此函子可表，则它对应于剩余域为 \(k\) 的一个局部拓扑 \(\Lambda \) 代数 \(O\)，而 \(O\) 的形式谱称为 \(X_0\)（及 \(\Lambda \)）的 \emph{形式模概形}；一些细节见 \hyperref[fga2]{[2]}。\par{}这里，若想应用下降技术，\(X_0\) 的“有限”自同构无妨，因为它们不影响 \(A\) 预概形 \(X\) 的自同构在上述精确意义下的存在性；当 \(A\) 不仅仅是域时，\(X\) 没有非平凡 \(A\) 自同构的充要条件是 \[     \operatorname {H}^0(X_0,\mathfrak {G}_{X_0/k})     = 0   \]，其中 \(\mathfrak {G}_{X_0/k}\) 表示 \(X_0\) 的 \(k\) 导子层，即切层。至少在 \(X_0\) 在 \(k\) 上光滑时，容易证明 \[     F(I_A,\xi )     = \operatorname {H}^1(X_A,\mathfrak {G}_{X_A/A}).   \]。于是照例得到：\begin{proposition}{4.1}\setnumber{4.1}\label{fga3.ii-c.4-proposition-4.1}\skipsinglepar\par{}假设 \(\operatorname {H}^0(X_0,\mathfrak {G}_{X_0/k})=0\)。则 \(X_0\) 的形式模概形存在。若此外 \(X_0\) 在 \(k\) 上固有，则此形式模概形为 Noether 概形。\end{proposition}\begin{remarks}{}\label{fga3.ii-c.4-remarks}\skipsinglepar\begin{enumerate}\item{}若不假设 \(X_0\) 在 \(k\) 上光滑，则 \(F(I_A,\xi )\) 可等同于 \[         \operatorname {Ext}_{\mathscr {O}_{P_A}}^1(P_A;\mathscr {I}_{X_A},\mathscr {O}_{X_A})       \] 的一个 \(A\) 子模，其中置 \(P_A=X_A\times _A X_A\)，并通过对角态射 \(X_A\to  P_A\) 将 \(\mathscr {O}_{X_A}\) 视为 \(P_A\) 上的凝聚层，而 \(\mathscr {I}_{X_A}\) 表示由对角态射定义的 \(P_A\) 上的凝聚理想层。更确切地说，将 Hochschild 理论作一个容易的整体化，即可知上述 \(\operatorname {Ext}^1\) 可等同于如下对象的同构类集：定义在 \(X_A\) 上、作为 \(I_A\) 代数平坦并赋有增广同构 \(\mathscr {O}\otimes _{I_A}A\to \mathscr {O}_{X_A}\) 的 \(I_A\) 代数层 \(\mathscr {O}\)；回顾我们曾置 \(I_A=At/(t^2)\)。子模 \(F(I_A,\xi )\) 对应于 \emph{交换} 代数层。因此，光滑性假设在模理论中并非本质，尽管 \hyperref[fga2]{[2]} 曾给人相反的印象。

    \item{}回顾（\emph{同前引处}），特别地，\(k\) 上的每个 \emph{光滑} 的 \emph{固有代数曲线} \(X_0\) 都有在 \(\Lambda \) 上光滑的形式模概形；当亏格 \(g\) 满足 \(\geqslant 2\) 时，其相对维数为 \(3g-3\)，当 \(g=0,1\) 时则为 \(g\)。后两种情形已不直接包含于 \hyperref[fga3.ii-c.4-proposition-4.1]{命题~\ref*{fga3.ii-c.4-proposition-4.1}} 中。不过，在椭圆曲线的情形，即 \(g=1\)，借助下面的说明可以恢复这些结论。\end{enumerate}\end{remarks}\par{}当然，可以通过考虑 \(k\) 上赋有各种结构的概形系，\emph{任意地} 改变 \hyperref[fga3.ii-c.4-proposition-4.1]{命题~\ref*{fga3.ii-c.4-proposition-4.1}}。例如，设 \(X_0\) 为 \(k\) 上带有标定原点的一个 \emph{abelian 概形}，即将 \(X_0\) 视为 \(k\) 上的 \emph{群} 概形；令 \(F(A)\) 为 \(A\) 上的 \emph{abelian} 概形的同构类集，即在 \(A\) 上固有且光滑的群概形，并赋有形如 \hyperref[fga3.ii-c.4-equation-star]{式~\ref*{fga3.ii-c.4-equation-star}} 的 abelian 概形同构。可以证明，规定乘法结构，甚至仅规定一个“单位截面”，就消除了无穷小自同构，因而存在一个形式模概形，对应于完备 Noether 局部环 \(O\)。还可以证明，若 \(X\) 是局部 Noether 预概形 \(S\) 上的固有光滑概形，且纤维“绝对连通”，则 \(X\) 上任何有单位截面的乘法结构都必然结合且交换，只要它在 \emph{一} 个纤维上结合且交换，并且 \(S\) 连通；此外，此结构由单位截面唯一确定。进一步，假设 \(S\) 为剩余域为 \(k\) 的 Artin 局部环 \(A\) 的谱，\(X\) 在 \(A\) 上固有且赋有截面 \(s\)，最后 \(X\otimes _A k\) 赋有 \(k\) 上的 abelian 概形结构，其零元为对应于 \(s\) 的 \(X\otimes _A k\) 的点。一个容易的障碍计算，加上 Serre 的一个论证，就能证明 \(X\) 上存在以截面 \(s\) 为单位截面的乘法结构。由此，利用 \hyperref[fga2]{[2]} 的“存在定理”过渡到 \(A\) 为完备 Noether 局部环的情形，再对一般情形应用 \hyperref[fga3.i]{[3]} 的下降技术，就能对所有局部 Noether 连通的 \(S\) 证明类似论断。这表明，此处考虑的函子 \(F(A)\) 同构于本节开头定义的、忘掉 \(X_0\) 的乘法结构所得的类似函子。因此，特别地，若 \(\mathfrak {m}\) 为 \(O\) 的极大理想，则 \(\mathfrak {m}/\mathfrak {m}^2\) 典范同构于 \(\operatorname {H}^1(X_0,\mathfrak {G}_{X_0/k})\) 的对偶，因而维数为 \(n^2\)，其中 \(n=\dim  X_0\)。判定 \(O\) 是否确实在 \(\Lambda \) 上 \emph{光滑}，即是否同构于 \(\Lambda \) 上 \(n^2\) 个变元的形式幂级数代数，将非常有趣。现在，\hyperref[fga3.ii-a.5-proposition-5.2]{§A, 命题 5.2} 允许将此问题等价地表述为一个 \emph{可约化为给定 abelian 概形的 abelian 概形的存在性问题}。无论如何，若 \(k\) 的特征为 \(0\)，可用超越方法看到答案是肯定的。在特征 \(p\neq 0\) 下，显然只需考虑 \(\Lambda \) 为代数闭域 \(k\) 上构造的 Witt 向量环的情形。这也许正是“\emph{Greenberg 函子}”显示作用的时候……\begin{remarks}{}\label{fga3.ii-c.3-remarks-ii}\skipsinglepar\par{}\emph{［编者补注］} 我最近证明，域上的 abelian 簇的形式模概形在 Witt 环上确实光滑；换言之，任意局部 Artin 环 \(A\) 上的每个 abelian 概形，若 \(A\) 是另一个局部 Artin 环 \(B\) 的商，则它都由 \(B\) 上的某个 abelian 概形约化而来。证明只用到 第 182 次报告第 12 页引入的提升障碍类的变性。还要回顾，Mumford 已构造了亏格为 \(g\) 的曲线及极化 abelian 概形的模概形，参见 Harvard University 的 \emph{Séminaire Mumford–Tate}（1961—62）。\end{remarks}\section{覆盖的延拓}\label{fga3.ii-c.5}\par{}设 \(\mathfrak {X}\) 为 \emph{形式} Noether 预概形\hyperref[fga2]{[2]}，\(U\) 为 \(\mathfrak {X}\) 的一个开子集，局部由 \(\mathscr {O}_\mathfrak {X}\) 的一个非零因子截面的“非零”条件定义；因此它足够大，使得 \(\mathscr {O}_\mathfrak {X}\) 在开子集 \(V\) 上的任何截面，只要在 \(U\cap  V\) 上为零，就必为零。（\emph{［编者补注］} 还须假设定义 \(U\) 的截面不仅在 \(\mathfrak {X}\) 上是非零因子，而且在每个 \(X_n\) 上也如此。）设 \(\mathfrak {J}\) 为 \(\mathfrak {X}\) 的一个“定义理想”，置 \(X_n=(\mathfrak {X},\mathscr {O}_\mathfrak {X}/\mathfrak {J}^{n+1})\)，于是它是通常的 Noether 预概形。若 \(\mathfrak {X}'\) 和 \(\mathfrak {X}''\) 是 \(\mathfrak {X}\) 的 \emph{平坦覆盖}，即由作为模层凝聚且局部自由的代数层定义的 \(\mathfrak {X}\) 上预概形，且它们 \emph{在 \(U\) 上非分歧}，则显然的映射 \[     \operatorname {Hom}_\mathfrak {X}(\mathfrak {X}',\mathfrak {X}'')     \to  \operatorname {Hom}_{X_0}(X'_0,X''_0)   \] 为 \emph{单射}；特别地，\emph{\(\mathfrak {X}'\) 的一个在 \(X'_0\) 上诱导恒等的自同构必为恒等}。这允许我们对此情形应用下降技术。具体地，从 \(X_0\) 的一个在 \(U_0\) 上非分歧的平坦覆盖 \(X'_0\) 出发，令 \(G(\mathfrak {X})\) 为 \(\mathfrak {X}\) 的那些在 \(X_0\) 上诱导 \(X'_0\) 的平坦覆盖 \(\mathfrak {X}'\) 的同构类集，所用同构在 \(X'_0\) 上诱导恒等；这些覆盖必然在 \(U\) 上非分歧。类似地，对 \(\mathfrak {X}\) 的每个开子集 \(V\) 定义 \(G(V)\)；更一般地，对 \(\mathfrak {X}\) 上的每个形式预概形 \(\mathfrak {Y}\) 定义 \(G(\mathfrak {Y})\)。于是，由 \hyperref[fga2]{[2]} 与 \hyperref[fga3.i]{[3]} 的结果首先得到如下结论：\begin{enumerate}
    
 \item[a.]
      \par{}当 \(V\) 遍历 \(\mathfrak {X}\) 的开子集时，各 \(G(V)\) 构成 \(\mathfrak {X}\) 上的一个 \emph{层}，记为 \(\mathscr {G}_\mathfrak {X}=\mathscr {G}\)。此层在 \(U\) 上的限制是茎只有一个元素的 \emph{常值} 层。
      \par{}更一般地，描述 \(\mathscr {G}_\mathfrak {X}\) 的茎，可以精确地归结为完备局部环的问题：
    

      
    
 \item[b.]对每个 \(x\in \mathfrak {X}\)，有 \[         \mathscr {G}_x = G(\operatorname {Spec}(\mathscr {O}_{\mathfrak {X},x})) \subset  G(\operatorname {Spec}(\widehat {\mathscr {O}}_{\mathfrak {X},x}))       \]。也就是说，考虑 \(\widehat {\mathscr {O}}_{\mathfrak {X},x}\) 上的有限自由代数 \(B\) 的同构类，并赋有从 \(B\otimes _{\widehat {\mathscr {O}}_{\mathfrak {X},x}}\mathscr {O}_{X_0,x}\) 到 \((\widehat {\mathscr {O}}'_0)_x\) 的同构，其中 \(\mathscr {O}'_0\) 是 \(X_0\) 上定义 \(X'_0\) 的代数层。


    
 \item[c.]有典范同构 \(\mathscr {G}_\mathfrak {X}=\varprojlim \mathscr {G}_{\mathfrak {X}_n}\)；换言之，对 \(\mathfrak {X}\) 的每个开子集 \(V\)，有 \(G(V)=\varprojlim  G(V_n)\)。


    
 \item[d.]假设 \(\mathfrak {X}\) 来自完备 Noether 局部环 \(\Lambda \) 上的一个通常的 \emph{固有} 概形 \(X\)：该环的定义理想为 \(\mathfrak {m}\)，对 \(\mathscr {O}_X\) 作 \(\mathscr {J}\) 进完备化，其中 \(\mathscr {J}=\mathfrak {m}\cdot \mathscr {O}_X\)。则 \(G(\mathfrak {X})\) 典范同构于 \emph{通常的} 概形 \(X\) 的那些“可约化为 \(X'_0\)”的平坦覆盖的类集。

  \end{enumerate} 形象地说，（a）与（b）建立了问题的局部方面与整体方面之间的基本关系；（c）给出“有限”方面与“无穷小”方面之间的关系；最后，（d）在精确条件下重申了“形式”方面与“代数”方面的一致性。\par{}现在设 \(\mathfrak {X}\) 由完备 Noether 局部环 \(\Lambda \) 定义，且 \(\mathscr {J}=\mathfrak {m}\cdot \mathscr {O}_X\)，所以 \(X_0\) 为 \({\Lambda }/\mathfrak {m}\) 上的预概形。对每个在 \(\Lambda \) 上有限的代数 \(A\)，置 \[     F(A)     = G(\mathfrak {X}\times _\Lambda  A).   \]。这是关于 \(A\) 的取值于集合范畴的协变函子；由（c），此函子完全由它在 Artin 代数 \(A\) 上的作用确定。说此函子 pro 可表，即形如 \[     F(A)     = \operatorname {Hom}_{\text {拓扑 }\Lambda \text { 代数}}(\mathscr {O},A)   \]，其中 \(\mathscr {O}\) 是 \cref{fga3.ii-a.5} 所述类型的拓扑 \(\Lambda \) 代数，等价于说限制到 Artin 代数 \(A\) 后此断言成立。于是，结合 \hyperref[fga3.ii-b-theorem-1]{定理~\ref*{fga3.ii-b-theorem-1}} 与 \hyperref[fga3.ii-b-theorem-2]{定理~\ref*{fga3.ii-b-theorem-2}} 确实得到：\begin{proposition}{5.1}\setnumber{5.1}\label{fga3.ii-c.5-proposition-5.1}\skipsinglepar\par{}上述函子 pro 可表。\end{proposition}\par{}当然，由（a），若 \(U=\mathfrak {X}\)，则对 \(\mathfrak {X}\) 上的每个 \(\mathfrak {Y}\)，\(G(\mathfrak {Y})\) 都只有一个元素，此时函子 \(F\) 意义不大，因为有 \(\mathscr {O}=\Lambda \)。看来，几乎在其他一切情形下，拓扑局部环 \(\mathscr {O}\) 都 \emph{不是 Noether 环}。不过，它的存在性鲜明地显示了解集 \(G(\mathfrak {X})\) 的 \emph{“连续”性质}；直观上，这对应于延拓 \(X'_0\) 时分歧如何扩散存在“连续”的选择。可以将此结果与 J.-P. Serre \hyperref[Ser1958]{[4]} 通过局部类域论所采取的观点比较；还应注意，“几何”局部域的极大 abelian 扩张的拓扑 Galois 群具有 \emph{连续} 性质，其 Pontrjagin 意义下的对偶群表现为代数群、或至少拟代数群的归纳极限；这里，扩张的分类同样由无限维“簇”给出。上述讨论中，也可以取 \(\mathfrak {X}\) 为某个完备局部环的形式谱，例如 \(\Lambda \) 是其 Cohen 子环；我们可以希望将本节结果用于维数 \(>1\) 的完备局部环的扩张研究。在局部和整体两种情形下，这些结果都有望使我们能够在高阶分歧现象与特征 \(0\) 的现象之间建立精确关系，后者可用超越方法处理。无论如何，正是 \hyperref[fga3.ii-c.5-proposition-5.1]{命题~\ref*{fga3.ii-c.5-proposition-5.1}} 的预备分析，使得 \hyperref[fga2]{[2]} 所述研究基本群的方法能够推广到“温分歧”情形，并用超越方法解决“三点问题”。\par{}最后指出，若 \(X_0\) 的维数为 \(1\)，情况会简化；此时由（a）和（b），\(G(\mathfrak {X})\) 可等同于 \(\prod _i G(\operatorname {Spec}(\mathscr {O}_{\mathfrak {X},x_i}))\)，其中 \(x_i\) 为 \(X_0\setminus  U\) 的各点：\emph{在分歧点处可以任意选取“局部”延拓}。进一步，若 \(X_0\) 正规，则可指出，\hyperref[fga3.ii-c.5-proposition-5.1]{命题~\ref*{fga3.ii-c.5-proposition-5.1}} 所保证的形式模概形在 \(\operatorname {Spec}(\Lambda )\) 上 \emph{光滑}。
\section*{原报告参考文献}
\begin{enumerate}
\item[1.]\setnumber{1}\label{fga3.ii-bib-1} GROTHENDIECK (Alexander). \emph{Sur quelques points d'algèbre homologique}，\emph{Tohoku Math. J.}，9 (1957)，119--221。
\item[2.]\setnumber{2}\label{fga3.ii-bib-2} GROTHENDIECK (Alexander). \emph{Géométrie formelle et géométrie algébrique}，\emph{Séminaire Bourbaki}，11 (1958/59)，分册 3，第 182 次报告。
\item[3.]\setnumber{3}\label{fga3.ii-bib-3} GROTHENDIECK (Alexander). \emph{Technique de descente et théorèmes d'existence en géométrie algébrique, I : Généralités. Descente par morphismes fidèlement plats}，\emph{Séminaire Bourbaki}，12 (1959/60)，分册 1，第 190 次报告。
\item[4.]\setnumber{4}\label{fga3.ii-bib-4} SERRE (Jean-Pierre). \emph{Corps locaux et isogénies}，\emph{Séminaire Bourbaki}，11 (1958/59)，分册 3，第 185 次报告。
\end{enumerate}
\section*{第 182 次报告勘误（1959 年 5 月）}
\label{fga3.ii-erratum-fga2}
\noindent\emph{Géométrie formelle et géométrie algébrique}：第 23 页第 12 行（推论 5 的第一行），加入：“\(X\) 或 \(Y\) 在 \(k\) 上固有”。

% END chapters/fga3.ii.tex

\clearpage
\setcounter{chapter}{-1}
\renewcommand\thechapter{\arabic{chapter}}
\renewcommand\fgaAnchorPrefix{fga3.iii}

% BEGIN chapters/fga3.iii.tex
\part{商预概形}\label{fga3.iii}\thispagestyle{fancy}
\begin{original}{}\label{fga3.iii-original-citation}\skipsinglepar
\par{}A. Grothendieck. "Techniques de construction et théorèmes d'existence en géométrie algébrique, III: Préschemas quotients". \emph{Séminaire Bourbaki} \textbf{13} (1960–61), 第212讲, 1961年2月. \url{http://www.numdam.org/book-part/SB_1960-1961__6__99_0/}
\end{original}
\chapter{引言}\thispagestyle{fancy}\label{fga3.iii-introduction}
\begin{remark}{}\label{fga3.iii-introduction-remark}\skipsinglepar
\par{}\emph{[Comp.]} 我们指出, \hyperref[fga3.v]{FGA~\ref*{fga3.v}} ("Picard概形: 存在性定理") 中对这里所建立理论的应用, 也可以用适当地运用Hilbert概形来替代 (参见\emph{Séminaire Mumford–Tate}, Harvard University (1961–62)). 如\cref{fga3.iii-8}所述, 这里的理论最重要的缺口是缺乏关于非固有等价关系的商的存在性判据, 例如射影群的某些作用所产生的等价关系. Mumford \cite{Mum1961} 已经在这一方向得到一个重要定理. 关于他的结果的精化及这一理论的各种应用, 见\emph{Séminaire Mumford–Tate}, Harvard University (1961–62).
\end{remark}
\par{}本讲所讨论的问题与前两讲不同: 我们试图表示的是以变动概形为自变量的某些协变函子, 而不再是反变函子. 然而, 在代数几何的许多构造问题中, 包括\hyperref[fga3.i]{[1]}和\hyperref[fga3.ii]{[2]}中的问题, 取商的过程都不可缺少. 事实上, 对于忠实平坦且拟紧的态射\(T\to  S\), 一个\(T\)-预概形\(X\)上的\emph{下降数据的有效性}问题, 等价于\(X\)关于该下降数据在\(X\)上所定义的平坦等价关系的商的存在性问题, 其中要求这个商具有下面将研究的一些合理性质; \hyperref[fga3.i-a.2.c]{[1], A, §2.c}中提出的问题, 或许可以与本讲\cref{fga3.iii-2}所提出的问题一并解决. 同样, 一个\(S\)-概形\(X\)的\emph{Picard概形} (定义见\hyperref[fga3.ii-c.3]{[2], C, §3}) 可以用多种方式定义为某些其他概形 (正除子概形, 或射影嵌入概形) 关于平坦等价关系的商, 而这些辅助概形的定义与构造也较为简单: 它们基本上是\(\operatorname {Hom}_S(X,Y)\)型的概形以及\hyperref[fga3.ii-c.2]{[2], C, §2}中定义的变体; 下一讲将讨论它们的构造, 其中须加上适当的射影性假设. 因而, 将本讲的结果与下一讲的结果结合, 就能在适当假设下构造Picard概形.
\par{}预概形中的取商问题仍有尚未解决的疑问. 最重要的一个将在\cref{fga3.iii-8}中说明. 目前, 它仍是构造\emph{任意亏格曲线在整数环上的模概形}、\emph{极化阿贝尔簇的模概形}等的唯一障碍. 因此, 它的解决值得代数群专家为之努力.
\chapter{等价关系与有效等价关系}\thispagestyle{fancy}\label{fga3.iii-1}
\par{}设\(\mathcal {C}\)是一个范畴, \(X\)是\(\mathcal {C}\)的一个对象. 如果对于\(\mathcal {C}\)的每个对象\(T\), 相应的映射\[     p_1(T),p_2(T)\colon  R(T)\rightrightarrows  X(T)   \] (这里对\(\mathcal {C}\)的任意对象\(Y\), 置\(Y(T)=\operatorname {Hom}(T,Y)\)) 所定义的映射\[     R(T)\to  X(T)\times  X(T),\] 都给出\(R(T)\)到集合\(X(T)\)上某个等价关系的图的双射, 则称态射对\[     p_1,p_2\colon  R\rightrightarrows  X,   \]是\(\mathcal {C}\)中以\(X\)为\emph{目标}、以\(R\)为\emph{来源}的"\emph{等价对}". 在以\(X\)为目标的等价对之间有一个显然的等价关系; 它的一个等价类称为\(X\)上的一个\emph{\(\mathcal {C}\)-等价关系}, 或在不会混淆时简称等价关系.
\par{}如果\(X\times  X\)存在, 则给出\(X\)上的一个等价关系, 等价于给出\(X\times  X\)的一个子对象\(R\), 使得对于\(\mathcal {C}\)的每个对象\(T\), 与\(R(T)\)相应的\((X\times  X)(T)=X(T)\times  X(T)\)的子集是\(X(T)\)上一个等价关系的图. 将投影\(\mathrm {pr}_1\)、\(\mathrm {pr}_2\)分别诱导的从\(R\)到\(X\)的态射记作\(p_1\)、\(p_2\), 上述条件就是说\((p_1,p_2)\)是一个等价对. 按照\hyperref[fga3.ii-a.1]{[2], A, §1}中阐述的一般原则, 也可以在\(\mathcal {C}\)中以交换图表示\(X(T)\)中的\(R(T)\)所满足的集合论等价关系公理 (这里假设\(X\times  X\)和纤维积\((R,p_2)\times _X(R,p_1)\)均存在). 我们不需要这样做.
\par{}只要有同来源\(R\)、同目标\(X\)的一对态射\((p_1,p_2)\), 就可以将这一对的\emph{余核}定义为\(\mathcal {C}\)的一个对象\(Y\), 它表示关于\(Z\)的协变函子\[     \operatorname {Hom}_{p_1,p_2}(X,Z);\] 这里的值是所有满足\(up_1=up_2\)的从\(X\)到\(Z\)的态射\(u\)所组成的集合. 如果\(Y\)存在, 则它在相差唯一同构的意义下确定. 我们将它记作\(Y/(p_1,p_2)\), 或滥用记号记作\(Y/R\); 后一个记号主要在\((p_1,p_2)\)为等价对时使用, 此时习惯于在记号上把该对定义的等价关系与\(R\)混同. 注意, 若将\(Y\)视为\(X\)的商, 那么它确实只依赖于态射对\((p_1,p_2)\)定义的等价关系.
\par{}现在从一个态射\[     f\colon  X\to  Y   \]出发, 它使我们能把\(X\)看成"\(Y\)之上的对象", 并假设纤维积\[     \mathcal {R}(f)     = X\times _Y X   \]存在. 设\(p_1\)、\(p_2\)是它的两个投影. 则\((p_1,p_2)\)是一个等价对, 称为与态射\(f\)\emph{相伴}的等价对. 它因而定义了一个等价关系, 称为与态射\(f\)\emph{相伴}的等价关系.
\par{}称以\(X\)为目标、\(R\)为来源的态射对\((p_1,p_2)\)为\emph{有效等价对}, 如果满足下列条件: 
\begin{enumerate}
    
 \item[i.]余核\(Y=X/(p_1,p_2)\)存在;

    
 \item[ii.]纤维积\(X\times _Y X\)存在;

    
 \item[iii.]以\(p_1\)和\(p_2\)为分量的态射\(R\to  X\times _Y X\)为同构.

  
\end{enumerate} 此时\((p_1,p_2)\)确实是等价对. 我们也称它定义的等价关系为\emph{有效等价关系}.
\par{}称态射\(f\colon  X\to  Y\)为\emph{有效满态射}, 如果满足下列条件: 
\begin{enumerate}
    
 \item[i.]纤维积\(R=X\times _Y X\)存在;

    
 \item[ii.]商\(X/(p_1,p_2)\)存在, 这里\(p_1\)、\(p_2\)是从\(R\)到\(X\)的两个投影;

    
 \item[iii.]由\(f\)诱导的态射\(X/(p_1,p_2)\to  Y\)为同构.

  
\end{enumerate} 此时\(f\)确实是满态射, 甚至是严格满态射 (参见\hyperref[fga3.i-a.2.c]{[1], A, §2.c}); 如果纤维积\(X\times _Y X\)存在, 则反之亦然. 我们也称满态射\(f\)所定义的\(X\)的商对象为\(X\)的一个\emph{有效商}.
\par{}上述定义蕴涵下述"\emph{Galois对应}":
\begin{proposition}{1.1}\setnumber{1.1}\label{fga3.iii-1-proposition-1.1}\skipsinglepar
\par{}\(X\)上有效等价关系\(R\)的集合与\(X\)的有效商\(Y\)的集合之间, 存在一个保持自然序的双射对应: \(R\)对应于有效商\(X/R\), 而\(Y\)对应于典范投影\(X\to  Y\)所定义的有效等价关系, 即由带有两个投影的纤维积\(X\times _Y X\)所定义的关系.
\end{proposition}
\par{}在性质很好的范畴 (集合、集合层等) 中, 每个商都是有效的, 每个等价关系也都是有效的. 对于给定预概形\(S\)之上的预概形范畴等范畴, 这就不再成立, 即使\(S\)是一个域的谱, 甚至只考虑\(S\)上的有限概形, 也不成立. 有效性问题, 乃至商的存在性问题 (在\(S\)上的预概形不有限时), 往往都很微妙.
\chapter{例: \(S\)上的有限预概形}\thispagestyle{fancy}\label{fga3.iii-2}
\par{}设\(\mathcal {C}\)是\(S\)上的有限预概形范畴, 其中假设基底局部Noether. 那么\(\mathcal {C}\)等价于\(S\)上凝聚交换代数层范畴的反范畴; 如果\(S\)是环为\(A\)的仿射概形, 则它等价于在\(A\)上有限的\(A\)-代数 (即在\(A\)上作为模为有限型的代数) 范畴的反范畴. 因而立即得到, \(\mathcal {C}\)中的有限投射极限与有限归纳极限均存在. 对前者, 这是熟知的, 而且不需任何有限性假设: 预概形\(X\)和\(Y\)在\(S\)上的纤维积对应于相应代数的张量积\(B\otimes _A C\), 由两个\(A\)-代数同态\(u,v\colon  C\rightrightarrows  B\)定义的两个态射\(X\rightrightarrows  Y\)的核, 对应于\(B\)模掉由\(u(c)-v(c)\)生成的理想所得的商, 等等. 对于有限归纳极限, 只需一方面考虑有限和, 它们对应于\(A\)-代数的通常直积; 另一方面考虑态射对\(X\rightrightarrows  Y\)的余核, 它对应于\(C\)中使两个同态\(u,v\colon  C\rightrightarrows  B\)取值相同的元素组成的子环, 这一点可以直接看出 (由Noether假设, 此子环在\(A\)上有限). 还应指出, 利用Noether假设, 可以证明在\(S\)上的有限预概形范畴\(\mathcal {C}\)中如此构造的有限归纳极限, 特别是商, 实际上也是\emph{所有}预概形范畴中的商.
\par{}如\hyperref[fga3.i]{[1]}所述, \emph{\(\mathcal {C}\)中存在非有效的满态射} (也即非严格的满态射; 因为纤维积存在, 这两者相同). 如果不加平坦性假设, \emph{我不知道等价关系是否总是有效的}. 在这一方向, 我只得到一些很局部的肯定结果, 而这些结果对于形式模理论基本定理的证明不可缺少 (参见\hyperref[fga3.ii-b-theorem-1]{[2], B, 定理1}). 注意, 在所提问题中, 很容易归结到\(S\)是剩余域代数闭的Artin局部环的谱的情形. 但即使\(A\)是一个域, 答案也还不知道.
\par{}也可以考虑\(S\)上的预概形\(X\)不再假设在\(S\)上有限, 而\(X\)上的一个等价关系\(R\)满足\(p_1\colon  R\to  X\)是有限态射的情形. 此时称\(R\)为\emph{有限等价关系}. 为简单起见, 假设\(S\)和\(X\)仿射 (于是\(R\)也仿射, 从而问题归结为纯交换代数的情形), \emph{即使在这种情形, 我们也不知道商\(X/R=Y\)是否存在, 以及典范态射\(X\to  Y\)是否有限}. (最简单的情形是假设\(S\)是域\(k\)的谱, 而\(X\)是\(k[t]\)的谱, 即仿射直线.) 当然, 如果上述两个问题的答案都为肯定, 那么在所述情形就能断定\(R\)是有效的. 注意, 若不是\(X\)中的等价关系图, 而只有\cref{fga3.iii-4}意义下\(X\)中的一个预等价关系图, 则商\(Y\)的\emph{存在性}问题和\(f\colon  X\to  Y\)的\emph{有限性}问题仍可用完全相同的方式提出.
\par{}沿某些闭子预概形"粘合"给定的预概形\(X_i\), 以构造预概形时, 会遇到关于一般有限等价关系的取商问题; 粘合法则恰好由\(X_i\)之和构成的预概形\(X\)上的有限等价关系表示. 也可以预期, 这些问题及其许多变体的解决, 将是沿\hyperref[fga3.ii]{[2]}开创的方向建立一般非射影构造方法的先决条件.
\par{}作者所知道的唯一一般性的肯定事实如下:
\begin{proposition}{2.1}\setnumber{2.1}\label{fga3.iii-2-proposition-2.1}\skipsinglepar
\par{}设\(S\)是局部Noether预概形, \(s\)是\(S\)的一点, \(\Omega \)是\(k(s)\)的一个代数闭扩张. 考虑相应的"纤维函子"\(F\), 它把每个在\(S\)上有限的\(S\)-概形\(X\)对应于\(X/S\)的所有\(\Omega \)值点组成的集合. 此函子显然左正合, 而且是\emph{右正合}的, 即它与有限归纳极限交换, 特别是与态射对的余核交换.
\end{proposition}
\par{}对\(S\)的所有"几何点"应用这个结果, 就得到\(\mathcal {C}\)的"商"范畴\(\mathcal {C}'\) (即"模掉满的根态射"来推理, 也就是形式地给这类态射添上逆) 是一个"几何"范畴, 即满足与集合范畴相同的"有限性质". 特别是每个等价关系都有效. 由此得出: 若\(R\)是\(X\)上的等价关系, 且\(X\)在\(S\)上有限, 则典范态射\(R\to  X\times _Y X\) (其中\(Y=X/R\)) 是\emph{根且满}的; 事实上, 因为它是单态射, 所以是满的闭浸入.
\chapter{群作用的情形}\thispagestyle{fancy}\label{fga3.iii-3}
\par{}现在重新假设\(\mathcal {C}\)是任意范畴. 设\(G\)、\(X\)是\(\mathcal {C}\)的对象, 且\(G\)是作用于对象\(X\)的\(\mathcal {C}\)-群. 这就是说 (参见\hyperref[fga3.ii-a.1]{[2], A, §1}), 对\(\mathcal {C}\)的每个对象\(T\), 在\(G(T)\)上给定一个群结构, 并在\(X(T)\)上给定以\(G(T)\)为算子群的集合结构, 使得当\(T\)变动时, 这些结构"随\(T\)函子地变动". 若\(G\times  G\)和\(G\times  X\)在\(\mathcal {C}\)中存在, 则这种结构也可以定义为一对态射\[     
\begin {gathered}       G\times  G\to  G     \\\pi \colon  G\times  X\to  X     
\end {gathered},\] 满足对于\(\mathcal {C}\)的每个对象\(T\), 集合\(G(T)\)和\(X(T)\)上的相应合成律使\(G(T)\)成为作用于\(X(T)\)的群. 将这个公理翻译为\(\mathcal {C}\)中若干图的交换性是容易的, 但很繁琐, 而且在我所知的所有情形中实际上完全没有必要.
\par{}假设\(G\times  X\)存在, 考虑两个态射\[     p_1,p_2\colon  G\times  X\rightrightarrows  X,\] 其中\[     
\begin {aligned}       p_1 &= \mathrm {pr}_1     \\p_2 &= \pi .     
\end {aligned}.\] 立即可见, \((p_1,p_2)\)是等价对当且仅当对于\(\mathcal {C}\)的每个对象\(T\), 该对定义的映射\[     G(T)\times  X(T)     \sim  (G\times  X)(T) \to  X(T)\times  X(T)   \]是单射, 即群\(G(T)\)在集合\(X(T)\)上\emph{自由地}作用; 换言之, 若\(g\in  G(T)\)、\(x\in  X(T)\)且\(g\cdot  x=x\), 则\(g\)是群\(G(T)\)的单位元. 此时称\(G\)在\(X\)上\emph{自由作用} (或称\(X\)是\emph{\(G\)下的主\(\mathcal {C}\)-空间}). 于是与\((p_1,p_2)\)相伴的等价关系称为在\(X\)上自由作用的\emph{群\(G\)所定义的等价关系}. 如果\(X\times  X\)也存在, 并考虑由\((p_1,p_2)\)定义的态射\[     p\colon  G\times  X\to  X\times  X,\] 那么\(G\)自由作用的条件就是说\(p\)是一个\emph{单态射}.
\par{}当然, 即使\(G\)在\(X\)上的作用不是自由的, 我们仍然希望有\(X\)关于\(G\)的商的存在性判据, 即上述态射对\((p_1,p_2)\)的余核的存在性判据.
\par{}这个余核通常记作\(X/G\); 若\(G\)在左边作用, 则更适合记作\(G\backslash  X\) (前一个记号留给\(G\)在右边作用的情形). 应当指出, 即使\(G\times  X\)在\(p\)下的"像"存在, 记作\(R\) (例如将此像定义为\(X\times  X\)中能使\(p\)经由它分解的最小子对象), 它通常也不是\(X\)上的等价关系. 若此时试图直接关于\(R\)取商 (更准确地说, 关于由两个投影\(\mathrm {pr}_i\)诱导的从\(R\)到\(X\)的态射对取商), 就会丢失原态射对\((p_1,p_2)\)的特有性质. 因此有必要推广等价关系的概念, 使它能直接适用于由一个带作用的\(\mathcal {C}\)-群定义的态射对.
\chapter{预等价关系}\thispagestyle{fancy}\label{fga3.iii-4}
\par{}回顾一下, \emph{群胚}是指所有态射均为同构的范畴. 另一方面, 一个范畴由两个基本集合\(X\)、\(R\) (前者是\emph{对象}集, 后者是\emph{箭头}集) 以及下列结构组成: 
\begin{enumerate}
    
 \item[i.]一对映射\[         p_1,p_2\colon  R\rightrightarrows  X,\] 称为\emph{来源映射}和\emph{目标映射};

    
 \item[ii.]一个映射\[         \pi \colon (R,p_2)\times _X(R,p_1) \to  R,\] 称为\emph{合成映射}.

  
\end{enumerate} 这些数据应满足熟知的公理, 这里不再重复. 公理可以表示为某些图的交换性, 再加上一个使另外两个图交换的必然唯一的映射\(D\colon  X\to  R\)的存在性; \(D\)把对象对应于其恒等映射, 并满足\[     p_1\circ  D = p_2\circ  D = \operatorname {id}_X.\] 说这个范畴是群胚, 就意味着存在一个必然唯一的映射\[     s\colon  R\to  R,\] 称为\(R\)的\emph{对称}, 它把每个箭头变成其逆箭头. 这可以表示为由\(s\)、\(\Delta \)及上述数据构成的另外四个图的交换性, 其中前两个可写成\[     
\begin {aligned}       p_1\circ  s &= p_2     \\p_2\circ  s &= p_1.     
\end {aligned}.\]
\par{}回顾这些概念之后, \hyperref[fga3.ii-a.1]{[2], A, §1}中的一般定义特别说明了, 对于任意范畴\(\mathcal {C}\)的一对对象\((X,R)\), 所谓"\emph{\(\mathcal {C}\)-范畴}结构" (分别为\emph{\(\mathcal {C}\)-群胚}结构) 应是什么意思: 按定义, 就是对\(\mathcal {C}\)的每个对象\(T\), 在以\(X(T)\)为对象集、以\(R(T)\)为箭头集的集合对上给定一个集合论意义的范畴 (分别为群胚) 结构, 并使这些结构"随\(T\)函子地变动". 于是得到两个态射\[     p_1,p_2\colon  R\rightrightarrows  X,\] 称为\emph{来源态射}和\emph{目标态射}, 并且若有关的纤维积存在, 则还有一个态射\[     \pi \colon  (R,p_2)\times _X(R,p_1) \to  R,\] 称为\emph{合成态射}. 这三个态射就足以确定\((X,R)\)上的范畴 (分别为群胚) 结构, 对它们所加的条件是: 对于每个\(T\), 与\(X(T)\)和\(R(T)\)相应的三个映射在集合对\((X(T),R(T))\)上定义一个范畴 (分别为群胚) 结构. 必要时, 也可将此表示为某些图的交换性, 其中涉及一个完全确定的态射\[     D\colon  X\to  R,\] 以及在群胚情形中的一个完全确定的态射\[     s\colon  R\to  R;\] 这些图与"集合论"情形相同. 幸而实际上不需要把公理作如此繁琐的解释. 能够用某些纤维积之间的态射及态射等式来表示这些数据和公理, 在理论上的唯一意义是: 若有一个左正合函子\(F\colon \mathcal {C}\to \mathcal {C}'\) (即与有限直积和纤维积交换的函子), 则它把\(\mathcal {C}\)-范畴 (分别为\(\mathcal {C}\)-群胚) 变成\(\mathcal {C}'\)-范畴 (分别为\(\mathcal {C}'\)-群胚), 条件是有限直积和纤维积在\(\mathcal {C}\)中存在.
\par{}实际中重要的是, 应懂得如何把态射\(p_1\)、\(p_2\)、\(\pi \)、\(D\)和\(s\)理解为\(\mathcal {C}\)中一个适当的半单纯或单纯对象的\emph{单纯运算} (至少在\(\mathcal {C}\)中纤维积存在时如此). 为固定术语, 引入\emph{单纯形型}范畴\(\mathcal {S}\): 它的对象为形如\[     \Delta _n     = [0,n]   \]的有限集合, 其中\(n\in \mathbb {Z}_{\geqslant 0}\), 而\([0,n]\)表示从\(0\)到\(n\)的整数闭区间; 它的态射为这些有限集合之间的\emph{任意映射}. 注意\(\mathcal {S}\)等价于\emph{非空}有限集合及其间所有映射构成的范畴. \(\mathcal {S}\)中显然存在\emph{非空}有限族对象的和, 也存在两个对象沿第三个对象的粘合和 (它是纤维积的对偶运算). 以\(\mathcal {S}'\)表示\(\mathcal {S}\)的具有相同对象的子范畴, 但态射只取\(\Delta _n\)之间的\emph{递增映射}. 此范畴等价于非空有限全序集的范畴. 在这个范畴中, 两个对象的和从不存在, 而两个对象\(A\)、\(B\)沿第三个对象\(C\)的粘合和一般也不存在 (例如取\(C=\Delta _0\)、\(A=B=\Delta _1\), 且两个结构映射\(u\colon  C\to  A\)与\(v\colon  C\to  B\)相同). 不过, 在某些情形粘合和\emph{确实}存在; 考虑\[     
\begin {gathered}       A = \Delta _m       \qquad  B = \Delta _n       \qquad  C = \Delta _0     \\u(0) = m       \qquad  v(0) = 0     
\end {gathered},\] 此时有\[     A\coprod _C B     = \Delta _{m+n}.\]
\par{}范畴\(\mathcal {C}\)中的\emph{单纯对象} (分别为\emph{半单纯对象}) 定义为从\(\mathcal {S}\) (分别为\(\mathcal {S}'\)) 到\(\mathcal {C}\)的一个反变函子\(K\). 因此单纯对象通过限制定义一个半单纯对象, 但前者与后者的本质区别在于\(K_n=K(\Delta _n)\)中还有\emph{对称运算}, 它们对应于\(n+1\)个元素的对称群的元素在函子\(K\)下的像 (此群视为\(\mathcal {S}\)中\(\Delta _n\)的自同构群).
\par{}按上述约定, 对每个\(n\), 以\(\Delta '_n\) (分别为\(\Delta ''_n\)) 表示对象集为\(\Delta _n\)、箭头集由\(\Delta _n\)上的"混沌序"关系 (分别为自然全序关系) 定义的有限范畴 (即箭头集是这一关系的图). 显然, \(\Delta '_n\) (分别为\(\Delta ''_n\)) 函子地依赖于\(\mathcal {S}\) (分别为\(\mathcal {S}'\)) 的对象\(\Delta _n\). 因而若\(Z\)是一个范畴, 则当\(\Delta _n\)变动时, \(\operatorname {Hom}(\Delta '_n,Z)\) (分别为\(\operatorname {Hom}(\Delta ''_n,Z)\)) 是从范畴\(\mathcal {S}\) (分别为\(\mathcal {S}'\)) 到集合范畴的一个函子, 即所谓与\emph{范畴\(Z\)相伴}的\emph{单纯集} (分别为\emph{半单纯集}), 记作\(Z'\) (分别为\(Z''\)). 此外, 从与\(Z'\)相应的半单纯集到\(Z''\)有一个显然的自然同态, 它是同构当且仅当\(Z\)是群胚. 于是:
\begin{proposition}{4.1}\setnumber{4.1}\label{fga3.iii-4-proposition-4.1}\skipsinglepar
\par{}从范畴的范畴到半单纯集范畴的函子\(Z\mapsto  Z''\)完全忠实, 并在\emph{范畴}的范畴与满足以下条件的半单纯集范畴之间定义等价: 即从\(\mathcal {S}'\)到\(\mathtt {Set}\)的反变函子\(K\)\emph{把上述类型的粘合和\(A\coprod _C B\)变成集合的纤维积}.
\par{}同样, 从群胚范畴到单纯集范畴的函子\(Z\mapsto  Z'\)完全忠实, 并在\emph{群胚}的范畴与满足以下条件的单纯集范畴之间定义等价: 即从\(\mathcal {S}\)到\(\mathtt {Set}\)的反变函子\(K\)\emph{把粘合和变成纤维积}.
\end{proposition}
\par{}因此可以把范畴看成半单纯集的特殊情形, 把群胚看成单纯集的特殊情形; 当然须在"相差同构"的意义下这样说, 这正是在用一种结构解释另一种结构时的严格说法. 通常的归结到集合论情形的办法于是给出:
\begin{corollary}{4.2}\setnumber{4.2}\label{fga3.iii-4-corollary-4.2}\skipsinglepar
\par{}将范畴、群胚和单纯集分别换成\(\mathcal {C}\)-范畴、\(\mathcal {C}\)-群胚和\(\mathcal {C}\)-单纯对象, 上述断言仍然成立, \emph{条件是}\(\mathcal {C}\)中存在纤维积.
\end{corollary}
\par{}与\(\mathcal {C}\)中的范畴\((X,R,\ldots )\)相伴的\(\mathcal {C}\)中的半单纯对象\(K\), 可以具体写出: 将\(K\)的分量\(K_n=K(\Delta _n)\)视为\((R,p_1)\)在\(X\)上的\((n+1)\)重纤维积, 或者更方便地用递推公式\[     
\begin {aligned}       K_0 &= R     \\K_n &= (K_{n-1},p_n^{(n-1)})\times _X(R,p_1)     
\end {aligned}   \]来定义; 这里\(p_i^{(n-1)}\) (其中\(0<i<n-1\)) 是从\(K_{n-1}\)到\(X\)的自然投影 (也可递推定义). 如此, \(p_1\)、\(p_2\)、\(\pi \)、\(D\)和\(s\)便可分别理解为与\(\mathcal {S}\)中下列态射相应的单纯运算: \(\Delta _1\)的\(0\)面、\(\Delta _1\)的\(1\)面、\(\Delta _2\)的\((0,2)\)面、退化\(\Delta _1\to \Delta _0\)以及\(\Delta _1\)的对称. 其余所有半单纯 (分别为单纯) 运算均可由上述四个 (分别为五个) 运算通过合成和纤维积形式地得到.
\par{}现在将一个范畴的对象\(X\)上的\emph{预等价关系}定义为以\(X\)为对象之对象的一个群胚的数据. 这样的数据, 除其他部分以外, 还给出对象\(R\)及两个态射\[     p_1,p_2\colon  R\rightrightarrows  X.\] 但应注意, 与等价对的情形相反, 仅这些数据并不能确定所讨论的结构. 本讲之所以研究这一概念, 是为了得到能够取商的判据, 即能够作出态射对\((p_1,p_2)\)的余核的判据. 因而这个问题的陈述不涉及群胚固有的其他数据. 不过, 在下面各结果的证明中, 我们将用到这些额外数据, 特别是直到\(3\)维的单纯运算 (包括对称运算), 其中将出现\(R\)在\(X\)上的四重纤维积.
\par{}\(\mathcal {C}\)的对象\(X\)上的一个等价关系定义一个预等价关系: 只需在集合论情形说明这一点, 此时对集合\(X\)上的一个等价关系, 配以对象集为\(X\)、箭头集为该等价关系的图的集合的群胚即可.
\par{}作用于\(\mathcal {C}\)的对象\(X\)的\(\mathcal {C}\)-幺半群\(G\), 定义一个以\(R=G\times  X\)和\(X\)为基本对象的\(\mathcal {C}\)-范畴 (条件是\(G\times  X\)存在); 而该范畴是\(\mathcal {C}\)-群胚当且仅当\(G\)是群. 仍只须在集合论情形证明. 此时把箭头\((g,a)\)与\((g',g\cdot  a)\)的合成定义为\[     (g',g\cdot  a) \circ  (g,a)     = (g'g,a);\] 也就是说, 若\(a,b\in  X\), 则按定义, \(\operatorname {Hom}(a,b)\)是把\(a\)变到\(b\)的搬运子, 而态射依靠\(G\)中元素的合成而合成.
\begin{remark}{}\label{fga3.iii-4-remark}\skipsinglepar
\par{}只要暗中假设所有涉及的对象均属于某个固定的"宇宙" (它本身是一个集合), 就可避开\hyperref[fga3.iii-4-proposition-4.1]{命题~\ref*{fga3.iii-4-proposition-4.1}}这类陈述所引起的逻辑困难.
\end{remark}
\chapter{关于有限平坦等价关系的商}\thispagestyle{fancy}\label{fga3.iii-5}
\begin{theorem}{5.1}\setnumber{5.1}\label{fga3.iii-5-theorem-5.1}\skipsinglepar
\par{}设\(X=\operatorname {Spec}(B)\)是仿射概形, \(\mathcal {R}\)是\(X\)上的一个预等价关系, 且其分量\(R_1\)仿射, 记为\(R_1=\operatorname {Spec}(C)\). 假设第一投影\(p_1\colon  R_1\to  X\)是有限局部自由态射, 即相应的环同态\(p'_1\colon  B\to  C\)使\(C\)成为有限型投射\(B\)-模. 设\(A\)是\(B\)中由同态对\(p'_1,p'_2\colon  B\rightrightarrows  C\)的核给出的子环 (即满足\(p'_1(b)=p'_2(b)\)的元素\(b\)的集合). 置\(Y=\operatorname {Spec}(A)\), 以\(f\colon  X\to  Y\)表示由\(A\)嵌入\(B\)定义的态射. 此时有: 
\begin{enumerate}
      
 \item[i.]\(B\)在\(A\)上整, 即\(f\)是整态射.

      
 \item[ii.]态射\(f\)是满的, 其纤维是\(X\)模\(\mathcal {R}\)的集合论等价类\(p_2(p_1^{-1}(x))\), 而\(Y\)的拓扑是\(X\)的拓扑的商拓扑.

      
 \item[iii.]\(Y\)是预概形范畴中\(X\)关于\(\mathcal {R}\)的商.

      
 \item[iv.]若\(\mathcal {R}\)来自一个等价\emph{关系}, 则态射\(f\colon  X\to  Y\)有限局部自由 (即\(B\)是有限型投射\(A\)-模), 且该等价关系有效, 即\(R_1\to  X\times _Y X\)是同构.

    
\end{enumerate}
\end{theorem}
\par{}这个定理推广了有限群\(G\)以自同构作用于环\(B\)、\(A\)为不变元环时的熟知定理, 其证明与已知的证明类似. 还可如下精化(iii):
\begin{corollary}{5.2}\setnumber{5.2}\label{fga3.iii-5-corollary-5.2}\skipsinglepar
\par{}典范态射\(R_1\to  X\times _Y X\)是\emph{满}的.
\end{corollary}
\par{}仍设\(\mathcal {R}\)是\(X\)上一个"有限局部自由"的预等价关系, 但现在\(X\)是任意预概形. 假设能找到预概形\(Y\)及态射\(f\colon  X\to  Y\), 使\(fp_1=fp_2\), 并使\(Y\)上的环层同态序列\[     \mathcal {O}_Y \to  f_*(\mathcal {O}_X) \rightrightarrows  g_*(\mathcal {O}_R)   \]正合 (其中\(g=fp_i\)). 由上述定理可得与该定理(i)至(iv)相应的结论; 特别由(iii), \(Y\)是\(X\)关于\(\mathcal {R}\)的商, 因而在相差唯一同构的意义下确定. 在这些条件下, 称\(X\)上的预等价关系\(\mathcal {R}\)为\emph{容许}的. 按此定义:
\begin{theorem}{5.3}\setnumber{5.3}\label{fga3.iii-5-theorem-5.3}\skipsinglepar
\par{}设\(X\)是预概形, \(\mathcal {R}\)是\(X\)上的预等价关系, 且\(p_1\colon  R_1\to  X\)有限局部自由. \(\mathcal {R}\)容许的充要条件是\(X\)模\(\mathcal {R}\)的每个集合论等价类\(p_2(p_1^{-1}(x))\)都包含在一个仿射开子集中 (如果\(X\)的每个有限子集都包含于一个仿射开子集, 这个条件总能满足, 例如\(X\)在某个仿射概形上拟射影时).
\end{theorem}
\par{}事实上容易证明, \(X\)模\(\mathcal {R}\)的每个等价类此时均包含于一个在\(\mathcal {R}\)下\emph{稳定}的仿射开子集, 于是将应用\hyperref[fga3.iii-5-theorem-5.1]{定理~\ref*{fga3.iii-5-theorem-5.1}}得到的各块粘合起来, 就构造出商\(Y\).
\begin{corollary}{5.4}\setnumber{5.4}\label{fga3.iii-5-corollary-5.4}\skipsinglepar
\par{}假设这个条件满足, 并且\(\mathcal {R}\)来自一个等价关系. 则该等价关系有效, 即\(R_1\to  X\times _Y X\)是同构, 且\(f\colon  X\to  Y\)是有限局部自由态射.
\end{corollary}
\par{}于是由下降立即得到:
\begin{corollary}{5.5}\setnumber{5.5}\label{fga3.iii-5-corollary-5.5}\skipsinglepar
\par{}在\hyperref[fga3.iii-5-corollary-5.4]{推论~\ref*{fga3.iii-5-corollary-5.4}}的条件下, \(X\)在\(Y\)上处处秩为\(n\)的充要条件是\((R_1,p_1)\)在\(X\)上处处秩为\(n\). 若\(X\)和\(R_1\)是\(Z\)-预概形, 且\(p_1\)、\(p_2\)是\(Z\)-态射 (因而\(Y\)是\(Z\)-预概形), 则\(X\)在\(Z\)上平坦当且仅当\(Y\)在\(Z\)上平坦.
\end{corollary}
\par{}归纳如下:
\begin{scholium}{}\label{fga3.iii-5-scholium}\skipsinglepar
\par{}给出预概形间一个有限、局部自由且满的态射\(f\colon  X\to  Y\), 等价于给出一个预概形\(X\)及其上的等价关系\(R\), 满足\(p_1\colon  R\to  X\)有限局部自由, 且每个等价类\(p_2(p_1^{-1}(x))\)包含于一个仿射开子集.
\end{scholium}
\begin{remarks}{5.6}\setnumber{5.6}\label{fga3.iii-5-remarks-5.6}\skipsinglepar
\begin{enumerate}\item{}我们不需要作任何Noether性假设.

    \item{}这种取商方式包含Cartier的"不可分下降"作为特例: 它相应于确定有限局部自由态射\(f\colon  X\to  Y\), 使得\(f_*(\mathcal {O}_X)\)相对于\(\mathcal {O}_Y\)具有\(p\)-基 (这里\(X\)是给定预概形, 其层\(\mathcal {O}_X\)被素数\(p>0\)零化). 注意, 不对局部环作任何正则性假设, 也不假设\(X\)是域上的代数概形, 也很容易表述这一结果. 取\(X\)为特征\(p\)的一个域的谱, 就得到Jacobson–Bourbaki理论.

    \item{}Gabriel在域\(k\)上有限交换群的取商理论中已经得到\hyperref[fga3.iii-5-theorem-5.3]{定理~\ref*{fga3.iii-5-theorem-5.3}}的一个特例. (比较\hyperref[fga3.iii-7-corollary-7.4]{推论~\ref*{fga3.iii-7-corollary-7.4}}.)
\end{enumerate}
\end{remarks}
\chapter{关于固有平坦等价关系的商}\thispagestyle{fancy}\label{fga3.iii-6}
\begin{theorem}{6.1}\setnumber{6.1}\label{fga3.iii-6-theorem-6.1}\skipsinglepar
\par{}设\(S\)是\emph{局部Noether}预概形, \(X\)是\emph{拟射影}的\(S\)-概形, \(\mathcal {R}\)是\(X\)上的预等价关系, 满足: 
\begin{enumerate}
      
 \item[a.]\(p_1\colon  R_1\to  X\)固有且平坦;

      
 \item[b.]\(R_1\to  X\times _S X\)是有限态射 (由(a), 等价地说, 是纤维有限的态射; 若\(\mathcal {R}\)来自等价关系, 则此条件自动满足).

    
\end{enumerate} 在这些条件下: 
\begin{enumerate}
      
 \item[i.]\(Y=X/\mathcal {R}\)存在, 且当\(S\)为Noether时, 它在\(S\)上拟射影. (\emph{[Trans.] \emph{[Comp.]} 关于\(Y=X/\mathcal {R}\)在\(S\)上拟射影的事实, 目前只在\(\mathcal {R}\)来自等价关系的情形得到证明.})

      
 \item[ii.]典范态射\(f\colon  X\to  Y\)满、固有且开, 其纤维是\(X\)模\(\mathcal {R}\)的等价类\(p_2(p_1^{-1}(x))\), 因而\(Y\)可与\(X\)关于\(\mathcal {R}\)所定义的集合论等价关系的拓扑商空间认同. 最后, \(R_1\to  X\times _Y X\)是满的.

      
 \item[iii.]若\(\mathcal {R}\)来自等价关系, 则该等价关系有效, 即\(R_1\to  X\times _Y X\)是同构; 而且\(f\colon  X\to  Y\)平坦, 从而忠实平坦.

    
\end{enumerate}
\begin{proof}
\par{}证明时, 可考虑\(X\)关于\(\mathcal {R}\)的适当拟截面, 将问题归结到\hyperref[fga3.iii-5-theorem-5.1]{定理~\ref*{fga3.iii-5-theorem-5.1}}, 其证明类似于\emph{Séminaire Chevalley}中代数商群的构造.
\end{proof}
\end{theorem}
\par{}归纳如下:
\begin{scholium}{}\label{fga3.iii-6-scholium}\skipsinglepar
\par{}设\(X\)在\(S\)上拟射影, 且\(S\)局部Noether. 则给出从\(X\)到某个\(S\)-预概形\(Y\)的一个固有、忠实平坦且满的态射\(f\colon  X\to  Y\), 等价于给出\(X\)上的一个等价关系\(R\), 使得\(p_1\colon  R\to  X\)固有且平坦.
\end{scholium}
\par{}同一方法还给出以下结果:
\begin{theorem}{6.2}\setnumber{6.2}\label{fga3.iii-6-theorem-6.2}\skipsinglepar
\par{}设\(S\)是Noether预概形, \(X\)是在\(S\)上有限型的预概形, \(\mathcal {R}\)是\(S\)-预概形\(X\)上的预等价关系. 假设
\begin{enumerate}
      
 \item[a.]\(p_1\colon  R_1\to  X\)平坦且有限型;

      
 \item[b.]态射\(R_1\to  X\times _S X\)拟有限, 即其纤维有限.

    
\end{enumerate}. 则\(X\)中存在一个关于\(\mathcal {R}\)\emph{饱和}的\emph{稠密}开子集\(U\), 使得: 
\begin{enumerate}
      
 \item[i.]若\(\mathcal {R}_U\)是\(\mathcal {R}\)在\(U\)上诱导的预等价关系, 则\(U/\mathcal {R}_U\)存在, 并在\(S\)上有限型.

      
 \item[ii.]典范态射\(U\to  U/\mathcal {R}_U\)满且开, 其纤维是\(\mathcal {R}_U\)的集合论等价类 (因此\(U/\mathcal {R}_U\)是\(U\)关于\(\mathcal {R}_U\)所定义的集合论等价关系的拓扑商空间). 最后, 态射\((\mathcal {R}_U)_1\to  U\times _{U/\mathcal {R}_U}U\)是满的.

      
 \item[iii.]若\(\mathcal {R}\)来自等价关系, 则可以假设\(U\to  U/\mathcal {R}_U\)忠实平坦, 且\(\mathcal {R}_U\)有效.

    
\end{enumerate}
\end{theorem}
\par{}这是一个本质上具有"双有理"性质的结果.
\begin{remarks}{6.3}\setnumber{6.3}\label{fga3.iii-6-remarks-6.3}\skipsinglepar
\begin{enumerate}\item{}我不知道在\hyperref[fga3.iii-6-theorem-6.1]{定理~\ref*{fga3.iii-6-theorem-6.1}}和\hyperref[fga3.iii-6-theorem-6.2]{定理~\ref*{fga3.iii-6-theorem-6.2}}中, 假设(b)是否多余. 在实际的群的取商中, 这个假设迫使我们限于所有稳定子都是有限群的情形.

    \item{}可以问, 在没有任何平坦性假设时, 是否也有类似于\hyperref[fga3.iii-6-theorem-6.1]{定理~\ref*{fga3.iii-6-theorem-6.1}}和\hyperref[fga3.iii-6-theorem-6.2]{定理~\ref*{fga3.iii-6-theorem-6.2}}的结果. 我在这一方向还没有反例. 但是, 即使保留平坦性假设, 限于\(p_1\colon  R\to  X\)平坦拟有限 (但不有限) 且\(X\)仿射的等价关系, \(R\)仍可能无效: 只需取覆盖Nagata簇的仿射开子集上所诱导的等价关系 (在该簇上, 一个二元群以"非容许"的方式作用).
\end{enumerate}
\end{remarks}
\chapter{应用}\thispagestyle{fancy}\label{fga3.iii-7}
\par{}如引言所述, \hyperref[fga3.iii-6-theorem-6.1]{定理~\ref*{fga3.iii-6-theorem-6.1}}最重要的应用是构造Picard概形以及解决其他各种"模"问题; 我们以后还会回到这些问题.
\par{}我们得到Shimura的下述结果的一个简单证明:
\begin{proposition}{7.1}\setnumber{7.1}\label{fga3.iii-7-proposition-7.1}\skipsinglepar
\par{}设\(A\)是在分式域为\(K\)的离散赋值环\(V\)上定义的阿贝尔概形. 则每个与\(A\otimes _V K\)的某个商同源的\(K\)上阿贝尔概形\(B'\), 都"在\(V\)处有好约化", 即它同构于某个\(B\otimes _V A\), 其中\(B\)是\(V\)上的阿贝尔概形 (回顾一下, 它本质上唯一).
\begin{proof}
\par{}可以假设\(B'\)是\(A_K\)关于一个子群概形\(C'\)的商. (注意, \(C'\)一般不"约化", 即其局部环可能有幂零元.) 考虑\(A\)中由\(C'\)的"闭包"给出的闭子概形\(C\), 即使\(C_K\)包含\(C'\)的\(A\)的最小闭子概形. 则\(C_K=C'\); 又因\(V\)是离散赋值环, 容易推出\(C\)是\(A\)在\(V\)上的子群概形. 由于\(A\)在\(\operatorname {Spec}(V)=S\)上固有, \(C\)也固有. 此外, \(A\)在\(S\)上射影. 于是可以应用\hyperref[fga3.iii-6-theorem-6.1]{定理~\ref*{fga3.iii-6-theorem-6.1}}构造\(A/C=B\), 这就是所求的\(B\).
\end{proof}
\end{proposition}
\par{}最后, 用基本上已知的论证可从\hyperref[fga3.iii-6-theorem-6.2]{定理~\ref*{fga3.iii-6-theorem-6.2}}得出以下结果:
\begin{theorem}{7.2}\setnumber{7.2}\label{fga3.iii-7-theorem-7.2}\skipsinglepar
\par{}设\(S\)是Artin环的谱, \(F\)和\(G\)是在\(S\)上有限型的群概形, \(u\colon  F\to  G\)是\(S\)上的群概形同态. 假设
\begin{enumerate}
      
 \item[i.]\(F\)在\(S\)上平坦;

      
 \item[ii.]\(u\)的核有限.

    
\end{enumerate}. 在这些条件下, 商概形\(G/F\)存在; 典范态射\(G\to  G/F\)满且开, 其纤维是\(F\)在\(G\)上的右作用定义的集合论等价类. 最后, 若\(u\)是单态射, 则态射\(G\to  G/F\)平坦, 而态射\(G\times  F\to  G\times _{(G/F)}G\)是同构; 换言之, \(G\)是\(G/F\)上以\(F\) (右作用) 为结构群的主齐性空间, 更准确地说, 结构群是视为\(G/F\)上群概形的\(F\times _S(G/F)\) (参见\hyperref[fga3.i-b.6]{[1], B, §6}).
\end{theorem}
\begin{corollary}{7.3}\setnumber{7.3}\label{fga3.iii-7-corollary-7.3}\skipsinglepar
\par{}在这些条件下, \(G\)在\(S\)上平坦的充要条件是\(G/F\)在\(S\)上平坦. 若此条件满足, 则关于\(F\)取商与基底\(S\)的任意扩张交换; 若\(F\)是\(G\)的不变子群, 则\(G/F\)具有\(G\)关于\(F\)的\emph{商群}结构.
\end{corollary}
\par{}当\(S\)是域的谱时情形尤为简单, 因为此时每个\(S\)-预概形自动在\(S\)上平坦. 我们得到:
\begin{corollary}{7.4}\setnumber{7.4}\label{fga3.iii-7-corollary-7.4}\skipsinglepar
\par{}设\(F\)和\(G\)是域\(k\)上的有限型群概形, \(u\colon  F\to  G\)是\(k\)-群同态. 则\(u\)可分解为\(F\to  F'\to  G\), 其中\(F\to  F'\)是关于\(F\)的闭子群\(\operatorname {Ker} u\)取商所得的同态, 而\(F'\to  G\)是一个闭浸入的群同态. 商\(G/F=G/F'\)存在. 通常的形式运算 (例如Noether同构定理中的运算) 在\(k\)上的代数群之间成立.
\end{corollary}
\par{}这个结果使我们可以统一处理经典意义的代数群 (即在\(k\)上不可约且在\(k\)上光滑的群) 的取商, 以及Cartier所考虑的关于"无穷小"子群的取商. 将Cartier在Dieudonné关于形式群的工作之后引入的"超代数"视为\(k\)上形式概形范畴中的群, 并在适当情形 (当它们对应于\(k\)上秩有限的超代数时) 视为\(k\)上有限的代数群, 是有利的.
\chapter{一个猜想}\thispagestyle{fancy}\label{fga3.iii-8}
\begin{remark}{}\label{fga3.iii-8-remark-i}\skipsinglepar
\par{}\emph{[Comp.]} 现在看来, 本节的猜想即使对于特征\(0\)的域上的非奇异簇也不成立: 无论商的存在性还是其拟射影性都可能失败, 甚至在\(G\)的作用具有闭图时也是如此.
\end{remark}
\par{}这里的猜想源于一个需要: 要知道如何对作用在"Hilbert概形"的某些子概形上的射影群取商 (在概形理论中, 这些"Hilbert概形"用来替代Chow簇).
\par{}设\(S\)是预概形, \(n\)是整数. 对每个\(S\)上的预概形\(S'\), 配以元素取值于\(\mathcal {O}_{S'}\)的截面环的可逆\((n\times  n)\)矩阵群\(\operatorname {GL}(n,\Gamma (S',\mathcal {O}_{S'}))\). 这样得到一个关于\(S'\)的反变函子, 很容易证明它可表示; 因而它对应于\(S\)上的一个群概形 (并且在\(S\)上仿射), 记作\(\operatorname {GL}(n)_S\). 它的构造与基底变换相容, 所以实际上所有这些均来自\(\mathbb {Z}\)上的一个群概形, 记作\(\operatorname {GL}(n)\). 群\(\operatorname {GL}(1)\)称为\emph{乘法群}, 常记作\(\operatorname {G_m}\), 它对应于函子\(S\mapsto \Gamma (S,\mathcal {O}_S)^*\), 后者是\(S\)上的"单位"群. 有一个显然的同态\(\operatorname {GL}(1)\to \operatorname {GL}(n)\), 且很容易构造其商群, 记作\(\operatorname {GP}(n-1)\), 称为\emph{\(\mathbb {Z}\)上次数为\((n-1)\)的射影群}. 它表示把\(S\)对应于层\(\operatorname {\mathscr {G}\kern  -2pt\mathscr {L}}(n)_S/\operatorname {\mathscr {G}\kern  -2pt\mathscr {L}}(1)_S\)的截面群的函子; 这里\(\operatorname {\mathscr {G}\kern  -2pt\mathscr {L}}(n)_S\)表示\(\operatorname {GL}(n)_S\)在\(S\)上的截面芽层. (注意, \(\operatorname {GP}(n-1)_S\)在\(S\)上的截面一般不来自\(\operatorname {GL}(n)_S\)在\(S\)上的截面!) 还可证明, 至少当\(S\)为Noether时, \(\operatorname {GP}(n-1)\)也表示函子\(S\mapsto \operatorname {Aut}_S(\mathbb {P}_S^{n-1})\) (其中\(\mathbb {P}_S^{n-1}\)是在\(S\)上相对维数为\((n-1)\)的射影型概形). 它正是这样出现在模理论中的.
\par{}设\(S\)是一个Noether概形, 也可以假设它仿射; 设\(X\)是一个拟射影的\(S\)-预概形, 带有关于\(S\)非常丰沛的可逆层\(\mathscr {L}\). 假设群\(G=\operatorname {GP}(n)_S\)同时作用于\(X\)及\(\mathscr {L}\) (其作用与\(X\)上的作用相容), 并在\(S\)上\emph{自由地}作用.
\begin{conjecture}{8.1}\setnumber{8.1}\label{fga3.iii-8-conjecture-8.1}\skipsinglepar
\par{}在上述条件下: 
\begin{enumerate}\item{}\(G\)定义的等价关系有效, 商\(Y=X/G\)在\(S\)上有限型, 且典范态射\(f\colon  X\to  Y\)平坦且满 (因而\(X\)成为\(Y\)上的齐性主丛, 群为\(G\times _S Y=\operatorname {GP}(n)_Y\)).

      \item{}设\(\mathscr {L}'\)是通过\(f\)下的"忠实平坦下降"从\(\mathscr {L}\)在\(Y\)上诱导的可逆层 (参见\hyperref[fga3.i-b.1-theorem-1]{[1], B, 定理1}). 则\(\mathscr {L}'\)在\(Y\)上关于\(S\)"预丰沛", 即存在一个整数\(m\)和从\(Y\)到某个适当射影型概形\(\mathbb {P}_S^N\)的拟有限态射, 使\((\mathscr {L}')^{\otimes  m}\)同构于\(\mathcal {O}_{\mathbb {P}_S^N}(1)\)的逆像.
\end{enumerate}
\end{conjecture}
\par{}应注意, 即使\(X\)在\(S\)上分离, \(Y\)也可能不在\(S\)上分离 (这种情形出现在完全不病态的"模问题"中). 如果(1)成立, 则\(Y\)分离当且仅当\(G\)所定义的等价关系具有闭图, 即\(G\times  X\to  X\times  X\)的像闭 (从而它是闭浸入). 若\(Y\)分离, 则\(\mathscr {L}'\)在\(Y\)上关于\(S\)预丰沛当且仅当它丰沛, 即某个适当张量幂定义一个射影浸入. 在引言提到的模问题中, 可以证明最终得到的等价关系确实具有闭图.
\begin{remarks}{8.1}\setnumber{8.1}\label{fga3.iii-8-remarks-8.2}\skipsinglepar
\par{}我们假设\(G=\operatorname {GP}(n)_S\), 是为了给出具体例子, 也因为这在目前实际应用中是最重要的情形. 对\(G\)作的合理假设, 似乎应是\(G\)为某个Tohokû群在\(S\)上的一种"形式" (Chevalley也已给出这些群在整数环上的构造). 关于上述猜想, 我所知道的唯一肯定事实如下: \emph{设\(X\)是特征\(0\)的域\(k\)上的仿射概形, 群\(\operatorname {GL}(n)_k\)或\(\operatorname {GP}(n-1)_k\)在其上自由作用. 则\(G\)定义的等价关系有效, 商\(X/G\)仿射, 且态射\(X\to  X/G\)平坦且满.} 证明使用下面的事实 (目前仅在特征\(0\)时得到证明): 让\(G\)作用于其仿射环\(A\) (此环为\(G\)的仿射环, 视为\(k\)上的向量空间), 则\(G\)的平凡表示仅出现一次 (在\(G\)下稳定的某个\(k\)上有限维向量子空间的合成列中). 系统地运用\(G\)的线性表示理论, 似乎可能证明这个猜想, 至少在基底是域时如此. 基底不再是域时, 作者一无所知.
\end{remarks}
\begin{remark}{}\label{fga3.iii-8-remark-ii}\skipsinglepar
\par{}\emph{[Comp.]} 正如下一讲末尾所指出的, 上述猜想确定是错误的. 上一注中的"肯定事实"似乎由不同作者同时证明 (Nagata、Rosenlicht、Grothendieck等).
\end{remark}

\section*{参考文献}

\begin{enumerate}
\item[{[1]}]\setnumber{1}\label{fga3.iii-bib-1} GROTHENDIECK (Alexander). --- Techniques de descente et théorèmes d'existence en géométrie algébrique, I: Généralités, \emph{Séminaire Bourbaki}, 1959/60, n\textsuperscript{o} 190, 29页.
\item[{[2]}]\setnumber{2}\label{fga3.iii-bib-2} GROTHENDIECK (Alexander). --- Techniques de descente et théorèmes d'existence en géométrie algébrique, II: Le théorème d'existence en théorie formelle des modules, \emph{Séminaire Bourbaki}, 1959/60, n\textsuperscript{o} 195, 22页.
\item[{[3]}]\setnumber{3}\label{fga3.iii-bib-3} GROTHENDIECK (A.) 与 DIEUDONNÉ (J.). --- \emph{Éléments de géométrie algébrique}, I: Le langage des schémas. Paris, Presses universitaires de France, 1960 (Institut des hautes Études scientifiques, Publications mathématiques, 4).

\end{enumerate}

% END chapters/fga3.iii.tex

\clearpage
\setcounter{chapter}{-1}
\renewcommand\fgaAnchorPrefix{fga3.iv}

% BEGIN chapters/fga3.iv.tex
\part{Hilbert概形}\label{fga3.iv}\thispagestyle{fancy}
\begin{original}{}\label{fga3.iv-original-citation}\skipsinglepar
\par{}A. Grothendieck. "Techniques de construction et théorèmes d'existence en géométrie algébrique, IV: Les schémas de Hilbert". \emph{Séminaire Bourbaki} \textbf{13} (1960–61), 第221讲, 1961年5月. \url{http://www.numdam.org/book-part/SB_1960-1961__6__249_0/}
\end{original}
\chapter{引言}\thispagestyle{fancy}\label{fga3.iv-introduction}
\par{}\hyperref[fga3.i]{[2, I]}和\hyperref[fga3.ii]{[2, II]}所述的方法, 大体上不依赖于对所讨论概形作任何射影性假设. 可惜它们至今尚未使我们解决\hyperref[fga3.ii]{[2, II]}所提出的存在性问题. 在本文和下一篇文章中, 我们将通过加上射影性假设来解决这些问题. 所用方法具有典型的射影性质, 实际上几乎不使用\hyperref[fga3.i]{[2, I]}和\hyperref[fga3.ii]{[2, II]}的结果. 这里将构造"Hilbert概形", 用来替代Chow坐标, 这在\hyperref[fga3.ii-c.2]{[2], II, C, §2}中已经提到. 在下一篇文章中, 将\hyperref[fga3.iii]{[2, III]}所建立的概形取商理论与Hilbert概形理论结合, 就能例如在相当一般的条件下构造Picard概形 (定义见\hyperref[fga3.ii-c.3]{[2], II, C, §3}).
\par{}总的说来, 现在我们已有一种大体令人满意的射影构造方法, 只是仍缺少\begingroup\renewcommand{\thefootnote}{*}\footnote{参见本讲末尾的增补。}\endgroup{}对"无不动点地"作用的射影群等群的取商理论 (参见\hyperref[fga3.iii-8]{[2], III, §8}). 解析几何的情形似乎还稍好一些 (若限于研究给定解析空间上的"射影"解析空间), 因为在解析空间中, 关于作用良好的群取商的困难已不再存在. 无论如何, 在代数几何和解析几何中, 都仍有待建立不需要任何射影性假设的构造方法.
\chapter{层的有界集合: 不变性质}\thispagestyle{fancy}\label{fga3.iv-1}
\par{}设\(k\)是域, \(X\)是\(k\)-预概形 (为简单起见, 假设它有限型). 每个扩张\(K/k\)给出一个\(K\)-预概形\(X_K=X\otimes _k K\). 若\(\mathscr {F}\)是\(X_K\)上的凝聚层, \(K'\)是\(K\)的一个扩张, 则\(\mathscr {F}\otimes _K K'=\mathscr {F}_{K'}\)是\(X_K\otimes _KK'=X_{K'}\)上的拟凝聚层. 因而若\(K\)、\(K'\)是\(k\)的任意扩张, \(\mathscr {F}\)是\(X_K\)上的拟凝聚层, \(\mathscr {F}'\)是\(X_{K'}\)上的拟凝聚层, 并且存在扩张\(K''/k\)和\(k\)-同态\(K\to  K''\)、\(K'\to  K''\), 使\(\mathscr {F}_{K''}\)与\(\mathscr {F}'_{K''}\)在\(X_{K''}\)上同构, 则称\(\mathscr {F}\)与\(\mathscr {F}'\)\emph{等价}. 这样定义了一个等价关系; 我们关心层在此关系下的等价类及这些等价类的集合. 注意, 若\(X_0\)在\(k\)上有限型, 则每个凝聚层的类均可由某个\(X_K\)上的凝聚层定义, 其中\(K\)是\(k\)的一个\emph{有限型}扩张. 因此, 在定义凝聚层的类时, 可以只考虑\(k\)的\emph{代数闭}扩张, 甚至可以只考虑\(k\)的一个固定的、具有无限超越次数的代数闭扩张\(\Omega \); 此时\(X_\Omega \)上两个凝聚层\(\mathscr {F}\)与\(\mathscr {F}'\)等价, 当且仅当存在\(\Omega \)的一个\(K\)-自同构\(\sigma \), 使\(\mathscr {F}\otimes _K(\Omega ,\sigma )\)同构于\(\mathscr {F}'\). 注意第一种定义下的凝聚层的类与第二种定义下的类之间有一个双射对应.
\par{}设\(E\)和\(E'\)是\(X\)上的两组凝聚层的类. 考虑所有形如\(\mathscr {F}\otimes \mathscr {F}'\)的层的类, 其中\(\mathscr {F}\)和\(\mathscr {F}'\)是\emph{同一个}的\(X_K\)上的凝聚层, 且\(\mathscr {F}\)的类属于\(E\), \(\mathscr {F}'\)的类属于\(E'\). 这样定义一个凝聚层的类的集合, 记作\(E\otimes  E'\). 同样可定义\(\mathscr {T}\kern  -2.75pt\mathit {or}_i(E,E')\)等. 一般地, 若一个运算\(\mathscr {U}\)把同一个\(X_K\)上的\(n\)个凝聚层组成的序列\(\mathscr {F}_1,\ldots ,\mathscr {F}_n\), 对应于\(X_K\)上的一个凝聚层集合\(\mathscr {U}(\mathscr {F}_1,\ldots ,\mathscr {F}_n)\), 并具有与层的同构及基底变换下的逆像相容这一显然的性质, 则给它配以一个仍记作\(\mathscr {U}\)的运算, 把\(n\)个凝聚层的类的集合所组成的序列\(E_1,\ldots ,E_n\)对应于一个凝聚层的类的集合\(\mathscr {U}(E_1,\ldots ,E_n)\).
\par{}本节的目的是定义某些称为\emph{有界}的层的类的集合, 并证明最常见的运算\(\mathscr {U}\)施于有界集合时, 所得集合仍然有界. 设\(X\)是在\(S\)上有限型的预概形, 且\(S\)为\emph{Noether}. 对每个\(s\in  S\), 纤维\(X_s\)是在\(k(s)\)上有限型的预概形; 我们考虑上述意义下\(X_s\)上的凝聚层的类. 这样,"\(X/S\)的一个纤维上的凝聚层的类"及类似说法便有了意义. 同样, 在各个纤维上分别进行上述步骤, 仍可考虑\(E\otimes  E'\)等运算: 它们把\(X/S\)的纤维上的凝聚层的类的集合族对应于\(X/S\)的纤维上的凝聚层的类的集合.
\begin{definition}{1.1}\setnumber{1.1}\label{fga3.iv-1-definition-1.1}\skipsinglepar
\par{}设\(E\)是\(X/S\)的纤维上的凝聚层的类的一个集合. 若存在\(S\)上的有限型预概形\(S'\)和\(X'=X\times _S S'\)上的凝聚层\(\mathscr {F}'\), 使\(E\)包含于\(\mathscr {F}'\)所定义的\(X/S\)的纤维上的层的类的集合中, 就称\(E\)是\emph{有界}的.
\end{definition}
\par{}按定义, 这个构造给\(s\in  S\)配以各层\(\mathscr {F}'\otimes _{S'}k(s')\)的类, 其中\(s'\)遍历\(S'\)中位于\(s\)之上的点 (于是\(k(s')\)是\(k(s)\)的扩张, 而\(X\otimes  (X_s)_{k(s')}\)可与\(X'\)在\(s'\)处的纤维\(X'\otimes _{S'}k(s')=X'_{s'}\)认同). 也可以说, 有界集合就是那些包含于某个凝聚层的\emph{代数族}中的集合, 此代数族由\(S\)上有限型的某个\(S'\)参数化.
\par{}有限多个有界集合的并有界 (取定义各个界定代数族的参数预概形\(S_i\)的和即可). 基底变换\(T\to  S\)把关于\(X/S\)有界的族变为关于\(X_T/T\)有界的族; 若\(T\to  S\)是满的, 则逆命题也成立 (更一般地, 只需它的像包含\(X/S\)的给定族\(E\)中出现的各个\(s\)). 因此在理论上, 只需在\(S\)为整数环\(\mathbb {Z}\)上有限型代数的谱的情形确定有界族.
\par{}若\(E\)和\(E'\)是关于\(X/S\)的层的类的有界族, 则\(E\otimes  E'\)也有界: 事实上, 若\(E\) (分别为\(E'\)) 被由\(T\to  S\)及\(X_T\)上的\(\mathscr {F}\) (分别为\(T'\to  S\)及\(X_{T'}\)上的\(\mathscr {F}'\)) 定义的代数族所界定, 则\(E\otimes  E'\)被下述代数族界定: 取\(T''=T\times _S T'\to  S\), 并在\(X_{T''}\)上取\(\mathscr {F}\)和\(\mathscr {F}'\)在\(X_{T''}\)上的逆像的张量积为层\(\mathscr {F}''\). 此论证成立, 正是因为函子\(\mathscr {F}\otimes \mathscr {F}'\)对\(\mathscr {F}\)和\(\mathscr {F}'\)均右正合, 因而与基底扩张交换, 特别与取纤维交换. 它不能原样用于\(\mathscr {T}\kern  -2.75pt\mathit {or}_i(E,E')\)、\(\mathscr {H}\kern  -2.5pt\mathit {om}(E,E')\)、\(\mathscr {E}\kern  -1.5pt\mathit {xt}^i(E,E')\)等局部运算. 不过, 按照\(E\otimes  E'\)的处理方式, 再利用下面这类结果 (均见[\hyperref[Gro1960b]{3}, IV 6.11]), 就可证明这些运算也把有界集合变为有界集合: 有界族\(E\)总能被某个\(X_T\)上的凝聚层\(\mathscr {F}\)所定义的代数族界定, 其中\(T\)在\(S\)上有限型, 且该层关于\(T\)\emph{平坦}. (将原来的参数空间"分成若干块".) 对适当的层要求这种平坦性, 确能保证\(\mathscr {T}\kern  -2.75pt\mathit {or}_i(\mathscr {F},\mathscr {F}')\)等运算与任意基底变换交换. 同一方法也适用于全局性质的运算: 固有态射下凝聚层的正像和导出正像、关于固有态射的全局\(\operatorname {Ext}\) (参见[\hyperref[GD1960]{5}, III §6]) 等; 所有这些运算都把层的有界族变为层的有界族 (注意: 这些运算可能改变承载所讨论各层的预概形).
\par{}下面两个断言可以用\emph{本质上}相同的平坦性方法证明; 关于有限型态射的纤维上的准素分解, 特别参见[\hyperref[GD1960]{5}, IV].
\begin{proposition}{1.2}\setnumber{1.2}\label{fga3.iv-1-proposition-1.2}\skipsinglepar
\par{}设\(E\)和\(E'\)是\(X/S\)的纤维上的层的类的有界集合, 并假设\(X\)在\(S\)上固有. 则
\begin{enumerate}
      
 \item[i.]同态\(\mathscr {F}\to \mathscr {F}'\)的核、余核和像所成的族有界 (其中\(\mathscr {F}\)的类属于\(E\), \(\mathscr {F}'\)的类属于\(E'\));

      
 \item[ii.]\(\mathscr {F}\)由\(\mathscr {F}'\)的扩张\(\mathscr {F}''\)所成的族有界 (其中\(\mathscr {F}\)的类属于\(E\), \(\mathscr {F}'\)的类属于\(E'\)).

    
\end{enumerate}
\begin{proof}
\par{}必要时作适当的基底变换, 可以假设\(E\)和\(E'\)分别由某个\(X_T/T\)上的凝聚层\(\mathscr {G}\)和\(\mathscr {G}'\)定义, 其中\(T\)在\(S\)上有限型. 还可以假设满足某些平坦性条件, 从而层\(f_T(\mathscr {H}\kern  -2.5pt\mathit {om}_{\mathcal {O}_{X_T}}(\mathscr {G},\mathscr {G}'))\)及\(\mathscr {E}\kern  -1.5pt\mathit {xt}_{f_T}^1(\mathscr {G},\mathscr {G}')\)的构造与任意态射\(T'\to  T\)所给出的基底变换交换. 又可假设上述凝聚层在\(T\)上局部自由. 设\(T_0\)、\(T_1\)为\(T\)上的向量丛, 它们的截面芽层分别是上述两层. 此时可以典范地定义\(X_{T_0}\)上的凝聚层同态\(\mathscr {G}_{T_0}\to \mathscr {G}'_{T_0}\), 以及\(X_{T_1}\)上的凝聚层扩张\[         0 \to  \mathscr {G}_{T_1} \to  \mathscr {G}'' \to  \mathscr {G}'_{T_1} \to  0,\] 后者具有显然的泛性质. 这个扩张中的层定义了一个代数族, 它界定(ii)中的族. 上述同态的核、余核和像也有同样性质; 考虑它们就证明(i) (这里须假设余核关于\(T_0\)平坦, 而这仍可归结到把\(T_0\)分块的办法...).
\end{proof}
\end{proposition}
\begin{proposition}{1.3}\setnumber{1.3}\label{fga3.iv-1-proposition-1.3}\skipsinglepar
\par{}设\(E\)是\(X/S\)的纤维上的层的类的有界族. 则各个\((\operatorname {supp}\mathscr {F})_\mathrm {red}\)的结构层的类组成一个有界族, 其中\(\mathscr {F}\)是某个\(X_K\)上的凝聚层, \(K\)代数闭, 且其类属于\(E\).
\end{proposition}
\par{}这里\((\operatorname {supp}\mathscr {F})_\mathrm {red}\)表示\(\mathscr {F}\)的支集, 赋以诱导的约化结构; 也就是说, 其结构层是\(\mathcal {O}_{X_K}\)模掉定义\(\operatorname {supp}\mathscr {F}\)的最大理想层所得的商. 对由准素分解理论从\(\mathscr {F}\)典范地导出的层, 可以证明与\hyperref[fga3.iv-1-proposition-1.3]{命题~\ref*{fga3.iv-1-proposition-1.3}}相类似的结果; 例如\(\mathscr {F}/\mathscr {F}_i\), 其中\(\mathscr {F}_i\)是\(\mathscr {F}\)中对应于\(\mathscr {F}\)的支集的各分支、并且关于此性质极小的准素子层; 又例如\(\mathcal {O}_{X_K}/\mathfrak {p}\), 其中\(\mathfrak {p}\)是\(\mathscr {F}\)的一个相伴素理想层; 或\(\mathcal {O}_{X_K}/\mathfrak {q}\), 其中\(\mathfrak {q}\)是对应于\(\mathscr {F}\)的支集的某个分支的准素理想层 (所取的基域为代数闭域).
\chapter{有界族与Hilbert多项式}\thispagestyle{fancy}\label{fga3.iv-2}
\par{}以下假设\(X\)在\(S\)上射影, 并带有一个非常丰沛层, 记作\(\mathcal {O}_X(1)\). 对\(S\)的一点\(s\)的剩余域\(k(s)\)的每个扩张\(K\), 考虑\(X_K\)上相应的层\(\mathcal {O}_{X_K}(1)\), 它仍然非常丰沛.
\par{}对\(X_K\)上每个凝聚层\(\mathscr {F}\), 配以函数\[     P_\mathscr {F}(n)     = \text {层 }\mathscr {F}(n)\text {在 }X_K\text { 上的Euler–Poincaré特征数}.\] 它是关于整数\(n\)的多项式, 称为\(\mathscr {F}\)的\emph{Hilbert多项式}. 对充分大的\(n\), \(P(n)\)恰为\(\operatorname {H}^0(X_K,\mathscr {F}(n))\)在\(K\)上的维数, 因为对\(i>0\)和充分大的\(n\), 各\(\operatorname {H}^i(X_K,\mathscr {F}(n))\)为零.
\par{}若\(\mathscr {F}\)是\(X\)上关于\(S\)\emph{平坦}的凝聚层, 则在纤维\(X_s\)上诱导的层\(\mathscr {F}_s\)的Hilbert多项式 (对于\(S\)的同一个连通分支) 均相同 [\hyperref[GD1960]{5}, III, §7]. 因而对\(X\)上的每个凝聚层\(\mathscr {F}\), 各层\(\mathscr {F}_s\) (\(s\in  S\)) 的Hilbert多项式的集合为有限集, 这里不需要平坦性假设.
\par{}还须回顾, \(X\)上的任意凝聚层\(\mathscr {F}\)都同构于某个形如\(\mathcal {O}_X(-n)^N\)的层的商层, 其中\(n,N\)充分大. 因而纤维上诱导的层\(\mathscr {F}_s\)也都是纤维上相应的层\(\mathcal {O}(-n)\)的商.
\par{}从这两点得到下述定理中条件的必要性:
\begin{theorem}{2.1}\setnumber{2.1}\label{fga3.iv-2-theorem-2.1}\skipsinglepar
\par{}设\(X\)在\(S\)上射影, \(S\)为Noether, 且\(\mathcal {O}_X(1)\)是\(X\)上关于\(S\)非常丰沛的层. 设\(E\)是\(X/S\)的纤维上的层的类的一个集合. \(E\)有界的充要条件是满足下列条件: 
\begin{enumerate}
      
 \item[a.]\(X\)上存在一个凝聚层\(\mathscr {L}\) (可以假设形如\(\mathcal {O}_X(-n)^N\)), 使得\(E\)包含于形如\(\mathscr {L}_K\)的层的凝聚商层的类所组成的族;

      
 \item[b.]类属于\(E\)的各层\(\mathscr {F}\)的Hilbert多项式\(P_\mathscr {F}\), 都属于同一个\emph{有限}多项式集合.

    
\end{enumerate}
\end{theorem}
\par{}尚须证明条件的充分性, 这将是一个更精确结果的特例. 对域\(K\)上有限型预概形上的每个凝聚模\(\mathscr {F}\), 以及每个整数\(r\), 以\(N_r\)表示\(\mathscr {F}\)的一个子模: 在任一开子集上的截面, 是\(\mathscr {F}\)在该开子集上支集维数\(<r\)的截面. 于是对\(r>\dim \operatorname {supp}\mathscr {F}\)有\(N_r=\mathscr {F}\), 对\(r\leqslant 0\)有\(N_r=0\), 从而得到\(\mathscr {F}\)的一个有限递增滤过, 其因子\(N_r/N_{r+1}\)的相伴素闭链, 恰为\(\mathscr {F}\)的维数为\(r\)的相伴素闭链. 置\[     \mathscr {F}_{(r)}     = \mathscr {F}/N_r,\] 则\(\mathscr {F}_{(r)}\)的相伴素闭链恰为\(\mathscr {F}\)的维数\(\geqslant  r\)的相伴素闭链; 特别地, \(\mathscr {F}_{(r)}\)等于\(\mathscr {F}\)当且仅当\(\mathscr {F}\)的相伴素闭链的维数均\(\geqslant  r\). 此时有:
\begin{theorem}{2.2}\setnumber{2.2}\label{fga3.iv-2-theorem-2.2}\skipsinglepar
\par{}在\hyperref[fga3.iv-2-theorem-2.1]{定理~\ref*{fga3.iv-2-theorem-2.1}}的条件下, 设\(s\)是整数, 并假设\(E\)满足条件(a)以及(b)的下述弱化形式: 
\begin{enumerate}
      
 \item[bs.]类属于\(E\)的各层\(\mathscr {F}\)的Poincaré多项式\(P_\mathscr {F}\), 其\emph{次数\(\leqslant  (s-1)\)的各项}的系数有界.

    
\end{enumerate} 此时, 各层\(\mathscr {F}_{(s)}\) (其中\(\mathscr {F}\)的类属于\(E\)) 组成一个有界族. 此外, \(P_\mathscr {F}\)的\((s-2)\)次项系数均有下界.
\end{theorem}
\par{}于是:
\begin{corollary}{2.3}\setnumber{2.3}\label{fga3.iv-2-corollary-2.3}\skipsinglepar
\par{}假设类属于\(E\)的各层\(\mathscr {F}\), 其所有相伴素闭链的维数\(d\)均满足\(s\leqslant  d\leqslant  r\). 那么在\hyperref[fga3.iv-2-theorem-2.1]{定理~\ref*{fga3.iv-2-theorem-2.1}}的条件(b)中, 只需考虑\(P_\mathscr {F}\)的从\((s-1)\)次到\(r\)次 (两端包括在内) 的系数.
\end{corollary}
\par{}本节的其余部分用来概述\hyperref[fga3.iv-2-theorem-2.2]{定理~\ref*{fga3.iv-2-theorem-2.2}}的证明. 关键是下面两个引理, 其中第一个是熟知的 (它概括了Chow坐标所包含的有效数学内容):
\begin{lemma}[Chow]{2.4}\setnumber{2.4}\label{fga3.iv-2-lemma-2.4}\skipsinglepar
\par{}考虑纤维\(X_K\)的子概形\(Y\)的结构层, 其中\(K\)是\(S\)的某个剩余域的代数闭扩张, \(Y\)约化且所有分支具有相同维数\(r\) (并把\(\mathcal {O}_Y\)视为\(\mathcal {O}_X\)的商层). 若\(Y\)的次数有界, 则这些\(Y\)组成一个有界族.
\end{lemma}
\par{}这里\(Y\)的次数\(a\)最方便地定义为\(P_{\mathcal {O}_Y}=an^r/r!+\ldots \)中最高次项的系数.
\begin{lemma}{2.5}\setnumber{2.5}\label{fga3.iv-2-lemma-2.5}\skipsinglepar
\par{}设\(\mathscr {L}\)是\(X\)上的凝聚层, \(E\)是层\(\mathscr {L}_K\)的商层\(\mathscr {F}\)的类所组成的一个集合 (其中\(K\)是\(S\)的一个剩余域的扩张). 假设\(X\)在\(S\)上的各纤维的维数\(\leqslant  r\), 并置\[       P_\mathscr {F}(n)       = a_\mathscr {F}n^r/r! + b_\mathscr {F}n^{r-1}/(r-1)! + \text {次数}<r-1\text { 的项}.\] 则系数\(a_\mathscr {F}\)有 (上) 界, 而\(b_\mathscr {F}\)有下界. 若\(b_\mathscr {F}\)有界, 则族\(\mathscr {F}_{(r)}\)有界.
\begin{proof}
\par{}将\(S\)替换为覆盖\(S\)的\(S\)的若干子概形的并, 可以假设存在\emph{有限}态射\(f\colon  X\to \mathbb {P}_S^r\), 使\(\mathcal {O}_X(1)\)同构于\(\mathcal {O}_{\mathbb {P}_S^r}(1)\)的逆像; 因而对于\(X\)上的每个凝聚层\(\mathscr {F}\), 均有\(P_\mathscr {F}=P_{f_*(\mathscr {F})}\). 利用上一节的方法, 也容易证明\(X\)上的层\(\mathscr {F}\)的集合有界, 当且仅当各\(f_*(\mathscr {F})\)的集合有界. 最后, 有\[         f_*(\mathscr {F})_{(r)}         = f_*(\mathscr {F}_{(r)}).\] 因此可归结到\(X=\mathbb {P}_S^r\)的情形. 还可对适当的\(k\)和\(s\)假设\(\mathscr {L}=\mathcal {O}_{\mathbb {P}_S^r}(k)^s\). 系数\(a_\mathscr {F}\)满足\[         0\leqslant  a_\mathscr {F} \leqslant  s,\] 故有界. 据此, 说\(P_\mathscr {F}(n)\)的\(n^{r-1}\)次项系数有下界 (分别为有界), 等价于对\(P_\mathscr {F}(n-k)=P_{\mathscr {F}(-k)}(n)\)作同样断言. 这样就归结到\[         \mathscr {L}         = \mathcal {O}_{\mathbb {P}_S^r}^s.       \]的情形.
\par{}考虑正合列\[         0 \to  N_r \to  \mathscr {F} \to  \mathscr {F}_{(r)} \to  0,\] 由此得\[         P_\mathscr {F}         = P_{\mathscr {F}_{(r)}} + P_{N_r}.\] 又因\(P_{N_r}\)的\(n^{r-1}\)次项系数非负 (因为\(\dim \operatorname {supp} N_r\leqslant  r-1\)), 所以\[         b_{\mathscr {F}_{(r)}} \leqslant  b_\mathscr {F}.\] 因此在证明引理时, 可以用\(\mathscr {F}_{(r)}\)替代\(\mathscr {F}\), 即假设所讨论的\(\mathscr {L}\)的商\(\mathscr {F}\)无挠.
\par{}由于\(\mathbb {P}_K^r\)正规, 所以\(\mathscr {F}\)在某个开集\(U=\mathbb {P}_K^r\setminus  Y\)上是秩为\(a=a_\mathscr {F}\)的局部自由层, 其中\(Y\)的余维\(\geqslant 2\). 因而\(\bigwedge ^a\mathscr {F}\)是\(\mathbb {P}_K^r\)上的一个层, 其在\(U\)上的限制可逆, 从而同构于\(\mathbb {P}_K^r\)上的一个在相差同构意义下确定的可逆层的限制 (因为\(\mathbb {P}_K^r\)正则, 且\(Y\)的余维\(\geqslant 2\)). 后一个层形如\(\mathcal {O}_{\mathbb {P}_K^r}(d)\), 其中\(d\)是完全确定的整数. 由于\(\bigwedge ^a\mathscr {F}\)是\(\bigwedge ^a\mathcal {O}_{\mathbb {P}_K^r}^n\simeq \mathcal {O}_{\mathbb {P}_K^r}^N\)的一个商, 其中\(N=\binom {n}{a}\), 所以它有\(N\)个典范截面. 这些截面给出\(\mathcal {O}_{\mathbb {P}_K^r}(d)\)在\(U\)上的截面, 后者是\(\mathcal {O}_{\mathbb {P}_K^r}(d)\)的截面\(s_i\) (\(1\leqslant  i\leqslant  N\)) 的限制 (因为\(\mathbb {P}_K^r\)正规, 且\(Y\)的余维\(\geqslant 2\)). 这些\(s_i\)在\(U\)的各点生成\(\mathcal {O}_{\mathbb {P}_K^r}(d)\), 因此不全为零, 从而\(d\geqslant 0\). 一个容易的计算还给出\[         b_\mathscr {F}         = a_\mathscr {F}(r+1)/2 + d.\] 特别有\(b_\mathscr {F}\geqslant 0\), 所以\(b_\mathscr {F}\)有下界. 它有界当且仅当\(d\)有界; 我们将证明此时\(\mathscr {F}\)仍属于一个有界族. 可以固定\(a_\mathscr {F}\)和\(b_\mathscr {F}\), 记为\(a\)和\(b\) (从而\(d\)也固定). 给定\(\mathcal {O}_{\mathbb {P}_K^r}(d)\)的\(N\)个截面\(s_i\), 即给定一个同态\(s\colon \bigwedge ^a\mathscr {L}_K\to \mathcal {O}_{\mathbb {P}_K^r}(d)\), 就能把\(\mathscr {F}\)恢复为相应复合同态\[         \mathscr {L}_K         \to  \mathscr {H}\kern  -2.5pt\mathit {om}\left ( \bigwedge ^{a-1}\mathscr {L}_{K'}, \bigwedge ^a\mathscr {L}_K \right )         \to  \mathscr {H}\kern  -2.5pt\mathit {om}\left ( \bigwedge ^{a-1}\mathscr {L}_{K'}, \mathcal {O}_{\mathbb {P}_K^r}(d) \right )       \]的余像; 其中第一个箭头是外积给出的典范同态, 第二个箭头来自\(s\). 再由\hyperref[fga3.iv-1-proposition-1.2]{命题~\ref*{fga3.iv-1-proposition-1.2}}的(i)即得结论.
\end{proof}
\end{lemma}
\par{}合用上述两个引理, 可以证明:
\begin{lemma}{2.6}\setnumber{2.6}\label{fga3.iv-2-lemma-2.6}\skipsinglepar
\par{}在\hyperref[fga3.iv-2-theorem-2.1]{定理~\ref*{fga3.iv-2-theorem-2.1}}的初始条件下, 假设对所有\(\mathscr {F}\)都有\[       P_\mathscr {F}(n)       = a_\mathscr {F} n^r/r! + b_\mathscr {F} n^{r-1}/(r-1)! + \text {次数}<r-1\text { 的项},\] 且系数\(a_\mathscr {F}\)有界. 则系数\(b_\mathscr {F}\)有下界. 此外, 若各\(b_\mathscr {F}\)有界, 则各\(\mathscr {F}_{(r)}\)有界.
\begin{proof}
\par{}可以假设各层\(\mathscr {F}\)的基域\(K\)代数闭. 给每个\(\operatorname {supp}\mathscr {F}_{(r)}\) (即\(r\)维分支的并) 赋以诱导的约化结构. 则各\(\operatorname {supp}\mathscr {F}_{(r)}\)的次数以\(a\)为上界, 所以由\hyperref[fga3.iv-2-lemma-2.4]{引理~\ref*{fga3.iv-2-lemma-2.4}}, 各\(\operatorname {supp}\mathscr {F}_{(r)}\)组成有界集合. 又对于\(\operatorname {supp}\mathscr {F}_{(r)}\)的每个分支, \(\mathscr {F}_{(r)}\)在该分支处的长度\(\leqslant  a\); 因而若\(\mathscr {I}_\mathscr {F}\)是定义\(\operatorname {supp}\mathscr {F}_{(r)}\)的理想, 则可将\(\mathscr {F}_{(r)}\)视为\(X\)中由\(\mathscr {I}_\mathscr {F}^a\)定义的子概形\(Y_\mathscr {F}\)上的模. 如上一引理, 还可归结到\(\mathscr {F}=\mathscr {F}_{(r)}\)的情形, 从而\(\mathscr {F}\)来自\(Y_\mathscr {F}\)上的一个模. 各\(Y_\mathscr {F}\)对应于\(\mathcal {O}_{X_K}\)的商模的一个有界族, 因而来自某个概形\(X\times _S T\)的闭子概形\(Y\). 于是对\(Y/T\)及\(\mathscr {L}\otimes _X Y\)应用\hyperref[fga3.iv-2-lemma-2.5]{引理~\ref*{fga3.iv-2-lemma-2.5}}, 即得结论.
\end{proof}
\end{lemma}
\par{}现在对\(\dim \operatorname {supp}\mathscr {F}\)的上界\(r\)作归纳, 证明\hyperref[fga3.iv-2-theorem-2.2]{定理~\ref*{fga3.iv-2-theorem-2.2}}. 当\(r<0\)时断言显然成立; 假设\(r\geqslant 0\), 且对\(r'<r\)已证得断言. 由\hyperref[fga3.iv-2-lemma-2.6]{引理~\ref*{fga3.iv-2-lemma-2.6}}, 各\(\mathscr {F}_{(r)}\)组成一个有界族; 再由\hyperref[fga3.iv-1-proposition-1.2]{命题~\ref*{fga3.iv-1-proposition-1.2}}的(i), 各同态\(\mathscr {L}_K\to \mathscr {F}_{(r)}\)的核也组成有界族. 因而\(X\)上存在一个凝聚模\(\mathscr {L}'\), 使这些核, 进而也使各\(N_r(\mathscr {F})=\operatorname {Ker}(\mathscr {F}\to \mathscr {F}_{(r)})\), 都是模\(\mathscr {L}'_K\)的商. 因为各\(\mathscr {F}_{(r)}\)有界, 所以各\(P_{\mathscr {F}_{(r)}}\)有界, 公式\[     P_\mathscr {F}     = P_{\mathscr {F}_{(r)}} + P_{N_r}   \]便说明各\(P_{N_r}\)满足与各\(P_\mathscr {F}\)相同的条件(bs). 因而各\(N_r\)满足条件(a)和(bs), 由归纳假设, 各\((N_r)_{(s)}\)有界. 而\(\mathscr {F}_{(s)}\)是\(\mathscr {F}_{(r)}\)由\((N_r)_{(s)}\)的扩张, 所以由\hyperref[fga3.iv-1-proposition-1.2]{命题~\ref*{fga3.iv-1-proposition-1.2}}的(ii), 各\(\mathscr {F}_{(s)}\)有界. 对于\hyperref[fga3.iv-2-theorem-2.2]{定理~\ref*{fga3.iv-2-theorem-2.2}}的最后一个断言, 注意由\hyperref[fga3.iv-1-proposition-1.2]{命题~\ref*{fga3.iv-1-proposition-1.2}}的(i), 各映射\(\mathscr {F}\to \mathscr {F}_{(s)}\)的核\(N_s\)有界, 且\(P_{N_s}\)的\(n^{s-1}\)次项系数有界; 此时\hyperref[fga3.iv-2-lemma-2.6]{引理~\ref*{fga3.iv-2-lemma-2.6}}说明下一项的系数有下界. \hyperref[fga3.iv-2-theorem-2.2]{定理~\ref*{fga3.iv-2-theorem-2.2}}的证明至此完成.
\chapter{Hilbert概形: 定义与存在性定理}\thispagestyle{fancy}\label{fga3.iv-3}
\par{}设\(X\)是预概形\(S\)之上的预概形, \(\mathscr {F}\)是\(X\)上的拟凝聚模. 用\[     \operatorname {Quot}({\mathscr {F}/X/S})   \]表示\(\mathscr {F}\)的各个在\(S\)上平坦的拟凝聚商模所组成的集合. 现设\(S'\to  S\)为基底变换态射, 置\(X'=X\times _S S'\)和\(\mathscr {F}'=\mathscr {F}\otimes _{\mathcal {O}_S}\mathcal {O}_{S'}\); 则\(X'\)是\(S'\)上的预概形, 带有拟凝聚模\(\mathscr {F}'\), 因而可以考虑\(\operatorname {Quot}({\mathscr {F}'/X'/S'})\). 置\[     \mathscr {Q}\kern  -.75pt\mathit {uot}_{{\mathscr {F}/X/S}} (S')     = \operatorname {Quot}({\mathscr {F}'/X'/S'})   \] (其中仍如上置\(X'=X\times _S S'\)).
\par{}若\(S''\to  S'\)是一个\(S\)-态射, 则\(X''=X\times _S S''\)同构于\(X'\times _{S'}S''\), \(\mathscr {F}''\)同构于\(\mathscr {F}'\otimes _{\mathcal {O}_{S'}}\mathcal {O}_{S''}\). 因为逆像函子\(\mathscr {G}'\mapsto \mathscr {G}'\otimes _{\mathcal {O}_{S'}}\mathcal {O}_{S''}\)右正合, 取值于\(X''\)上的拟凝聚模范畴, 并把\(S'\)-平坦模变成\(S''\)-平坦模, 所以得到自然映射\[     \operatorname {Quot}({\mathscr {F}'/X'/S'}) \to  \operatorname {Quot}(\mathscr {F}''/X''/S''),\] 从而\(\mathscr {Q}\kern  -.75pt\mathit {uot}_{{\mathscr {F}/X/S}} (S')\)是一个取值于集合范畴的\emph{关于\(S'\)的反变函子} (其中\(S'\)是\(S\)上的预概形). 以下假设\(X\)在\emph{Noether}的\(S\)上\emph{射影}, 且\(\mathscr {F}\)为\emph{凝聚}的; 为简单起见, 只考虑在\(S\)上的\emph{局部Noether}的\(S'\).
\begin{theorem}{3.1}\setnumber{3.1}\label{fga3.iv-3-theorem-3.1}\skipsinglepar
\par{}在这些条件下, 局部Noether的\(S\)-预概形范畴上的反变函子\(\mathscr {Q}\kern  -.75pt\mathit {uot}_{{\mathscr {F}/X/S}}\), 可由一个\(S\)-预概形\(\underline {\operatorname {Quot}}_{{\mathscr {F}/X/S}}\)\emph{表示}, 这个预概形是一个射影\(S\)-概形序列的和 (因而\(\underline {\operatorname {Quot}}_{{\mathscr {F}/X/S}}\)在\(S\)上局部有限型).
\end{theorem}
\par{}我们将按以下方式得到这样的分解. 设\(\mathcal {O}_X(1)\)是\(X\)上关于\(S\)非常丰沛的可逆层. 对每个有理系数多项式\(P(n)\), 以\(\operatorname {Quot}^P({\mathscr {F}/X/S})\)表示\(\operatorname {Quot}({\mathscr {F}/X/S})\)的一个子集, 其元素为\(\mathscr {F}\)的凝聚商\(\mathscr {G}\), 它在\(S\)上平坦, 且在每个\(s\in  S\)处的Hilbert多项式均为\(P\). 再置\[     \mathscr {Q}\kern  -.75pt\mathit {uot}_{\mathscr {F}/X/S}^P(S')     = \operatorname {Quot}^P({\mathscr {F}'/X'/S'}),\] 就得到\(\mathscr {Q}\kern  -.75pt\mathit {uot}_{\mathscr {F}/X/S}\)的一个子函子. §2所回顾的Hilbert多项式的不变性质给出: \emph{\(\mathscr {Q}\kern  -.75pt\mathit {uot}_{\mathscr {F}/X/S}\)可表示的充要条件是各\(\mathscr {Q}\kern  -.75pt\mathit {uot}_{\mathscr {F}/X/S}^P\)可表示; 此时表示它的\(S\)-预概形\(\underline {\operatorname {Quot}}_{{\mathscr {F}/X/S}}\)同构于表示各函子\(\mathscr {Q}\kern  -.75pt\mathit {uot}_{\mathscr {F}/X/S}^P\)的预概形\(\underline {\operatorname {Quot}}_{{\mathscr {F}/X/S}}^P\)之和.} 据此, \hyperref[fga3.iv-3-theorem-3.1]{定理~\ref*{fga3.iv-3-theorem-3.1}}将由下述定理推出:
\begin{theorem}{3.2}\setnumber{3.2}\label{fga3.iv-3-theorem-3.2}\skipsinglepar
\par{}在上述记号下, 函子\(\mathscr {Q}\kern  -.75pt\mathit {uot}_{\mathscr {F}/X/S}^P\)可由一个射影\(S\)-预概形\(\operatorname {Quot}_{\mathscr {F}/X/S}^P\)表示.
\end{theorem}
\par{}本节其余部分证明\hyperref[fga3.iv-3-theorem-3.2]{定理~\ref*{fga3.iv-3-theorem-3.2}}.
\par{}设\(\nu \)是整数. 对每个\(S\)上的\(S'\), 以\(A_\nu (S')\)表示\(\mathscr {F}'=\mathscr {F}\otimes _{\mathcal {O}_S}\mathcal {O}_{S'}\)的各个凝聚商\(\mathscr {G}=\mathscr {F}'/\mathscr {H}\)所组成的集合, 其中要求这些商在\(S'\)上平坦, 且满足下列条件: 
\begin{enumerate}
    
 \item[a.]当\(i>0\)且\(n\geqslant \nu \)时, \(\operatorname {R}^i f'_*(\mathscr {G} (n))=0\);

    
 \item[b.]当\(i>0\)且\(n\geqslant \nu \)时, \(\operatorname {R}^i f'_*(\mathscr {H} (n))=0\);

    
 \item[c.]当\(k\geqslant 0\)时, \(f'_*(\mathscr {H} (\nu +k))=\mathcal {S}'_k f'_*(\mathscr {H} (\nu ))\).

  
\end{enumerate} 对最后一个条件, 假设\(X\)写成\(S\)上一个由\(\mathcal {S}_1\)生成的拟凝聚正分次代数\(\mathcal {S}_\bullet \)的齐次素谱, 并置\(\mathcal {S}'=\mathcal {S}\otimes _{\mathcal {O}_S}\mathcal {O}_{S'}\), 于是\(X'\)是\(\mathcal {S}'\)的齐次素谱. 为证明\hyperref[fga3.iv-3-theorem-3.2]{定理~\ref*{fga3.iv-3-theorem-3.2}}, 很容易归结到\(X=\mathbb {P}_S^r\)的情形 (因为\(S\)是开子集\(U\)的并, 在这些开子集上\(X|U\)是\(\mathbb {P}_U^r\)的闭子概形, 且\(\mathcal {O}_X(1)\)由\(\mathcal {O}_{\mathbb {P}_U^r} (1)\)诱导), 并归结到\(\mathscr {F}\)形如\((\mathcal {O}_{\mathbb {P}_S^r})^N\), 因而在\(S\)上平坦的情形. 此时, 上述各层\(\mathscr {H}\)也在\(S'\)上平坦. 于是由Künneth关系 [\hyperref[GD1960]{5}, III, §7] 以及(b), 条件(a)和(b)在基底变换下稳定, 并且蕴涵: 当\(n\geqslant \nu \)时, \(f'_*(\mathscr {G} (n))\)和\(F'_*(\mathscr {H} (n))\)的构造与基底扩张交换 (同上), 因而(c)也在基底扩张下稳定. 换言之, \emph{\(A_\nu (S')\)是关于\(S'\)的反变函子, 并且在准确的意义下是\(A(S')=\mathscr {Q}\kern  -.75pt\mathit {uot}_{{\mathscr {F}/X/S}}^P(S')\)的子函子}. 令\(\nu \)变动, 就得到\(A(S')\)的一列递增子集\(A_\nu (S')\); 由Serre的一个熟知定理 [\hyperref[GD1960]{5}, III, §2], 其并为\(A(S')\). 现注意, 若\(\mathscr {G}\)是\(\mathscr {F}'\)的凝聚商, 在\(S\)上平坦, 而\(S\)的一个元素\(s\)使得基底变换\(\operatorname {Spec}(k(s))\to  S\)所得的\(\mathscr {F}_s\)的商\(\mathscr {G}_s\)满足(a)、(b)、(c), 即属于\(A_\nu \operatorname {Spec}(k(s))\), 则\(s\)有一个开邻域\(U\), 使\(\mathscr {G}|(f')^{-1} (U)\)也满足同样条件, 即这个商属于\(A_\nu (U)\); 对(a)和(b), 这事实上由"全纯函数定理" [\hyperref[GD1960]{5}, III, §7] 得到; 对(c), 则由Nakayama引理以及对充分大的\(n\)和\(k\geqslant 0\)总有\(f'_*(\mathscr {H} (n+k))=S'_kf'_*(\mathscr {H} (n))\)这一事实 [\hyperref[GD1960]{5}, III, §2] 得到. 由这些说明可得以下结论 (比较[\hyperref[Gro1960a]{4}, IV]): \emph{函子\(A\)可表示的充要条件是各函子\(A_\nu \)可表示; 此时表示\(A\)的\(S\)-预概形\(Q\), 是表示各\(A_\nu \)的开子集\(Q_\nu \)的递增并}
\par{}置\[     M_\bullet      = \sum _{n\geqslant 0} f_*(\mathscr {F} (n))     = \mathcal {S}_\bullet ^N,\] 因而有\[     M'_\bullet      = M_\bullet \otimes _{\mathcal {O}_S}\mathcal {O}_{S'}     = \sum _{n\geqslant 0} f'_*(\mathscr {F} (n))     = {\mathcal {S}'_\bullet }^N.\]
\par{}由(a)有
\begin{enumerate}
    
 \item[a'.]当\(n\geqslant \nu \)时, \(f'_*(\mathscr {G} (n))\)是秩为\(P(n)\)的局部自由层 [\hyperref[GD1960]{5}, III, §7]

  
\end{enumerate}, 而由(b)在\(i=1\)时有
\begin{enumerate}
    
 \item[a''.]\(f'_*(\mathscr {G} (n))\)是\(M'_n\)的商模.

  
\end{enumerate}. 又对\(n=\nu \)知道这个商模, 就由(c)知道\(M_n\)在\(n\geqslant \nu \)时的各个子模\(f'_*(\mathscr {H} (n))\), 从而知道\(\mathscr {H}\), 进而知道\(\mathscr {G}\). 这样得到一个\emph{单射}\[     A_\nu (S') \to  \mathscr {G}\kern  -1.25pt\mathit {rass}_{P(\nu )} (M'_\nu ),\] 从\(A_\nu (S')\)到\(M'\)的秩为\(P(\nu )\)的局部自由商模的集合; 因而得到一个\emph{函子性}的同态\[     i_\nu \colon  A_\nu (S') \to  \mathscr {G}\kern  -1.25pt\mathit {rass}_{P(\nu )} (M_\nu )(S').\] 右侧函子可由Grassmann概形\(\operatorname {Grass}_{P(\nu )} (M_\nu )\)表示 (比较[\hyperref[Gro1960a]{4}, V]), 而此概形在\(S\)上射影. 于是有:
\begin{lemma}{3.3}\setnumber{3.3}\label{fga3.iv-3-lemma-3.3}\skipsinglepar
\par{}\(A_\nu (S')\)是可表示函子, 且表示同态\(i_\nu \)的态射\(Q_\nu \to \operatorname {Grass}_{P(\nu )} (M_\nu )\)是浸入 (这蕴涵\(Q_\nu \)在\(S\)上拟射影).
\end{lemma}
\par{}这个断言等价于以下断言 (比较[\hyperref[Gro1960a]{4}, IV]): \emph{若给定\(M'_\nu \)的一个局部自由秩为\(P(\nu )\)的商模\(N\), 则\(S'\)有一个子预概形\(Z\), 使得对于每个\(S'\)上的局部Noether预概形\(T'\), \(N\)在\(T'\)上的逆像属于\(\Im  A_\nu (T')\)的充要条件是\(T'\to  S'\)经由子预概形\(Z\)分解.} 改换记号, 可以假设\(S'=S\), 即有\(M_\nu \)模掉子模\(R_\nu \)得到的商\(N_\nu \). 它来自\(A(s)\)的一个元素的充要条件是满足以下两个条件: 
\begin{enumerate}
    
 \item[i.]对\(k\geqslant 0\), \(M_{\nu +k}/\mathcal {S}_k R_\nu \)局部自由且秩为\(P(\nu +k)\).

    
 \item[ii.]由\(M_\bullet \)的分次子模\(R_\bullet =\sum _{k\geqslant 0}\mathscr {S}_kR_\nu \)定义的\(\mathscr {F}\)的子层\(\mathscr {H}\) (参见[\hyperref[GD1960]{5}, II, §3]) 及商层\(\mathscr {G}=\mathscr {F}/\mathscr {H}\)均满足上述条件(a)和(b).

  
\end{enumerate} 这些条件显然必要; 若它们满足, 则上述(ii)定义的层\(\mathscr {G}\)同构于分次\(\mathscr {S}_\bullet \)-模\(N_\bullet \) (即各\(M_{n+k}/\mathscr {S}_kR\)之和) 所对应的层, 所以在\(S\)上\emph{平坦} (因为它的各茎是\(N\)关于\(\mathcal {S}_\bullet \)的齐次素理想的局部化的直和因子), 而且由(i), 对应于Hilbert多项式\(P\). 再考虑(ii), 就看到(a')和(a'')满足; 因而当\(n\geqslant \nu \)时, 自然同态\(N_n\to  f_*(\mathscr {G} (n))\)是等秩局部自由模间的\emph{满}同态, 故为同构, 从而对所有\(k\geqslant 0\)有\(f_*(\mathscr {H} (\nu +k))=\mathcal {S}_kR_\nu \). 这就证明\(\mathscr {G}\in  A_\nu (S)\), 且\(R_\nu \)是由\(\mathscr {G}\)所定义的\(\operatorname {Grass}_{P(\nu )} (M_\nu )\)的元素.
\par{}上述判据(i)和(ii)同样适用于基底变换\(S'\to  S\)后的情形. 首先证明, 基底变换\(S'\to  S\)后条件(i)满足, 可以表示为\(S'\to  S\)经由\(S\)的某个子预概形\(Z\)分解. 一旦证明这一点, 用\(Z\)替代\(S\), 就归结到(i)已在\(S\)上满足的情形; 因为它在基底变换下稳定, 所以只须表述条件(ii). 但此时以\(U\)表示这样的\(s\in  S\)的集合: \(\mathscr {G} (n)\)与\(\mathscr {H} (n)\)在纤维\(X_s\)上诱导的层, 对\(n\geqslant \nu \)在所有\(>0\)维的上同调均为零. 我们已经证明\(U\)是开集, 而条件(ii)在基底变换\(S'\to  S\)后满足, 当且仅当\(S'\to  S\)经由\(U\)分解; 这就证明\hyperref[fga3.iv-3-lemma-3.3]{引理~\ref*{fga3.iv-3-lemma-3.3}}. 因此尚须证明以下引理:
\begin{lemma}{3.4}\setnumber{3.4}\label{fga3.iv-3-lemma-3.4}\skipsinglepar
\par{}设\(S\)是一个局部Noether预概形, 其上有一个由\(\mathcal {S}_1\)生成的拟凝聚正分次代数\(\mathcal {S}_\bullet \); 设\(M_\bullet \)是有限型拟凝聚分次\(\mathcal {S}\)-模, \(P\)是有理系数多项式, \(\nu \)是整数. 则\(S\)存在一个显然唯一的子预概形\(Z\), 具有如下性质: 对\(S\)上的每个预概形\(S'\), \(M_n\otimes _{\mathcal {O}_S}\mathcal {O}_{S'}\)在所有\(n\geqslant \nu \)时均局部自由且秩为\(P(n)\)的充要条件是\(S'\to  S\)经由\(Z\)分解.
\end{lemma}
\par{}显然可以假设\(S\)仿射, 因而为Noether. 于是:
\begin{lemma}{3.5}\setnumber{3.5}\label{fga3.iv-3-lemma-3.5}\skipsinglepar
\par{}对每个整数\(N\geqslant \nu \), 设\(U_N\)是\(S\)中由满足下列条件的\(s\in  S\)组成的开子集: 对所有\(\nu \leqslant  n\leqslant  N\), 均有\(\operatorname {rank}_{k(s)}M_{ns}\otimes _{\mathcal {O}_{S,s}} k(s)\leqslant  P(n)\). 则这列递减开子集\(U_N\)最终稳定.
\begin{proof}
\par{}已知 [\hyperref[Gro1960b]{3}, IV], \(S\)可以有限地分割成约化子概形\(S_i\), 使每个\(M\otimes _{\mathcal {O}_S}\mathcal {O}_{S_i}\)在\(S\)上\emph{平坦}. 因而可假设\(M\)平坦, 从而各\(M_n\)平坦. 最后, 显然可假设\(S\)连通. 但此时 [\hyperref[GD1960]{5}, III, §7], 存在一个整数\(n_0\)和一个多项式\(Q\), 使\[         \operatorname {rank}_{k(s)}M_{ns}\otimes _{\mathcal {O}_{S,s}} k(s)         = Q(n)         \qquad \text {当 } n\geqslant  n_0.\] 首先假设\(P<Q\), 从而对充分大的\(n\)有\(P(n)\neq  Q(n)\). 则显然对充分大的\(N\)有\(U_N=\varnothing \), 因而各\(U_N\)组成的序列当然稳定. 在相反情形, 对充分大的\(n\)有\(P(n)\geqslant  Q(n)\), 所以当\(N\geqslant  n_0\)时\(U_N=U_{n_0}\), 各\(U_N\)组成的序列同样稳定.
\end{proof}
\end{lemma}
\par{}特别地, 满足
\begin{eqn}{*}\setnumber{*}\label{fga3.iv-3-equation-star}\skipsinglepar\[       \operatorname {rank}_{k(s)}M_{ns}\otimes _{\mathcal {O}_{S,s}} k(s) \leqslant  P(n)       \qquad \text {对所有 } n\geqslant \nu      \tag{*}     \]
\end{eqn}的\(s\in  S\)所组成的集合\(U_\infty \)是开集, 因为它是各\(U_N\)之交. 因此在证明\hyperref[fga3.iv-3-lemma-3.4]{引理~\ref*{fga3.iv-3-lemma-3.4}}时, 可以用开子集\(U\)替代\(S\), 这就归结到\hyperref[fga3.iv-3-equation-star]{公式~\ref*{fga3.iv-3-equation-star}}中的不等式在所有\(s\in  S\)处成立的情形.
\begin{lemma}{3.6}\setnumber{3.6}\label{fga3.iv-3-lemma-3.6}\skipsinglepar
\par{}设\(M\)是局部Noether预概形\(S\)上的模, \(r\)是整数. 则\(S\)有一个显然唯一的子预概形\(Z\), 具有如下性质: 对每个\(S\)上的\(S'\), \(M\otimes _{\mathcal {O}_S}\mathcal {O}_{S'}\)局部自由且秩为\(r\)的充要条件是\(S'\to  S\)经由\(Z\)分解. 若对所有\(s\)都有\(\operatorname {rank}_{k(s)}M_s\otimes _{\mathcal {O}_{S,s}} k(s)\leqslant  r\), 则\(Z\)是\(S\)的闭子预概形 (假设\(M\)凝聚).
\begin{proof}
\par{}事实上, 上述推理把我们归结到对所有\(s\in  S\)都有不等式\hyperref[fga3.iv-3-equation-star]{公式~\ref*{fga3.iv-3-equation-star}}的情形 (必要时, 用使此不等式成立的\(s\)组成的开子集替代\(S\)). 此时可假设\(M\)纳入一个正合列\[         \mathcal {O}_S^q \to  \mathcal {O}_S^r \to  M \to  0;\] 而对于\(S\)上的\(S'\), 所要求的条件就是说在相应正合列\(\mathcal {O}_{S'}^q\to \mathcal {O}_{S'}^r\to  M'\to  0\)中, 第二个箭头为同构, 即第一个箭头为零. 于是可见, 由定义同态\(\mathcal {O}_S^q\to \mathcal {O}_S^r\)的矩阵的系数生成的理想所定义的\(S\)的闭子预概形\(Z\), 满足所求条件.
\end{proof}
\end{lemma}
\par{}回到\hyperref[fga3.iv-3-lemma-3.4]{引理~\ref*{fga3.iv-3-lemma-3.4}}的证明中刚才停下之处, 以\(Z_n\)表示由\hyperref[fga3.iv-3-lemma-3.6]{引理~\ref*{fga3.iv-3-lemma-3.6}}给模\(M_n\)和整数\(r=P(n)\)配以的\(S\)的\emph{闭}子预概形; 以\(Z'_N\)表示\(\nu \leqslant  n\leqslant  N\)时各\(Z_n\)的下确界. 则各\(Z_N\)组成\(Z\)的一列递减闭子预概形, 因而必然最终稳定. 设\(Z\)为充分大的\(N\)时各\(Z_W\)的恒定值. 这就是\hyperref[fga3.iv-3-lemma-3.4]{引理~\ref*{fga3.iv-3-lemma-3.4}}所要求的\(Z\). 至此证明了\hyperref[fga3.iv-3-lemma-3.4]{引理~\ref*{fga3.iv-3-lemma-3.4}}, 因而也证明了\hyperref[fga3.iv-3-lemma-3.3]{引理~\ref*{fga3.iv-3-lemma-3.3}}.
\par{}这样, 我们已证明\emph{\(\mathscr {Q}\kern  -.75pt\mathit {uot}_{{\mathscr {F}/X/S}}^P\)可由一个\(S\)-预概形\(Q\)表示, 后者是\(S\)上的拟射影开子预概形\(Q_\nu \)的递增并}. 为继续推进, 需要引用\hyperref[fga3.iv-2-theorem-2.1]{定理~\ref*{fga3.iv-2-theorem-2.1}}, 很容易从中推出\(Q\)是\emph{拟紧}的 (因为它是一个在\(S\)上有限型的预概形\(S'\)的像, 后者参数化各\(\mathscr {F}_K\)的Hilbert多项式为\(P\)的商层族). 因此\(Q\)等于某个\(Q_\nu \), 从而在\(S\)上拟射影. 为证明它在\(S\)上射影, 只须证明它在\(S\)上\emph{固有}; 为此只需应用[\hyperref[GD1960]{5}, II, 7.3.8]所给形式的固有性赋值判据. 只须验证:
\begin{lemma}{3.7}\setnumber{3.7}\label{fga3.iv-3-lemma-3.7}\skipsinglepar
\par{}设\(S\)是离散赋值环的谱, \(s\)是其一般点, \(X\)是\(S\)上的预概形, \(\mathscr {F}\)是\(X\)上的拟凝聚模, \(\mathscr {G}_s\)是\(X_s\)上\(\mathscr {F}_s=\mathscr {F}\otimes _{\mathcal {O}_S} k(s)\)的拟凝聚商模. 则存在唯一的\(\mathscr {F}\)的拟凝聚商模\(\mathscr {G}\), 它在\(S\)上平坦, 且其在\(X_s\)上的限制为\(\mathscr {G}_s\).
\begin{proof}
\par{}事实上, 若\(\mathscr {G}_s=\mathscr {F}_s/\mathscr {H}_s\), 只需考虑在一般纤维诱导\(\mathscr {H}_s\)的\(\mathscr {F}\)的最大子层\(\mathscr {H}\) ([\hyperref[GD1960]{5}, I, 9.4.2]), 并取\(\mathscr {G}=\mathscr {F}/\mathscr {H}\). 容易验证这个层满足要求.
\end{proof}
\end{lemma}
\par{}至此, \hyperref[fga3.iv-3-theorem-3.2]{定理~\ref*{fga3.iv-3-theorem-3.2}}以及\hyperref[fga3.iv-3-theorem-3.1]{定理~\ref*{fga3.iv-3-theorem-3.1}}均已完全证明.
\par{}这个证明同时还表明:
\begin{proposition}{3.8}\setnumber{3.8}\label{fga3.iv-3-proposition-3.8}\skipsinglepar
\par{}在\hyperref[fga3.iv-3-theorem-3.2]{定理~\ref*{fga3.iv-3-theorem-3.2}}的条件下, 置\(Q=\underline {\operatorname {Quot}}_{{\mathscr {F}/X/S}}^P\)、\(X_Q=X\times _S Q\)、\(\mathscr {F}_Q=\mathscr {F}\otimes _{\mathcal {O}_S}\mathcal {O}_Q\), 并设\(\mathscr {G}\)是\(\mathscr {F}_Q\)的一个在\(Q\)上平坦、相对Hilbert多项式为\(P\)的凝聚商, 使得\((Q,\mathscr {G})\)表示函子\(\mathscr {Q}\kern  -.75pt\mathit {uot}_{{\mathscr {F}/X/S}}^P\). 则存在一个整数\(\nu \), 使得当\(n\geqslant \nu \)时, \((f_Q)_*(\mathscr {G} (n))\)是\(Q\)上秩为\(P(n)\)的局部自由模, 并关于\(S\)非常丰沛, 即定义\(Q\)到\(S\)上某个Grassmann概形\(\underline {\operatorname {Grass}}_{P(n)} (M)\)的一个浸入. 因而当\(n\geqslant \nu \)时, \(Q\)上的层\(\bigwedge ^{P(n)} (f_Q)_*(\mathscr {G} (n))\)可逆, 并关于\(S\)非常丰沛.
\begin{proof}
\par{}事实上, 如\hyperref[fga3.iv-3-theorem-3.2]{定理~\ref*{fga3.iv-3-theorem-3.2}}中一样可归结到\(\mathscr {F}\)在\(S\)上平坦的情形, 此时只需取整数\(\nu \), 使\(A_\nu =A\) (记号同上).
\end{proof}
\end{proposition}
\par{}\hyperref[fga3.iv-3-theorem-3.2]{定理~\ref*{fga3.iv-3-theorem-3.2}}最重要的应用是\(\mathscr {F}=\mathcal {O}_X\)的情形. 此时写作\[     
\begin {aligned}       \underline {\operatorname {Quot}}_{\mathcal {O}_X/X/S}       &= \underline {\operatorname {Hilb}}_{X/S}     \\\underline {\operatorname {Quot}}_{\mathcal {O}_X/X/S}^P       &= \underline {\operatorname {Hilb}}_{X/S}^P     
\end {aligned},\] 因而有分解\[     \underline {\operatorname {Hilb}}_{X/S}     = \coprod _P \underline {\operatorname {Hilb}}_{X/S}^P.\] 按定义, \(\underline {\operatorname {Hilb}}_{X/S}\)表示函子\(\mathscr {H}\kern  -1.75pt\mathit {ilb}_{X/S} (S')\), 其值为\(X'=X\times _S S'\)的各个在\(S\)上平坦的闭子预概形的集合; \(\underline {\operatorname {Hilb}}_{X/S}^P\)则表示与具有给定Hilbert多项式\(P\)的闭子预概形相应的子函子. 这两个预概形也分别称为\(X\)在\(S\)上的\emph{Hilbert预概形}和\emph{指标为\(P\)的Hilbert预概形}. 这一术语来自Hilbert多项式在理论中所起的作用. 它们与经典Chow簇 (其目的是参数化闭链, 而非簇) 的性质差别, 类似于一个簇的闭链类的\emph{Chow环}与该簇的\emph{层的类环} (如Riemann–Roch定理\hyperref[BS1958]{[1]}中所引入的) 的差别. 注意当\(X=\mathbb {P}_S^r\)且\(S\)是一个域的谱时, 知道\(X\)上一个凝聚模\(\mathscr {F}\)的Hilbert多项式, 等价于知道\(\mathscr {F}\)的Chern类, 甚至等价于知道\(\mathscr {F}\)在\(X\)上凝聚层的类环中的类.
\begin{remarks}{3.9}\setnumber{3.9}\label{fga3.iv-3-remarks-3.9}\skipsinglepar
\par{}还应指出, \(\underline {\operatorname {Quot}}_{{\mathscr {F}/X/S}}\)和\(\underline {\operatorname {Quot}}_{{\mathscr {F}/X/S}}^P\)的构造已归结到\(X=\mathbb {P}_S^r\)且\(\mathscr {F}=\mathcal {O}_X^N\)、\(\mathcal {O}_X(1)\)为通常的非常丰沛层的情形; 更准确地说, 一般的\(\underline {\operatorname {Quot}}_{{\mathscr {F}/X/S}}\)都是上述预概形的闭子预概形. 由于各\(\underline {\operatorname {Quot}}_{{\mathscr {F}/X/S}}^P\)的构造显然与基底变换\(S'\to  S\)相容, 所以还可进一步归结到\(S=\operatorname {Spec}(\mathbb {Z})\)的情形: \[       \mathscr {Q}_{r,N}^P       = \underline {\operatorname {Quot}}_{(\mathcal {O}_{\mathbb {P}_{\mathbb {Z}}^r})^N/\mathbb {P}_{\mathbb {Z}}^r/\operatorname {Spec}(\mathbb {Z})},\] 特别地, 得到绝对Hilbert概形\[       \underline {\operatorname {Hilb}}_r^P       = \mathscr {Q}_{r,1}^P.\] 更详细地研究这些概形将很有意思, 首先是确定它们的连通分支 (它们是否连通?) 和不可约分支 (由Serre \hyperref[Ser1961]{[7]}, 可能存在完全位于某个素数\(p\neq 0\)之上的不可约分支). 回顾Weil提出的问题: \(\underline {\operatorname {Hilb}}_r^P\)在各\(s\in \operatorname {Spec}(\mathbb {Z})\)之上的纤维的不可约分支, 是否对应于素域的"正则"扩张, 即是否"相对连通". 对Hilbert概形, 这些问题或许比对"Chow簇"更容易处理. \emph{[Comp.]} Hartshorne研究了代数闭域上Hilbert概形的连通分支, 证明各\(\underline {\operatorname {Hilb}}_r^P\)连通, 并确定了使\(\underline {\operatorname {Hilb}}_r^P\neq \varnothing \)的数对\((r,P)\) \hyperref[Har1966]{[Har1966]}.
\end{remarks}
\chapter{变体}\thispagestyle{fancy}\label{fga3.iv-4}
\begin{enumerate}

  
 \item[a.]
    
\par{}\label{fga3.iv-4.a} 在\hyperref[fga3.iv-3-theorem-3.1]{定理~\ref*{fga3.iv-3-theorem-3.1}}的条件下, 设\(U\)是\(X\)的开子集; 以\(A'\)表示\(A=\mathscr {Q}\kern  -.75pt\mathit {uot}_{{\mathscr {F}/X/S}}\)的一个子函子, 使\(A'(S')\)为\(\mathscr {F}'\)的所有在\(S'\)上平坦且支集包含于\(U'\)的商模\(\mathscr {G}\)所组成的集合. 立即可见\(A'\)可由表示\(A\)的预概形\(\underline {\operatorname {Quot}}_{{\mathscr {F}/X/S}}\)的一个开子集表示. 因此, 若把\(X\)在\(S\)上射影改为在\(S\)上\emph{拟射影}, 并把结论中的"射影"一词改为"拟射影", 再以\(\mathscr {Q}\kern  -.75pt\mathit {uot}_{{\mathscr {F}/X/S}}(S')\)表示\(\mathscr {F}'\)的所有在\(S'\)上平坦且\emph{其支集在\(S'\)上固有}的凝聚商\(\mathscr {G}\)的集合, 则\hyperref[fga3.iv-3-theorem-3.1]{定理~\ref*{fga3.iv-3-theorem-3.1}}和\hyperref[fga3.iv-3-theorem-3.2]{定理~\ref*{fga3.iv-3-theorem-3.2}}仍然成立.
  

  
 \item[b.]
    
\par{}\label{fga3.iv-4.b} 一般地, 对\(\mathscr {F}'=\mathscr {F}\otimes _{\mathcal {O}_S}\mathcal {O}_{S'}\)的在\(S'\)上平坦的商\(\mathscr {G}\), 可以再加上各种在基底变换下稳定的自然条件, 从而得到我们想要表示的\(\mathscr {Q}\kern  -.75pt\mathit {uot}_{{\mathscr {F}/X/S}}\)的各种子函子. 在许多情形中, 通常的判据可以证明这些函子仍然由\(\underline {\operatorname {Quot}}_{{\mathscr {F}/X/S}}\)的开子集表示. 特别地, 加上下列某一性质时就是如此: 
\begin{enumerate}\item{}在各纤维\(X'_{s'}\)上诱导的模\(\mathscr {G}_{s'}\) (\(s'\in  S'\)) 的相伴素闭链的维数属于给定的整数集合.
        \item{}(在\(\mathscr {F}=\mathcal {O}_X\)的情形, 此时\(\mathscr {G}\)对应于\(X'\)的一个闭子预概形\(Y\):) \(Y\)是在\(S\)上\emph{光滑}的预概形 [\hyperref[Gro1960b]{3}, IV], 或者在\(S\)上\emph{正规} (即各纤维\(Y_{s'}\)在"\(k(s)\)上"正规, 也就是任意基域扩张后均正规), 或者 (若\(X\)在\(S\)上平坦) 是\(X\)中关于\(S\)的局部完全\(k\)-交 (即各纤维\(Y_{s'}\)是各\(X_{s'}\)中的局部完全交).
\end{enumerate} 其他条件可涉及在各\(X'_{s'}\)上诱导的模\(\mathscr {G}_{s'}\)的上同调性质, 等等. 当然, 如果每个条件都由\(\underline {\operatorname {Quot}}_{{\mathscr {F}/X/S}}\)的一个开集\(U_i\)表示, 那么同时满足这些条件就由这些开集之交表示. 例如, 对\(S\)上的每个\(S'\), 考虑\(X'=X\times _S S'\)的闭子预概形\(Y\)所组成的集合, 其中要求它们是在\(S'\)上具有给定秩\(r\)的平展覆盖 [\hyperref[Gro1960b]{3}, I]; 这样就得到一个关于\(S'\)的可表示反变函子.
  

  
 \item[c.]
    
\par{}\label{fga3.iv-4.c} \hyperref[fga3.ii-c.2]{[2], II, C, §2}所定义的预概形\(\underline {\operatorname {Hom}}_S(X,Y)\)、\(\prod _{X/S}Z/S\)和\(\underline {\operatorname {Isom}}_S(X,Y)\), 在适当射影性假设下存在, 并可实现为适当的Hilbert预概形的开子预概形. 因为\(\underline {\operatorname {Hom}}_S(X,Y)=\prod _{X/S}((X\times  Y)/X)\), 所以\(\operatorname {Hom}_S(X,Y)\)的情形归结到\(\prod _{X/S}(Z/X)\)的情形. 此时注意, 对\(S\)上的每个\(S'\), \(Z'=Z\times _S S'\)在\(X'=X\times _S S'\)上的截面集合, 与\(Z\)的满足下述条件的子预概形\(\Gamma \)的集合双射对应 (若\(Z\)在\(X\)上分离, 则这些子预概形必然闭): 由\(Z'\to  X'\)诱导的态射\(\Gamma \to  X'\)是同构. 因而, \emph{若\(X\)在\(S\)上平坦且固有, \(Z\)在\(S\)上拟射影, 则\(\prod _{X/S}(Z/X)\)存在, 并实现为\(\underline {\operatorname {Hilb}}_{Z/S}\)的开子预概形}. 于是, \emph{若\(X\)在\(S\)上射影且平坦, \(Y\)在\(S\)上拟射影, 则\(\underline {\operatorname {Hom}}_S(X,Y)\)存在, 并实现为\(\underline {\operatorname {Hilb}}_{(X\times _S Y)/S}\)的开子预概形}. 若\(X\)和\(Y\)都在\(S\)上射影, 则立即得到\(\underline {\operatorname {Isom}}_S(X,Y)\)也存在, 并由\(\underline {\operatorname {Hom}}_S(X,Y)\)的一个开子集表示. 同样, 若\(X\)在\(S\)上平坦且射影, \(Y\)在\(S\)上拟射影, 则\(S\)-预概形\(\underline {\operatorname {Imm}}_S(X/Y)\)也可由\(\underline {\operatorname {Hom}}_S(X,Y)\)的开子集表示; 这个预概形对应于\(\underline {\operatorname {Hom}}_S(X,Y)\)所表示的函子的一个子函子, 后者对应于那些为浸入的\(S'\)-态射\(X'\to  Y'\).

    
\par{}设\(\mathscr {L}\) (分别为\(\mathscr {M}\)) 是\(X\) (分别为\(Y\)) 上关于\(S\)非常丰沛的可逆层, 从而在\(X\times _S Y\)上得到关于\(S\)非常丰沛的层\(\mathscr {L}\otimes _{\mathcal {O}_S}\mathscr {M}\). 对任意有理系数多项式\(P\), \(\underline {\operatorname {Hilb}}_{(X\times _S Y)/S}^P\)有定义, 且为\(S\)上的拟射影预概形. 它因而在\(\underline {\operatorname {Hom}}_S(X,Y)\)上诱导一个既开又闭、并在\(S\)上拟射影的子集, 记作\(\operatorname {Hom}_S(X,Y)^P\). 于是\(\underline {\operatorname {Hom}}_S(X,Y)^P\)在\(S\)上的截面, 就是那些使任意整数\(n\)均满足\[         \chi \Big ((\mathscr {L}\otimes _{\mathcal {O}_X}g^*(\mathscr {M}))^{\otimes  n}\Big )         = P(n).       \]的\(S\)-态射\(g\colon  X\to  Y\). 这样得到Matsusaka定理的推广: 一个"极化"射影簇的自同构组成一个代数群. 这个断言在这里显然具有更精确的意义, 因为该群被定义为一个泛问题的解. 还应注意, 在代数闭域上, 过去所考虑的自同构群, 是将这里定义的"真正"自同构群模掉幂零元所得的群; 这说明为何那些历史上的构造难以在非完美基域上进行, 因为基域扩张后出现的幂零元理想未必"在\(k\)上有定义". 同样的说明也适用于过去的大多数构造.
  

\end{enumerate}
\chapter{Hilbert概形的微分研究}\thispagestyle{fancy}\label{fga3.iv-5}
\par{}有以下结果:
\begin{proposition}{5.1}\setnumber{5.1}\label{fga3.iv-5-proposition-5.1}\skipsinglepar
\par{}设\(S\)是预概形, \(S_0\)是由平方为零的拟凝聚理想\(\mathscr {I}\)定义的子预概形, \(X\)是\(S\)-预概形, \(\mathscr {F}\)是\(X\)上的拟凝聚模. 置\(X_0=\mathscr {F}\times _S S_0\)及\(\mathscr {F}_0=\mathscr {F}\otimes _{\mathcal {O}_S}\mathcal {O}_{S_0}\). 最后, 设\(\mathscr {G}_0=\mathscr {F}_0/\mathscr {H}_0\)是\(\mathscr {F}_0\)的一个在\(S_0\)上平坦的拟凝聚商模. 对\(X\)的每个开子集\(U\), 以\(\mathscr {E}(U)\)表示\(\mathscr {F}|U\)的所有在\(S\)上平坦并满足\(\mathscr {G}\otimes _{\mathcal {O}_S}\mathcal {O}_{S_0}=\mathscr {G}_0\)的拟凝聚商模\(\mathscr {G}\)所组成的集合; 当\(U\)变动时, 各\(\mathscr {E}(U)\)是\(U\)上一个层\(\mathscr {E}\)的截面. 此时群层\[       \mathscr {A}       = \underline {\operatorname {Hom}}_{\mathcal {O}_{X_0}}(\mathscr {H}_0,\mathscr {G}_0\otimes _{\mathcal {O}_{S_0}}\mathscr {I})     \]自然作用于\(\mathscr {E}\), 使它成为一个"形式地为\(\mathscr {A}\)-主齐性"的层 (即对于\(X\)的每个开子集\(U\), \(\mathscr {E}(U)\)或者为空, 或者是一个\(\mathscr {A}(U)\)-主齐性集).
\end{proposition}
\par{}因此得到:
\begin{corollary}{5.2}\setnumber{5.2}\label{fga3.iv-5-corollary-5.2}\skipsinglepar
\par{}假设在\(X\)上局部存在将\(\mathscr {G}_0\)延拓为\(\mathscr {F}\)的在\(S\)上平坦的商的延拓\(\mathscr {G}\) (即层\(\mathscr {E}\)的各茎非空). 则存在一个典范的障碍类\[       c(\mathscr {G}_0) \in  \operatorname {H}^1(X,\mathscr {A}),\] 它为零的充要条件是存在把\(\mathscr {G}_0\)延拓为\(\mathscr {F}\)的在\(S\)上平坦的商的整体延拓\(\mathscr {G}\). 若这个类为零, 则所有可能延拓组成的集合\(\mathscr {E}(X)\)是\(\mathscr {A}(X)=\operatorname {Hom}_{\mathcal {O}_X}(\mathscr {H}_0,\mathscr {G}_0\otimes _{\mathcal {O}_{S_0}}\mathscr {I})\)下的一个主齐性集.
\end{corollary}
\par{}特别地, 若\(\operatorname {H}^1(X,\mathscr {A})=0\), 就可保证整体延拓存在.
\begin{corollary}{5.3}\setnumber{5.3}\label{fga3.iv-5-corollary-5.3}\skipsinglepar
\par{}假设\(Q=\underline {\operatorname {Quot}}_{\mathscr {F}/X/S}\)存在 (参见\hyperref[fga3.iv-4.a]{§4.a}); 例如可假设\(X\)在某个局部Noether的\(S\)上拟射影, 且\(\mathscr {F}\)凝聚. 设\(x\in  Q\), 它对应于某个\(k(s)\)的剩余域扩张\(K=k(x)\) (其中\(s\in  S\)). 则\(x\)由\(K\)-预概形\(X_K\)上的模\(\mathscr {F}_0=F_K\)的一个凝聚商模\(\mathscr {G}_0=\mathscr {F}_0/\mathscr {H}_0\)定义. 设\(\mathscr {A}\)是\(X_K\)上由\[       \mathscr {A} = \underline {\operatorname {Hom}}_{\mathcal {O}_{X_0}}(\mathscr {H}_0,\mathscr {G}_0).     \]定义的凝聚层. 则纤维\(Q_s\)在点\(x\)处的Zariski切空间 (即\(\mathfrak {m}/\mathfrak {m}^2\)的\(K\)-对偶, 其中\(\mathfrak {m}\)是\(\mathcal {O}_{Q_K,x}\)的极大理想) 典范同构于\(\operatorname {H}^0(X_K,\mathscr {A})\).
\end{corollary}
\par{}这个关于Zariski切空间的结果可以推广: 对给定的\(S\)-态射\(g\colon  S'\to  Q\), 即\(Q'=Q\times _S S'\)在\(S'\)上的一个截面\(g'\), 可用以下公式刻画模\[     \Omega      = g^*(\Omega _{Q/S}^1)     = {g'}^*(\mathscr {J}/\mathscr {J}^2)   \] (其中\(\mathscr {J}\)是由\(Q'\)在\(S'\)上的截面\(g'\)所定义的\(Q'\)上的理想): \[     \operatorname {Hom}_{\mathcal {O}_{S'}}(\Omega ,\mathscr {M})     \simeq  \operatorname {H}^0(X',\mathscr {A});\] 这个公式关于\(S'\)上的凝聚模\(\mathscr {M}\)函子化, 其中\(\mathscr {A}\)仍为\(X'=X\times _S S'\)上由\[     \mathscr {A} = \underline {\operatorname {Hom}}_{\mathcal {O}_{X'}}(\mathscr {H},\mathscr {G}\otimes _{\mathcal {O}_S}\mathscr {M})   \]定义的模 (\(\mathscr {G}=\mathscr {F}'/\mathscr {H}\)是\(\mathscr {F}'=\mathscr {F}\otimes _{\mathcal {O}_S}\mathcal {O}_{S'}\)中与\(g\)相应的商模). 只需在\hyperref[fga3.iv-5-proposition-5.1]{命题~\ref*{fga3.iv-5-proposition-5.1}}中以\(S'\)替代\(S_0\), 并以预概形\(D(\mathscr {M})=(S',\mathcal {O}_{S'}+\mathscr {M})\)替代\(S\), 其中把\(\mathscr {M}\)视为平方为零的理想.
\par{}若在\hyperref[fga3.iv-5-proposition-5.1]{命题~\ref*{fga3.iv-5-proposition-5.1}}中\(\mathscr {F}=\mathcal {O}_X\), 则给定\(\mathscr {G}_0\)等价于给定\(X_0\)的一个在\(S_0\)上平坦的闭子预概形\(Y_0\), 其定义理想为\(\mathscr {J}_0=\mathscr {M}_0\), 而\((*)\)给出\[     \mathscr {A}     = \underline {\operatorname {Hom}}_{\mathcal {O}_{X_0}}(\mathscr {J}_0/\mathscr {J}_0^2,\mathcal {O}_{Y_0}\otimes _{\mathcal {O}_{S_0}}\mathscr {J}).\] 此时把\(\mathscr {J}/\mathscr {J}^2\)视为\(Y_0\)在\(X_0\)中的\emph{余法层}, 亦记作\(\mathscr {N}_{Y_0/X_0}\); 将\(\mathscr {A}\)看成\(Y_0\)上的模, 并在\(Y\)上计算\(\operatorname {H}^0\)和\(\operatorname {H}^1\), 是有利的. 若\(Y_0\)局部是\(X_0\)中的完全交, 且\(X\)在\(S\)上平坦, 则在\hyperref[fga3.iv-5-proposition-5.1]{命题~\ref*{fga3.iv-5-proposition-5.1}}中局部延拓的可能性可得到保证, 且\(\mathscr {J}/\mathscr {J}^2\)在\(Y_0\)上局部自由, 可以写作\[     \mathscr {A}     = \check {\mathscr {N}}_{X_0/Y_0}\otimes _{\mathcal {O}_{S_0}}\mathscr {J};\] 其中右侧第一因子是\(Y_0\)在\(X_0\)中的法丛. 应用光滑性的基本判据 [\hyperref[Gro1960b]{3}, III, 3.1], 例如得到:
\begin{corollary}{5.4}\setnumber{5.4}\label{fga3.iv-5-corollary-5.4}\skipsinglepar
\par{}在\hyperref[fga3.iv-5-corollary-5.3]{推论~\ref*{fga3.iv-5-corollary-5.3}}的条件下, 假设\(\mathscr {F}=\mathcal {O}_X\), \(X\)在\(S\)上平坦, 且与\(\mathscr {G}_0\)相应的\(X_0\)的闭子预概形\(Y_0\)局部为完全交. 则\(Q_s\)在\(x\)处的Zariski切空间典范同构于\(\operatorname {H}^0(Y_0,\check {\mathscr {N}}_{X_0/Y_0})\). 若\(\operatorname {H}^1(Y_0,\check {\mathscr {N}_{X_0/Y_0}})=0\), 则Hilbert预概形\(X\)在点\(x\)处关于\(S\)光滑 (其中\(\check {\mathscr {N}}_{X_0/Y_0}\)是\(Y_0\)在\(X_0\)中的法层).
\end{corollary}
\begin{remark}{5.5}\setnumber{5.5}\label{fga3.iv-5-remark-5.5}\skipsinglepar
\par{}这个断言特别适用于\(Y_0\)是\(X_0\)中由\emph{一}个方程定义的完全交的情形, 即它是一个正"Cartier除子". 此时\(\check {\mathscr {N}}_{X_0/Y_0}\)同构于由除子\(Y_0\)所定义的\(X_0\)上的可逆层\(\mathscr {J}^{-1}\)在\(Y_0\)上诱导的层. 在研究非奇异射影簇\(X_0\)上的正除子族时, 特别会遇到这种情形. \(Q\) (或若愿意, \(Q\)中对应于除子的开集\(D\)) 在点\(x\)处的Zariski切空间与\(\mathscr {H}^0(Y_0,\check {\mathscr {N}}_{X_0/Y_0})\)的同构, 在经典代数几何中称为"特征同态" (从前者到后者). 过去只在\(x\)是某个除子的"完全连续族"的参数簇\(T\)的光滑点时定义它; 从我们的观点看, 此参数簇就是概形\(D\)的某个不可约分支, 赋以诱导的\emph{约化}结构. \(T\)在\(x\)处的切空间于是为\(D\)在\(x\)处切空间的一个\emph{子空间}, 所以旧式特征同态确实单射, 但只有加上额外条件才满射, 例如\(D\)在\(x\)处整时. 事实上, Zappa \hyperref[Zap1945]{[8]}构造过一个例子 (其中\(X\)是\emph{复数域}上的非奇异射影曲面), 使特征同态即使在\(T\)的一般点处也不满射. \emph{因此\(D\)即使在所讨论不可约分支的一般点处也不是整的} 这格外鲜明地表明, 为理解曲面理论中最经典的现象, 必须容许带幂零元的簇.
\par{}\emph{[Comp.]} 关于Zappa的例子, 注意Mumford甚至已构造\(\mathbb {P}_\mathbb {C}^3\)的Hilbert概形的一个不可约分支 (其一般点表示次数\(14\)、亏格\(24\)的非奇异曲线), 它在其一般点处不约化. 沿所得曲线作爆破, 他还得到一个\(\mathbb {C}\)上的\(3\)维正则射影概形, 其形式模概形在一般点处不约化; 等价地, 它在\(\mathbb {C}\)上按解析几何意义定义的局部模簇 (见\emph{Séminaire Cartan} \textbf{13}, 1960/61) 在所有点处均不约化.
\end{remark}
\begin{remark}{5.6}\setnumber{5.6}\label{fga3.iv-5-remark-5.6}\skipsinglepar
\par{}我们在\hyperref[fga3.iv-5-corollary-5.4]{推论~\ref*{fga3.iv-5-corollary-5.4}}中给出了光滑性判据, 特别适用于除子概形. Kodaira在\hyperref[Kod1956]{[6]}中给出另一判据, 即要求\(\operatorname {H}^1(X_0,\mathscr {L})\)为零, 其中\(\mathscr {L}=\mathscr {J}_0^{-1}\)是由除子\(Y_0\)定义的\(X_0\)上的可逆层; 这一判据适用于\(S\)是特征\(0\)的域的谱的情形, 而\hyperref[Kod1956]{[6]}在基域为\(\mathbb {C}\)时用超越方法证明了它. 这里指出, 一般地, 对任意\(S\), Kodaira的条件足以保证从除子预概形到\(X/S\)的Picard预概形的典范态射\(D\to \underline {\operatorname {Pic}}_{X/S}\), 在所讨论的点\(x\)处光滑 (一旦知道\(\underline {\operatorname {Pic}}_{X/S}\)存在, 用通常的光滑性判据就容易验证). 此外, 若\(\underline {\operatorname {Pic}}_{X/S}\)在\(x\)的像点处关于\(S\)光滑 (例如\(\underline {\operatorname {Pic}}_{X/S}\)在\(S\)上光滑), 则\(D\)在\(x\)处关于\(S\)光滑. 另一方面, Cartier已证明在\emph{特征\(0\)的域\(k\)}上局部有限型的每个群预概形都在\(k\)上光滑. 合用这两个结果, 就恢复了Kodaira的结果. 由这些说明还可见, 在特征\(p>0\)的域\(K\)上, 若\(\underline {\operatorname {Pic}}_{X/S}\)不在\(k\)上光滑 (例如\(X\)是Igusa曲面时), 则条件\(\operatorname {H}^1(X_0,\mathscr {L})=0\)反而蕴涵\(D\)在\(x\)处不光滑; 若\(K\)代数闭, 则它甚至在\(x\)处不约化.
\end{remark}
\par{}最后给出下面的结果, 它在纤维空间的微分研究中起重要作用:
\begin{proposition}{5.7}\setnumber{5.7}\label{fga3.iv-5-proposition-5.7}\skipsinglepar
\par{}设\(S\)局部 Noether，\(X\)是在\(S\)上有限且平坦的预概形, \(Z\)是\(S\)上的预概形, 并使\(\prod _{X/S}(Z/X)\)存在 (若\(Z\)在\(X\)上拟射影就如此). 若\(Z\)在\(X\)上光滑, 则\(\prod _{X/S}(Z/X)\)在\(S\)上光滑.
\begin{proof}
\par{}这由定义以及通常的光滑性判据 [\hyperref[Gro1960b]{3}, III, §3.1] 立即得到.
\end{proof}
\end{proposition}
\par{}注意, 若\(X\)在\(S\)上有限且平坦, 则\(\prod _{X/S}(Z/X)\)的存在性问题可用很初等的方式处理, 无须使用Hilbert概形理论. 例如可得: 若\(X\)在\(S\)上为根态射, 则无须对\(Z\)作任何限制, \(\prod _{X/S}(Z/X)\)就存在. 举例说, 设\(T\)是\(S\)-预概形, \(T_n\)是\(T\times _S T\)中对角线\(T\times _S T\)的"\(n\)阶无穷小邻域", 带有由两个投影诱导的态射\(p_1,p_2\colon  T_n\to  T\). 借助\(p_1\), 可把\(T_n\)视为\(T\)上的有限预概形; 并进一步假设\(T_n\)在\(T\)上平坦 (当\(T\)在\(S\)上光滑时即如此). 对每个\(T\)上的预概形\(X\), 置\[     (X/T/S)^{(n)}     = \prod _{T_n/S}(p_2^*(X/T)/T_n);\] 这是\(T\)上的一个预概形, 称为\emph{\(X\)在\(T\)上的\(n\)阶截面芽丛} (相对于\(S\)). 它函子地依赖于\(X\), 且若\(X\)光滑, 则它也在\(T\)上光滑.
\chapter{与范数及对称积概念的关系}\thispagestyle{fancy}\label{fga3.iv-6}
\par{}设\(S\)是预概形, \(X\)、\(Y\)是\(S\)-预概形, \[     u\colon (X/S)^n\to  Y   \]是从\(X/S\)的\(n\)重笛卡尔幂到\(Y\)的一个\emph{对称}的\(S\)-态射. 为简单起见, 假设\(S\)局部Noether, \(X\)和\(Y\)在\(S\)上有限型. 此时对于\(X\)上每个支集在\(S\)上有限、并且在\(S\)上平坦且相对\(S\)的秩等于\(n\)的凝聚模\(\mathscr {F}\) (即\(f_*(\mathscr {F})\)是\(S\)上秩为\(n\)的局部自由模), 可以自然地配以\(Y\)在\(S\)上的一个截面: \[     \mathscr {N}_{X/S}^u(\mathscr {F}) \in  \Gamma (Y/S).\] 这里不详述定义, 只指出所得的形式运算是通常范数及迹的形式运算的自然推广. 当\(X\)在\(S\)上的\(n\)重对称幂存在时 (例如对称群\(\sigma _n\)作用于\((X/S)^n\)的各轨道均包含于仿射开集时), 可取\(Y\)为这个对称幂\(\operatorname {Symm}_S^n(X)\), 从而得到一个典范元素\[     \mathscr {N}_{X/S}(\mathscr {F}) \in  \Gamma (\operatorname {Symm}_S^n(X)/S),\] 并可由它恢复各\(\mathscr {N}_{X/S}^u(\mathscr {F})\). 另一个重要情形是\(X\)为\(S\)上的交换幺半群, \(X=Y\), 且态射\(u\)来自\(X\)的合成律. 此时对于\(X\)上的模\(\mathscr {F}\)所对应的\(X\)在\(S\)上的截面, 简记为\(\mathscr {N}(\mathscr {F})\).
\par{}现在假设\(X\)上给定一个凝聚模\(\mathscr {F}\), 使\(\underline {\operatorname {Quot}}_{{\mathscr {F}/X/S}}\)存在; 或者至少使下述函子\(\mathscr {Q}\kern  -.75pt\mathit {uot}_{{\mathscr {F}/X/S}}^n\)可由一个\(S\)-预概形\(\underline {\operatorname {Quot}}_{{\mathscr {F}/X/S}}^n\)表示: 它给\(S\)上的每个\(S'\)配以\(\mathscr {F}'=\mathscr {F}\otimes _{\mathcal {O}_S}\mathcal {O}_{S'}\)的所有在\(S\)上平坦、相对秩为\(n\)的凝聚商层\(\mathscr {M}\)所组成的集合. (若\(X\)在\(S\)上拟射影, 则\(\underline {\operatorname {Quot}}_{{\mathscr {F}/X/S}}^n\)确实存在; 按\cref{fga3.iv-3}的记号, 它恰为\(\underline {\operatorname {Quot}}_{{\mathscr {F}/X/S}}^P\), 其中\(P\)是常数项为\(n\)的常值多项式.) 由于各\(\mathscr {N}_{X/S}^u(\mathscr {M})\)的构造与基底变换相容, 就得到一个典范态射\[     \mathscr {N}_{X/S}^u\colon  \underline {\operatorname {Quot}}_{{\mathscr {F}/X/S}}^n \to  Y;\] 特别地, 若\(X\)在\(S\)上的\(n\)重对称幂存在, 则有\[     \mathscr {N}_{X/S}\colon  \underline {\operatorname {Quot}}_{{\mathscr {F}/X/S}}^n \to  \operatorname {Symm}_S^n(X).\] 最重要的情形是\(\mathscr {F}=\mathcal {O}_X\), 此时得到态射\[     \mathscr {N}_{X/S}\colon  \underline {\operatorname {Hilb}}_{X/S}^n \to  \operatorname {Symm}_S^n(X).\] 显然在\(n=0\)和\(n=1\)时, 它为同构. 但对\(n\geqslant 1\), 即使\(S\)是域\(k\)的谱, 且\(X\)在\(S\)上光滑, 它一般也不是同构, 甚至不是单射态射; 因为仅知道\(X\)的\(0\)维子概形所定义的闭链, 并不能确定该子概形 (例如它对应于\(X\)的闭点\(x\)处局部环\(\mathcal {O}_{X,x}\)中关于极大理想的准素理想\(\mathscr {I}\); 准确地说, 只知道\(\mathscr {I}\)在\(\mathcal {O}_{X,x}\)中作为\(k\)-向量空间的余维). 只能作如下断言 (这里\(S\)又是任意的): 
\begin{enumerate}
    
 \item[a.]若\(X\)在\(S\)上光滑, 则范数态射在\(\underline {\operatorname {Hilb}}_{X/S}^n\)中对应于\(X\)内秩为\(n\)的平展覆盖的分类的开集 (参见\hyperref[fga3.iv-4.b]{§4.b}) 与\(\operatorname {Symm}_S^n(X)\)中对应于无重分支的\(n\)-闭链的开集之间定义一个同构.

    
 \item[b.]若此外\(X\)在\(S\)上的相对维数为\(1\), 则范数态射甚至定义从\(\underline {\operatorname {Hilb}}_{X/S}^n\)到\(\operatorname {Symm}_{X/S}^n\)的一个同构.

  
\end{enumerate} 第二点来自以下事实: 非奇异代数曲线的\(0\)维子概形, 只要知道其相应的闭链就能确定. 同样的说明更一般地适用于非奇异代数概形上的正Cartier除子 (而在这个很特殊的情形, 也不排除Chow簇与Hilbert簇给出相同结果的可能).
\chapter{补充与问题}\thispagestyle{fancy}\label{fga3.iv-7}
\par{}如J.-P. Serre所指出, 由Nagata的一个熟知例子可知, 可以找到一个为域\(k\)的谱的概形\(S\)、一个为\(k\)的二次扩张\(k'\)的谱的\(S\)-概形\(S'\), 以及一个\(3\)维的光滑固有 (但非射影) \(S'\)-概形\(X\), 使\(\prod _{S'/S}(X/S)\)不存在. 这当然更蕴涵Hilbert概形\(\underline {\operatorname {Hilb}}_{X/S}^2\)不存在 (甚至表示\(X\)内\(S\)的秩为\(2\)的平展覆盖的\(k\)-概形也不存在; 从而\(X\)的对称平方也不存在, 参见\cref{fga3.iv-6}). 这对代数几何中进行非射影构造的可能性构成了严重限制. (不过可以相信, 解析几何中不会出现这样的限制, 正如形式几何中不出现一样, 参见\hyperref[fga3.ii]{[2, II]}.) 然而, 若\(X\)是域\(k\)的谱\(S\)上的固有概形, 且\(Z\)在\(X\)上拟射影, 则\(\prod _{X/S}(Z/X)\)存在, 并且是由\(S\)上一列拟射影概形的和构成的概形 (如同射影情形\hyperref[fga3.iv-3-theorem-3.1]{定理~\ref*{fga3.iv-3-theorem-3.1}}). 为证明这一点, 可以取一个射影\(S\)-概形\(X'\)支配\(X\), 从而归结到\(X\)本身射影的情形; 这里不详述证明, 它还使用\hyperref[fga3.i-a.2.b]{[2], I, A, §2.b}中有限态射的分解结果. 这一方法成功的关键在于, 当\(S\)是域的谱时, Chow引理中出现的\(X'\)自动在\(S\)上平坦. 若对\(S\)不加假设, 而仅假设\(X\)在\(S\)上固有平坦, \(Z\)在\(X\)上拟射影, 我不知道结论是否仍成立. 应用中一个重要情形是\(Z\)为\(X\)的闭子概形; 此时若\(\prod _{X/S}(Z/X)\)存在, 它必为\(S\)的闭子概形. 只要\(X\)在\(S\)上射影, 就可以用相当简单的办法直接构造它, 无须用Hilbert概形理论; 所用方法还更一般地表明, 若\(Z\)在\(X\)上仿射, 则\(\prod _{X/S}(Z/X)\)存在且在\(S\)上仿射. 它也表明, 若\(X\)在\(S\)上固有平坦 (但不一定在\(S\)上射影), 则对\(X\)上每个局部平凡向量丛\(Z\), \(\prod _{X/S}(Z/X)\)存在且为\(S\)上的向量丛. 希望能重新研究并统一这些结果.

\section*{参考文献}

\begin{enumerate}
\item[{[1]}]\setnumber{1}\label{fga3.iv-bib-1} BOREL (A.) 与 SERRE (J.-P.). --- Le théorème de Riemann--Roch, \emph{Bull. Soc. math. France}, t. 86, 1958, p. 97--136.
\item[{[2]}]\setnumber{2}\label{fga3.iv-bib-2} GROTHENDIECK (Alexander). --- Technique de construction et théorèmes d'existence en géométrie algébrique, I, II, III, \emph{Séminaire Bourbaki}, t. 13, 1959/60, n\textsuperscript{o} 190, 29页, n\textsuperscript{o} 195, 22页; t. 14, 1960/61, n\textsuperscript{o} 212, 20页.
\item[{[3]}]\setnumber{3}\label{fga3.iv-bib-3} GROTHENDIECK (Alexander). --- \emph{Séminaire de Géométrie algébrique}, I, II, III, IV. Paris, Institut des hautes Études scientifiques, 1960/61 (油印本).
\item[{[4]}]\setnumber{4}\label{fga3.iv-bib-4} GROTHENDIECK (Alexander). --- Technique de construction en géométrie analytique, IV, V, \emph{Séminaire Cartan}, t. 13, 1960/61, n\textsuperscript{o} 11 与 12.
\item[{[5]}]\setnumber{5}\label{fga3.iv-bib-5} GROTHENDIECK (A.) 与 DIEUDONNÉ (J.). --- \emph{Éléments de géométrie algébrique}, I: Le langage des schémas. Paris, Presses universitaires de France, 1960 (Institut des hautes Études scientifiques, Publications mathématiques, 4); II, III, IV (待出版).
\item[{[6]}]\setnumber{6}\label{fga3.iv-bib-6} KODAIRA (K.). --- Characteristic linear systems of complete continuous systems, \emph{Amer. J. of Math.}, t. 78, 1956, p. 716--744.
\item[{[7]}]\setnumber{7}\label{fga3.iv-bib-7} SERRE (Jean-Pierre). --- Exemples de variétés projectives en caractéristique \(p\) non relevables en caractéristique \(0\), \emph{Proc. Nat. Acad. Sc. U. S. A.}, t. 47, 1961, p. 108--109.
\item[{[8]}]\setnumber{8}\label{fga3.iv-bib-8} ZAPPA (G.). --- Sull'esistenza, sopra la superficie algebriche, di sistemi continui completi infiniti, la cui curva e a serie caratteristica incompleta, \emph{Pont. Acad. Sc. Acta}, t. 9, 1945, p. 91--93.

\end{enumerate}

\section*{增补}
\label{fga3.iv-additif}

\begin{center}[校样校正时增补]
\end{center}

\par{}现在看来, 文献[2]的III, §8中的各个猜想都是错误的, 即使对于特征\(0\)的域上的非奇异簇, 无论商的存在性还是商的拟射影性, 都是如此; 即使该群的作用具有闭图, 也仍然如此.

% END chapters/fga3.iv.tex

\clearpage
\setcounter{chapter}{0}
\renewcommand\fgaAnchorPrefix{fga3.v}

% BEGIN chapters/fga3.v.tex
\part{Picard 概形：存在定理}\label{fga3.v}\thispagestyle{fancy}\begin{original}{}\label{fga3.v-original-citation}\skipsinglepar\par{}A. Grothendieck．“Technique de descente et théorèmes d'existence en géométrie algébrique, V: Les schémas de Picard: Théorèmes d'existence”．\emph{Séminaire Bourbaki} \textbf{14}（1961–62），第 232 次报告，1962 年 2 月。\url{http://www.numdam.org/book-part/SB_1961-1962__7__143_0/}\end{original}\setcounter{chapter}{0}\chapter{相对 Picard 群与函子}\thispagestyle{fancy}\label{fga3.v-1}\par{}对于每个预概形（更一般地，每个赋环空间）\(X\)，定义 \(X\) 的（\emph{绝对}）\emph{Picard 群}，记为 \(\operatorname {Pic}(X)\)，为 \(X\) 上可逆模（即局部同构于 \(\mathcal {O}_X\) 的模）的同构类所成的群。因此有典范同构 \begin{eqn}{1.1}\setnumber{1.1}\label{fga3.v-1-equation-1.1}\skipsinglepar\[       \operatorname {Pic}(X) \xrightarrow {\sim } \operatorname {H}^1(X,\mathcal {O}_X^\times )     \tag{1.1}     \]\end{eqn}，其中 \(\mathcal {O}_X^\times \) 表示 \(\mathcal {O}_X\) 的可逆元层（可将它等同于可逆模 \(\mathcal {O}_X\) 的自同构层）。注意，\(X\mapsto \operatorname {Pic}(X)\) 显然是关于 \(X\) 的反变函子，且同构 \hyperref[fga3.v-1-equation-1.1]{式~\ref*{fga3.v-1-equation-1.1}} 是函子性的。\par{}设 \(X\) 是预概形 \(S\) 上的预概形。由以上讨论，当 \(S'\) 在 \(S\) 上预概形的范畴 \(\mathtt {Sch}_{/S}\) 中变化时，就得到反变函子 \(S'\mapsto \operatorname {Pic}(X\times _S S')\)。这个函子不可能“可表”（\hyperref[fga3.ii-a.1]{[4]，II，A}），因为我们拟加以分类的可逆模具有 \emph{非平凡自同构}，因而此函子并不具有“\emph{局部性质}”（[\hyperref[Gro1960a]{5}，IV，5.4]）。因此，有必要使其“局部化”：对每个相对预概形 \(X/S\)，引入一个相对性质的群 \begin{eqn}{1.2}\setnumber{1.2}\label{fga3.v-1-equation-1.2}\skipsinglepar\[       \operatorname {Pic}'(X/S)       = \operatorname {H}^0(S,\operatorname {R}^1f_*(\mathcal {O}_X^\times ))     \tag{1.2}     \]\end{eqn}（其中 \(f\colon  X\to  S\) 是结构态射；参见 \hyperref[fga3.ii-c.3]{[4]，II，C.3}）。在 \emph{同上} 中，这个群称为相对 Picard 群；不过，出于下文将说明的原因，这里最好称之为 \(X/S\) 的 \emph{受限相对 Picard 群}。当 \(S'\) 在 \(\mathtt {Sch}_{/S}\) 中变化时，\(S'\mapsto \operatorname {Pic}'(X\times _S S'/S')\) 是关于 \(S'\) 的反变函子，也记为 \(\mathscr {P}\kern  -1pt\mathit {ic}'_{X/S}\)，其本质上由公式 \begin{eqn}{1.3}\setnumber{1.3}\label{fga3.v-1-equation-1.3}\skipsinglepar\[       \mathscr {P}\kern  -1pt\mathit {ic}'_{X/S}(S')       = \operatorname {Pic}(X\times _S S'/S').     \tag{1.3}     \]\end{eqn} 给出。现在，这个函子具有“局部性质”，因为我们已经作了为此所需的处理。直观地说，\hyperref[fga3.v-1-equation-1.3]{式~\ref*{fga3.v-1-equation-1.3}} 的右端可理解为以参数预概形 \(S'/S\) 为指标的、\(X/S\)（的纤维）上可逆层的类的“代数族”所成的集合。当函子 \(\mathscr {P}\kern  -1pt\mathit {ic}'\) 可表时，将表示它的 \(S\) 上的预概形记为 \(\underline {\operatorname {Pic}}_{X/S}\)，称为 \(X\) 在 \(S\) 上的 \emph{Picard 预概形}；这时有 \begin{eqn}{1.4}\setnumber{1.4}\label{fga3.v-1-equation-1.4}\skipsinglepar\[       \operatorname {Hom}_S(S',\underline {\operatorname {Pic}}_{X/S})       \cong  \mathscr {P}\kern  -1pt\mathit {ic}'_{X/S}(S')       = \operatorname {Pic}'(X\times _S S'/S').     \tag{1.4}     \]\end{eqn}。然而，在某些重要情形中，\(\mathscr {P}\kern  -1pt\mathit {ic}'_{X/S}\) 不可表（例如域 \(k\) 上没有 \(k\)-有理点的“Brauer–Severi”簇），但相对 Picard 预概形仍有自然的定义。这是因为，从绝对 Picard 群 \(\operatorname {Pic}(X\times _S S'/S')\) 定义函子 \(\mathscr {P}\kern  -1pt\mathit {ic}'\) 时，我们所作的局部化还不够；更确切地说，\(\mathscr {P}\kern  -1pt\mathit {ic}'\) 一般并不“与忠实平坦下降相容”。下面详细说明。\par{}令 \((\mathscr {M})\) 为所有 \emph{忠实平坦且拟紧} 的预概形态射所成的集合；此集合在基变换及复合下稳定。设 \(P\) 是从 \(\mathtt {Sch}_{/S}\) 到集合范畴的反变函子。对每个满足 \(u\in (\mathscr {M})\) 的 \(S\)-态射 \(u\colon  T'\to  T\)，考虑图表 \begin{eqn}{1.5}\setnumber{1.5}\label{fga3.v-1-equation-1.5}\skipsinglepar\[       P(T)       \to  P(T')       \rightrightarrows  P(T'\times _T T')     \tag{1.5}     \]\end{eqn}；它由对图表 \[     T     \leftarrow  T'     \underset {\mathrm {pr}_2}{\overset {\mathrm {pr}_1}{\leftleftarrows }} T'\times _T T'.   \] 施以 \(P\) 得到。若 \(P\) 可表，则由下降理论（\hyperref[fga3.i-b.1-theorem-2]{[4]，I，B，定理 2}）可知，对所有 \(u\in (\mathscr {M})\)，图表 \hyperref[fga3.v-1-equation-1.5]{式~\ref*{fga3.v-1-equation-1.5}} 都正合。我们称这一事实为 \(P\) 与 \((\mathscr {M})\) 相容；就此处而言，也就是说 \(P\)“与忠实平坦下降相容”，或说 \(\mathtt {Sch}_{/S}\) 上的 \emph{“预层”} \(P\) 对集合 \((\mathscr {M})\) 所给出的局部化概念而言是一个 \emph{“层”}。对于任意 \(P\)，用一个在通常拓扑局部化情形中熟知的标准过程，可以给它配备一个“层”\(\mathscr {P}\) 及函子同态 \(P\to \mathscr {P}\)，后者具有显然的泛性质。\(\mathscr {P}\) 的构造可明确写成如下形式：为定义 \(\mathscr {P}(T)\)，对每个满足态射 \(u\colon  T'\to  T\) 属于 \((\mathscr {M})\) 的、\(T\) 上的 \(T'\)，用 \(\overline {\operatorname {H}}^0(T'/T,P)\) 表示 \(P(T')\) 的如下子集：其元素 \(\xi \) 在 \(P(T'\times _T T')\) 中的像 \(\xi _1,\xi _2\) 满足，存在 \((\mathscr {M})\) 中的态射 \(v\colon  T''\to  T'\times _T T'\)，使 \(\xi _1\) 与 \(\xi _2\) 在 \(P(T'')\) 中有相同的像。\par{}（注意：如此定义的集合 \(\overline {\operatorname {H}}^0\) 比 \hyperref[fga3.i-a.4.a]{[4]，I，A.4(a)} 中引入的集合 \(\operatorname {H}^0(T'/T,P)\) 大。）让 \(T'\) 在固定的 \(T\) 上变化（始终要求 \(u\in (\mathscr {M})\)），各个 \(\overline {\operatorname {H}}^0(T'/T,P)\) 构成一个归纳系（在所有 \(T'\) 构成的集合上，以支配关系给出预序），并令 \begin{eqn}{1.6}\setnumber{1.6}\label{fga3.v-1-equation-1.6}\skipsinglepar\[       \mathscr {P}(T)       = \varinjlim _{T'} \overline {\operatorname {H}}^0(T'/T,P).     \tag{1.6}     \]\end{eqn}。这个表达式关于 \(T\) 的函子性是显然的。\par{}当 \[     P(T)     = \operatorname {Pic}(X\times _S T)   \] 时，称由 \hyperref[fga3.v-1-equation-1.6]{式~\ref*{fga3.v-1-equation-1.6}} 定义的 \(\mathtt {Sch}_{/S}\) 上的反变函子为 \(X\) 在 \(S\) 上的 \emph{相对 Picard 函子}，记为 \(\mathscr {P}\kern  -1pt\mathit {ic}_{X/S}\)；并将群 \(\mathscr {P}\kern  -1pt\mathit {ic}_{X/S}(S)\) 定义为 \(X\) 在 \(S\) 上的 \emph{相对 Picard 群}，记为 \(\operatorname {Pic}(X/S)\)。于是有显然的双射 \begin{eqn}{1.7}\setnumber{1.7}\label{fga3.v-1-equation-1.7}\skipsinglepar\[       \mathscr {P}\kern  -1pt\mathit {ic}_{X/S}(T)       \xrightarrow {\sim } \operatorname {Pic}(X\times _S T/T).     \tag{1.7}     \]\end{eqn}。因此，\(\operatorname {Pic}(X/S)\) 的一个元素由群 \(\operatorname {Pic}(X\times _S S')\) 的一个元素 \(\xi '\) 给出（其中 \(S'\to  S\) 忠实平坦且拟紧），并要求能够找到忠实平坦拟紧态射 \(S''\to  S'\times _S S'\)，使 \(\xi '\) 在 \(\operatorname {Pic}(X\times _S S'')\) 中的两个逆像相同。满足刚才条件的元素 \(\xi '\)（属于 \(\operatorname {Pic}(X\times _S S')\)）和 \(\xi _1\)（属于 \(\operatorname {Pic}(X\times _S S_1)\)）定义 \(\operatorname {Pic}(X/S)\) 的同一元素，当且仅当存在忠实平坦拟紧态射 \(S'_1\to  S'\times _S S_1\)，使这两个元素在 \(\operatorname {Pic}(X\times _S S'_1)\) 中的像相等。实际运算时，往往改用上面引入的函子 \(P'=\mathscr {P}\kern  -1pt\mathit {ic}'_{X/S}\) 更为方便。立刻可见，典范态射 \(P\to  P'\) 定义了一个 \emph{同构} \begin{eqn}{1.8}\setnumber{1.8}\label{fga3.v-1-equation-1.8}\skipsinglepar\[       \mathscr {P}       \xrightarrow {\sim } \mathscr {P}'     \tag{1.8}     \]\end{eqn}，从而用 \(\mathscr {P}\kern  -1pt\mathit {ic}'_{X/S}=P'\) 对 \(\mathscr {P}\kern  -1pt\mathit {ic}_{X/S}\) 给出了一种通常更方便的描述。由下面的 \hyperref[fga3.v-2-corollary-2.3]{推论~\ref*{fga3.v-2-corollary-2.3}}，若将刚才对 \(\operatorname {Pic}(X/S)\) 的描述中的 \(P\) 换为 \(P'\)，则至少在 \emph{同上} 所给的条件下，可以取 \(S''=S'\times _S S'\) 及 \(S'_1=S'\times _S S_1\)。\par{}若函子 \(\mathscr {P}\kern  -1pt\mathit {ic}_{X/S}\) 可表，就说 \(X/S\) 有 Picard 预概形，并称表示此函子的 \(S\) 上的预概形为 \(X\) 在 \(S\) 上的 \emph{Picard 预概形}，记为 \(\underline {\operatorname {Pic}}_{X/S}\)。显然，为此只须 \(P'=\mathscr {P}\kern  -1pt\mathit {ic}_{X/S}\) 可表，因为此时 \(P'\) 已是一个“层”，而等式 \hyperref[fga3.v-1-equation-1.8]{式~\ref*{fga3.v-1-equation-1.8}} 表明，态射 \(P'\to \mathscr {P}'\) 可等同于典范态射 \begin{eqn}{1.9}\setnumber{1.9}\label{fga3.v-1-equation-1.9}\skipsinglepar\[       \mathscr {P}\kern  -1pt\mathit {ic}'_{X/S}       \to  \mathscr {P}\kern  -1pt\mathit {ic}_{X/S}     \tag{1.9}     \]\end{eqn}，后者于是是同构。这说明，我们的术语与上面在 \hyperref[fga3.v-1-equation-1.4]{式~\ref*{fga3.v-1-equation-1.4}} 中引入的术语相容。一般地，只要 \(\underline {\operatorname {Pic}}_{X/S}\) 存在，它就由函子性同构 \begin{eqn}{1.10}\setnumber{1.10}\label{fga3.v-1-equation-1.10}\skipsinglepar\[       \operatorname {Hom}_S(S',\underline {\operatorname {Pic}}_{X/S})       \xrightarrow {\sim } \operatorname {Pic}(X\times _S S'/S').     \tag{1.10}     \]\end{eqn} 所定义。\chapter{各类相对与绝对 Picard 群之间的关系}\thispagestyle{fancy}\label{fga3.v-2}\begin{proposition}{2.1}\setnumber{2.1}\label{fga3.v-2-proposition-2.1}\skipsinglepar\par{}设 \(f\colon  X\to  S\) 是满足 \(\mathcal {O}_S\xrightarrow {\sim } f_*(\mathcal {O}_X)\) 的态射。则有正合列 \[       0       \to  \operatorname {Pic}(S)       \to  \operatorname {Pic}(X)       \to  \operatorname {Pic}'(X/S).     \]。若 \(X\) 在 \(S\) 上有截面，则最后一个态射满射，即有同构 \[       \operatorname {Pic}'(X/S)       \xrightarrow {\sim } \operatorname {Pic}(X)/\operatorname {Pic}(S).     \]。\begin{proof}\par{}此正合列可视为对应于 \(f\) 和 \(\mathcal {O}_X\) 的 Leray 谱序列所给出的正合列的低次部分。第二个断言同样是形式的。\end{proof}\end{proposition}\begin{proposition}{2.2}\setnumber{2.2}\label{fga3.v-2-proposition-2.2}\skipsinglepar\par{}设 \(f\colon  X\to  S\) 是满足 \(\mathcal {O}_S\xrightarrow {\sim } f_*(\mathcal {O}_X)\) 的拟紧分离态射，\(S'\to  S\) 是忠实平坦拟紧态射。那么：\begin{enumerate}
      
 \item[i.]\(\operatorname {Pic}'(X/S)\to \operatorname {Pic}'(X\times _S S'/S')\) 是单射；

      
 \item[ii.]若 \(X\) 在 \(S\) 上局部有截面（即每个 \(s\in  S\) 都有开邻域 \(U\)，使得 \(X|U\) 在 \(U\) 上有截面），则图表 \[           \operatorname {Pic}'(X/S)           \to  \operatorname {Pic}'(X\times _S S'/S')           \rightrightarrows  \operatorname {Pic}'(X\times _S S''/S'')         \]（其中 \(S''=S'\times _S S'\)）正合。

    \end{enumerate}\begin{proof}\par{}借助忠实平坦下降的初等性质，第一个断言由下面的一般注记得到。若 \(f\colon  X\to  S\) 是满足 \(\mathcal {O}_S\xrightarrow {\sim } f_*(\mathcal {O}_X)\) 的态射，则从 \(S\) 上有限型局部自由模的范畴到 \(X\) 上有限型局部自由模的范畴的函子 \(\mathscr {F}\mapsto  f^*(\mathscr {F})\) 是全忠实的；其本质像是 \(X\) 上满足如下条件的模 \(\mathscr {G}\)：\(f_*(\mathscr {G})\) 局部自由，且典范同态 \[         f^*f_*(\mathscr {G})         \to  \mathscr {G}       \] 为同构。第二个断言已在 \hyperref[fga3.i-b.4]{[4]，I，B.4} 中借助下降理论证明。\end{proof}\end{proposition}\par{}\hyperref[fga3.v-2-proposition-2.2]{命题~\ref*{fga3.v-2-proposition-2.2}} 的结果也可表述如下：\begin{corollary}{2.3}\setnumber{2.3}\label{fga3.v-2-corollary-2.3}\skipsinglepar\par{}在 \hyperref[fga3.v-2-proposition-2.2]{命题~\ref*{fga3.v-2-proposition-2.2}} 的条件下，典范同态 \hyperref[fga3.v-1-equation-1.9]{式~\ref*{fga3.v-1-equation-1.9}} \(\mathscr {P}\kern  -1pt\mathit {ic}'_{X/S}\to \mathscr {P}\kern  -1pt\mathit {ic}_{X/S}\) 是单射；若 \(X\) 局部地在 \(S\) 上有截面，则它还是双射。（后一情形中，相对 Picard 群 \(\operatorname {Pic}(X/S)\) 等同于受限相对 Picard 群 \(\operatorname {Pic}'(X/S)\)。）\end{corollary}\par{}将它与 \hyperref[fga3.v-2-proposition-2.1]{命题~\ref*{fga3.v-2-proposition-2.1}} 结合，便得到：\begin{corollary}{2.4}\setnumber{2.4}\label{fga3.v-2-corollary-2.4}\skipsinglepar\par{}在 \hyperref[fga3.v-2-proposition-2.2]{命题~\ref*{fga3.v-2-proposition-2.2}} 的条件下，有正合列 \[       0       \to  \operatorname {Pic}(S)       \to  \operatorname {Pic}(X)       \to  \operatorname {Pic}(X/S).     \]。若 \(X\) 在 \(S\) 上有截面，则最后一个同态满射，即有同构 \[       \operatorname {Pic}(X/S)       \xrightarrow {\sim } \operatorname {Pic}(X)/\operatorname {Pic}(S).     \]。\end{corollary}\begin{remark}{2.5}\setnumber{2.5}\label{fga3.v-2-remark-2.5}\skipsinglepar\par{}设 \(f\colon  X\to  S\) 是满足 \(\mathcal {O}_S\xrightarrow {\sim } f_*(\mathcal {O}_X)\) 的态射，\(g\) 是 \(X\) 在 \(S\) 上的一个截面。设 \(\mathscr {L}\) 是 \(X\) 上的可逆模。将 \(\mathscr {L}\) 的 \emph{\(g\)-刚性化} 定义为一个同构 \(\mathcal {O}_S\xrightarrow {\sim } g^*(\mathscr {L})\)，并将 \emph{\(g\)-刚性化可逆模} 定义为 \(X\) 上配有一个 \(g\)-刚性化的可逆模 \(\mathscr {L}\)。这样一个结构的每个自同构都是平凡的，而 \(\operatorname {Pic}'(X/S)\) 可等同于 \(S\) 上 \(g\)-刚性化可逆模的同构类所成的集合。（正是这一事实使我们能利用下降理论证明 \hyperref[fga3.v-2-proposition-2.2]{命题~\ref*{fga3.v-2-proposition-2.2}} 的 (ii)。）这给出了对 \(\operatorname {Pic}(X/S)\) 的新解释，至少当 \(f\) 还拟紧且分离时是如此；因为这时由 \hyperref[fga3.v-2-corollary-2.3]{推论~\ref*{fga3.v-2-corollary-2.3}} 可知 \(\operatorname {Pic}(X/S)\xrightarrow {\sim }\operatorname {Pic}'(X/S)\)。\end{remark}\begin{remark}{2.6}\setnumber{2.6}\label{fga3.v-2-remark-2.6}\skipsinglepar\par{}设 \(f\colon  X\to  S\) 是 \hyperref[fga3.v-2-proposition-2.2]{命题~\ref*{fga3.v-2-proposition-2.2}} 中的态射，\(S'\to  S\) 是忠实平坦拟紧态射，使得存在一个 \(S\)-态射 \(S'\to  X\)，即 \(X'=X\times _S S'\) 在 \(S'\) 上有截面。令 \(S''=S'\times _S S'\)、\(X''=X\times _S S''\)，并考虑正合列 \[       \operatorname {Pic}(X/S)       \to  \operatorname {Pic}(X'/S')       \rightrightarrows  \operatorname {Pic}(X''/S'').     \]。施以 \hyperref[fga3.v-2-corollary-2.4]{推论~\ref*{fga3.v-2-corollary-2.4}}，得到正合列 \[       \operatorname {Pic}(X/S)       \to  \operatorname {Pic}(X')/\operatorname {Pic}(S')       \rightrightarrows  \operatorname {Pic}(X'')/\operatorname {Pic}(S'').     \]。特别地，相对 Picard 群的每个元素都已经“来自”\(\operatorname {Pic}(X')\) 的某个元素。这大大简化了上一节对相对 Picard 群的描述，甚至也简化了 \(X\) 在 \(S\) 上的 Picard 函子的描述，因为对每个 \(S\) 上的 \(T\)，都可以将上述讨论用于 \(X\times _S T/T\) 及态射 \(T'=S'\times _S T\to  T\)。例如，若 \(f\) 本身忠实平坦，就可取 \(S'=X\)；于是，只要 \(f\) 还为有限型（分别为光滑态射，等等），在描述相对 Picard 函子 \(T\mapsto \operatorname {Pic}(X\times _S T/T)\) 时，就可以只考虑有限型（分别为光滑态射，等等）的基变换 \(T'\to  T\)。若 \(f\) 射影且平坦，\(S\) 局部 Noether，则可以证明，上面可取这样的 \(S'\to  S\)，使 \(S'\) 是 \(S\) 中覆盖 \(S\) 的若干开集 \(S_i\) 的 \emph{平坦覆盖} \(S'_i\) 的直和；若 \(f\) 还可分，则可取 \(S'_i\) 在 \(S_i\) 上 étale。\end{remark}\chapter{主要存在定理：陈述}\thispagestyle{fancy}\label{fga3.v-3}\par{}关于 Picard 预概形的存在性，目前还没有一个涵盖所有已知情形的陈述，甚至连猜想也没有。如果可以这样说，一个“实际上必要”的条件是 \(f\colon  X\to  S\) 为 \emph{固有}（以保证基本的有限性性质）并且 \emph{平坦}。这些条件并不充分，即使 \(S\) 是域 \(k\)（譬如复数域 \(\mathbb {C}\)）上的对偶数代数 \(k[t]/(t^2)\) 的谱，且 \(X\) 的维数为 \(1\)，也是如此。在撰写本报告时，关于 Picard 预概形的最重要的存在定理都由下列定理推出：\begin{theorem}{3.1}\setnumber{3.1}\label{fga3.v-3-theorem-3.1}\skipsinglepar\par{}设 \(f\colon  X\to  S\) 是局部 Noether 预概形之间的态射。假设：\begin{enumerate}
      
 \item[i.]\(f\) 射影；

      
 \item[ii.]\(f\) 平坦；

      
 \item[iii.]\(f\) 的几何纤维整。

    \end{enumerate}在这些条件下，\(\underline {\operatorname {Pic}}_{X/S}\) 存在。\end{theorem}\par{}下面两节将给出的证明梗概也将同时说明下述事实。设 \(\xi \) 是 \(\underline {\operatorname {Pic}}_{X/S}\) 的一个截面，对应于 \(X/S\) 上的一个甚丰沛层 \(\mathcal {O}_X(1)\)（即由射影嵌入 \(X\to \mathbb {P}(\mathscr {E})\) 诱导的层）；那么，\(\underline {\operatorname {Pic}}_{X/S}\) 中存在开子集 \(U\)，它是 \(S\) 上拟射影开子集的不交并，使得 \(U\) 在由 \(\xi \) 给出的平移下稳定，且 \(\underline {\operatorname {Pic}}_{X/S}\) 是开集 \(U-n\xi \) 的递增并（每个开集都同构于 \(U\)）。特别地，由此可得，在 \hyperref[fga3.v-3-theorem-3.1]{定理~\ref*{fga3.v-3-theorem-3.1}} 的条件下，\(\underline {\operatorname {Pic}}_{X/S}\) 在 \(S\) 上 \emph{分离}。\begin{remark}{3.2}\setnumber{3.2}\label{fga3.v-3-remark-3.2}\skipsinglepar\par{}一些例子（例如取 \(S\) 为离散赋值环的谱，\(X\) 在 \(S\) 上的相对维数为 \(1\)）表明，若删去 \hyperref[fga3.v-3-theorem-3.1]{定理~\ref*{fga3.v-3-theorem-3.1}} 中的假设 (iii)，代之以较弱的假设：对每个 \(s\in  S\)，同态 \(k(s)\to \operatorname {H}^0(X_s,\mathcal {O}_{X_s})\) 均为同构，那么 \(\underline {\operatorname {Pic}}_{X/S}\) 在 \(S\) 上未必分离。这既可能出现在 \(f\) 的几何纤维约化、但一个整的泛几何纤维经特殊化“分裂”为两个不可约分支时，也可能出现在 \(f\) 的几何纤维不可约、但一个整的泛几何纤维特殊化为一个“多重纤维”时。第一种情形的例子是一条圆锥曲线退化为两条相交直线；D. Mumford 曾给我看过第二种情形的一个例子：一条椭圆曲线退化为一条二重椭圆曲线。这些例子在任意特征下都成立。\end{remark}\begin{remark}{3.3}\setnumber{3.3}\label{fga3.v-3-remark-3.3}\skipsinglepar\par{}在 \hyperref[fga3.v-3-theorem-3.1]{定理~\ref*{fga3.v-3-theorem-3.1}} 的条件下，我不知道 \(\underline {\operatorname {Pic}}_{X/S}\) 是否是若干开集的不交并，其中每个开集在 \(S\) 上为有限型，因而拟射影。注意，与 Hilbert 概形的情形（\hyperref[fga3.iv]{[4]，IV}）相同，对 Hilbert 多项式 \(Q\in \mathbb {Q}[t]\) 的研究可以将 \(\underline {\operatorname {Pic}}_{X/S}\) 分解为开集 \(\underline {\operatorname {Pic}}_{X/S}^Q\) 的不交和；这些开集在 \(S\) 上似乎应当拟射影。至少当 \(f\) 为光滑态射时，下一次报告将证明这一点。请注意，若将假设 (i) 换为“\(X\) 在 \(S\) 上 \emph{局部地} 射影”（这足以证明 \hyperref[fga3.v-3-theorem-3.1]{定理~\ref*{fga3.v-3-theorem-3.1}}，因为 \(\underline {\operatorname {Pic}}_{X/S}\) 的存在问题显然在 \(S\) 上是局部的），则很容易给出满足 \emph{\(\underline {\operatorname {Pic}}_{X/S}\) 含有在 \(S\) 上并非有限型的连通分支} 的例子。例如，设 \(X_0\) 是代数闭域 \(k\) 上的非奇异射影代数簇，配有自同构 \(u\) 及 \(X_0\) 的 Néron–Severi 群中的元素 \(\xi \)，使得各 \(u^n(\xi )\) 两两不同。譬如，可取 \(X_0\) 为一条椭圆曲线 \(E\) 与自身的乘积，取 \(u\) 为 \(E\times  E\) 的自同构 \((x,y)\mapsto (x,y+x)\)。设 \(S\) 是两条非奇异不可约曲线的并，它们交于两点 \(a\) 和 \(b\)。在 \(S\) 上存在群 \(\mathbb {Z}\) 的连通主覆盖 \(P\)；利用 \(u\) 所定义的 \(\mathbb {Z}\) 在 \(X_0\) 上的作用，便得到 \(S\) 上纤维为 \(X_0\) 的伴随纤维丛（在 \(S\setminus \{a\}\) 和 \(S\setminus \{b\}\) 上平凡）；在这里所考虑的特殊情形中，它实际上是 \(S\) 上的一个 \emph{Abel} 概形。容易看出，\(\underline {\operatorname {Pic}}_{X/S}\) 也是由 \(P\) 及 \(\mathbb {Z}\) 通过 \(u\) 在 \(\underline {\operatorname {Pic}}_{X_0/k}\) 上的作用所得到的伴随纤维丛；它含有一个同构于 \(P\times \underline {\operatorname {Pic}}_{X_0/k}^0\) 的连通分支（其中 \(\underline {\operatorname {Pic}}^0\) 表示 \(\underline {\operatorname {Pic}}\) 中单位元的连通分支），这个分支在 \(S\) 上并非有限型。（在 \hyperref[fga3.v-3-remark-3.3]{注~\ref*{fga3.v-3-remark-3.3}} 所述的 \(S\) 上非分离 Picard 预概形的各种情形中，也可以构造类似现象。）\emph{[编者]} 这里提出的问题已由 Mumford 给出肯定回答（见 \hyperref[fga3.vi]{FGA~\ref*{fga3.vi}}）。\end{remark}\chapter{相对 Cartier 除子与射影丛}\thispagestyle{fancy}\label{fga3.v-4}\par{}我们只需用到正除子，因此本节以下省略“正”这一限定词。\par{}设 \(X\) 是一个预概形。\(X\) 上的 \emph{Cartier 除子}，简称除子，是 \(X\) 的一个闭子预概形 \(D\)，其定义理想 \(\mathscr {J}\) 是一个 \emph{可逆} 模，即局部由 \(\mathcal {O}_X\) 的一个 \emph{非零因子} 截面生成。给 \(D\) 配上可逆模 \[     \mathscr {L}(D)     = \mathscr {J}^{-1}   \]，则典范单射 \(\mathscr {J}\to \mathcal {O}_X\) 给出典范同态 \[     s_D\colon  \mathcal {O}_X     \to  \mathscr {J}^{-1}     = \mathscr {L}(D)   \]，即 \(s_D\in \Gamma (X,\mathscr {L}(D))\)。给定一个除子，本质上 \emph{等价} 于给定 \(X\) 上的一个可逆模 \(\mathscr {L}\) 及其处处不是零因子的截面 \(s\)：给这样一对 \((\mathscr {L},s)\) 配上 \(s\) 的“除子”，记为 \(\operatorname {div}(s)\)。对 \(X\) 上给定的可逆模 \(\mathscr {L}\)，定义 \(\mathscr {L}\) 的除子 \(D\) 所成的集合，与商集 \(\Gamma (X,\mathscr {L})^{\times }/\Gamma (X,\mathcal {O}_X^\times )\) 一一对应；其中 \(\Gamma (X,\mathscr {L})^\times \) 表示 \(\Gamma (X,\mathscr {L})\) 中处处不是零因子的截面所成的子集。\par{}现在设有一个局部有限型态射 \(f\colon  X\to  S\)；为简单起见，假设 \(S\) 局部 Noether。设 \(\mathscr {J}\) 是 \(X\) 上的凝聚理想，\(D\) 是它在 \(X\) 中定义的子概形，并令 \(x\in  X\)、\(s=f(x)\)。我们将证明下列条件等价：\begin{enumerate}
    
 \item[i.]\(\mathscr {J}\) 在 \(x\) 处可逆（即 \(\mathscr {J}_x\) 由 \(\mathcal {O}_{X,x}\) 的一个正则元生成），且 \(D\) 在 \(x\) 处在 \(S\) 上平坦。

    
 \item[ii.]\(X\) 和 \(D\) 在 \(x\) 处在 \(S\) 上平坦，且 \(D_s\) 在点 \(x\) 处是纤维 \(X_s\) 上的 Cartier 除子。

    
 \item[iii.]\(X\) 在 \(x\) 处在 \(S\) 上平坦，且 \(\mathscr {J}_x\) 由一个元素 \(f_x\) 生成，该元素在 \(X_s\) 上诱导一个非零因子芽。

  \end{enumerate}此时称 \(D\) 在所考虑的点处为 \(X/S\) 上的一个 \emph{相对 Cartier 除子}，或简称相对除子。注意，在 (i) 中，\(D\) 在 \(x\) 的某个邻域内的各点处也都是相对除子。因此，若 \(X\) 和 \(D\) 在 \(S\) 上平坦，且 \(D\) 在 \(S\) 上固有，那么，使 \(D_s\) 成为 \(X_s\) 中的 Cartier 除子的 \(s\in  S\)（即使 \(D\) 在 \(X_s\) 的各点处都是相对 Cartier 除子的那些点）所成的集合是 \(S\) 的开子集。上面的定义也已作了必要安排，以保证相对 Cartier 除子的概念在任意基变换 \(S'\to  S\) 下稳定。于是考虑 \(X/S\) 上相对除子的集合 \(\operatorname {Div}(X/S)\)，以及由 \[     \mathscr {D}\kern  -1pt\mathit {iv}_{X/S}(S')     = \operatorname {Div}(X\times _S S'/S').   \] 定义的、关于在 \(S\) 中变化的 \(S'\) 的反变函子。假设 \(X\) 在 \(S\) 上平坦且固有。由相对 Cartier 除子的刻画 (ii)，可将 \(\mathscr {D}\kern  -1pt\mathit {iv}_{X/S}\) 看作 \hyperref[fga3.iv]{[4]，IV} 中定义的函子 \(\mathscr {H}\kern  -1.75pt\mathit {ilb}_{X/S}\) 的子函子；由上述注记，包含态射 \[     \mathscr {D}\kern  -1pt\mathit {iv}_{X/S}     \to  \mathscr {H}\kern  -1.75pt\mathit {ilb}_{X/S}   \] “可由开浸入表示”（参见 [\hyperref[Gro1960a]{5}，IV，3.13]）。利用 \hyperref[fga3.iv]{[4]，IV} 的主要存在定理，得到：\begin{proposition}{4.1}\setnumber{4.1}\label{fga3.v-4-proposition-4.1}\skipsinglepar\par{}假设 \(f\colon  X\to  S\) 射影且平坦。那么函子 \(\mathscr {D}\kern  -1pt\mathit {iv}_{X/S}\) 可表；更确切地说，它由 \(\underline {\operatorname {Hilb}}_{X/S}\) 的一个开集表示。\end{proposition}\par{}给定 \(X/S\) 上的甚丰沛层 \(\mathcal {O}_X(1)\)，利用 \(\underline {\operatorname {Hilb}}_{X/S}\) 按 Hilbert 多项式 \(Q\in \mathbb {Q}[t]\) 典范分解为开集 \(\underline {\operatorname {Hilb}}_{X/S}^Q\) 之和的分解，得到类似的分解 \[     \underline {\operatorname {Div}}_{X/S}     = \sqcup _{Q\in \mathbb {Q}[t]} \underline {\operatorname {Div}}_{X/S}^Q   \]，它是若干两两不交、在 \(S\) 上 \emph{拟射影} 的开集之和。\par{}利用映射 \(D\mapsto \mathscr {L}(D)\)，得到一个函子性同态 \[     \mathscr {D}\kern  -1pt\mathit {iv}_{X/S}     \to  \mathscr {P}\kern  -1pt\mathit {ic}_{X/S}   \tag{+}   \]；我们拟研究它。在相当一般的条件下，它相对可表（[\hyperref[Gro1960a]{5}，IV，3]）。因此，从 \(\mathscr {P}\kern  -1pt\mathit {ic}_{X/S}(S')\) 的一个元素 \(\xi \) 出发，并为简化记号而假设 \(S'=S\)，我们将证明 \(\mathscr {D}\kern  -1pt\mathit {iv}_{X/S}\) 的相应子函子可表。首先考虑 \(\xi \) 由 \(X\) 上的一个可逆模 \(\mathscr {L}\) 定义的情形。假设 \(X\) 在 \(S\) 上固有且平坦，\(X\) 在 \(S\) 上的几何纤维都是整的；这还蕴含（[\hyperref[GD1960]{2}，III，§7]）\(\mathcal {O}_S\xrightarrow {\sim } f_*(\mathcal {O}_X)\)，而且这一等式在任意基变换 \(S'\to  S\) 后仍成立。于是，\(X/S\) 上满足 \(\mathscr {L}(D)\) 与 \(\mathscr {L}\) 定义 \(\operatorname {Pic}(X/S)=\mathscr {P}\kern  -1pt\mathit {ic}_{X/S}(S)\) 中同一元素的相对 Cartier 除子 \(D\)，也就是由 \hyperref[fga3.v-2-corollary-2.4]{推论~\ref*{fga3.v-2-corollary-2.4}} 可知满足 \(\mathscr {L}(D)\) 与 \(\mathscr {L}\) 在 \(S\) 上局部同构的相对 Cartier 除子，与商层 \(f_*(\mathscr {L})^{\times }/\mathscr {O}_S^\times \) 的截面一一对应。此对应与基变换相容。\emph{同上} 及 \hyperref[fga3.v-bib-6]{[6]} 中的“Künneth”型一般论证表明，\(X/S\) 的固有性及 \(\mathscr {L}\) 在 \(S\) 上的平坦性，蕴含 \(S\) 上存在一个凝聚模 \(\mathscr {Q}\)（在唯一同构意义下确定）及层的同构 \[     f_*(\mathscr {L})     \xrightarrow {\sim } \mathscr {H}\kern  -2.5pt\mathit {om}_{\mathcal {O}_X}(\mathscr {Q},\mathcal {O}_S)   \]；此外，\(\mathscr {Q}\) 的构造与基变换相容。这里，\(f_*(\mathscr {L})^\times \) 表示 \(f_*(\mathscr {L})\) 的集合子层，其在 \(U\) 上的截面是 \(\mathscr {L}\) 在 \(f^{-1}(U)\) 上那些在 \(f^{-1}(U)/U\) 上定义相对 Cartier 除子的截面，即在各 \(X_s\) 上（\(s\in  U\)）诱导出非零因子截面的截面。利用纤维 \(X_s\) 整的假设，这一条件等价于在纤维 \(X_s\) 上诱导的截面不恒为零；用局部同态 \(\mathscr {Q}\to \mathcal {O}_S\) 来说，便是这些同态满射（Nakayama）。这说明，\(f_*(\mathscr {L})^{\times }/\mathcal {O}_S^\times \) 的截面所成的集合，与 \emph{\(\mathscr {Q}\) 的可逆商模} 所成的集合一一对应；或者，按照凝聚模 \(\mathscr {Q}\) 所伴随的射影丛 \(\mathbb {P}(\mathscr {Q})\) 的定义（参见 [\hyperref[Gro1960a]{5}，V，§2]），与 \(\mathbb {P}(\mathscr {Q})\) 在 \(S\) 上的截面所成的集合一一对应。此描述与取逆像相容，因而得到下述定理。\begin{theorem}{4.3}\setnumber{4.3}\label{fga3.v-4-theorem-4.3}\skipsinglepar\par{}设 \(f\colon  X\to  S\) 是几何纤维整的平坦固有态射，\(S\) 局部 Noether，\(\mathscr {L}\) 是 \(X\) 上的可逆模。对每个 \(S\) 上的 \(S'\)，以 \(T(S')\) 表示 \(X\times _S S'/S'\) 上满足 \(\mathscr {L}(D)\) 在 \(S'\) 上局部同构于 \(\mathscr {L}\otimes _{\mathcal {O}_S}\mathcal {O}_{S'}\) 的相对除子 \(D\) 所成的集合（即要求 \(\mathscr {L}(D)\) 与 \(\mathscr {L}\otimes _{\mathcal {O}_S}\mathcal {O}_{S'}\) 定义 \(\operatorname {Pic}(X\times _S S'/S')\) 的同一元素）。那么，\(S\) 上存在一个在唯一同构意义下确定的凝聚模 \(\mathscr {Q}\)，使得函子 \(T\) 由射影丛 \(\mathbb {P}(\mathscr {Q})\) 表示。\end{theorem}\begin{corollary}{4.4}\setnumber{4.4}\label{fga3.v-4-corollary-4.4}\skipsinglepar\par{}若假设 \(f\) 射影，则函子性同态 \(\mathscr {D}\kern  -1pt\mathit {iv}_{X/S}\to \mathscr {P}\kern  -1pt\mathit {ic}_{X/S}\) 可由射影态射表示。\end{corollary}\par{}若 \(X\) 在 \(S\) 上有截面（分别地，局部有截面），则由 \hyperref[fga3.v-4-theorem-4.3]{定理~\ref*{fga3.v-4-theorem-4.3}} 及 \hyperref[fga3.v-2-proposition-2.1]{命题~\ref*{fga3.v-2-proposition-2.1}}，上述同态可由射影丛（分别地，局部射影丛）表示。当 \(f\) 拟射影时，利用 \(X\) 在 \(S\) 上的有限平坦局部拟截面，通过下降方法很容易归约到上述情形。\begin{remark}{4.5}\setnumber{4.5}\label{fga3.v-4-remark-4.5}\skipsinglepar\par{}在 \hyperref[fga3.v-4-theorem-4.3]{定理~\ref*{fga3.v-4-theorem-4.3}} 的条件下，\(S\) 上的模 \(\mathscr {Q}\) 一般并不局部自由；这是因为 \(\mathscr {Q}\) 的约化纤维的维数，即当 \(s\in  S\) 变化时 \(\operatorname {H}^0(X_s,\mathscr {L}_s)\) 的维数，可以 \emph{跳跃}。给定局部 Noether 预概形 \(S\) 上的凝聚模 \(\mathscr {Q}\)，容易证明：对给定的任意 \(s\in  S\)，\(\mathscr {Q}\) 在 \(s\) 处自由，当且仅当 \(\mathbb {P}(\mathscr {Q})\) 在 \(s\) 上方各点处在 \(S\) 上平坦（此时它甚至在 \(s\) 上方各点处在 \(S\) 上是光滑的）。当这一条件满足，且 \(\mathscr {Q}\) 如上由 \(\mathscr {L}\) 定义时，这也等价于直像 \(f_*(\mathscr {L})\) 的构造 \emph{“与 \(s\) 的某个邻域上的基变换交换”}，或蕴含 \(f_*(\mathscr {L})_s\to \operatorname {H}^0(X_s,\mathscr {L}_s)\) 满射。例如，若 \(\operatorname {H}^1(X_s,\mathscr {L}_s)=0\)，就属于这一情形。只要所讨论的预概形存在，这些判据尤其适用于泛情形 \(\underline {\operatorname {Div}}_{X/S}\to \underline {\operatorname {Pic}}_{X/S}\)，并给出此态射在 \(\underline {\operatorname {Div}}_{X/S}\) 的给定点处光滑的充要条件（分别地，充分条件）。\end{remark}\chapter{主要存在定理的证明}\thispagestyle{fancy}\label{fga3.v-5}\par{}在 \hyperref[fga3.v-3-theorem-3.1]{定理~\ref*{fga3.v-3-theorem-3.1}} 的条件下，选取一个在 \(X/S\) 上甚丰沛的模 \(\mathcal {O}_X(1)\)，以 \(\xi \) 表示 \(\operatorname {Pic}(X/S)\) 中的相应元素。为简便起见，令 \(\mathscr {P}(S')=\operatorname {Pic}(X\times _S S'/S')\)，并假设 \(X/S\) 有截面。令 \(\mathscr {P}^+(S')\) 为 \(\mathscr {P}(S')\) 的如下子集：其元素是 \(X\times _S S'\) 上满足 \[     \begin {aligned}       \operatorname {R}^i f'_*(\mathscr {L}(n)) &= 0       \qquad \text {当 }i>0\text { 且任意 }n\geqslant 0     \\f'_*(\mathscr {L}(n)) &\neq 0       \qquad \text {对所有 }n\geqslant 0.     \end {aligned}   \] 的可逆模 \(\mathscr {L}\) 的类。这些条件在基变换下稳定，因而定义了 \(\mathscr {P}\) 的子函子 \(\mathscr {P}^+\)；它显然在由 \(\xi \) 给出的平移下稳定。利用 Serre 的“定理 A 和 B”（[\hyperref[GD1960]{2}，III，§2]）及一般理论（[\hyperref[Gro1960a]{5}，IV，5]），容易看出，\(\mathscr {P}\) 可表当且仅当 \(\mathscr {P}^+\) 可表；于是，\(\mathscr {P}^+\) 将由表示 \(\mathscr {P}\) 的预概形 \(\underline {\operatorname {Pic}}_{X/S}\) 中的一个开集 \(U\) 表示，而后者将是开集 \(U-n\xi \) 的递增并。\par{}为简便起见，令 \(\mathscr {D}=\mathscr {D}\kern  -1pt\mathit {iv}_{X/S}\)，并令 \(\mathscr {D}^+\) 为 \(\mathscr {P}^+\) 在典范态射 \(\mathscr {D}\to \mathscr {P}\) 下的逆像。于是有态射 \begin{eqn}{+}\setnumber{+}\label{fga3.v-5-equation-plus}\skipsinglepar\[       \mathscr {D}^+       \to  \mathscr {P}^+     \tag{+}     \]\end{eqn}；我们已经知道，\(\mathscr {D}^+\) 可由预概形 \(D=\underline {\operatorname {Div}}_{X/S}\) 中的开集 \(D^+\) 表示（其存在性由\hyperref[fga3.v-4-proposition-4.1]{命题 4.1} 保证）。由\hyperref[fga3.v-4-theorem-4.3]{定理 4.3} 容易得到，态射 (+) 可由光滑、射影且满射的态射表示（更确切地说，可由处处非零的\emph{局部自由模所伴随的射影丛}表示）。这是因为，若 \(X\times _S S'\) 上的 \(\mathscr {L}\) 满足本节开头的条件，则 \(f'_*(\mathscr {L})\) 是非零的 \emph{局部自由} 模，其构造与基变换交换；采用 \hyperref[fga3.v-4-theorem-4.3]{定理~\ref*{fga3.v-4-theorem-4.3}} 的记号，此时 \(\mathscr {Q}\) 同构于 \(f'_*(\mathscr {L})\) 的对偶。因此，利用一般判据（[\hyperref[Gro1960a]{5}，IV，4.7]），便可断定 \(P^+\) 可表。在 \emph{同上} 中，取 \(\mathcal {S}\) 为所有忠实平坦拟紧预概形态射所成的集合（由 \hyperref[fga3.i-b.1]{[4]，I，B}，它们确为有效满态射）。正如我们刚才所见，\emph{同上} 的条件 (a)，即 \hyperref[fga3.v-5-equation-plus]{式~\ref*{fga3.v-5-equation-plus}} 可由属于 \(\mathcal {S}\) 的态射表示，已得到满足；条件 (b) 是说函子 \(\mathscr {P}^+\) 与忠实平坦拟紧下降相容，这一点是直接的。只须再证明 \emph{同上} 的条件 (c)：预概形 \(D^+\) 中由 \hyperref[fga3.v-5-equation-plus]{式~\ref*{fga3.v-5-equation-plus}} 内的 \(\mathcal {S}\)-可表态射诱导的等价关系 \(R\) 是 \(\mathcal {S}\)-有效的，即它有效且 \(D^+\to  D^+/R\) 属于 \(\mathcal {S}\)。为此首先注意，\(D^+\) 中对应于虚 Hilbert 多项式 \(Q\in \mathbb {Q}[t]\) 的开集 \({D^+}^Q\) 在 \(R\) 下稳定（因为 \(R\) 的纤维连通）；于是归约为对每个 \(Q\) 证明 \({D^+}^Q\) 上诱导的等价关系 \(R^Q\) 是 \(\mathcal {S}\)-有效的。但现在 \({D^+}^Q\) 拟射影，且等价关系 \(R^Q\) 射影且平坦。因此，\hyperref[fga3.iii-6-theorem-6.1]{[4]，III，6.1} 的假设得到满足，它给出所需结论。\par{}在 \(X/S\) 未必有截面的一般情形中，利用下降技术很容易归约到上述情形；也可以对 \(\mathscr {P}^+\) 的定义作必要的修改，再重复上述证明。\begin{remarks}{5.1}\setnumber{5.1}\label{fga3.v-5-remarks-5.1}\skipsinglepar\par{}我们采用的方法，本质上是 Matsusaka 用来射影地构造 Picard 簇的方法。我们从 \hyperref[fga3.iii]{[4]，III} 引用的、使我们能够取有效商的结果，也很容易由 Hilbert 概形的存在定理推出（例如参见 \hyperref[fga3.v-bib-6]{[6]}）。（在经典理论中，这些商是利用 Chow 坐标构造的。）注意，\(\underline {\operatorname {Pic}}_{X/S}\) 中开集 \(\underline {\operatorname {Pic}}_{X/S}^+\) 的构造，以及按照定义所讨论的可逆模的 \emph{除子} 的 Hilbert 多项式、将其分解为在 \(S\) 上拟射影的开集 \({\underline {\operatorname {Pic}}_{X/S}^+}^Q\) 的分解，都与基变换相容（这使我们能够使用下降技术）。\end{remarks}\begin{remarks}{5.2}\setnumber{5.2}\label{fga3.v-5-remarks-5.2}\skipsinglepar\par{}只要 \(f\colon  X\to  S\) 固有且平坦，且同态 \(k(s)\to \operatorname {H}^0(X_s,\mathcal {O}_{X_s})\)（\(s\in  S\)）都是同构，就有 \(\underline {\operatorname {Pic}}_{X/S}\) 存在，这并非不可能（后一条件还蕴含 \(\mathcal {O}_S\xrightarrow {\sim } f_*(\mathcal {O}_X)\)，且在任意基变换 \(S'\to  S\) 后仍然成立）。至少，在解析空间的范畴中，若进一步假定 \(f\) 射影，就可以用 [\hyperref[Gro1960a]{5}，IX，3.1] 中阐明的一种不同的方法来证明这一点（若我没有记错，这是 Chow 的方法）。在这种方法中，不再对除子概形的开集取一个 \emph{固有} 且平坦的等价关系的商，而是在 \(X\) 到 \(\mathbb {P}_S^r\) 的浸入概形中对射影群取商。利用 Mumford 关于对射影群 \hyperref[Mum1961]{[8]} 取商的结果，这种方法很可能可以推广到概形的情形；目前尚无成文的证明，除了 \(X\) 在 \(S\) 上有“许多局部截面”的情形，例如 \(X\) 在剩余域代数闭的完备局部环上可分的情形。原则上，这种方法适用范围更广，因为它也给出非分离 Picard 预概形的存在性，而在这种情形中，第一种方法必然失效。（从技术上说，困难在于：当 \(f\) 的几何纤维并非整的时，\hyperref[fga3.v-4-theorem-4.3]{定理~\ref*{fga3.v-4-theorem-4.3}} 中考虑的函子不再由射影丛 \(\mathbb {P}(\mathscr {Q})\) 本身表示，而由它的一个 \emph{开集} 表示；这就引出了对平坦但不固有的等价关系取商这一微妙问题。）\par{}\emph{[编者]} 正如我在下一次报告开头所指出的，这里提出的存在性猜想是错误的；但 Mumford 用他的方法证明了一个稍弱的定理。\end{remarks}\begin{remark}{5.3}\setnumber{5.3}\label{fga3.v-5-remark-5.3}\skipsinglepar\par{}注意，这里的证明既不需要预先构造曲线或曲线族的 Jacobian，也不需要 Abel 簇或 Abel 概形的理论。因此，它与沿 Weil 所勾画的道路展开的传统论述——例如 Lang 的书 \hyperref[Lan1959]{[7]} 或 Chevalley 的文章 \hyperref[Che1960]{[1]}——有本质区别。即使对于代数闭域（譬如复数域）上非奇异曲线的 Jacobian，这里给出的构造也是唯一已知的、具有我们在 \cref{fga3.v-1} 中用作定义的极强性质的构造（这些性质本质上是 Chevalley 的性质，但允许“参数簇”带有幂零元）。Picard 概形的构造理应先于而非后于 Abel 簇理论，这一点先验地就很清楚：一般的 Picard 概形既不是 Abel 概形，也不能归约为 Abel 概形。代数闭域上奇异曲线的情形已经表明这一点；这里出现 Rosenlicht 的“广义 Jacobian”，它们并非 Abel 簇。此外，一旦有了一般的 Picard 概形理论，Abel 簇、更一般的 Abel 概形的理论就会大为简化。特别地，Abel 概形的对偶理论，尤其是 Cartier 型结果，此时几乎成为形式的推论（例如参见 \hyperref[fga3.v-bib-10]{[10]}）。\end{remark}\begin{remarks}{5.4}\setnumber{5.4}\label{fga3.v-5-remarks-5.4}\skipsinglepar\par{}对于一条退化为奇异曲线的曲线的 Jacobian，Igusa 的“相容性原理”只有被理解为相对曲线概形 \(X/S\) 的 Picard 概形的存在定理时，才能得到恰当理解；这里的曲线在 \(S\) 上未必光滑。因此，当特殊纤维曲线整时（即用经典术语说，不可约且重数为 \(1\) 时），这是主要存在定理 \hyperref[fga3.v-3-theorem-3.1]{定理~\ref*{fga3.v-3-theorem-3.1}} 的一个特例。注意，目前对于可约的特殊纤维曲线（即使各分支的重数均为 \(1\)，即特殊纤维曲线在剩余域上可分），已知存在定理仍未涵盖，除非底环是剩余域代数闭的 \emph{完备} 离散赋值环；参见 \hyperref[fga3.v-5-remarks-5.2]{注~\ref*{fga3.v-5-remarks-5.2}}。在概形理论中，要对亏格为 \(g\) 的曲线的模概形作 Baily–Satake“紧化”的几何构造，必然会遇到这一存在性问题。（目前，由 Igusa 的工作，只在 \(g=1\) 时知道这种紧化。）\end{remarks}\chapter{相对存在定理}\thispagestyle{fancy}\label{fga3.v-6}\par{}这里将概述一些有用的情形：某些 Picard 概形的存在性蕴含另一些 Picard 概形的存在性。这使我们能由主要存在定理 \hyperref[fga3.v-3-theorem-3.1]{定理~\ref*{fga3.v-3-theorem-3.1}} 推出各种其他存在定理。\begin{proposition}{6.1}\setnumber{6.1}\label{fga3.v-6-proposition-6.1}\skipsinglepar\par{}设 \(f\colon  X\to  S\) 是平坦射影态射，使其 Stein 分解 \(f=f''f'\) 中的态射 \(f'\colon  X\to  S'\) 为 \emph{平坦} 且几何纤维整（从而满足 \hyperref[fga3.v-3-theorem-3.1]{定理~\ref*{fga3.v-3-theorem-3.1}} 的假设），并使有限态射 \(f''\colon  S'\to S\) 平坦。那么 \(\underline {\operatorname {Pic}}_{X/S}\) 存在，而且（用 \hyperref[fga3.ii-c.2]{[4]，II，C.2} 中引入的记号）有典范同构 \[       \underline {\operatorname {Pic}}_{X/S}       \xrightarrow {\sim } \prod _{S'/S} \underline {\operatorname {Pic}}_{X/S'}.     \]。\begin{proof}\par{}为证明这一点，先建立函子同构 \[         \mathscr {P}\kern  -1pt\mathit {ic}_{X/S}         \xrightarrow {\sim } \prod _{S'/S} \mathscr {P}\kern  -1pt\mathit {ic}_{X/S'}       \]，再利用 \hyperref[fga3.v-3-theorem-3.1]{定理~\ref*{fga3.v-3-theorem-3.1}}，由此得到 \(\mathscr {P}\kern  -1pt\mathit {ic}_{X/S'}\) 可表。为保证 \(\prod _{S'/S}\underline {\operatorname {Pic}}_{X/S'}\) 存在，使用 \cref{fga3.v-3} 中所说明的 \(\underline {\operatorname {Pic}}_{X/S'}\) 的结构（它蕴含 \(\underline {\operatorname {Pic}}_{X/S'}\) 在 \(S\) 上的任一纤维的有限子集都包含在某个仿射开集中）。\end{proof}\end{proposition}\par{}例如，若概形 \(X\) 是 \(S\) 上满足 \hyperref[fga3.v-3-theorem-3.1]{定理~\ref*{fga3.v-3-theorem-3.1}} 的条件的概形 \(X_i\) 的 \emph{和}，则 \hyperref[fga3.v-6-proposition-6.1]{命题~\ref*{fga3.v-6-proposition-6.1}} 的陈述归结为平凡的陈述 \[     \underline {\operatorname {Pic}}_{X/S}     \xrightarrow {\sim } \prod _i \underline {\operatorname {Pic}}_{X_i/S}.   \]。\begin{corollary}{6.2}\setnumber{6.2}\label{fga3.v-6-corollary-6.2}\skipsinglepar\par{}设 \(f\colon  X\to  S\) 是几何纤维局部整的射影平坦态射（例如射影且正规的态射）。那么 \(\underline {\operatorname {Pic}}_{X/S}\) 存在。\begin{proof}\par{}这时，\(S'\) 是 \(S\) 的 étale 覆盖（只要 \(f\) 可分，即平坦且几何纤维约化，这一点就成立）；由 \(\underline {\operatorname {Pic}}_{X/S'}\) 的类似结构可见，\cref{fga3.v-3} 中为 \(\underline {\operatorname {Pic}}_{X/S}\) 陈述的结构定理仍然成立。\end{proof}\end{corollary}\par{}应用下降过程，可以得到一个相对存在定理；其适用范围取决于 \hyperref[fga3.i-a.2.c]{[4]，I，A(c)} 中提出的非平坦下降问题能否解决。这里我们仅说明一个特殊情形：\begin{proposition}{6.3}\setnumber{6.3}\label{fga3.v-6-proposition-6.3}\skipsinglepar\par{}设 \(f\colon  X\to  S\) 是固有态射，\(X_1\) 与 \(X_2\) 是 \(X\) 的子预概形，均在 \(S\) 上平坦，分别由凝聚理想 \(\mathscr {J}_1\) 与 \(\mathscr {J}_2\) 定义，使得 \(\mathscr {J}_1\cap \mathscr {J}_2=(0)\)，且 \(\mathcal {O}_X/(\mathscr {J}_1+\mathscr {J}_2)\) 在 \(S\) 上平坦（即 \(X\) 中作为 \(X_1\) 与 \(X_2\) 的上确界的子预概形就是 \(X\) 本身，而它们的下确界 \(Z\) 在 \(X\) 上平坦）。进一步假设，对所有 \(s\in  S\)，当 \(i=1,2\) 时同态 \(k(s)\to \operatorname {H}^0(X_{i_s},\mathcal {O}_{X_{i_s}})\) 均为双射。那么自然的函子同态 \[       \mathscr {P}\kern  -1pt\mathit {ic}_{X/S}       \to  \mathscr {P}\kern  -1pt\mathit {ic}_{X_1/S} \times  \mathscr {P}\kern  -1pt\mathit {ic}_{X_2/S}     \] 可由仿射态射表示。因此，若 \(\underline {\operatorname {Pic}}_{X_1/S}\) 与 \(\underline {\operatorname {Pic}}_{X_2/S}\) 存在，则 \(\underline {\operatorname {Pic}}_{X/S}\) 也存在，且典范态射 \[       \underline {\operatorname {Pic}}_{X/S}       \to  \underline {\operatorname {Pic}}_{X_1/S} \times  \underline {\operatorname {Pic}}_{X_2/S}     \] 是仿射的。\begin{proof}\par{}借助忠实平坦下降，可归约到 \(Z\) 在 \(S\) 上有截面的情形。这个截面在 \(S\) 上定义 \(X\)、\(X_1\) 及 \(X_2\) 的截面，并使我们能如 \hyperref[fga3.v-2-remark-2.5]{注~\ref*{fga3.v-2-remark-2.5}} 中所述消除所考虑结构的自同构。此时，证明在于注意到：给定 \(X\) 上的“刚性化”可逆模 \(\mathscr {L}\)，等价于给定三元组 \((\mathscr {L}_1,\mathscr {L}_2,u)\)，其中 \(\mathscr {L}_i\) 是 \(X_i\) 上的“刚性化”模，而 \(u\) 是从 \(\mathscr {L}_1|Z\) 到 \(\mathscr {L}_2|Z\) 的、与刚性化相容的同构。只须再验证，当 \(\mathscr {L}_1\) 与 \(\mathscr {L}_2\) 固定时，给定 \(u\) 等价于给定 \(S\) 上某个适当的、在 \(S\) 上仿射的概形的一个截面；这一点很容易。\end{proof}\end{proposition}\par{}由 \hyperref[fga3.v-6-proposition-6.3]{命题~\ref*{fga3.v-6-proposition-6.3}} 容易得出：\begin{corollary}{6.4}\setnumber{6.4}\label{fga3.v-6-corollary-6.4}\skipsinglepar\par{}设 \(X\) 是域 \(k\) 上的固有可分概形，\(X_i\) 是 \(X\) 的不可约分支。若各 \(\underline {\operatorname {Pic}}_{X_i/k}\) 存在，则 \(\underline {\operatorname {Pic}}_{X/k}\) 也存在，且典范态射 \[       \underline {\operatorname {Pic}}_{X/k}       \to  \prod _i \underline {\operatorname {Pic}}_{X_i/k}     \] 是仿射的。\end{corollary}\par{}与 \hyperref[fga3.v-6-corollary-6.2]{推论~\ref*{fga3.v-6-corollary-6.2}} 结合，这表明，例如，只要 \(X\) 是域 \(k\) 上的射影可分概形，\(\underline {\operatorname {Pic}}_{X/k}\) 就存在。即使 \(X\) 在 \(k\) 上不再可分，利用 Oort 的论证 \hyperref[Oor1962]{[11]}，同样也有归约结果。这一方法也适用于任意底上的概形（例如，在下一次报告中证明 \hyperref[fga3.v-3-remark-3.3]{注~\ref*{fga3.v-3-remark-3.3}} 所陈述的有限性结果时，这种情形就很有用）。为避免过于技术性的陈述，我们只考虑底为域的情形：\begin{proposition}{6.5}\setnumber{6.5}\label{fga3.v-6-proposition-6.5}\skipsinglepar\par{}设 \(X\) 是域 \(k\) 上的固有概形，\(X_0\) 是一个具有相同底层集合的子概形（因而由 \(X\) 上的幂零理想定义）。那么函子性态射 \(\mathscr {P}\kern  -1pt\mathit {ic}_{X/k}\to \mathscr {P}\kern  -1pt\mathit {ic}_{X_0/k}\) 可由仿射态射表示。特别地，若 \(\underline {\operatorname {Pic}}_{X_0/k}\) 存在，则 \(\underline {\operatorname {Pic}}_{X/k}\) 也存在，且态射 \(\underline {\operatorname {Pic}}_{X/k}\to \underline {\operatorname {Pic}}_{X_0/k}\) 是仿射的。\end{proposition}\par{}将它与 \hyperref[fga3.v-6-corollary-6.4]{推论~\ref*{fga3.v-6-corollary-6.4}} 结合，容易得出：\begin{corollary}{6.6}\setnumber{6.6}\label{fga3.v-6-corollary-6.6}\skipsinglepar\par{}设 \(X\) 是域 \(k\) 上的射影预概形。那么 \(\underline {\operatorname {Pic}}_{X/k}\) 存在。\end{corollary}\begin{remark}{6.6}\setnumber{6.6}\label{fga3.v-6-remark-6.6}\skipsinglepar\par{}对于域 \(k\) 上的每个 \emph{固有} 概形 \(X\)，\(\underline {\operatorname {Pic}}_{X/k}\) 应当存在，这是极为可信的。由上述结果，这个问题可以归约到 \(k\) 代数闭、\(X\) 整的情形。此时知道，存在一个在 \(k\) 上 \emph{射影} 的整概形 \(X'\) 及支配态射 \(g\colon  X'\to  X\)（Chow 引理）。因此，只须证明相应的函子性态射 \(\mathscr {P}\kern  -1pt\mathit {ic}_{X/k}\to \mathscr {P}\kern  -1pt\mathit {ic}_{X'/k}\) 可表（很可能还可由仿射态射表示），因为已经知道 \(\mathscr {P}\kern  -1pt\mathit {ic}_{X'/k}\) 可表。这引出了目前尚无法回答的非平坦下降问题。注意，若只考虑函子 \(\mathscr {P}\kern  -1pt\mathit {ic}_{X/k}\) 在约化预概形上的限制（其中 \(X\) 是代数闭域 \(k\) 上的固有整概形），则确实得到一个可表函子：在 \(X\) 正规的情形中，Chevalley 在 \hyperref[Che1960]{[1]} 中证明了这一点；在一般情形中，Seshadri 在 \hyperref[Ses1962]{[12]} 中用下降方法证明了这一点。不过，按我们的记号，这些作者所构造的概形并非 \(\underline {\operatorname {Pic}}_{X/k}\)，而是 \((\underline {\operatorname {Pic}}_{X/k})_\mathrm {red}\)，即与 \(\underline {\operatorname {Pic}}_{X/k}\) 对应的约化概形。\par{}\emph{[编者]} 正如我们在 \hyperref[fga3.vi]{下一次报告} 开头所指出的，Murre 最近已经对这里提出的问题给出了肯定回答。\end{remark}\begin{remark}{6.7}\setnumber{6.7}\label{fga3.v-6-remark-6.7}\skipsinglepar\par{}更一般地，设 \(f\colon  X'\to  X\) 是 \(k\) 上固有预概形之间的满射态射。考虑非平坦下降，就会猜想 \(\underline {\operatorname {Pic}}_{X/k}\to \underline {\operatorname {Pic}}_{X'/k}\) 是 \emph{仿射} 态射；特别地，这将蕴含（除去单位元的连通分支后）相应的 Néron–Severi 群同态模去挠元后为单射。当 \(X\) 和 \(X'\) 非奇异时，利用相交理论可验证这一点。在其他情形中，答案似乎尚未知道。\par{}\emph{[编者]} 这里提出的问题已得到肯定回答（参见下一次报告的评论的最后一段）。\end{remark}\begin{remark}{6.8}\setnumber{6.8}\label{fga3.v-6-remark-6.8}\skipsinglepar\par{}考虑带有幂零元的代数概形的 Picard 概形，对于许多问题是有用的，甚至不可或缺；这或许出乎意料。例如，若 \(X\) 是 \(k\) 上的光滑射影概形，\(Y\) 是一个超平面截面，就可以考虑 \(Y\) 的各阶“无穷小邻域”\(X_n\) 及其 Picard 概形 \(\underline {\operatorname {Pic}}_{X_n/k}\)。当 \(X\) 不可约且维数为 \(\geqslant 4\)（分别地，\(\geqslant 3\)）时，典范态射 \[       \underline {\operatorname {Pic}}_{X/k}       \to \underline {\operatorname {Pic}}_{X_n/k}       \qquad \text {对充分大的 }n     \] 是同构（分别地，在 Néron–Severi 群的挠子群的逆像之间诱导同构）；下一次报告对 Picard 概形作定性研究时将用到这一结果。同样，考虑某些带有幂零元的曲线的 Picard 概形，并利用形式几何的基本定理 \hyperref[fga3.v-bib-3]{[3]}，就能在等特征情形中证明 Mumford 的一个猜想：对于每个维数为 \(2\) 的完备正规 Noether 局部环 \(A\)，\(A\) 的除子类群可以看成剩余域 \(k\) 上的一个代数群 \(G\) 的 \(k\)-有理点所成的集合（一旦选定 \(A\) 中的代表域，\(G\) 就典范地确定）。当 \(A\) 的维数任意时，似乎应当存在一个 \(k\) 上的代数 pro-群，起到上述 \(G\) 的作用；当可以对 \(\operatorname {Spec}(A)\) 作“解奇点”时，它可构造成 \(k\) 上适当的（带有幂零元的）射影概形的 Picard 概形的投射极限。\end{remark}
\section*{原报告参考文献}
\begin{enumerate}
\item[{[1]}]\setnumber{1}\label{fga3.v-bib-1} Chevalley (Claude). \emph{Sur la théorie de Picard}. Amer. J. of Math., t.~82, 1960, pp.~435--490.
\item[{[2]}]\setnumber{2}\label{fga3.v-bib-2} Grothendieck (A.) et Dieudonné (J.). \emph{Éléments de géométrie algébrique. Chapitre I et suivants}. Paris, Presses universitaires de France, 1960--1961 (Institut des Hautes Études Scientifiques, Publications mathématiques 4, 8, \ldots).
\item[{[3]}]\setnumber{3}\label{fga3.v-bib-3} Grothendieck (Alexander). \emph{Géométrie formelle et géométrie algébrique}. Séminaire Bourbaki, t.~11, 1958/59, n\textsuperscript{o}~182, 28~页。
\item[{[4]}]\setnumber{4}\label{fga3.v-bib-4} Grothendieck (Alexander). \emph{Techniques de descente et théorèmes d'existence en géométrie algébrique, I--IV}. Séminaire Bourbaki, t.~12, 1959/60, n\textsuperscript{o}~190, 29~页；n\textsuperscript{o}~195, 22~页；t.~13, 1960/61, n\textsuperscript{o}~212, 20~页；n\textsuperscript{o}~221, 28~页。
\item[{[5]}]\setnumber{5}\label{fga3.v-bib-5} Grothendieck (Alexander). \emph{Techniques de construction en géométrie analytique, I--X}. Séminaire Cartan, t.~13, 1960/61: \emph{Déformation des structures analytiques complexes}, n\textsuperscript{o}~7--17.
\item[{[6]}]\setnumber{6}\label{fga3.v-bib-6} Grothendieck (Alexander). \emph{Séminaire de géométrie algébrique}，由 Lichtenbaum 撰写，Harvard University, 1961（待出版）。
\item[{[7]}]\setnumber{7}\label{fga3.v-bib-7} Lang (Serge). \emph{Abelian varieties}. New York, Interscience Publishers, 1959 (Interscience Tracts in Pure and Applied Mathematics, 7).
\item[{[8]}]\setnumber{8}\label{fga3.v-bib-8} Mumford (D.). \emph{An elementary theorem in geometric invariant theory}. Bull. Amer. Math. Soc., t.~67, 1961, pp.~483--486.
\item[{[9]}]\setnumber{9}\label{fga3.v-bib-9} Mumford (D.). \emph{The topology of normal singularities of an algebraic surface and a criterion for simplicity}. Paris, Presses universitaires de France, 1961 (Institut des Hautes Études Scientifiques, Publications mathématiques, 9).
\item[{[10]}]\setnumber{10}\label{fga3.v-bib-10} Mumford (D.) and Tate (J.). \emph{Séminaire de géométrie algébrique}. Harvard University, Spring Term 1962（待出版）。
\item[{[11]}]\setnumber{11}\label{fga3.v-bib-11} Oort (F.). \emph{Sur le schéma de Picard}. Bull. Soc. Math. France, t.~90, 1962（待出版）。
\item[{[12]}]\setnumber{12}\label{fga3.v-bib-12} Seshadri (C. S.). 学位论文，待发表于 \emph{Annali di Matematica}。
\end{enumerate}

% END chapters/fga3.v.tex

\clearpage
\setcounter{chapter}{-1}
\renewcommand\fgaAnchorPrefix{fga3.vi}

% BEGIN chapters/fga3.vi.tex
\part{Picard 概形：一般性质}\label{fga3.vi}\thispagestyle{fancy}\begin{original}{}\label{fga3.vi-original-citation}\skipsinglepar\par{}A. Grothendieck．“Technique de descente et théorèmes d'existence en géométrie algébrique, VI: Les schémas de Picard: Propriétés générales”．\emph{Séminaire Bourbaki} \textbf{14}（1961–62），第 236 次报告，1962 年 5 月。\url{http://www.numdam.org/book-part/SB_1961-1962__7__221_0/}\end{original}\setcounter{chapter}{-1}\chapter[{对上次报告（[3]，V）的补充}]{对上次报告（[3]，V）的补充\textsuperscript{(1)}}\thispagestyle{fancy}\label{fga3.vi-0}\par{}\emph{［译校注］原刊标题保留了脚注标记(1)，但本库法文汇编扫描及 Numdam 原刊扫描中均未定位到相应脚注正文。此处保留标记，未以推测文字补造原注。}\par{}对于 \hyperref[fga3.v]{[3]，V} 中提出的 Picard 预概形的存在性问题，已有一些进展：\begin{enumerate}
    
 \item[a.]（Mumford。）若 \(f\colon  X\to  S\) 是射影可分态射（即平坦且纤维可分），则 Picard 预概形 \(\underline {\operatorname {Pic}}_{X/S}\) 存在，这在一般情形中并不成立；即使 \(f\) 的纤维维数为 \(1\)，且 \(S\) 是完备离散赋值环的谱，也是如此。一个反例如下：取 \(S=\operatorname {Spec}\mathbb {R}[{[t]}]\)，取 \(X\) 为 \(\mathbb {P}_S^2\) 中由方程 \(x^2+y^2=tz^2\) 定义的子概形（齐次变量为 \(x,y,z\)）。它表示一条退化为几何上相交的两条直线的圆锥曲线，但在域 \(\mathbb {R}\) 上的特殊纤维却不可约（它在 \(\mathbb {R}\) 中由方程 \(x^2+y^2=0\) 给出）。容易看出，作 étale 扩张 \(S'\to  S\)（其中 \(S'=\operatorname {Spec}\mathbb {C}[{[t]}]\)）后，\(X'/S'\) 的 Picard 预概形存在，并由此将其明确描述为若干 \(\widetilde {S}\) 的副本之和，其中 \(\widetilde {S}\) 是将 \(S\) 的原点复制无穷多次所得。容易看到，\(\underline {\operatorname {Pic}}_{X'/S'}\) 上相对于 \(S'\to  S\) 的下降数据（这里由 \(S'\) 在 \(S\) 上的 Galois 群 \(\mathbb {Z}/2\mathbb {Z}\) 的作用给出）并不有效，因为这个群置换某些被复制的点（因而有些轨道不包含在任何仿射开集中）。不过，Mumford 已证明：若 \(f\colon  X\to  S\) 是可分射影态射，且对所有 \(s\in  S\)，\(X_s\) 的不可约分支相对于 \(k(s)\) 都几何不可约，则 \(\underline {\operatorname {Pic}}_{X/S}\) 存在。证明基于他关于取商定理的一个精细化；参见 \hyperref[fga3.vi-bib-7]{[7]}。还要注意，即使不对纤维 \(X_s\) 的不可约分支作任何假设，下面将引入的概形 \(\underline {\operatorname {Pic}}_{X/S}^\tau \) 仍有可能总是存在。


    
 \item[b.]（Murre。）若 \(X\) 是一个域上的固有概形，则 \(\underline {\operatorname {Pic}}_{X/k}\) 存在。证明部分地利用了 Chevalley 的证明 \hyperref[Che1960]{[1]}，而根本上依靠 Picard 函子的群结构。

  \end{enumerate}关于 Picard 概形理论的其他一些评注，特别是它与 Abel 概形的关系，建议参阅 \hyperref[fga3.vi-bib-7]{[7]}。最后，本报告一个明显的不足是缺少“等价判据”，即能够比较射影概形及其超平面截面的 Picard 概形的判据。建立这些判据所需的关键定理可见 \hyperref[fga3.vi-bib-5]{[5]}，并须将它们与 Picard 概形的存在定理结合。\chapter{交换群预概形的拓扑性质}\thispagestyle{fancy}\label{fga3.vi-1}\par{}设 \(k\) 是域，\(G\) 是 \(k\) 上的群预概形。单位元 \(e\) 是 \(k\)-有理点，因而必为闭点，立刻可得 \(G\times _k G\) 的对角线闭，从而 \(G\) 是 \emph{分离} 的：\emph{域上的每个群预概形都是分离的}。用 \(G^0\) 表示单位元 \(e\) 的连通分支。因为 \(e\) 是 \(k\)-有理点，\(G^0\) 实际上是 \emph{几何连通} 的，即 \(G\mapsto  G^0\) 与到另一个域的基变换相容。还可得出，\(G^0\) 在乘法下稳定（集合意义下）；若 \(G\) 局部 Noether，则 \(G^0\) 为开集，于是可将 \(G^0\) 视为 \(G\) 的一个 \emph{开子群}。以下假设 \(G\) 在 \(k\) 上局部有限型；则 \(G^0\) 是 \emph{几何不可约且在 \(k\) 上有限型} 的。事实上，可以假设 \(k\) 代数闭，进而假设 \(G\) 约化（此时 \(G_\mathrm {red}\) 是 \(G\) 的子群，因为 \(G_\mathrm {red}\times _k G_\mathrm {red}\) 此时也约化）。于是它在某个非空开集上在 \(k\) 上光滑；平移此开集，便知它处处光滑。但这样一来，\(G\) 局部不可约，所以其不可约分支就是连通分支，因而 \(G^0\) 不可约。现在设 \(U\) 是 \(G^0\) 中 \(e\) 的一个仿射邻域；利用 \(G^0\) 不可约，立刻可知 \(U\cdot  U=G^0\)，这就证明 \(G^0\) 拟紧，因而在 \(k\) 上有限型。\par{}为简单起见，假设 \(G\) 交换。对每个整数 \(n>0\)，以 \(G^{(n)}\) 表示 \(G^0\) 在\(G\) 中的 \(n\) 次幂同态 \(\varphi _n\)下的逆像，于是 \(G^{(n)}\) 是 \(G\) 的开子群。令 \[     \begin {aligned}       G^\tau        &= \bigcup _{n>0} G^{(n)}     \\G^\sigma        &= \bigcup _{(n,p)=1} G^{(n)}     \\G^\rho        &= \bigcup _{h>0} G^{(p^h)}     \end {aligned}   \]，其中 \(p\) 是域 \(k\) 的特征指数。这些都是 \(G\) 的开子群，且满足 \[     \begin {aligned}       G^\sigma  \cap  G^\rho        &= G^0     \\G^\sigma  \cdot  G^\rho        &= G^\tau .     \end {aligned}   \]。\begin{remark}{}\label{fga3.vi-1-remark}\skipsinglepar\par{}可以构造商群概形 \(G/G^0=N\)（参见 \hyperref[fga3.iv]{[3]，IV}），并将 \(G^\tau \)、\(G^\sigma \)、\(G^\rho \) 分别定义为 \(N\) 的挠子群（分别地，其 \(p\)-准素分量，以及其 \(p\)-准素分量的自然补群，即对满足 \(q\neq  p\) 的素数 \(q\) 取各 \(q\)-准素分量之和）在 \(G\) 中的逆像。注意，\(N\) 是 \(k\) 上可分的离散群概形，因此（一旦选定 \(k\) 的代数闭包 \(\overline {k}\)，从而得到 Galois 群 \(\pi \)）可将它等同于一个通常的离散群，其上有 \(\pi \) 通过自同构的作用。挠子群的构造及其分解为各 \(q\)-准素分量，可以用这一方式作显然的解释。若 \(G\) 是 \(k\) 上固有概形 \(X\) 的 Picard 概形，则可以将 \(N\) 称为 \(X\) 在 \(k\) 上的（约化）\emph{Néron–Severi 概形}。若 \(G_\mathrm {red}^0\) 是 \(G^0\) 的群子概形——特别地，当 \(k\) 为完美域，或 \(G^0\) 在 \(k\) 上固有时（例如 \(X\) 几何正规时）即是如此——则还可以引入商 \(N'=G/G_\mathrm {red}^0\)。从特殊化的观点，即当 \(X\) 在代数概形族中变化时看，它的表现往往比 \(N\) 更好。\end{remark}\par{}现在设 \(S\) 是局部 Noether 预概形，\(G\) 是 \(S\) 上的群预概形，且在 \(S\) 上局部有限型。不假设 \(G\) 在 \(S\) 上有限型，也不假设它在 \(S\) 上分离。令 \[     G^0     = \bigcup _{s\in  S} (G_s)^0   \]；若 \(G\) 交换，再令 \[     \begin {aligned}       G^\tau        &= \bigcup _{s\in  S} (G_s)^\tau      \\G^\sigma        &= \bigcup _{s\in  S} (G_s)^\sigma      \\G^\rho        &= \bigcup _{s\in  S} (G_s)^\rho .     \end {aligned}   \]。这些是 \(G\) 的子集，在 \(G\) 的乘法下稳定，但这并不显然蕴含它们可以由 \(G\) 的子群预概形定义。特别地，似乎 \emph{一般并不存在底层集合为 \(G^0\) 的 \(G\) 的子群预概形}。当然，若其中某个集合是开集，赋予诱导结构后，它就是 \(G\) 的开子群预概形。我们将看到，对于 \(G^\tau \)，情形总是如此；因此，从可表函子的观点看，特别有 \emph{从“特殊化”的观点看，数值等价比代数等价表现得更令人满意}。以上所定义的集合的主要一般性质如下：\begin{theorem}{1.1}\setnumber{1.1}\label{fga3.vi-1-theorem-1.1}\skipsinglepar\par{}\(G^0\)、\(G^\tau \)、\(G^\sigma \) 和 \(G^\rho \) 都是 \emph{局部可构造} 的。此外：\begin{enumerate}
      
 \item[i.]\(G^0\) 在 \(S\) 上拟紧。若各 \(G_s^0\) 固有，且 \(G\) 在 \(S\) 上分离，则 \(G^0\) 在 \(S\) 上固有，因而在 \(G\) 中闭。

      
 \item[ii.]\(G^\tau \) 为开集。若 \(G^0\) 为闭集，则 \(G^\tau \) 亦然。

      
 \item[iii.]若 \(G^0\) 为闭集，则只要处于等特征情形，即 \(S\) 的所有剩余域有相同特征，\(G^\sigma \) 也为闭集。若 \(G^0\) 为闭集，且 \(G\to  S\) 在 \(G^\sigma \) 的各点处普遍开（参见下面的 \hyperref[fga3.vi-1-corollary-1.5]{推论~\ref*{fga3.vi-1-corollary-1.5}}），则 \(G^\sigma \to  S\) 普遍开。

      
 \item[iv.]若 \(G^0\) 为闭集，则 \(G^\rho \) 亦然。若处于等特征情形，且对每个满足 \((n,p)=1\) 的整数 \(n>0\)，\(G\) 中的 \(n\) 次幂同态为开映射，则 \(G^\rho \) 为开集。

    \end{enumerate}\end{theorem}\par{}现在给出证明的一些提示。\(G^0\) 局部可构造且在 \(S\) 上拟紧，这一事实包含在以下引理中：\begin{lemma}{1.2}\setnumber{1.2}\label{fga3.vi-1-lemma-1.2}\skipsinglepar\par{}若 \(S\neq \varnothing \)，则 \(S\) 中存在非空开集 \(U\)、\(U\) 上纤维连通的有限型群概形 \(H\)，以及一个群概形同态 \(H\to  G|U\)，后者是像为 \(G^0|U\) 的开浸入。\begin{proof}\par{}为证明此引理，可假设 \(S\) 不可约；设 \(\eta \) 为其泛点。作基变换 \(S'=\operatorname {Spec}(\mathcal {O}_{S,\eta })\to  S\)，便得到局部 Artin 环 \(\mathcal {O}_{S,s}=A\) 上的群概形 \(G'\)，其中有一个 \emph{开} 群子概形 \({G'}^0\)，它 \emph{在 \(A\) 上有限型}，如上所述。因此，它来自 \(\eta \) 的某个开邻域 \(U\) 上的有限型群概形 \(H\)，而典范浸入 \({G'}^0\to  G'\) 来自一个开浸入 \(H\to  G|U\)；当 \(U\) 足够小时，后者是群概形同态。若 \(U\) 足够小，\(H\) 的纤维连通；若 \(U\) 足够小，这些纤维均有相同维数，即 \(G\) 的纤维的维数。因此，当 \(U\) 足够小时，对所有 \(s\in  U\)，\(H_s\) 在 \(G_s\) 中的像恰为 \(G_s^0\)。引理得证。\end{proof}\end{lemma}\par{}\hyperref[fga3.vi-1-theorem-1.1]{定理~\ref*{fga3.vi-1-theorem-1.1}} 的 (i) 中的第二个断言包含在下列引理中（将此引理用于 \(G\) 中 \(G^0\) 的一个拟紧开邻域）：\begin{lemma}{1.3}\setnumber{1.3}\label{fga3.vi-1-lemma-1.3}\skipsinglepar\par{}设 \(X\) 是 \(S\) 上的有限型分离预概形，\(S\) 局部 Noether；设 \(g\) 是 \(X\) 在 \(S\) 上的截面，\(X^0\) 是各 \(g(s)\) 在各 \(X_s\) 中的连通分支之并。设 \(s\in  S\) 使得 \(X_s^0\) 在 \(k(s)\) 上固有。那么，\(s\) 有一个开邻域 \(U\)，使 \(X^0|U\) 在 \(U\) 上固有，从而也在 \(X|U\) 中闭。\begin{proof}\par{}对底作忠实平坦下降，可归约到 \(S\) 是完备局部环的谱、\(s\) 为其闭点的情形。应用 [\hyperref[GD1960]{2}，III，§5.5.1]，可见 \(X\) 分解为两个不交开集 \(X'\) 和 \(X''\) 的和，其中 \(X'\) 在 \(S\) 上固有，并且 \(X'_s=X_s^0\)。这使我们可归约到 \(X=X'\) 的情形，即 \(X\) 在 \(S\) 上固有的情形。此时可运用标准证明，利用子集固有性的赋值判据（[\hyperref[GD1960]{2}，II] 中遗漏了这一判据）。\end{proof}\end{lemma}\par{}现在证明 \(G^\tau \) 是开集；考虑到 \(G^\tau \) 的构造（与 \(G^0\)、\(G^\sigma \)、\(G^\rho \) 的构造一样）与基扩张交换，这等价于证明：对于 \(G\) 在 \(S\) 上的每个截面 \(g\)，\(g^{-1}(G^\tau )\) 是开集。这要求证明两点：\begin{enumerate}
    
 \item[a.]设 \(y\in  S\) 满足 \(g(y)\in  G^\tau \)。那么，对 \(y\) 的所有邻近点 \(y'\in \overline {y}\)，都有 \(g(y')\in  G^\tau \)。

    
 \item[b.]设 \(y'\in  S\) 满足 \(g(y')\in  G^\tau \)。那么，对 \(y'\) 的每个泛化 \(y\)，都有 \(g(y)\in  G^\tau \)。

  \end{enumerate}对于 (a)，注意存在整数 \(n>0\) 使 \(g^n(y)\in  G^0\)；由于 \(G^0\) 可构造，\((g^n)^{-1}G^0\) 也可构造；于是对 \(y\) 的所有邻近点 \(y'\in \overline {y}\)，都有 \(g^n(y')\in  G^0\)。对于 (b)，注意各 \(g(y')^n=g^n(y')\) 都留在 \(G_{y'}\) 的某个拟紧开集中（因为它们包含在模 \(G_{y'}^0\) 的有限多个类中）；因此 \(G\) 中存在一个拟紧开集 \(U\)，包含各 \(g^n(y')\)，从而也包含其泛化 \(g^n(y)\)。所以，\(g(y)\) 的各次幂都留在 \(G_y\) 的某个拟紧开集中，容易由此推出 \(g(y)\in  G_y^\tau \)。\par{}假设 \(G^0\) 为闭集；我们将证明 \(G^\tau \) 也是闭集。已知 \(G^\tau \) 为开集，因而局部可构造，所以只须证明它在特殊化下稳定。这是因为它是一族闭子集的并，这些闭子集就是闭子集 \(G^0\) 在各 \(n\) 次幂同态 \(\varphi _n\) 下的逆像。\par{}同一论证表明，只要 \(G^0\) 为闭集，\(G^\sigma \) 和 \(G^\rho \) 也都是闭集（对前者还须附加等特征假设）；只需先证明 \(G^\sigma \) 和 \(G^\rho \) 局部可构造。但对于 \(x\in  G^\tau \)，令 \(\nu (x)\) 为使 \(n\) 次幂同态 \(\varphi _n\) 将 \(x\) 映入 \(G\) 的最小整数 \(n>0\)。则 \(G^\sigma \)（分别地，\(G^\rho \)）由满足 \(\nu (x)\) 与 \(p\) 互素（分别地，为 \(p\) 的某个幂）的那些 \(x\in  X\) 组成。因而，所断言的可构造性由下面更精确的引理得到：\begin{lemma}{1.4}\setnumber{1.4}\label{fga3.vi-1-lemma-1.4}\skipsinglepar\par{}\(G^\tau \) 上的函数 \(\nu \) 局部可构造。\begin{proof}\par{}这就是说，对每个整数 \(n>0\)，满足 \(\nu (x)=n\) 的 \(x\in  G^\tau \) 所成的集合局部可构造。但这个集合是 \(\varphi _n^{-1}(G^0)\) 与各 \(\varphi _d^{-1}(G^0)\) 之并的差集，其中 \(d\) 遍历 \(n\) 的真因子。由于 \(G^0\) 局部可构造，各 \(\varphi _i^{-1}(G^0)\) 也都局部可构造，故上述差集亦然。\end{proof}\end{lemma}\par{}假设 \(G^0\) 闭，且 \(G\to  S\) 在 \(G^\sigma \) 的各点处普遍开；我们将证明 \(G^\sigma \to  S\) 为开映射，即将任意 \(x\in  G^\sigma \) 在 \(G^\sigma \) 中的邻域映为 \(y=f(x)\) 的邻域。由于 \(G^\sigma \) 局部可构造，只须证明：对 \(y\) 的每个泛化 \(y'\)，在 \(y'\) 上方存在 \(x\) 在 \(G\) 中的一个泛化 \(x\)'。作基变换，可归约到 \(S\) 是离散赋值环的谱，\(y\) 和 \(y'\) 分别为闭点和泛点的情形。利用 \(G\to  S\) 在 \(x\) 处为开映射（因而在 \(y'\) 上方存在 \(x\) 在 \(G\) 中的泛化 \(x_1\)），还可以（再作一次基变换后）假设存在 \(G\) 在 \(S\) 上的截面 \(g\)，使 \(x=g(y)\)。若 \(k(y')\) 的特征为 \(0\)，只须任取 \(x\) 在 \(G\) 中位于 \(y'\) 上方的泛化 \(x'\)；由于 \(G^\tau \) 为开集，它属于 \(G^\tau \)，进而由 \(G_{y'}^\sigma =G_{y'}^\tau \) 可知它属于 \(G^\sigma \)。若 \(k(y')\) 的特征为 \(p>0\)，则令 \[     \nu (g(y'))     = p^hm     \qquad \text {其中 }(m,p)=1   \]，取满足 \(ap^h+bm=1\) 的整数 \(a\) 和 \(b\)，并令 \[     \begin {aligned}       g_1       &= g^{ap^h}     \\g_2       &= g^{bm}     \end {aligned}   \]，于是 \(g=g_1g_2\)。由构造，有 \(g_1(y')\in  G^\sigma \) 和 \(g_2(y')\in  G^\rho \)。由于 \(G^0\) 闭，推出 \(g_1(y)\in  G^\sigma \) 和 \(g_2(y)\in  G^\rho \)；又因 \[     g(y)     = g_1(y)g_2(y) \in  G^\sigma    \]，还有 \(g_2(y)\in  G^\sigma \)，所以 \[     g_2(y)\in  G_y^\sigma \cap  G_y^\rho      = G_y^0.   \]。然而，由假设及 \(S\) 为离散赋值环的谱这一事实，\(G\setminus (G_y\setminus  G_y^0)\) 是 \(G\) 的一个开集，\(G\to  S\) 在其上诱导开态射，因而在 \(G_y^0\) 的每一点处给出“拟截面”。所以，必要时再扩张底 \(S\)，可以假设存在 \(G^0\) 在 \(S\) 上的截面 \(g'_2\)，使 \(g'_2(y)=g_2(y)\)。令 \(g'=g_1g'_2\)；由构造，\(g'(y)=g(y)=x\) 且 \(g'(y')=g_1(y')g'_2(y')\in  G^\sigma \)，所以 \(g(y')=x'\) 是 \(x\) 在 \(G^\sigma \) 中位于 \(y'\) 上方的泛化。这就证明了 (iii)。\par{}最后证明断言 (iv)。只须证明，若 \(x\in  G^\rho \)，则 \(x\) 的每个泛化 \(x'\) 都属于 \(G^\rho \)。（对适当的 \(h\)，取在 \(\varphi _{n^h}\) 下的像后，）可以假设甚至有 \(x\in  G^0\)。于是，对于每个与特征互素的整数 \(n>0\)，\(x\) 都在 \(\varphi _n\) 的像中（因为在特征与 \(n\) 互素的域上的连通有限型群中，\(n\) 次幂映射是满射）。由于 \(\varphi _n\) 为开映射，\(x'\) 也在 \(\varphi _n\) 的像中。更确切地说，取 \(G^\tau \) 中包含 \(G_y^0\) 的拟紧开集 \(U\)，则对所有与 \(p\) 互素的 \(n\)，都有 \(x'\in \varphi _n(U)\)。取 \(n\) 为各 \(\nu (z)\)（\(z\in  U\)）的与 \(p\) 互素的因子的公倍数，便得 \(x'\in  G^\rho \)。\par{}\hyperref[fga3.vi-1-theorem-1.1]{定理~\ref*{fga3.vi-1-theorem-1.1}} 的断言 (iii) 还有如下补充：\begin{corollary}{1.5}\setnumber{1.5}\label{fga3.vi-1-corollary-1.5}\skipsinglepar\par{}设 \(n>1\) 是一个整数，使 \(n\) 次幂同态 \(\varphi _n\colon  G\to  G\) 普遍开（例如 étale）。令 \(G^{(n)}=\varphi _n^{-1}(G^0)\)，并假设连通纤维 \(G_s^0\) “不含加法分量”（即将 \(k(s)\) 扩张到代数闭包所得的群不含同构于 \(\operatorname {G_a}\) 的子群）。那么，\(G\to  S\) 在 \(G^{(n^h)}\) 的各点处普遍开。特别地，若各 \(G_s^0\) 不含加法分量，且对每个整数 \(n>1\)，同态 \(\varphi _n\) 在剩余特征与 \(p\) 互素的各点处普遍开，则 \(G\to  S\) 在 \(G^\sigma \) 的各点处普遍开；从而（\hyperref[fga3.vi-1-theorem-1.1]{定理~\ref*{fga3.vi-1-theorem-1.1}} 的 (iii)），若 \(G^0\) 为闭集，则 \(G^\sigma \to  S\) 普遍开。在这种情形中，\(s\mapsto \dim  G_s\) 是 \(S\) 上的局部常值函数。\begin{proof}\par{}由假设，\(\varphi _n\) 的核 \({}_nG\) 在 \(S\) 上普遍开，因此 \(G\to  S\) 在 \({}_nG\) 的各点处普遍开，进而（将 \(n\) 换为 \(n^h\)）在 \({}_{(n^h)}G\) 的各点处也普遍开。但是，关于纤维 \(G_s^0\) 的假设恰好说明，\(G_s\) 中阶为 \(n\) 的某个幂的点在各 \(G^{(n^h)}\) 的并中稠密。容易由此得到，\(G\to  S\) 在此并集的所有点处普遍开，特别地，在单位截面沿线普遍开；又容易由此得到 \(s\mapsto \dim  G_s\) 为局部常值函数。\end{proof}\end{corollary}\begin{remark}{1.6}\setnumber{1.6}\label{fga3.vi-1-remark-1.6}\skipsinglepar\par{}回忆，态射 \(X\to  S\) 称为普遍开的，若它将每个开集映为开集，且任意基变换 \(S'\to  S\) 后仍保有这一性质。这也等价于下述条件（当 \(S\) 局部 Noether、\(X\to  S\) 局部有限型时），\(X\) 的每个不可约分支都支配 \(S\)，且此性质在任意基变换 \(S'\to  S\) 后保持。在这两个断言中，由于上述有限性假设，只须对这样的基变换 \(S'\to  S\) 验证：\(S'\) 是离散赋值环的谱（愿意的话，还可要求完备且剩余域代数闭……）。此定义显然推广到 \(X\) 的子集 \(Z\)（例如 \(G\) 的子集 \(G^\sigma \)）。即使从几何纤维连通的光滑射影态射 \(X\to  S\) 出发（例如 \(S=\operatorname {Spec}\mathbb {C}[j]\) 上椭圆曲线模族的纤维平方），\(\underline {\operatorname {Pic}}_{X/S}\) 在 \(S\) 上也可能并不普遍开，即 \(\underline {\operatorname {Pic}}_{X/S}\) 可能有完全位于 \(S\) 的某一个点上方的不可约分支；这是完全正常的现象。它与 \(X/S\) 的纤维的 Néron–Severi 群的秩可能向上跳跃这一事实有关（“复乘”现象）。不过，(1.5) 保证，在良好情形中，\(\underline {\operatorname {Pic}}_{X/S}^\sigma \)（而且通常看来甚至 \(\underline {\operatorname {Pic}}_{X/S}^\tau \)）在 \(S\) 上普遍开。\end{remark}\par{}最后给出一个有用的情形，在这一特殊情形下，\(G^0\) 的确是开集：\begin{corollary}{1.7}\setnumber{1.7}\label{fga3.vi-1-corollary-1.7}\skipsinglepar\par{}假设 \(G\to  S\) 在 \(G^0\) 的各点处普遍开（参见 \hyperref[fga3.vi-1-corollary-1.5]{推论~\ref*{fga3.vi-1-corollary-1.5}}），且纤维 \(G_s\) 可分，因而在 \(k(s)\) 上光滑（由 Cartier 的一个结果，在剩余特征为零时，后一个条件自动满足）。那么 \(G^0\) 在 \(G\) 中开。若 \(S\) 还约化，则 \(G^0\) 在 \(S\) 上光滑（特别地，平坦）。\begin{proof}\par{}第一个断言可进一步精确表述如下：对某个给定的 \(s\in  S\)，若 \(G\to  S\) 在 \(G_s^0\) 的各点处普遍开，且 \(G_s^0\) 在 \(k(s)\) 上可分，则 \(G^0\) 是 \(G_s^0\) 的一个邻域（而且对所有邻近点，关于 \(s\) 所作的假设仍然满足）。这一陈述的一个证明见 [\hyperref[GD1960]{2}，IV，§7]；此证明与群结构完全无关。\hyperref[fga3.vi-1-corollary-1.7]{推论~\ref*{fga3.vi-1-corollary-1.7}} 的最后一个断言也与群结构及连通分支问题无关；它是 [IV]，第5节 中一个平坦性判据的特例。更一般地，该判据给出：\end{proof}\end{corollary}\begin{corollary}{1.8}\setnumber{1.8}\label{fga3.vi-1-corollary-1.8}\skipsinglepar\par{}设 \(U\) 是纤维 \(G_s\) 的一个开子集，使得 \(G\to  S\) 在 \(U\) 的各点处普遍开（参见 (1.5)）。若 \(G_s\) 在 \(k(s)\) 上可分，且 \(S\) 在 \(s\) 处约化，则 \(G\) 在 \(U\) 的各点处在 \(S\) 上平坦（从而在 \(U\) 的各点处也在 \(S\) 上光滑）。\end{corollary}\begin{remarks}{1.9}\setnumber{1.9}\label{fga3.vi-1-remarks-1.9}\skipsinglepar\par{}若 \(G\) 在 \(S\) 上分离，我不知道 \(G^0\) 或 \(G^\tau \) 是否总在 \(G\) 中闭；这似乎不大可能。无论如何，去掉分离假设就有明显的反例，例如取具有可数无穷多个原点副本的仿射直线，便得到仿射直线上的群概形 \(G\)，其所有纤维都是平凡群，唯有一个纤维是 \(\mathbb {Z}\)。此群概形是一个 Picard 预概形的开群子概形，亦即单位截面的闭包；这里所说的 Picard 预概形对应于 \(S\)-概形 \(X\)，后者是一族退化为两条相交直线的圆锥曲线。\par{}即使从离散赋值环 \(V\) 的谱 \(S\) 上的有限平坦群概形 \(G\) 出发，\(G^0\) 化为单位截面，因而为闭集，只要去掉某些假设，各种结论仍会失效。假设 \(V\) 等特征且特征为 \(p>0\)，令 \(G\) 为函子同态 \(f\mapsto  f^p-tf\) 所给出的同态 \(\operatorname {G_a}\to \operatorname {G_a}\) 的核，其中 \(t\) 是 \(V\) 的一个素元。（回忆，按照定义，\(S\) 上的“加法群”\(\operatorname {G_a}\) 表示函子 \(S'\mapsto \Gamma (S',\mathcal {O}_{S'})\)。）于是，作为 \(S\)-概形，\(G=\operatorname {Spec} V[x]/(x^p-tx)\) 是 \(p\) 条相交“直线”的并，给出一个退化为无穷小加法型群的群 \(\mathbb {Z}/p\mathbb {Z}\)。在这个例子中，\(G^0=G^\sigma \)（这是 \(p\) 条直线中的一条），所以 \(G^0\) 在 \(G\) 中不开。在不等特征、剩余特征为 \(p>0\) 的情形中，从群概形 \(H=\mu _p\) 出发，它由 \(\operatorname {G_m}\) 中 \(p\) 次幂映射的核给出，故 \(H=\operatorname {Spec} V[x]/(x^p-1)\)；这同样是 \(p\) 条相交直线的并，给出群 \(\mathbb {Z}/p\mathbb {Z}\)（在特征 \(0\) 时），它退化为特征 \(p\) 的无穷小乘法型群。令 \(H'\) 为通常的有限群 \(\mathbb {Z}/p\mathbb {Z}\) 所定义的“常值群概形”，它是 \(\mathbb {Z}/p\mathbb {Z}\) 个 \(S\) 的副本的不交和，也就是 \(\operatorname {Spec} V^{\mathbb {Z}/p\mathbb {Z}}\)。于是，\(G=H\times _S H'\) 描述了特征 \(0\) 下的群 \((\mathbb {Z}/p\mathbb {Z})^2\)，它退化为特征 \(p\) 下的一个无穷小群与 \(\mathbb {Z}/p\mathbb {Z}\) 的乘积。这里，\(G^\rho \) 是单位截面与特殊纤维之并，故 \(G^\rho \) 不开；而 \(G^\sigma \) 是单位截面与泛纤维之并，故不闭。这与 \hyperref[fga3.vi-1-theorem-1.1]{定理~\ref*{fga3.vi-1-theorem-1.1}} 的 (iii)、(iv) 中等特征时成立的结论不同。当然，这些都是与特征 \(p>0\) 有关的现象。上述结果给出：\end{remarks}\begin{corollary}{1.10}\setnumber{1.10}\label{fga3.vi-1-corollary-1.10}\skipsinglepar\par{}假设 \(S\) 的剩余特征均为 \(0\)，于是 \(G^\tau =G^\sigma \) 且 \(G^0=G^\rho \)。那么 \(G^\tau =G^\sigma \) 为开集；若 \(G\) 在 \(S\) 上分离，且各 \(G_s^0\) 固有，则它甚至既开又闭。在相同假设下，\(G^0=G^\rho \) 在 \(S\) 上固有，因而在 \(G\) 中闭。最后，若对每个整数 \(n>1\)，\(G\) 中的 \(n\) 次幂同态普遍开，则 \(G^0\) 为开集。若各 \(G_s^0\) 还没有加法分量（例如，像上面一样各 \(G_s^0\) 固有），则 \(G^\tau \to  S\) 普遍开；若 \(S\) 约化，它甚至是光滑的。\end{corollary}\par{}最后指出下述容易得到的结果：\begin{proposition}{1.11}\setnumber{1.11}\label{fga3.vi-1-proposition-1.11}\skipsinglepar\par{}\(S\) 中存在开子集 \(U\)，使得 \(G\) 中满足 \(G\) 在 \(S\) 上光滑（分别地，平坦）的点所成的集合 \(G'\)，是由 \(G|U\) 诱导的某个开群子概形的底层集合。此外，\(G\) 在 \(U\) 上的每个截面都是 \(G'\) 在 \(U\) 上的截面。\end{proposition}\begin{corollary}{1.12}\setnumber{1.12}\label{fga3.vi-1-corollary-1.12}\skipsinglepar\par{}若 \(G\) 在单位截面的各点处在 \(S\) 上光滑（分别地，平坦），则它在 \(G\) 在 \(S\) 上的每个截面的各点处，以及在 \(G^0\) 的所有点处，也都是光滑的（分别地，平坦的）。若进一步对每个整数 \(n>0\)，\(n\) 次幂同态 \(\varphi _n\colon  G\to  G\) 在特征与 \(n\) 互素的各点处 étale，则 \(G\) 在 \(G^\sigma \) 的所有点处在 \(S\) 上光滑（分别地，平坦）。\end{corollary}\chapter{对 Picard 概形的局部性质的应用}\thispagestyle{fancy}\label{fga3.vi-2}\begin{theorem}{2.1}\setnumber{2.1}\label{fga3.vi-2-theorem-2.1}\skipsinglepar\begin{enumerate}
    
 \item[i.]设 \(f\colon  X\to  S\) 是固有光滑态射，并假设 \(\underline {\operatorname {Pic}}_{X/S}\) 存在（例如 \(f\) 射影时）。那么，\(\underline {\operatorname {Pic}}_{X/S}\) 在 \(S\) 上分离；而且对 \(\underline {\operatorname {Pic}}_{X/S}\) 的每个在 \(S\) 上有限型的闭子集 \(Z\)，都有 \(Z\) 在 \(S\) 上固有。

    
 \item[ii.]设 \(X\) 是域 \(k\) 上固有且几何正规的预概形。那么 \(\underline {\operatorname {Pic}}_{X/S}^0\) 在 \(k\) 上固有。

  \end{enumerate}\begin{proof}\par{}对于 (i)，由 [\hyperref[GD1960]{2}，II，§7] 的赋值判据，只须证明：若 \(S\) 是完备离散赋值环的谱，\(U\) 是由 \(S\) 的泛点组成的开集，则 \(\underline {\operatorname {Pic}}_{X/S}\) 在 \(S\) 上的每个有理截面，即在 \(U\) 上的每个截面，都唯一延拓为 \(S\) 上的截面。考虑到 \(\underline {\operatorname {Pic}}_{X/S}\) 的定义，这等价于下述陈述：对 \(V=f^{-1}(U)\) 上的每个可逆模 \(\mathscr {L}\)，\(X\) 上存在一个延拓 \(\mathscr {L}\) 的可逆模，且在同构意义下唯一。用“Cartier”除子类来描述 \(V\) 上（分别地，\(X\) 上）的可逆模，很容易得到这一点：\(X\) 的局部环是正则的（因为 \(X\) 在正则的 \(S\) 上光滑），因而由 Auslander 的结果，它们都是析因环，所以 \(S\) 上的每个除子都是 Cartier 除子。事实上，\(V\) 上的每个除子，取其“闭包”即可延拓为 \(X\) 上的除子。\par{}对于 (ii)，利用 Chow 引理可归约到 \(X\) 射影、因而嵌入某个 \(\mathbb {P}_k^n\) 的情形；还可假设 \(X\) 连通。若 \(\dim  X=1\)，则 \(X\) 在 \(k\) 上光滑，可以应用 (i)。若 \(\dim  X\geqslant 2\)，则可以利用已知“等价判据”的一个变式：存在有限条在 \(X\) 上光滑的曲线 \(Y_i\)（由 \(X\) 与 \(\mathbb {P}_k^n\) 中适当线性子空间相交而得），使得 \(\underline {\operatorname {Pic}}_{X/k}^\tau \to \prod _i\underline {\operatorname {Pic}}_{Y_i/k}^\tau \) 为单态射，从而在连通分支上也诱导单态射。由上述讨论，靶在 \(k\) 上固有；这里的态射是群概形同态，必为闭浸入，因此 \(\underline {\operatorname {Pic}}_{X/k}^0\) 在 \(k\) 上也固有。利用代数闭域上交换代数群的结构（Chevalley–Borel 的结果），可以避免使用微妙的等价判据。此时只须证明，从去掉原点的仿射直线到 \(\underline {\operatorname {Pic}}_{X/S}^\tau \) 的每个态射都是常值态射；这等价于说，\(X[t,t^{-1}]\) 上的每个可逆模都来自 \(X\) 上的可逆模。这是一个熟知的初等结果，甚至不需要 \(X\) 在 \(k\) 上固有这一事实（\(X\) 正规的假设使我们可立即归约到 \(X\) 正则的情形）。\end{proof}\end{theorem}\begin{remark}{2.2}\setnumber{2.2}\label{fga3.vi-2-remark-2.2}\skipsinglepar\par{}即使仅假设 \(f\) 平坦，且其纤维 \(X_s\) 为局部完全交并在余维 \(\leqslant 2\) 处在 \(k(s)\) 上光滑，上面对 (i) 的证明仍然成立；这里利用了 \hyperref[fga3.vi-bib-5]{[5]} 中证明的下述事实：在余维 \(\leqslant 3\) 处正则的 Noether 完全交局部环是析因环（“Samuel 猜想”）。注意，若将“余维 \(\leqslant 2\)”换为“余维 \(\leqslant 1\)”，即换为“正规”的假设，结论就不再成立；考虑一族退化为二次锥面的非奇异二次曲面即可看出。\end{remark}\begin{corollary}{2.3}\setnumber{2.3}\label{fga3.vi-2-corollary-2.3}\skipsinglepar\par{}设 \(f\colon  X\to  S\) 是固有正规态射（即平坦且几何纤维正规），并假设 \(\underline {\operatorname {Pic}}_{X/S}\) 存在。那么 \(\underline {\operatorname {Pic}}_{X/S}^0\) 在 \(S\) 上固有，从而在 \(\underline {\operatorname {Pic}}_{X/S}\) 中闭；此外，\(\underline {\operatorname {Pic}}_{X/S}^\tau \) 与 \(\underline {\operatorname {Pic}}_{X/S}^\rho \) 都闭，在“等特征”情形下，\(\underline {\operatorname {Pic}}_{X/S}^\sigma \) 也闭。\begin{proof}\par{}应用 \hyperref[fga3.vi-1-theorem-1.1]{定理~\ref*{fga3.vi-1-theorem-1.1}} 及 \hyperref[fga3.vi-2-theorem-2.1]{定理~\ref*{fga3.vi-2-theorem-2.1}} 的 (ii)。\end{proof}\end{corollary}\begin{corollary}{2.4}\setnumber{2.4}\label{fga3.vi-2-corollary-2.4}\skipsinglepar\par{}设 \(f\colon  X\to  S\) 是固有光滑态射，且 \(\underline {\operatorname {Pic}}_{X/S}\) 存在，并为 \(S\) 上有限型概形 \(P^{(i)}\) 的和（参见 \hyperref[fga3.v-4-proposition-4.1]{[3]，V，4.1}）。那么每个 \(P^{(i)}\) 都在 \(S\) 上固有。\begin{proof}\par{}这由 \hyperref[fga3.vi-2-theorem-2.1]{定理~\ref*{fga3.vi-2-theorem-2.1}} 的 (i) 推出。\end{proof}\end{corollary}\par{}如 \hyperref[fga3.vi-2-remark-2.2]{注~\ref*{fga3.vi-2-remark-2.2}} 中所指出，减弱对 \(f\) 的纤维所作的假设，仍可推广上述结果；但若仅假设 \(f\) 正规，则结果不再成立。在这种情形中，即使假定 \(\underline {\operatorname {Pic}}_{X/S}^\tau \) 在 \(S\) 上有限型，我也不知道它是否仍在 \(S\) 上固有。\begin{theorem}{2.5}\setnumber{2.5}\label{fga3.vi-2-theorem-2.5}\skipsinglepar\par{}设 \(f\colon  X\to  S\) 是固有平坦态射，且 \(\underline {\operatorname {Pic}}_{X/S}\) 存在；对每个整数 \(n\)，以 \(\varphi _n\) 表示此群预概形中的 \(n\) 次幂同态。那么，\(\varphi _n\) 在剩余特征与 \(n\) 互素的所有点 \(x\in  X\) 处 étale。\end{theorem}\par{}由 étale 态射的无穷小刻画，上面的断言等价于下述陈述：\begin{lemma}{2.6}\setnumber{2.6}\label{fga3.vi-2-lemma-2.6}\skipsinglepar\par{}假设 \(S\) 是 Artin 局部环 \(A\) 的谱，其极大理想 \(\mathfrak {m}\) 的 \((\nu +1)\) 次幂为零，并令 \(A_{\nu -1}=A/\mathfrak {m}^\nu \)、\(X_{\nu -1}=X\otimes _A A_{\nu -1}\)。设 \(\mathscr {L}\) 是 \(X\) 上的可逆模，\(\mathscr {L}'_{\nu -1}\) 是 \(X_{\nu -1}\) 上的可逆模，其 \(n\) 次张量幂同构于 \(\mathscr {L}_{\nu -1}=\mathscr {L}\otimes _A A_{\nu -1}\)。那么，\(X\) 上存在可逆模 \(\mathscr {L}'\)，其 \(n\) 次张量幂同构于 \(\mathscr {L}\)（只要 \(n\) 与 \(k=k(\mathscr {L})\) 的剩余特征互素）。\begin{proof}\par{}令 \(V=\mathfrak {m}^\nu =\mathfrak {m}^{\nu }/\mathfrak {m}^{\nu +1}\)，它是 \(k=k(A)\) 上的向量空间。先将 \(\mathscr {L}'_{\nu -1}\) 任意延拓为 \(X\) 上的可逆模 \(\mathscr {L}'\)。这样做的障碍位于 \(\operatorname {H}^2(X_0,\mathcal {O}_{X_0})\otimes _k V\) 中；但由关于 \(\mathscr {L}'_{\nu -1}\) 的假设及 \(\mathscr {L}_{\nu -1}\) 可以延拓这一事实，此障碍与 \(n\) 的乘积为零；由于 \(n\) 与特征互素，障碍本身必为零。延拓的不唯一性位于 \(\operatorname {H}^1(X_0,\mathcal {O}_{X_0})\otimes _k V\) 中，而 \({\mathscr {L}'}^{\otimes  n}\) 与 \(\mathscr {L}\) 之间的偏差 \(\xi \) 也位于同一个模中。若想修正 \(\mathscr {L}'\) 使此偏差为零，就须在上述模中找到一个 \(\eta \)，使 \(n\eta =\xi \)。由于 \(n\) 与特征互素，这仍然可以做到。\end{proof}\end{lemma}\begin{corollary}{}\label{fga3.vi-2-corollary}\skipsinglepar\par{}在 \hyperref[fga3.vi-2-theorem-2.5]{定理~\ref*{fga3.vi-2-theorem-2.5}} 的条件下，进一步假设纤维 \(X_s\) 的 Picard 概形均不含加法分量（例如，各 \(X_s\) 几何正规时；参见 \hyperref[fga3.vi-2-theorem-2.1]{定理~\ref*{fga3.vi-2-theorem-2.1}} 的 (ii)）。那么 \(\underline {\operatorname {Pic}}_{X/S}\to  S\) 在 \(\underline {\operatorname {Pic}}_{X/S}^\sigma \) 的各点处普遍开。若 \(\underline {\operatorname {Pic}}_{X/S}^0\) 闭（例如，各 \(X_s\) 几何正规时；参见 \hyperref[fga3.vi-2-corollary-2.3]{推论~\ref*{fga3.vi-2-corollary-2.3}}），则 \(\underline {\operatorname {Pic}}_{X/S}^\sigma \) 本身在 \(S\) 上普遍开。最后，在等特征情形中，\(\underline {\operatorname {Pic}}_{X/S}^\rho \to  S\) 普遍开。\begin{proof}\par{}应用 \hyperref[fga3.vi-1-corollary-1.5]{推论~\ref*{fga3.vi-1-corollary-1.5}} 和 \hyperref[fga3.vi-1-theorem-1.1]{定理~\ref*{fga3.vi-1-theorem-1.1}}。\end{proof}\end{corollary}\begin{corollary}{2.7}\setnumber{2.7}\label{fga3.vi-2-corollary-2.7}\skipsinglepar\par{}设 \(f\colon  X\to  Y\) 是固有平坦态射，且 \(\underline {\operatorname {Pic}}_{X/S}\) 存在。那么，\(S\) 上的函数 \(s\mapsto \dim \underline {\operatorname {Pic}}_{X_s/k(s)}\) 上半连续（即只可能向上跳跃，不能向下跳跃）；若各 \(\underline {\operatorname {Pic}}_{X_s/k(s)}\) 不含加法分量，则它甚至连续（即局部常值）。\begin{proof}\par{}对于局部 Noether 底上任意局部有限型群预概形，第一个断言都是平凡的，或几乎平凡的，因为只须考察单位截面沿线的情形。第二个断言由 \hyperref[fga3.vi-2-theorem-2.5]{定理~\ref*{fga3.vi-2-theorem-2.5}} 得到。\end{proof}\end{corollary}\begin{remark}{2.8}\setnumber{2.8}\label{fga3.vi-2-remark-2.8}\skipsinglepar\par{}设 \(s,s'\in  S\)，且 \(s\) 是 \(s'\) 的特殊化。那么，\hyperref[fga3.vi-2-corollary-2.7]{推论~\ref*{fga3.vi-2-corollary-2.7}} 等价于 \(X_{s'}\) 及其“特殊化”\(X_s\) 的 Picard 概形的维数之间的一个不等式（分别地，等式）。甚至在 Picard 概形构造出来以前，Serre 就已注意到：对于 \emph{光滑} 态射 \(f\colon  X\to  S\)，各 \(X_s\) 的 Picard 簇的维数不变这一事实，是基本群特殊化理论（[\hyperref[Gro1960b]{4}，X]）以及 Picard 簇的有限阶点与基本群 Abel 化之间的经典 Kummer 型关系（[\hyperref[Gro1960b]{4}，XI]）的形式推论。用 \(\alpha \)、\(\mu \)、\(\lambda \) 分别表示 \(\underline {\operatorname {Pic}}_{X_s/k(s)}\) 的 Abel 部分、乘法部分和加法部分的维数，并类似定义 \(\alpha '\)、\(\mu '\)、\(\lambda '\)，则已知关系可写为下列不等式：\begin{eqn}{*}\setnumber{*}\label{fga3.vi-2-equation-star}\skipsinglepar\[         \alpha +\mu +\lambda          \geqslant  \alpha '+\mu '+\lambda '       \tag{*}       \]\end{eqn}（只要 \(\underline {\operatorname {Pic}}_{X/S}\) 存在就成立，因此很可能在所有情形中均成立）。当 \(\lambda =\lambda '=0\) 时，这个不等式化为在同样存在性假设下成立的等式：\[       \alpha +\mu        = \alpha '+\mu '.     \]。由 Serre 的论证，若各 \(X_s\) 可分（甚至不必假设 \(\underline {\operatorname {Pic}}_{X/S}^\tau \) 存在），或者各 \({}_n\underline {\operatorname {Pic}}_{X/S}\)（Picard 函子中各 \(\varphi _n\) 的核）在 \(S\) 上 \emph{分离}，则还有 \begin{eqn}{**}\setnumber{**}\label{fga3.vi-2-equation-star-star}\skipsinglepar\[         2\alpha +\mu          \leqslant  2\alpha '+\mu '       \tag{**}       \]\end{eqn}（这里利用了 \hyperref[fga3.vi-2-theorem-2.5]{定理~\ref*{fga3.vi-2-theorem-2.5}} 所保证的它们在 \(S\) 上 étale 这一事实）。我们倾向于猜想，\hyperref[fga3.vi-2-equation-star]{式~\ref*{fga3.vi-2-equation-star}} 在所有情形中均为等式，至少在各 \(X_s\) 可分时如此；还猜想有不等式 \begin{eqn}{***}\setnumber{***}\label{fga3.vi-2-equation-star-star-star}\skipsinglepar\[         \begin {aligned}           \alpha  &\leqslant  \alpha '         \\\lambda  &\geqslant  \lambda '         \end {aligned}       \tag{***}       \]\end{eqn}。只要所考虑的是局部 Noether 的 \(S\) 上的局部有限型群预概形，且纤维维数恒定，这些不等式就应当成立（这个方向的一个肯定结果见 \hyperref[fga3.vi-1-lemma-1.3]{引理~\ref*{fga3.vi-1-lemma-1.3}}）。\end{remark}\begin{remark}{2.9}\setnumber{2.9}\label{fga3.vi-2-remark-2.9}\skipsinglepar\par{}在所有已知情形中，\(\underline {\operatorname {Pic}}_{X/S}^\tau \) 在 \(S\) 上普遍开；但即使 \(f\colon  X\to  S\) 光滑，我们恐怕也不宜过于乐观。无论如何，Mumford 通过对 Igusa 曲面作无穷小变形，构造了 \(\underline {\operatorname {Pic}}_{X/S}^\tau \) 在 \(S\) 上不平坦的例子（诚然，其中 \(S\) 不约化，实际上 \(S\) 是 Artin 环的谱）。Mumford 所考察的点也属于 \(\underline {\operatorname {Pic}}_{X/S}^\rho \)；因此，（对光滑的 \(f\colon  X\to  S\)）\(\underline {\operatorname {Pic}}_{X/S}\) 在 \(\underline {\operatorname {Pic}}_{X/S}^\sigma \) 的各点处在 \(S\) 上平坦仍有可能成立。不过，报告人怀疑这一点并不总成立，即使只考虑 \(\underline {\operatorname {Pic}}_{X/S}^0\) 的各点，并假设 \(S\) 是离散赋值环的谱。这个问题与研究 Abel 概形在有限自同构群作用下的不动点有关；在这种情形中，我们似乎还缺少例子。看来，即使只考虑光滑射影的 \(f\colon  X\to  S\)，如果不对 \(f\) 的纤维的性质作更特殊的假设，本节关于 \(\underline {\operatorname {Pic}}_{X/S}\) 的局部正则性的结果以及 \hyperref[fga3.vi-2-remark-2.8]{注~\ref*{fga3.vi-2-remark-2.8}} 中提出的猜想，基本上已经穷尽了在这方面所能说的内容。不过，回忆，若 \(\underline {\operatorname {Pic}}_{X/S}\) 的几何纤维 \emph{约化} 且没有加法分量，则由 \hyperref[fga3.vi-1-corollary-1.8]{推论~\ref*{fga3.vi-1-corollary-1.8}} 和 \hyperref[fga3.vi-2-theorem-2.5]{定理~\ref*{fga3.vi-2-theorem-2.5}} 可知，只要 \(S\) 约化，\(\underline {\operatorname {Pic}}_{X/S}\) 就在 \(\underline {\operatorname {Pic}}_{X/S}^\sigma \) 的各点处在 \(S\) 上光滑。若 \(f\colon  X\to  S\) 正规，则即使不对 \(S\) 作任何假设，此结果仍成立；这将在 \hyperref[fga3.vi-3-theorem-3.5]{定理~\ref*{fga3.vi-3-theorem-3.5}} 中看到。就此指出：\end{remark}\begin{proposition}{2.10}\setnumber{2.10}\label{fga3.vi-2-proposition-2.10}\skipsinglepar\begin{enumerate}
    
 \item[i.]若 \(\underline {\operatorname {Pic}}_{X/S}\) 在单位截面的各点处在 \(S\) 上光滑（分别地，平坦），则它在 \(\underline {\operatorname {Pic}}_{X/S}^\sigma \) 的所有点处，以及 \(\underline {\operatorname {Pic}}_{X/S}\) 在 \(S\) 上的每个截面的各点处，也都是光滑的（分别地，平坦的）。

    
 \item[ii.]设 \(s\in  S\) 满足 \(\operatorname {H}^2(X_s,\mathcal {O}_{X_s})=0\)。那么 \(s\) 有一个开邻域 \(U\)，使 \(\underline {\operatorname {Pic}}_{X/S}|U\) 在 \(U\) 上光滑。

    
 \item[iii.]设 \(X\) 是一个域上的固有概形。那么有 \[         \dim \underline {\operatorname {Pic}}_{X/k}         \leqslant  \dim \operatorname {H}^1(X,\mathcal {O}_X)       \]，其中等号成立当且仅当 \(\underline {\operatorname {Pic}}_{X/k}\) 在 \(k\) 上光滑；若 \(k\) 的特征为 \(0\)，则总是如此。

  \end{enumerate}\begin{proof}\par{}断言 (i) 由 \hyperref[fga3.vi-2-theorem-2.5]{定理~\ref*{fga3.vi-2-theorem-2.5}} 和 \hyperref[fga3.vi-1-corollary-1.12]{推论~\ref*{fga3.vi-1-corollary-1.12}} 得到；(ii) 则由光滑态射的无穷小判据及一个熟知的障碍计算得到，这里利用了关于 \(s\) 的假设由“半连续性定理”对邻近点仍成立。最后，(iii) 由 \(\operatorname {H}^1(X,\mathcal {O}_X)\) 同构于 \(\underline {\operatorname {Pic}}_{X/k}\) 在单位元处的 Zariski 切空间这一事实得到；最后一个断言是 Cartier 定理的特例，该定理说，特征 \(0\) 下的一个“形式群”在 \(k\) 上形式光滑。\end{proof}\end{proposition}\chapter{\texorpdfstring{\(\underline {\operatorname {Pic}}_{X/S}\)}{Pic(X/S)} 的典范 Abel 子概形与 Albanese 概形}\thispagestyle{fancy}\label{fga3.vi-3}\begin{proposition}{3.1}\setnumber{3.1}\label{fga3.vi-3-proposition-3.1}\skipsinglepar\par{}设 \(k\) 是域，\(G\) 是 \(k\) 上交换且“不含加法分量”的有限型群概形。那么，\(G_\mathrm {red}^0\) 在 \(k\) 上可分，因而是 \(k\) 上的光滑群概形。\begin{proof}\par{}若 \(k\) 是完美域，断言显然成立；特别地，对 \(G_{\overline {k}}\) 也成立，其中 \(\overline {k}\) 是 \(k\) 的代数闭包。因此，只须证明 \((G_{\overline {k}}^0)_\mathrm {red}\) 来自 \(G\) 的一个子概形。由 \(G_{\overline {k}}\) 不含加法分量的假设，容易得到存在整数 \(m\)，使 \((G_{\overline {k}}^0)_\mathrm {red}\) 是 \(G_{\overline {k}}\) 中 \(m\) 次幂同态的“概形论”像。由于后一个同态来自 \(G\) 的类似同态，该同态的概形论像就是所需对象。\end{proof}\end{proposition}\begin{corollary}{3.2}\setnumber{3.2}\label{fga3.vi-3-corollary-3.2}\skipsinglepar\par{}设 \(X\) 是 \(k\) 上的正规固有概形，且 \(\underline {\operatorname {Pic}}_{X/k}\) 存在。那么 \(\underline {\operatorname {Pic}}_{X/k}\) 有一个底层集合为 \(\underline {\operatorname {Pic}}_{X/k}^0\) 的 Abel 子概形。\begin{proof}\par{}由 \hyperref[fga3.vi-2-theorem-2.1]{定理~\ref*{fga3.vi-2-theorem-2.1}} 的 (ii)，\(\underline {\operatorname {Pic}}_{X/k}^0\) 在 \(k\) 上固有，因而满足 \hyperref[fga3.vi-3-proposition-3.1]{命题~\ref*{fga3.vi-3-proposition-3.1}} 的条件。\end{proof}\end{corollary}\par{}上述结果表明，在某些情形中，经典的“Picard 簇”（在现今理论中即 \((\underline {\operatorname {Pic}}_{X/\overline {k}})_\mathrm {red}^0\)）“定义在 \(k\) 上”，无须假设域 \(k\) 完美。\par{}现在设 \(f\colon  X\to  S\) 是固有平坦的相对概形，为简单起见设 \(\mathcal {O}_S\xrightarrow {\sim } f_*(\mathcal {O}_X)\)，并假设 \(\underline {\operatorname {Pic}}_{X/S}\) 存在，且 \(\underline {\operatorname {Pic}}_{X/S}^0\) 在 \(S\) 上固有。对于 \hyperref[fga3.vi-3-theorem-3.3]{定理~\ref*{fga3.vi-3-theorem-3.3}} 的 (ii)，进一步假设 \(\underline {\operatorname {Pic}}_{X/S}\) 中存在一个包含 \(\underline {\operatorname {Pic}}_{X/S}^0\)、在 \(S\) 上拟射影的开集；如上所见，若 \(f\) 射影且几何纤维可分且不可约，此条件就得到满足。回忆，\(S\) 上的 \emph{Abel 概形} 是 \(S\) 上的群概形，它在 \(S\) 上固有且光滑，并且几何纤维连通。我们拟考察是否存在 \(\underline {\operatorname {Pic}}_{X/S}\) 的群子概形 \(A\)，使其为 Abel 概形且底层集合是 \(\underline {\operatorname {Pic}}_{X/S}^0\)。刚刚已看到，若 \(S\) 是域的谱，这样的 \(A\) 总是存在。在这里所考虑的一般情形中，目前可以说的是：\begin{theorem}{3.3}\setnumber{3.3}\label{fga3.vi-3-theorem-3.3}\skipsinglepar\par{}在上述条件下：\begin{enumerate}
      
 \item[i.]若 \(\underline {\operatorname {Pic}}_{X/S}\) 有一个底层集合为 \(\underline {\operatorname {Pic}}_{X/S}^0\) 的 Abel 子概形，则它唯一。因此，它的构造与基变换相容。

      
 \item[ii.]这样的 Abel 子概形存在的充要条件是：作每个 \(S'\) 为局部 Artin 概形的基变换 \(S'\to  S\) 后，它都存在。若 \(S\) 是局部环的谱，则甚至只须对满足 \(A_n=A/\mathfrak {m}^{n+1}\) 的各 \(S'=\operatorname {Spec}(A_n)\) 验证。若 \(S\) 约化，则也只须对各 \(S'\) 是离散赋值环的谱的情形验证（愿意的话，可取完备且剩余域代数闭）。

      
 \item[iii.]假设 \(A\) 存在，设 \(B=\operatorname {Alb}^0(X/S)\) 是其对偶 Abel 概形（即 \(B=\operatorname {Pic}_{A/S}^0\)；参见 \hyperref[fga3.vi-bib-7]{[7]}）。那么，可以典范地构造 \(B\) 的一个主齐性空间 \(P=\operatorname {Alb}^1(X/S)\)，以及一个 \(S\)-态射 \(X\to  P\)；它对于从 \(X\) 到准 Abel 概形（即 Abel 概形的主齐性空间）的 \(S\)-态射是泛的。\(\operatorname {Alb}^0(X/S)\)、\(\operatorname {Alb}^1(X/S)\)、\(X\to \operatorname {Alb}^1(X/S)\) 的构造均与基变换交换。

    \end{enumerate}\begin{proof}\par{}证明梗概如下：\begin{enumerate}\item[i.]这是交换群概形的 Abel 子概形所具有的一般刚性性质：若两个这样的子概形在点 \(s\in  S\) 处在集合意义下相同，则它们在 \(s\) 的整个连通分支上相同（\hyperref[fga3.vi-bib-7]{[7]}）。（这一结果推广了 Chow 的一个经典定理。）

      \item[ii.]利用 Hilbert 概形可见，将 \(S\) 上的每个 \(S'\) 映到 \((\underline {\operatorname {Pic}}_{X/S})\times _S S'\) 的典范 Abel 子概形所成的集合（其中至多一个元素）的函子，可由 \(S\) 上的有限型概形 \(T\) 表示。由 (i)，\(T\to  S\) 为单态射；由 \hyperref[fga3.vi-3-corollary-3.2]{推论~\ref*{fga3.vi-3-corollary-3.2}}，它是满射。说 \(\underline {\operatorname {Pic}}_{X/S}\) 有典范 Abel 子概形，等价于说 \(T\) 在 \(S\) 上有截面，也等价于说 \(T\to  S\) 为同构。这又等价于说 \(T\to  S\) 为 étale，或者，在 \(S\) 约化时，等价于说 \(T\to  S\) 固有。于是立刻得到 (ii)。

      \item[iii.]只用 \(\underline {\operatorname {Pic}}_{X/S}\) 的定义，就可注意到：对于 \(S\) 上的每个 Abel 概形 \(C\)，给定从 \(X\) 到 \(C\) 的一个主齐性空间的 \(S\)-态射，等价于给定群同态 \(C'\to \underline {\operatorname {Pic}}_{X/S}\)，其中 \(C'\) 是 \(C\) 的对偶 Abel 概形。但若 \(\underline {\operatorname {Pic}}_{X/S}\) 的典范 Abel 子概形 \(A\) 存在，则这些同态必通过 \(A\) 分解（再次利用有限阶点即可看出）。于是立刻得到 (iii)。\qedhere\end{enumerate}\end{proof}\end{theorem}\begin{remark}{3.4}\setnumber{3.4}\label{fga3.vi-3-remarks-3.4}\skipsinglepar\par{}若 \(\underline {\operatorname {Pic}}_{X/S}\) 的典范 Abel 子概形存在，将其记为 \(\underline {\operatorname {Pic}}_{X/S}^{00}\)。遗憾的是，它并不总是存在；对 Igusa 曲面作无穷小变形（一阶模变形）即可看到这一点。不过，至少在 \(S\) 约化时，\(\underline {\operatorname {Pic}}_{X/S}^{00}\) 可能总存在；由 (ii)，这等价于在 \(S\) 为离散赋值环的谱时存在。设 \(X_0\) 和 \(X_1\) 分别为 \(X\) 的特殊纤维和泛纤维，令 \(A_1=\underline {\operatorname {Pic}}_{X_1/K_1}\)，其中 \(K\) 是赋值环 \(V\) 的分式域。由 Koizumi 的结果，\(S\) 上存在本质唯一的 Abel 概形 \(A\)，其泛纤维为 \(A_1\)。如同 (2.1), (i) 的证明那样（从现在起假设 \(X\) 在 \(S\) 上光滑），容易看出，\(A_1\) 的恒等态射延拓为态射 \[       A \to  \underline {\operatorname {Pic}}_{X/S}.     \]由此得到同态 \begin{eqn}{*}\setnumber{*}\label{fga3.vi-3-equation-star}\skipsinglepar\[         A_0 \to  \underline {\operatorname {Pic}}_{X_0/k}^{00}       .\tag{*}       \]\end{eqn}容易证明，此同态为满射，核是一个有限 \(p\)-群，其中 \(p\) 是剩余域 \(k\) 的特征指数（仍利用有限阶点）。于是，关于 \(X/S\) 的下列条件等价：\begin{enumerate}
      
 \item[a.]上述同态 \hyperref[fga3.vi-3-equation-star]{式~\ref*{fga3.vi-3-equation-star}} 是同构（Shimura 将此表述为“Picard 簇”的构造“与特殊化相容”）。

      
 \item[b.]\(\underline {\operatorname {Pic}}_{X/S}^{00}\) 存在（此时它恰为 \(A\)）。

      
 \item[c.]（为便于回忆。）各 \(\underline {\operatorname {Pic}}_{X_n/S}^{00}\) 存在。

    \end{enumerate}\par{}由上面对 \hyperref[fga3.vi-3-equation-star]{式~\ref*{fga3.vi-3-equation-star}} 的核所作的注记，剩余特征为零时，条件 (a) 得到满足；不过，\hyperref[fga3.vi-3-theorem-3.5]{定理~\ref*{fga3.vi-3-theorem-3.5}} 将显著推广这一结果。\par{}当然，若 \(\underline {\operatorname {Pic}}_{X/S}\) 在 \(\underline {\operatorname {Pic}}_{X/S}^0\) 的各点处在 \(S\) 上光滑，则后者在 \(\underline {\operatorname {Pic}}_{X/S}\) 中开（参见 \hyperref[fga3.vi-1-corollary-1.7]{推论~\ref*{fga3.vi-1-corollary-1.7}}）；赋予诱导结构后，它是 \(\underline {\operatorname {Pic}}_{X/S}\) 的一个 Abel 子概形，因而等于 \(\underline {\operatorname {Pic}}_{X/S}^{00}\)；后者在此情形中存在。但我们有更强的结果：\end{remark}\begin{theorem}{3.5}\setnumber{3.5}\label{fga3.vi-3-theorem-3.5}\skipsinglepar\par{}在 \hyperref[fga3.vi-3-theorem-3.3]{定理~\ref*{fga3.vi-3-theorem-3.3}} 的条件下，设 \(s\in  S\) 使得 \(\underline {\operatorname {Pic}}_{X_s/k(s)}\) 在 \(k(s)\) 上光滑（或者等价地，使得 \(\dim \underline {\operatorname {Pic}}_{X_s/k(s)}=\dim \operatorname {H}^1(X_s,\mathcal {O}_{X_s})\)）。那么，\(s\) 有一个开邻域 \(U\)，使 \(\underline {\operatorname {Pic}}_{X/S}\) 在 \(\underline {\operatorname {Pic}}_{X/S}^0|U\) 的各点处在 \(S\) 上光滑，因而后者是 \(\underline {\operatorname {Pic}}_{X/S}|U\) 中的开 Abel 子概形。特别地，\(\underline {\operatorname {Pic}}_{X|U/U}^{00}\) 存在。\begin{proof}\par{}说明证明的思路。上述结果使我们可归约到 \(S\) 是 Artin 局部环 \(A\) 的谱的情形，对 \(A\) 的无穷小阶数作归纳。因此可以假设 \(\underline {\operatorname {Pic}}_{X_n/A_n}^0\) 在 \(A_n\) 上光滑，只须证明 \(\underline {\operatorname {Pic}}_{X_{n+1}/A_{n+1}}^0\) 在 \(A_{n+1}\) 上光滑。为此，注意只须构造 \(A_{n+1}\) 上的 Abel 概形 \(P_{n+1}\) 来延拓 \(P_n=\underline {\operatorname {Pic}}_{X_n/A_n}^0\)，并构造 \(X_{n+1}\times _{A_{n+1}}P_{n+1}\) 上的可逆模 \(\mathscr {L}_{n+1}\)，来延拓 \(X_n\times _{A_n}P_n\) 上的可逆模 \(\mathscr {L}_n\)；后一个模出现于将 Picard 概形 \(\underline {\operatorname {Pic}}_{X_n/A_n}\) 定义为泛问题的解时。（注意：可以假设 \(X\) 在 \(S\) 上配有截面……。）为进行这一构造，必须使用下列关键结果：\emph{定义在 Artin 局部环的商环上的每个 Abel 概形都可以延拓}（换言之，代数闭域上的 Abel 概形的绝对“形式模概形”（\hyperref[fga3.ii]{[3]，II}）在 \(k\) 上的 Witt 向量环上光滑）。这一结果可由提升障碍的一般形式性质及群运算得到。利用此结果，先将 \(P_n\) 任意延拓为 \(P_{n+1}\)；然后得到提升 \(\mathscr {L}_n\) 的障碍，位于 \(\operatorname {H}^2(X_0\times  P_0,\mathcal {O}_{X_0\times  P_0})\otimes _k V\) 中（其中 \(V=\mathfrak {m}^{n+1}/\mathfrak {m}^{n+2}\)），更确切地说，位于子空间 \(\operatorname {H}^1(X_0,\mathcal {O}_{X_0})\otimes \operatorname {H}^1(P_0,\mathcal {O}_{P_0})\otimes _k V\) 中（因为 \(\mathscr {L}_n\) 在两个因子 \(X_n\) 和 \(P_n\) 上的限制平凡）。但后一个空间恰为 \(\operatorname {H}^1(P_0,\mathscr {G}_{P_0/k})\otimes  V\)，其中 \(\mathscr {G}_{P_0/k}\) 是 \(P_0/k\) 的切层；因而它也是将 \(P_n\) 提升为 \(P_{n+1}\) 时描述不确定性的空间（\hyperref[fga3.ii]{[3]，II}）。所以，可以修正此提升（而且正如预期的那样，修正方式唯一），使提升 \(\mathscr {L}_n\) 的障碍消失。\end{proof}\end{theorem}\begin{corollary}{3.6}\setnumber{3.6}\label{fga3.vi-3-corollary-3.6}\skipsinglepar\par{}在 \hyperref[fga3.vi-3-theorem-3.5]{定理~\ref*{fga3.vi-3-theorem-3.5}} 的条件下，\(\operatorname {R}^1 f_*(\mathcal {O}_X)\) 是 \(S\) 上在 \(s\) 的某个邻域内局部自由的模，其构造与基变换交换。\begin{proof}\par{}此模恰好是 \(\underline {\operatorname {Pic}}_{X/S}\) 沿单位截面的切模。\end{proof}\end{corollary}\begin{remark}{3.7}\setnumber{3.7}\label{fga3.vi-3-remark-3.7}\skipsinglepar\par{}利用与 \hyperref[fga3.vi-3-theorem-3.5]{定理~\ref*{fga3.vi-3-theorem-3.5}} 相同的论证，可以证明：若 \(S'\) 是 \(S\) 中由幂零凝聚理想定义的子概形，且仅假设 \(\underline {\operatorname {Pic}}_{X'/S'}^0\) 存在并具有光滑纤维，则 \(\underline {\operatorname {Pic}}_{X/S}^0\) 必然存在，且为 \(S\) 上的 Abel 概形。这使我们在不作任何射影性假设时，也能在某些情形中构造 Picard 概形；例如，Artin 环上 Abel 概形的对偶 Abel 概形总存在。在 \(X\) 是 \(S\) 上 Abel 概形的情形中，利用 \hyperref[fga3.vi-3-corollary-3.6]{推论~\ref*{fga3.vi-3-corollary-3.6}}，以及 \(\operatorname {H}^\bullet (X_s,\mathcal {O}_{X_s})\) 是 \(\operatorname {H}^1(X_s,\mathcal {O}_{X_s})\) 上外代数的已知结构（Rosenlicht–Serre），便得 \(\operatorname {R}^if_*(\mathcal {O}_X)\) 对 \emph{所有} \(i\) 均局部自由；更确切地说，它同构于 \(\operatorname {R}^1f_*(\mathcal {O}_X)\) 的 \(i\) 次外幂。\end{remark}\begin{remark}{3.8}\setnumber{3.8}\label{fga3.vi-3-remark-3.8}\skipsinglepar\par{}对于光滑射影态射 \(f\colon  X\to  S\)，若 \(S\) 约化且所有剩余特征为零，则 \hyperref[fga3.vi-3-corollary-3.6]{推论~\ref*{fga3.vi-3-corollary-3.6}} 的结果早已通过超越方法作为 Hodge 理论的推论而知道。事实上，此时所有 \(\operatorname {R}^pf_*(\Omega _{X/S}^q)\) 都局部自由。不过，在不等特征情形中，Serre 簇（[\hyperref[Ser1958b]{8}，第 50 页]）给出“\(\operatorname {R}^1f_*(\mathcal {O}_X)\) 局部自由”的反例。等特征情形中似乎尚无反例。\end{remark}\begin{corollary}{3.8}\setnumber{3.8}\label{fga3.vi-3-corollary-3.8}\skipsinglepar\par{}在 \hyperref[fga3.vi-3-theorem-3.5]{定理~\ref*{fga3.vi-3-theorem-3.5}} 的条件下，\(\underline {\operatorname {Pic}}_{X/S}^\sigma |U\) 在 \(\underline {\operatorname {Pic}}_{X/S}^\sigma |U\) 的所有点处在 \(U\) 上光滑。\begin{proof}\par{}应用 \hyperref[fga3.vi-2-proposition-2.10]{命题~\ref*{fga3.vi-2-proposition-2.10}} 的 (i)。\end{proof}\end{corollary}\par{}特别地，考虑到 \hyperref[fga3.vi-2-proposition-2.10]{命题~\ref*{fga3.vi-2-proposition-2.10}} 的 (iii)，有：\begin{corollary}{3.9}\setnumber{3.9}\label{fga3.vi-3-corollary-3.9}\skipsinglepar\par{}假设 \(S\) 的所有剩余特征为零。那么，\(\underline {\operatorname {Pic}}_{X/S}^\tau \) 在 \(S\) 上光滑。\end{corollary}\begin{remark}{3.10}\setnumber{3.10}\label{fga3.vi-3-remark-3.10}\skipsinglepar\par{}例如，由此推出，若 \(\underline {\operatorname {Pic}}_{X/S}\) 还在 \(S\) 上 \emph{固有}，则 \(f\) 的几何纤维的 Néron–Severi 挠群在 \(S\) 的每个连通分支上恒定（当 \(f\) 光滑且射影时，这一点用超越方法也很显然）。注意，直接利用 \hyperref[fga3.vi-2-theorem-2.5]{定理~\ref*{fga3.vi-2-theorem-2.5}}，可以更一般地证明：若 \(\underline {\operatorname {Pic}}_{X/S}\) 在 \(S\) 上固有（例如 \(f\colon  X\to  S\) 光滑且射影），且 \(q\) 是一个不同于 \(S\) 的各剩余特征的素数，则 \(X/S\) 的几何纤维的 Néron–Severi 挠群的 \(q\)-准素分量在 \(S\) 的每个连通分支上恒定。不过，在特征 \(p>0\) 时，对于挠群的 \(p\)-准素分量，情形就不再如此。然而，\(\underline {\operatorname {Pic}}_{X_s/k(s)}/\underline {\operatorname {Pic}}_{X_s/k(s)}^{00}=T_{X_s/k(s)}\) 在各 \(k(s)\) 上的秩仍有可能局部常值。当 \(S\) 约化时，可以证明，这等价于 \(\underline {\operatorname {Pic}}_{X/S}^{00}\) 存在，且 \(\underline {\operatorname {Pic}}_{X/S}\) 在 \(S\) 上 \emph{平坦}；并且只须对 \(S\) 是离散赋值环的谱的情形验证。这一点在我考察过的少数例子中确实成立；但由于 \(S\) 为 Artin 环的谱时，相应陈述错误（参见 \hyperref[fga3.vi-2-remark-2.9]{注~\ref*{fga3.vi-2-remark-2.9}} 和 \hyperref[fga3.vi-3-remarks-3.4]{注~\ref*{fga3.vi-3-remarks-3.4}}），我们也不能过于乐观。\end{remark}\chapter{Picard 概形的有限性定理}\thispagestyle{fancy}\label{fga3.vi-4}\par{}设 \(f\colon  X\to  S\) 是射影平坦态射，且 \(\underline {\operatorname {Pic}}_{X/S}\) 存在。“Hilbert 多项式”\(Q\in \mathbb {Q}[t]\) 使我们可以将 \(\underline {\operatorname {Pic}}_{X/S}\) 分解为开集 \(\underline {\operatorname {Pic}}_{X/S}^Q\) 的和。若不加进一步假设，例如保证 \(\underline {\operatorname {Pic}}_{X/S}\) 在 \(S\) 上分离的假设，则这些开集在 \(S\) 上一般并非有限型；当 \(X\) 是“一条退化为两条直线的圆锥曲线”时就得到反例。不过，若 \(f\) 是 \emph{可分} 的，且具有 \emph{不可约几何纤维}，则这可能成立。这个问题与 \(\underline {\operatorname {Pic}}_{X/S}^\tau \) 是否在 \(S\) 上有限型有关；后者可能不对 \(X/S\) 的纤维作任何假设就成立。当 \(f\colon  X\to  S\) 光滑时，注意到 \(\underline {\operatorname {Pic}}_{X/S}^\tau \) 包含在某个 \(\underline {\operatorname {Pic}}_{X/S}^Q\) 中（因为在非奇异射影簇上，对于除子，“挠”等价细于数值等价——实际上由 Matsusaka 的结果，二者相同）；因此，若各 \(\underline {\operatorname {Pic}}_{X/S}^Q\) 在 \(S\) 上有限型，它也有限型。注意，刚才提出的这些有限性问题，即使不假设 \(\underline {\operatorname {Pic}}_{X/S}\) 存在，仍有意义，因为可以将它们表述为某些可逆模族在 \hyperref[fga3.iv]{[3]，IV} 的意义下“有界”。对于每个代数闭域 \(k\)，考虑 \(\mathbb {P}_k^r\) 中维数为 \(n\)、次数为 \(d\) 的整子概形，以及这些预概形上 Hilbert 多项式为 \(Q\) 的可逆模（其中 \(r\)、\(n\)、\(d\)、\(Q\) 均给定）；要证明这些模构成的族（视为各 \(\mathbb {P}_k^r\) 上的凝聚模）是 \emph{有界} 的，即可由 \(\mathbb {Z}\) 上的有限型概形参数化。\par{}利用 Matsusaka 的方法 \hyperref[Mat1957]{[6]}，一个相当技术性的证明（以适当形式使用“等价判据”）给出：若只考虑 \(\mathbb {P}_k^r\) 的非奇异子簇，答案为肯定。更确切地说，有下述结果：\begin{theorem}{4.1}\setnumber{4.1}\label{fga3.vi-4-theorem-4.1}\skipsinglepar\par{}设 \(f\colon  X\to  S\) 是几何纤维连通的光滑射影态射，\(S\) Noether。设 \(\mathcal {O}_X(1)\) 是 \(X\) 上相对于 \(S\) 的甚丰沛模，\(E\) 是 \(\underline {\operatorname {Pic}}_{X/S}\) 的子集，对应于 \(X/S\) 的几何纤维 \(\overline {X}_{s_i}\) 上的可逆模族 \((\mathscr {L}_i)\)；设 \(D_i\) 是 \(\overline {X}_{s_i}\) 上定义 \(\mathscr {L}_i\) 的除子（未必为正），\(a_n^{(i)}t^n+\ldots +a_0^{(i)}\) 是 \(\mathscr {L}_i\) 的 Hilbert 多项式，\(\xi =\xi _i\) 是定义 \(\mathcal {O}_X(1)\) 的除子，即一个超平面截面。\par{}那么，下列条件等价：\begin{enumerate}
    
 \item[a.]\(Q\) 拟紧，即包含在 \(S\) 上某个有限型开集中，也即族 \((\mathscr {L}_i)\) 有界。

    
 \item[b.]\(E\) 包含在有限多个集合 \(\underline {\operatorname {Pic}}_{X/S}^Q\)（\(Q\in \mathbb {Q}[t]\)）之并中，即各 \(\mathscr {L}_i\) 的 Hilbert 多项式构成的集合有限。

    
 \item[b'.]（若 \(X/S\) 的所有纤维具有相同维数 \(n\)。）各 \(\mathscr {L}_i\) 的 Hilbert 多项式的系数 \(a_{n-1}^{(i)}\) 和 \(a_{n-2}^{(i)}\) 包含在某个有限集中。

    
 \item[b''.]（若 \(X/S\) 的所有纤维具有相同维数 \(n\)。）整数 \(\xi ^{n-1}D_i\) 和 \(\xi ^{n-2}D_i^2\)（即 \(D_i\) 和 \(D_i^2\) 的射影次数）有上界。

  \end{enumerate}\end{theorem}\begin{corollary}{4.2}\setnumber{4.2}\label{fga3.vi-4-corollary-4.2}\skipsinglepar\par{}设 \(f\colon  X\to  S\) 是几何纤维连通的光滑射影态射。那么，概形 \(\underline {\operatorname {Pic}}_{X/S}^Q\)（\(Q\in \mathbb {Q}[t]\)）和 \(\underline {\operatorname {Pic}}_{X/S}^\tau \) 都在 \(S\) 上射影。\begin{proof}\par{}由 \hyperref[fga3.vi-4-theorem-4.1]{定理~\ref*{fga3.vi-4-theorem-4.1}}，它们在 \(S\) 上有限型，因此可应用 \hyperref[fga3.vi-2-corollary-2.4]{推论~\ref*{fga3.vi-2-corollary-2.4}}。\end{proof}\end{corollary}\begin{remark}{}\label{fga3.vi-4-remark-comp}\skipsinglepar\par{}\emph{[编者]} 自本报告编写以来，本节所提出的这一类有限性问题几乎已全部解决。这里指出目前已知的主要事实。为简化陈述，以下默认为各陈述中出现的所有 Picard 预概形都存在，尽管对陈述作显然的修改，就可以去掉任何存在性假设。以下 \(S\) 表示 Noether 概形，\(X\) 和 \(Y\) 是 \(S\) 上的固有概形。\begin{enumerate}
      
 \item[i.]设 \(f\colon  X\to  Y\) 是满射态射。那么 \(f^*\colon \underline {\operatorname {Pic}}_{Y/S}\to \underline {\operatorname {Pic}}_{X/S}\) 是有限型态射。


      
 \item[ii.]假设 \(Y\) 在 \(S\) 上射影，配有相对于 \(S\) 的丰沛可逆模，设 \(X\) 是此模的任意一个截面的零点概形。令 \(f\colon  X\to  Y\) 为典范浸入。最后，假设 \(Y/S\) 的纤维的不可约分支的维数为 \(\geqslant 3\)。那么 \(f^*\colon \underline {\operatorname {Pic}}_{Y/S}\to \underline {\operatorname {Pic}}_{X/S}\) 为有限型。


      
 \item[iii.]假设 \(X\) 在 \(S\) 上射影，且其所有几何纤维整，维数为 \(n\)。设 \(\mathcal {O}_X(1)\) 是 \(X\) 上的丰沛可逆模，用它来定义 Hilbert 多项式。设 \(M\) 是 \(\underline {\operatorname {Pic}}_{X/S}\) 的子集。那么，\(M\) 拟紧，当且仅当 \(M\) 的各元素的 Hilbert 多项式 \(a_0x^n+a_1x^{n-1}+a_2x^{n-2}+\ldots +a_n\) 中，系数 \(a_1\) 和 \(a_2\) 有界。


      
 \item[iv.]对每个整数 \(n\neq 0\)，群预概形 \(\underline {\operatorname {Pic}}_{X/S}\) 中的 \(n\) 次幂同态是有限型态射。

    \end{enumerate}\par{}注意，(i) 和 (ii) 还蕴含（在各自陈述的假设下）\emph{\(\underline {\operatorname {Pic}}_{Y/S}\) 的子集 \(M\) 拟紧，当且仅当它在 \(\underline {\operatorname {Pic}}_{X/S}\) 中的像拟紧}。因此，\(Y\) 上的可逆模 \(\mathscr {L}\) 与 \(0\) 是 \(\tau \)-等价的，当且仅当它在 \(X\) 中的逆像亦然；换言之，\(\underline {\operatorname {Pic}}_{Y/S}^\tau \) 是 \(\underline {\operatorname {Pic}}_{X/S}^\tau \) 的逆像。特别地，要证明第一个预概形在 \(S\) 上有限型，只须证明第二个如此，因为 \(\underline {\operatorname {Pic}}_{Y/S}\to \underline {\operatorname {Pic}}_{X/S}\) 为有限型。再利用 (i)、Chow 引理及 (iii)，得到：\begin{enumerate}
      
 \item[v.]\(\underline {\operatorname {Pic}}_{X/S}^\tau \) 在 \(S\) 上有限型。

    \end{enumerate}\par{}一般地，将有限态射的情形中的 (i) 与 (ii) 合用，对于大部分有限性问题，可以归约到 \(X/S\) 的几何纤维整、正规且维数为 \(\leqslant 2\) 的情形；往往甚至可以将满射但未必有限的态射的情形中的 (i)，与代数曲面的解奇点（在任意特征下由 Abhyankar 给出）合用，归约到 \(X/S\) 甚至光滑、从而具有维数为 \(2\) 的非奇异几何纤维的情形。例如，结合 (v) 及界定非奇异射影曲面的 Néron–Severi 群的秩的 Picard–Igusa 不等式，可证明 Néron 有限性定理的下述推广：\begin{enumerate}
      
 \item[vi.]设 \(X/S\) 在 \(S\) 上固有，除此之外任意。那么，\(X/S\) 的几何纤维 \(X_i/k_i\) 的 Néron–Severi 群 \(\underline {\operatorname {Pic}}_{X_i/k_i}/\underline {\operatorname {Pic}}_{X_i/k_i}^0\) 都是有限生成的，且它们的秩以及挠子群的阶有界。

    \end{enumerate}\par{}同样，归约到非奇异曲面的情形，再利用此情形中的已知定理（如 Néron 有限性定理及 Néron–Severi 群上的相交型非退化这一事实），便得到：\begin{enumerate}
      
 \item[vii.]设 \(X\) 是代数闭域上的固有概形。那么，\(X\) 中存在有限条整闭曲线 \(C_i\)（\(1\leqslant  i\leqslant  r\)），满足下列性质：对 \(\underline {\operatorname {Pic}}_{X/k}\) 的每个子集 \(M\)，\(M\) 拟紧，当且仅当整数 \(\deg \mathscr {L}_{C'_i}\)（\(\mathscr {L}\in  M\)）有界（其中 \(C'_i\) 表示 \(C_i\) 的正规化）。

    \end{enumerate}在上述陈述中，可以取 \(r\) 为 Néron–Severi 群的秩。一旦知道此群为有限型，(vii) 就归结为下述事实：由 \(X\) 中的曲线 \(C\) 定义的 Néron–Severi 群上的线性型，只有在 Néron–Severi 群的挠元上才同时为零。在 \(X\) 非奇异且射影的情形中，这一结果与 (v) 都由 Matsusaka 给出。利用 (vii)，容易得到 \(X\) 上与 \(0\) 为 \(\tau \)-等价的可逆模的下述刻画：\begin{enumerate}
      
 \item[viii.]
        \par{}设 \(X/k\) 是一个域上的固有概形，\(\mathscr {L}\) 是 \(X\) 上的可逆模。那么下列条件等价：\begin{enumerate}
            
 \item[a.]\(\mathscr {L}\) 与 \(0\) 是 \(\tau \)-等价的。

            
 \item[b.]对每个凝聚模 \(F\)，有 \(\chi (F\otimes \mathscr {L})=\chi (F)\)，其中 \(\chi \) 表示 Euler–Poincaré 特征。

            
 \item[b'.]与 (b) 相同，但取 \(F=\mathcal {O}_Y\)，其中 \(Y\) 是 \(X\) 中维数为 \(1\) 的整闭子概形。

            
 \item[c.]对上述每个 \(Y\)，以 \(Y'\) 表示其正规化曲线，有 \(\deg \mathscr {L}_{Y'}=0\)。

          \end{enumerate}若 \(X/k\) 射影且配有丰沛可逆模 \(\mathcal {O}_X(1)\)，上述条件还等价于：\begin{enumerate}
            
 \item[d.]对每个整数 \(m\)，\(\mathscr {L}^{\otimes  m}(1)\) 都丰沛。

            
 \item[e.]（若 \(X\) 整。）对每对整数 \((m,n)\)，有 \[                 \chi (\mathscr {L}^{\otimes  m}(n))                 = \chi (\mathcal {O}(n))               \]（即 (b) 对 \(F=\mathscr {L}^{\otimes  m}(n)\) 成立）。

          \end{enumerate}
      

    \end{enumerate}\par{}为证明最后一个条件的充分性，注意它蕴含各 \(\mathscr {L}^{\otimes  m}\) 的 Hilbert 多项式全都相等；于是由 Mumford 判据 (iii)，各 \(\mathscr {L}^{\otimes  m}\) 留在 \(\underline {\operatorname {Pic}}_{X/k}\) 的某个拟紧子集中，即满足 (a)。条件 (b)、(b') 和 (c) 应视为通常对非奇异射影簇定义的 \emph{数值等价} 概念在任意固有概形上的变式。对于这样的簇，(a) 与 (c) 的等价性显然早已知道（Matsusaka）。\par{}上述判据 (e) 还蕴含下列结果：\begin{enumerate}
      
 \item[ix.]设 \(f\colon  X\to  S\) 是几何纤维整的平坦射影态射。那么 \(\underline {\operatorname {Pic}}_{X/S}^\tau \) 在 \(\underline {\operatorname {Pic}}_{X/S}\) 中开 \emph{且闭}。

    \end{enumerate}\par{}我们只就关键结果 (i)、(ii)、(iii) 的证明作一些评注（结果 (iv) 与其他结果稍有不同，仅用根式满射有限态射的情形中的 (i)，更确切地说，仅用 Frobenius 态射的情形，就可证明）。对于 (i)，非平坦下降的思想（见 \hyperref[fga3.i-a.2.c]{TDTE I，第9页}）起着本质作用。可以发现，（若只考虑有限性结果，）缺少下降数据的有效性判据并无妨碍。Mumford 最近证明了 (iii) 的一个稍弱形式，其中判据用到 Hilbert 多项式的 \emph{全部} 系数。他的论证极为简单，受到 Nakai 的一个丰沛性判据的证明的启发（Nakai 对非奇异曲面陈述了该判据，Mumford 将它推广到任意射影态射）。在我看来，这一论证需要对 \(X/S\) 的纤维附加一个温和的限制（Serre 的（\(\text {S}_2\)）性质）；当几何纤维正规时，此限制得到满足。然后，在 (ii) 的证明中使用这一受限判据：判据 (i) 使我们可以归约到 \(Y/S\) 平坦、几何纤维整且正规的情形；应用 Mumford 判据，又容易归约到 \(X/S\) 满足相同条件的情形。由维数假设，\(Y\) 和 \(X\) 的几何纤维在其闭点处的深度为 \(\geqslant 2\)；这就可以按 [\hyperref[fga3.vi-bib-5]{5}，XII] 中给出的形式应用“等价判据”，从而完成 (ii) 的证明。有了 (i) 和 (ii) 以后，在 Mumford 判据中去掉纤维正规的假设，并证明 (iii) 所给的更强形式，就不困难了。最后注意，(i) 和 (ii) 的证明还表明，当 \(S\) 是域 \(k\) 的谱时，态射 \(\underline {\operatorname {Pic}}_{Y/k}\to \underline {\operatorname {Pic}}_{X/k}\) 是 \emph{仿射} 的（不仅仅是有限型的）。\end{remark}
\section*{原报告参考文献}
\begin{enumerate}
\item[{[1]}]\setnumber{1}\label{fga3.vi-bib-1} Chevalley (Claude). \emph{Sur la théorie de Picard}. Amer. J. of Math., t.~82, 1960, pp.~435--490.
\item[{[2]}]\setnumber{2}\label{fga3.vi-bib-2} Grothendieck (A.) et Dieudonné (J.). \emph{Éléments de géométrie algébrique. Chapitre I et suivants}. Paris, Presses universitaires de France, 1960--1961 (Institut des Hautes Études Scientifiques, Publications mathématiques 4, 8, 11, \ldots).
\item[{[3]}]\setnumber{3}\label{fga3.vi-bib-3} Grothendieck (Alexander). \emph{Technique de descente et théorèmes d'existence en géométrie algébrique, I--V}. Séminaire Bourbaki, t.~12, 1959/60, n\textsuperscript{o}~190, 29~页；n\textsuperscript{o}~195, 22~页；t.~13, 1960/61, n\textsuperscript{o}~212, 20~页；n\textsuperscript{o}~221, 28~页；t.~14, 1961/62, n\textsuperscript{o}~232, 19~页。
\item[{[4]}]\setnumber{4}\label{fga3.vi-bib-4} Grothendieck (Alexander). \emph{Séminaire de géométrie algébrique, 1960/61}. Paris, Institut des Hautes Études Scientifiques, 1961.
\item[{[5]}]\setnumber{5}\label{fga3.vi-bib-5} Grothendieck (Alexander). \emph{Séminaire de géométrie algébrique, 1961/62}，由一组听众撰写。Paris, Institut des Hautes Études Scientifiques（待出版）。
\item[{[6]}]\setnumber{6}\label{fga3.vi-bib-6} Matsusaka (T.). \emph{The criteria for algebraic equivalence and the torsion group}. Amer. J. Math., t.~79, 1957, pp.~53--66.
\item[{[7]}]\setnumber{7}\label{fga3.vi-bib-7} Mumford (D.) and Tate (J.). \emph{Séminaire de géométrie algébrique}. Harvard University, Spring Term 1962（待出版）。
\item[{[8]}]\setnumber{8}\label{fga3.vi-bib-8} Serre (Jean-Pierre). \emph{Sur la topologie des variétés algébriques en caractéristique $p$}. \emph{Symposium internacional de topologia algebraica} [1956, Mexico], pp.~24--53. Mexico, Universidad nacional autonoma, 1958.
\end{enumerate}

% END chapters/fga3.vi.tex

\backmatter
\pagestyle{plain}
\nocite{*}
\begingroup
\emergencystretch=1em
\printbibliography[heading=bibintoc,title=参考文献]
\endgroup
\end{document}
